% !TeX program = xelatex
% Standalone English volume. Use XeLaTeX until the references stabilize.
\documentclass[a4paper,11pt]{article}
\usepackage[T1]{fontenc}
\usepackage{lmodern}
\usepackage[english]{babel}
\usepackage[left=26mm,right=26mm,top=25mm,bottom=25mm,headheight=15pt,headsep=7mm,footskip=11mm]{geometry}
\usepackage{amsmath,amssymb,mathtools,amsthm}
\newtheorem{proposition}{Proposition}[section]
\newtheorem{lemma}[proposition]{Lemma}
\newtheorem{theorem}[proposition]{Theorem}
\newtheorem{corollary}[proposition]{Corollary}
\newtheorem{maintheorem}[proposition]{Main Theorem}
\theoremstyle{definition}
\newtheorem{definition}[proposition]{Definition}
\theoremstyle{remark}
\newtheorem{remark}[proposition]{Remark}
\makeatletter
\renewcommand*\l@section{\@dottedtocline{1}{0em}{2.4em}}
\renewcommand*\l@subsection{\@dottedtocline{2}{1.45em}{3.45em}}
\renewcommand*\l@subsubsection{\@dottedtocline{3}{3.8em}{4.55em}}
\makeatother
\usepackage{xcolor,fancyhdr,titlesec,needspace,xurl}
\usepackage{longtable,multicol}
\usepackage{hyperref}
\definecolor{ink}{HTML}{1A3542}
\definecolor{muted}{HTML}{4F606A}
\hypersetup{unicode=true,colorlinks=true,linkcolor=ink,filecolor=ink,urlcolor=ink,citecolor=ink,bookmarksnumbered=true,
 pdfdisplaydoctitle=true,pdflang={en},
 pdftitle={The generalized Riemann hypothesis for Dirichlet L-functions. Edition 930, Volume V},
 pdfsubject={A natural-language reconstruction of the proposed proof},
 pdfauthor={A mathematical reconstruction}}
\titleformat{\section}{\Large\bfseries\color{ink}}{\thesection}{.65em}{}
\titlespacing*{\section}{0pt}{.8em}{.65em}
\setlength{\parindent}{1.2em}
\linespread{1.045}
\setlength{\parskip}{.22em plus .05em minus .03em}
\setlength{\emergencystretch}{2em}
\setcounter{topnumber}{2}
\setcounter{tocdepth}{3}\setcounter{secnumdepth}{3}
\widowpenalty=4500
\clubpenalty=4500
\pagestyle{fancy}\fancyhf{}
\fancyhead[L]{\small\color{muted}Dirichlet GRH \textperiodcentered\ Volume I}
\fancyhead[R]{\small\color{muted}Dirichlet zeros and power intervals}
\fancyfoot[C]{\small\thepage}
\renewcommand{\headrulewidth}{.3pt}
\newcommand{\lead}[1]{\par\smallskip\noindent{\bfseries\color{ink}#1.}\enspace}
\newcommand{\sour}[2]{\unskip\penalty150\hspace{.35em}{\small\hyperlink{ref#1}{[#2]}}}
\titleformat{\subsection}{\large\bfseries}{\thesubsection}{.6em}{}
\titlespacing*{\subsection}{0pt}{1.15em plus .15em}{.4em}
\setlength{\abovedisplayskip}{9pt plus 2pt minus 3pt}
\setlength{\belowdisplayskip}{9pt plus 2pt minus 3pt}
\newcommand{\RR}{\mathbb R}\newcommand{\CC}{\mathbb C}\newcommand{\QQ}{\mathbb Q}
\newcommand{\ZZ}{\mathbb Z}\newcommand{\ord}{\operatorname{ord}}
\newcommand{\Frob}{\operatorname{Frob}}\newcommand{\Gal}{\operatorname{Gal}}
\newcommand{\len}{\operatorname{length}}\newcommand{\Res}{\operatorname*{Res}}
\newcommand{\wmean}{\operatorname{av}}\newcommand{\Ker}{\operatorname{ker}}
\newcommand{\Lh}{\mathcal L}\newcommand{\Gg}{\mathcal G}
\newcommand{\real}{\operatorname{Re}}\newcommand{\imag}{\operatorname{Im}}
\newcommand{\abs}[1]{\left|#1\right|}\newcommand{\eps}{\varepsilon}
\newcommand{\anchoritem}[2]{\par\Needspace{2.8\baselineskip}\hypertarget{ref#1}{}\noindent{\bfseries\color{ink}#2.}\enspace}
\newcommand{\srcloc}[2]{{\small\nolinkurl{#1.lean}}, line~#2}
\numberwithin{equation}{section}
\raggedbottom

\newcommand{\Sel}{\mathcal S}
\newcommand{\Cl}{\operatorname{Cl}}
\newcommand{\Aut}{\operatorname{Aut}}
\newcommand{\GL}{\operatorname{GL}}
\newcommand{\OO}{\mathcal O}
\hypersetup{bookmarksdepth=3}


\addto\captionsenglish{\renewcommand{\partname}{Volume}}


\makeatletter
\expandafter\def\csname r@cor:dyadiccyclic661\endcsname{{22.60}{110}{}{proposition.22.60}{Dirichlet-GRH-v930-Volume-IV.pdf}}
\expandafter\def\csname r@cor:levels113\endcsname{{11.23}{21}{}{proposition.11.23}{Dirichlet-GRH-v930-Volume-II.pdf}}
\expandafter\def\csname r@cor:tower27\endcsname{{19.11}{9}{}{proposition.19.11}{Dirichlet-GRH-v930-Volume-IV.pdf}}
\expandafter\def\csname r@cor:transfer100\endcsname{{11.18}{17}{}{proposition.11.18}{Dirichlet-GRH-v930-Volume-II.pdf}}
\expandafter\def\csname r@eq:budget\endcsname{{30.6}{12}{}{equation.30.6}{Dirichlet-GRH-v930-Volume-VI.pdf}}
\expandafter\def\csname r@eq:ideal\endcsname{{30.5}{11}{}{equation.30.5}{Dirichlet-GRH-v930-Volume-VI.pdf}}
\expandafter\def\csname r@eq:idele113\endcsname{{11.6}{20}{}{equation.11.6}{Dirichlet-GRH-v930-Volume-II.pdf}}
\expandafter\def\csname r@eq:qbound102\endcsname{{20.5}{27}{}{equation.20.5}{Dirichlet-GRH-v930-Volume-IV.pdf}}
\expandafter\def\csname r@eq:strict113\endcsname{{11.7}{21}{}{equation.11.7}{Dirichlet-GRH-v930-Volume-II.pdf}}
\expandafter\def\csname r@eq:tail114\endcsname{{19.15}{17}{}{equation.19.15}{Dirichlet-GRH-v930-Volume-IV.pdf}}
\expandafter\def\csname r@lem:fit13\endcsname{{11.2}{3}{}{proposition.11.2}{Dirichlet-GRH-v930-Volume-II.pdf}}
\expandafter\def\csname r@lem:hecke171\endcsname{{72.9}{53}{}{proposition.72.9}{Dirichlet-GRH-v930-Volume-XIV.pdf}}
\expandafter\def\csname r@lem:locallevel345\endcsname{{22.16}{75}{}{proposition.22.16}{Dirichlet-GRH-v930-Volume-IV.pdf}}
\expandafter\def\csname r@nav:I-26\endcsname{{6.4}{28}{}{subsection.6.4}{Dirichlet-GRH-v930-Volume-I.pdf}}
\expandafter\def\csname r@nav:I-28\endcsname{{6.6}{32}{}{subsection.6.6}{Dirichlet-GRH-v930-Volume-I.pdf}}
\expandafter\def\csname r@nav:II-00\endcsname{{11.1}{3}{}{subsection.11.1}{Dirichlet-GRH-v930-Volume-II.pdf}}
\expandafter\def\csname r@nav:II-08\endcsname{{11.2}{6}{}{subsection.11.2}{Dirichlet-GRH-v930-Volume-II.pdf}}
\expandafter\def\csname r@nav:III-01\endcsname{{30.5}{10}{}{subsection.30.5}{Dirichlet-GRH-v930-Volume-VI.pdf}}
\expandafter\def\csname r@nav:volume-IV\endcsname{{19}{3}{}{section.19}{Dirichlet-GRH-v930-Volume-IV.pdf}}
\expandafter\def\csname r@prop:example235\endcsname{{21.7}{44}{}{proposition.21.7}{Dirichlet-GRH-v930-Volume-IV.pdf}}
\expandafter\def\csname r@prop:finite15\endcsname{{11.4}{6}{}{proposition.11.4}{Dirichlet-GRH-v930-Volume-II.pdf}}
\expandafter\def\csname r@prop:moment267\endcsname{{9.5}{110}{}{proposition.9.5}{Dirichlet-GRH-v930-Volume-I.pdf}}
\expandafter\def\csname r@prop:qbound102\endcsname{{20.6}{27}{}{proposition.20.6}{Dirichlet-GRH-v930-Volume-IV.pdf}}
\expandafter\def\csname r@prop:realclass100\endcsname{{11.19}{18}{}{proposition.11.19}{Dirichlet-GRH-v930-Volume-II.pdf}}
\expandafter\def\csname r@prop:residue269\endcsname{{21.10}{48}{}{proposition.21.10}{Dirichlet-GRH-v930-Volume-IV.pdf}}
\expandafter\def\csname r@sec:beforecyclic661\endcsname{{22.27}{108}{}{subsection.22.27}{Dirichlet-GRH-v930-Volume-IV.pdf}}
\expandafter\def\csname r@sec:coho12\endcsname{{6.4}{28}{}{subsection.6.4}{Dirichlet-GRH-v930-Volume-I.pdf}}
\expandafter\def\csname r@sec:missing\endcsname{{6.6}{32}{}{subsection.6.6}{Dirichlet-GRH-v930-Volume-I.pdf}}
\expandafter\def\csname r@seriesguide\endcsname{{1.1}{4}{}{subsection.1.1}{Dirichlet-GRH-v930-Volume-I.pdf}}
\expandafter\def\csname r@thm:alltrunc329\endcsname{{9.16}{120}{}{proposition.9.16}{Dirichlet-GRH-v930-Volume-I.pdf}}
\expandafter\def\csname r@thm:common279\endcsname{{9.4}{109}{}{proposition.9.4}{Dirichlet-GRH-v930-Volume-I.pdf}}
\expandafter\def\csname r@thm:commonlength267\endcsname{{21.9}{46}{}{proposition.21.9}{Dirichlet-GRH-v930-Volume-IV.pdf}}
\expandafter\def\csname r@thm:conductor232\endcsname{{19.15}{12}{}{proposition.19.15}{Dirichlet-GRH-v930-Volume-IV.pdf}}
\expandafter\def\csname r@thm:critical13\endcsname{{11.1}{3}{}{proposition.11.1}{Dirichlet-GRH-v930-Volume-II.pdf}}
\expandafter\def\csname r@thm:cyclic278\endcsname{{19.24}{19}{}{proposition.19.24}{Dirichlet-GRH-v930-Volume-IV.pdf}}
\expandafter\def\csname r@thm:cyclicdescent661\endcsname{{22.59}{109}{}{proposition.22.59}{Dirichlet-GRH-v930-Volume-IV.pdf}}
\expandafter\def\csname r@thm:decomploss330\endcsname{{22.8}{67}{}{proposition.22.8}{Dirichlet-GRH-v930-Volume-IV.pdf}}
\expandafter\def\csname r@thm:deep114\endcsname{{19.19}{16}{}{proposition.19.19}{Dirichlet-GRH-v930-Volume-IV.pdf}}
\expandafter\def\csname r@thm:integralblock656\endcsname{{22.57}{107}{}{proposition.22.57}{Dirichlet-GRH-v930-Volume-IV.pdf}}
\expandafter\def\csname r@thm:intersection279\endcsname{{19.25}{21}{}{proposition.19.25}{Dirichlet-GRH-v930-Volume-IV.pdf}}
\expandafter\def\csname r@thm:levelstrict345\endcsname{{22.17}{75}{}{proposition.22.17}{Dirichlet-GRH-v930-Volume-IV.pdf}}
\expandafter\def\csname r@thm:logstrict319\endcsname{{22.5}{64}{}{proposition.22.5}{Dirichlet-GRH-v930-Volume-IV.pdf}}
\expandafter\def\csname r@thm:long232\endcsname{{5.6}{23}{}{proposition.5.6}{Dirichlet-GRH-v930-Volume-I.pdf}}
\expandafter\def\csname r@thm:norm91\endcsname{{94.9}{44}{}{proposition.94.9}{Dirichlet-GRH-v930-Volume-XVIII.pdf}}
\expandafter\def\csname r@thm:raylength656\endcsname{{22.55}{106}{}{proposition.22.55}{Dirichlet-GRH-v930-Volume-IV.pdf}}
\expandafter\def\csname r@thm:simultaneous307\endcsname{{34.22}{85}{}{proposition.34.22}{Dirichlet-GRH-v930-Volume-VI.pdf}}
\expandafter\def\csname r@thm:specializations280\endcsname{{19.26}{22}{}{proposition.19.26}{Dirichlet-GRH-v930-Volume-IV.pdf}}
\expandafter\def\csname r@thm:strictdigits329\endcsname{{22.7}{66}{}{proposition.22.7}{Dirichlet-GRH-v930-Volume-IV.pdf}}
\expandafter\def\csname r@thm:towerzero267\endcsname{{8.2}{100}{}{proposition.8.2}{Dirichlet-GRH-v930-Volume-I.pdf}}
\expandafter\def\csname r@thm:twolayers327\endcsname{{22.4}{62}{}{proposition.22.4}{Dirichlet-GRH-v930-Volume-IV.pdf}}
\expandafter\def\csname r@thm:wholeblock661\endcsname{{22.58}{108}{}{proposition.22.58}{Dirichlet-GRH-v930-Volume-IV.pdf}}
\makeatother
\begin{document}
\setcounter{part}{4}
\fancyhead[L]{\small\color{muted}Dirichlet GRH \textperiodcentered\ Volume V}
\fancyhead[R]{\small\color{muted}Class groups and special values}
\thispagestyle{plain}\begin{center}
{\LARGE\bfseries The generalized Riemann hypothesis\\[2pt]for Dirichlet \(L\)-functions}\\[6pt]
{\small 27 November 2026 \quad\textperiodcentered\quad Edition 930}\end{center}
\part{Class groups, special values and strict orders}\label{part:V}
\begin{abstract}\noindent We express strict orders in terms of Kummer classes, special values and norm maps. We use a single determinant to measure full strict orders at several twists, including classes that capitulate. For an extension ramified at two primes we prove that one decomposition group can detect every strict class. This gives full vanishing for specified character components across all twists with trivial residual character.\end{abstract}
\setcounter{section}{23}
\noindent The \hyperref[seriesguide]{guide in Volume~I} lists all twenty-seven volumes.
\begingroup\small\setlength{\parskip}{0pt}\tableofcontents\endgroup
\clearpage

\section{Kummer classes and real class groups}\label{nav:g-III-0}\label{nav:volume-V}

\subsection{Strict classes as triangular extensions}\label{nav:III-02}
We describe the arithmetic group whose length would have to be compared with \eqref{eq:budget}. Let \(E/\QQ\) be finite abelian with group \(G\), let \(\ell\nmid |G|\) be odd, and choose the valuation ring \(\OO\) of a finite unramified extension \(J/\QQ_\ell\) containing all values of the characters of \(G\). Such a ring exists because \(|G|\) is prime to \(\ell\). For a character \(\alpha:G\to\OO^\times\) put
\[
  e_\alpha=|G|^{-1}\sum_{g\in G}\alpha(g)^{-1}g,\qquad
  \vartheta=\kappa_\ell^{-r}\alpha^{-1},\qquad D=(J/\OO)(\vartheta).
\]
Here the character of \(G\) is inflated to the absolute Galois group \(G_\QQ\). Define a strict group over \(\QQ\) as the subgroup of \(H^1(G_\QQ,D)\) of classes trivial on inertia at every \(p\ne\ell\) and trivial on the full decomposition group at \(\ell\).

\begin{theorem}[Character descent]\label{thm:char}
There is an isomorphism of \(\OO\)-modules between this strict group and \(e_\alpha(\Sel_{\ell,r}(E)\otimes_{\ZZ_\ell}\OO)\). The \(\Sel_{\ell,r}(E)\) in this formula has the same unramified and strict local conditions as in Volume~II. For \(E=E_x\) the theorem applies to every odd \(\ell\nmid\operatorname{ord}(\psi)\). The component with \(\alpha(c)=(-1)^r\) corresponds exactly to an even character \(\vartheta\).
\end{theorem}
\begin{proof}
On \(G_E\) the finite character \(\alpha\) is trivial. Restriction on absolute cohomology identifies \(H^1(G_\QQ,D)\) with \(H^1(G_E,D)^G\), by inflation--restriction: the outside terms in degrees one and two for the finite quotient acting on \(D^{G_E}\) vanish by averaging. Restricting coefficients to \(G_E\) and then twisting the \(G\)-action back by \(\alpha^{-1}\) identifies this invariant subspace with the displayed idempotent on unrestricted cohomology of \(E\). Tensoring by \(\OO\) commutes with continuous cohomology here, as \(\OO\) is finite free over \(\ZZ_\ell\). Inflation from \(G_{E,S_\ell}\) identifies unramified classes outside \(\ell\) with the classes having trivial inertia there, since the coefficient is unramified away from \(\ell\) over \(E\).

For inertia over \(p\ne\ell\), the quotient of the inertia group over \(\QQ\) by inertia over \(E\) has order dividing \(|G|\). Inflation--restriction therefore gives injectivity on local \(H^1\) upon restriction, including when \(\alpha\) is ramified. The same holds for full decomposition groups over \(\ell\). The local conditions descend in both directions. Finally \([E_x:\QQ]\mid 4\,2^l\operatorname{ord}(\psi)\), and \(\vartheta(c)=(-1)^r\alpha(c)^{-1}\). This proves all assertions.
\end{proof}

Under this description a class is represented by a continuous function \(z:G_\QQ\to J/\OO\) satisfying
\[
 z(gh)=z(g)+\vartheta(g)z(h).
\]
Equivalently, it gives affine transformations \(y\mapsto\vartheta(g)y+z(g)\). Changing the origin adds a coboundary to \(z\). The strict local condition means that an origin can be chosen locally so that the translations disappear on the required local group. A continuous cocycle has finite image by compactness, hence takes values in \(\ell^{-m}\OO/\OO\) for some \(m\). After multiplication by \(\ell^m\) in representatives, identify this module with \(\OO/\ell^m\OO\) and write \(\bar z\) for its coordinate. The finite-level extensions have the matrices
\[
 g\longmapsto
 \begin{pmatrix}
  \vartheta(g)&\bar z(g)&0\\
  0&1&0\\
  0&0&\vartheta(g)^{-1}
 \end{pmatrix}\pmod{\ell^m}.
\]
Their multiplication law is precisely the cocycle law. Their trace is \(\vartheta(g)+1+\vartheta(g)^{-1}\) for every \(z\). At a good arithmetic Frobenius at \(p\), this is the reduction of \(d_\alpha\). Thus the upper-right entry carries extension information that these traces alone do not detect. For instance the nonzero homomorphism from the cyclic group of order \(\ell\) to the additive group of \(\mathbf F_\ell\) gives a nonsplit unipotent extension with the trace of the split extension at every element.

To deduce a bound for the strict groups from \eqref{eq:budget}, the proposed argument needs an integral deformation theorem that bounds their lengths using congruences with the characteristic-zero systems counted above. It also needs level and weight control for the systems that theorem produces. In the source archive, \nolinkurl{UnitaryThreeGroups/OneCharOf.lean} in the automorphic varieties directory, lines 102--138, uses a deformation ring, its integral model and automorphy of its points in deriving the single-character length inequality. Formula \eqref{eq:ideal} proves the congruence-ideal calculation for the Hecke traces. Theorem~\ref{thm:char} identifies the strict classes on the other side. A comparison of these two integral modules, with uniform length bounds, has not been established here.

\subsection{A class-group test for the complementary sign}\label{nav:III-03}
\label{sec:reflect}

Retain \(E=E_x\subseteq\QQ(\zeta_N)\), \(\alpha,\OO,J,\vartheta,D\) from Theorem~\ref{thm:char}. Suppose \(\ell\) is odd, \(\ell\nmid [E:\QQ]\), \(\alpha(c)=(-1)^r\), and the reduction \(\bar\vartheta\) is nontrivial on \(G_{\QQ_\ell}\). The last condition is automatic if \(\ell\nmid N\) and \(\ell-1\nmid r\): at such a prime \(\alpha\) is unramified and the mod-\(\ell\) cyclotomic character maps inertia onto \(\mathbf F_\ell^\times\). Denote the strict group over \(\QQ\) in Theorem~\ref{thm:char} by \(S_\alpha\). For \(m\ge1\) put
\[
 n_m=N\ell^m,\quad L_m=\QQ(\zeta_{n_m}),\quad
 \Gamma_m=\Gal(L_m/\QQ),\quad R_m=\OO/\ell^m\OO ,
 \qquad \beta_m=\kappa_{\ell,m}^{\,r+1}\alpha\bmod\ell^m .
\]
Here \(\kappa_{\ell,m}\) is the cyclotomic character modulo \(\ell^m\). For an \(R_m[\Gamma_m]\)-module \(U\), let \(U[\beta_m]\) be the subgroup on which every \(\gamma\) acts as \(\beta_m(\gamma)\). This definition also applies when \(\ell\mid |\Gamma_m|\).

\begin{theorem}[Reflection at every exponent]\label{thm:mirror}
After choosing a generator of \(\mu_{\ell^m}\), there is an \(R_m\)-linear injection
\begin{equation}
 S_\alpha[\ell^m]\ \hookrightarrow\
      \big(\Cl(\OO_{L_m})[\ell^m]\otimes_{\ZZ_\ell}\OO\big)[\beta_m].
 \label{eq:mirror}
\end{equation}
In particular the target lies in the minus part of the ideal class group.
\end{theorem}
\begin{proof}
Consider \(0\to D[\ell^m]\to D\xrightarrow{\ell^m}D\to0\). On a group \(U\) the kernel on \(H^1\) from finite to divisible coefficients is \(D^U/\ell^mD^U\). At \(U=G_\QQ\) and \(U=G_{\QQ_\ell}\) the invariants vanish: an element with \(\bar\vartheta\ne1\) acts by a scalar whose difference from one is a unit in \(\OO\). At inertia over \(p\ne\ell\), the cyclotomic character is unramified and the inertia action through \(\alpha\) has order prime to \(\ell\). If it is nontrivial its invariants vanish by the same argument. Otherwise \(D^U=D\) is divisible. In every case the finite-coefficient map is locally injective. Exactness globally now identifies \(S_\alpha[\ell^m]\) with the classes for \(R_m(\vartheta)\) which are unramified away from \(\ell\) and zero on the full decomposition group at \(\ell\).

The coefficient is trivial over \(L_m\). Restriction from \(G_\QQ\) to \(G_{L_m}\) on these degree-one classes is injective even when \(\ell\mid|\Gamma_m|\). Its kernel inflates from \(H^1(\Gamma_m,R_m(\vartheta))\). Choose \(\tau\in\Gamma_m\) with \(\vartheta(\tau)-1\) a unit. For a cocycle \(z\) on the abelian group \(\Gamma_m\), evaluate \(z(\tau\sigma)=z(\sigma\tau)\). The identity
\[
   (\vartheta(\tau)-1)z(\sigma)=(\vartheta(\sigma)-1)z(\tau)
\]
makes \(z\) the coboundary of \(z(\tau)/(\vartheta(\tau)-1)\). This proves injectivity.

Use the chosen root of unity to identify the trivial coefficient on \(G_{L_m}\) with \(\mu_{\ell^m}\otimes_{\ZZ/\ell^m} R_m\). Under this change the action on the restricted classes is by \(\beta_m=\kappa_{\ell,m}\vartheta^{-1}\). Indeed a restricted cocycle is a homomorphism \(b\) satisfying \(b(\sigma h\sigma^{-1})=\vartheta(\sigma)b(h)\). After identifying its values with roots of unity, conjugation on the values supplies the factor \(\kappa_{\ell,m}(\sigma)\); precomposition with \(h\mapsto\sigma^{-1}h\sigma\) supplies \(\vartheta(\sigma)^{-1}\). Kummer theory thus gives classes in \((L_m^\times/L_m^{\times\ell^m})\otimes R_m\) with this eigencharacter. Here tensoring over \(\ZZ/\ell^m\) is exact since \(R_m\) is finite free over that ring.

For these classes every finite valuation is a multiple of \(\ell^m\), coordinate by coordinate in a basis of \(R_m\). At places above \(\ell\), the local Kummer class is trivial. At the other places it becomes trivial over the maximal unramified extension, which has the same value group as the original local field. Taking the valuation of an \(\ell^m\)-th root there proves the assertion. If \((a)=\mathfrak a^{\ell^m}\), send its Kummer class to \([\mathfrak a]\). The ambiguity of multiplying \(a\) by an \(\ell^m\)-th power changes \(\mathfrak a\) by a principal ideal. The resulting equivariant map to \(\Cl(\OO_{L_m})[\ell^m]\) has kernel the global units modulo \(\ell^m\)-th powers: if \(\mathfrak a=(b)\), \(a/b^{\ell^m}\) is a unit. The same descriptions hold after tensoring.

The character \(\beta_m\) is odd because \(\vartheta\) is even. The odd part of the unit quotient is generated by roots of unity. In fact \(u/\bar u\) for any global unit has modulus one at every embedding and is an algebraic integer, hence is a root of unity by Kronecker's theorem. Since two is invertible, the differences \(1-c\) generate the odd part. On its root-of-unity generators the Galois action is \(\kappa_{\ell,m}\); prime-to-\(\ell\) torsion makes no contribution to this quotient. The character \(\beta_m\) differs from \(\kappa_{\ell,m}\) on the residue field by our choice of \(\tau\), so an element with both actions is zero. This proves injectivity in \eqref{eq:mirror}.
\end{proof}

\begin{theorem}[An explicit annihilator]\label{thm:stick}
For any positive integer \(b\) prime to \(N\ell\) define
\begin{equation}
 A_{b,m}(\alpha,r)=
   \sum_{\substack{1\le j\le n_m\\(j,n_m)=1}}
       \left\lfloor\frac{bj}{n_m}\right\rfloor
          \big(j^{r+1}\alpha(j)\big)^{-1}\quad\text{in }R_m .
 \label{eq:floor}
\end{equation}
This scalar kills \(S_\alpha[\ell^m]\). If it is nonzero and its valuation is \(v<m\), the entire finite group \(S_\alpha\) is killed by \(\ell^v\). In particular any nonzero residue at \(m=1\) proves \(S_\alpha=0\) at every exponent. The scalars in \eqref{eq:floor} are compatible under reduction as \(m\) increases.
\end{theorem}
\begin{proof}
Write \(\sigma_j(\zeta_{n_m})=\zeta_{n_m}^j\) and let
\[
   \Theta_n=\frac1n\sum_{1\le j\le n,\ (j,n)=1}j\,\sigma_j^{-1},
             \qquad n=n_m .
\]
The classical Stickelberger theorem states that the integral group-ring elements in \(\ZZ[\Gamma_m]\Theta_n\) annihilate the ideal class group of the full cyclotomic field. Apply it to \((b-\sigma_b)\Theta_n\). The coefficient of \(\sigma_j^{-1}\) here is \(bj/n-\{bj/n\}=\lfloor bj/n\rfloor\). Thus the element is integral. On the \(\beta_m\)-eigensubgroup in \eqref{eq:mirror} its action is multiplication by \eqref{eq:floor}. This proves the first assertion.~\cite{V-III-V-stick}

If the scalar has valuation \(v<m\), its unit factor is invertible in \(R_m\). Hence \(S_\alpha[\ell^m]\) is killed by \(\ell^v\). Any element of an \(\ell\)-primary group of order higher than \(\ell^v\) generates a cyclic subgroup containing an element of order exactly \(\ell^{v+1}\). That element already belongs to \(S_\alpha[\ell^m]\), a contradiction. This proves the assertion about the full group.

To prove compatibility put \(n=n_m\). Above each unit \(j\bmod n\) the units modulo \(n\ell\) are \(j+kn\), \(0\le k<\ell\), because \(\ell\mid n\). Their coefficients in \(\Theta_{n\ell}\) sum to \(j/n+(\ell-1)/2\). Thus projection to \(\QQ[\Gamma_m]\) takes \(\Theta_{n\ell}\) to \(\Theta_n+(\ell-1)\sum_{g\in\Gamma_m}g/2\). Multiplication by \(b-\sigma_b\) changes the difference into an integral multiple of \(\sum_g g\). Its value on \(\beta_m\) is zero, since \(g\) and \(cg\) have opposite values and \(\beta_m(c)=-1\). The character \(\beta_{m+1}\) reduces to \(\beta_m\), giving the claimed compatibility.
\end{proof}

\begin{corollary}[Application on the interval]\label{cor:kill}
For every \(x\) on the common intervals, every odd \(\ell\nmid N\operatorname{ord}(\psi)\), and \(1\le r<t(x)\) with \(\ell-1\nmid r\), Theorem~\ref{thm:stick} applies to each component of the sign \((-1)^r\) of the full group \(\Sel_{\ell,r}(E_x)\). It tests every twist in the sum if \(\ell>t(x)+1\). More explicitly, a nonzero test of valuation \(v\) bounds the corresponding length by
\[
   \len_{\OO}S_\alpha\ \le\
   v\,\dim_{\OO/\ell}\!\left(\big(\Cl(\OO_{L_1})[\ell]\otimes_{\mathbf F_\ell}(\OO/\ell)\big)[\beta_1]\right).
\]
\end{corollary}
\begin{proof}
The prime and twist conditions ensure nontrivial residual inertia at \(\ell\), as noted above. By Theorem~\ref{thm:char} the complementary sign is the direct sum of these components after extending to \(\OO\). Extension of scalars from \(\ZZ_\ell\) to the unramified \(\OO\) preserves length in their respective normalized units, by the structure theorem for finite modules. The number of cyclic summands of \(S_\alpha\) is \(\dim_{\OO/\ell} S_\alpha[\ell]\), bounded by Theorem~\ref{thm:mirror} with \(m=1\). Each summand has length at most \(v\).
\end{proof}

For example, take \(E=\QQ(i)\), \(N=4\), \(\ell=3\), and odd \(r\). The complementary character is the quadratic character \(\alpha\) modulo four. At \(m=1\) the units modulo twelve are \(1,5,7,11\); their values under \(\beta_1\) are \(1,1,-1,-1\), since \(r+1\) is even. With \(b=5\) the four floors are \(0,2,2,4\), giving \(A_{5,1}=-4\equiv2\bmod3\). Thus the complementary group vanishes at three for every such twist. This example illustrates a local test. In the fields \(E_x\) of the prime count the residue sums vary with \(N,\ell,\alpha,r\). Section~\ref{sec:all24} treats the other local types when the coefficient prime is prime to the field degree. The lengths for vanishing residues still require control. The divisibility theorem for Hecke congruences in Section~\ref{nav:III-01} and the independent tests here leave this uniform arithmetic estimate outstanding.

\subsection{Reflection without a local residual hypothesis}\label{nav:III-04}
\label{sec:all24}

Let \(E/\QQ\) be abelian and let \(\ell\nmid[E:\QQ]\) be odd. Use \(\alpha,\vartheta,D,S_\alpha\) as in Theorem~\ref{thm:char}, with \(\alpha(c)=(-1)^r\). Let \(k=\OO/\ell\OO\), write a bar for reduction to \(k\), and put \(f=\operatorname{cond}(\alpha)\) and \(f'=f/\ell^{v_\ell(f)}\). Since the order of \(\alpha\) is prime to \(\ell\), its local character at \(\ell\) kills the principal units, a pro-\(\ell\) group. Thus \(v_\ell(f)\le1\). For each \(m\ge1\) define
\[
 n=f'\ell^m,\qquad K_m=\QQ(\zeta_n),\qquad
 R_m=\OO/\ell^m\OO,\qquad \beta_m=\kappa_{\ell,m}\,(\vartheta\bmod\ell^m)^{-1}.
\]
Both \(\alpha\) and \(\kappa_{\ell,m}\) factor through \(\Gal(K_m/\QQ)\). Using this field allows the reflection test to discard auxiliary conductor factors belonging to other characters.

\begin{lemma}[Local valuation of a finite lift]\label{lem:localunit24}
Let \(F\) be a finite extension of \(\QQ_\ell\) containing \(\mu_{\ell^m}\), and fix a generator identifying \(\ZZ/\ell^m\) with \(\mu_{\ell^m}\). If a continuous homomorphism \(G_F\to\ZZ/\ell^m\) factors through \(\Gal(F(\mu_{\ell^\infty})/F)\), its Kummer class is represented by a unit of \(F\).
\end{lemma}
\begin{proof}
The cyclotomic image of \(G_F\) is open in \(1+\ell^m\ZZ_\ell\). The \(\ell\)-adic logarithm identifies this last group with \(\ell^m\ZZ_\ell\), so the image is precisely \(1+\ell^a\ZZ_\ell\) for some \(a\ge m\). A primitive \(\ell^a\)-th root \(u\) then belongs to \(F\). Choose \(v\) with \(v^{\ell^m}=u\) and of order \(\ell^{a+m}\). The Kummer homomorphism of \(u\) is \(g\mapsto g(v)/v=v^{\kappa_\ell(g)-1}\). It sends an element with cyclotomic value \(1+\ell^a\) to a primitive \(\ell^m\)-th root of unity. That element generates the cyclotomic image topologically, as its logarithm has valuation \(a\). Every homomorphism from this image to \(\mu_{\ell^m}\) is a multiple of the one for \(u\). All its Kummer classes are consequently represented by powers of the unit \(u\). After choosing a \(\ZZ/\ell^m\)-basis the same result applies to each coordinate for coefficients \(R_m\).
\end{proof}

\begin{theorem}[All complementary characters of prime-to-\(\ell\) order]\label{thm:all24}
If \(\bar\vartheta\) is trivial on \(G_\QQ\), then \(S_\alpha=0\). For every other \(\bar\vartheta\) and every \(m\ge1\), a choice of \(\ell^m\)-th root of unity of exact order \(\ell^m\) gives an injection
\begin{equation}
 S_\alpha[\ell^m]\ \lhook\joinrel\longrightarrow\
  \big(\Cl(\OO_{K_m})[\ell^m]\otimes_{\ZZ_\ell}\OO\big)[\beta_m].
 \label{eq:all24}
\end{equation}
This includes characters that become trivial modulo \(\ell\) on the entire decomposition group at \(\ell\).
\end{theorem}
\begin{proof}
First suppose \(\bar\vartheta=1\) globally. Let \(\omega\) be the Teichm\"uller lift of the mod-\(\ell\) cyclotomic character. Reduction is injective on roots of unity of order prime to \(\ell\) in \(\OO\). Thus \(\alpha=\omega^{-r}\) and \(\vartheta=(\kappa_\ell/\omega)^{-r}\). It is unramified at every finite prime other than \(\ell\). Let \(G_{\QQ,S_\ell}\) allow ramification at \(\ell\) and the infinite place. Every class in \(S_\alpha\) inflates from this quotient. Indeed, the inertia groups away from \(\ell\) act trivially on \(D\), so restriction to each of them is an actual homomorphism. The class condition makes it zero on those groups and hence on their closed normal closure.

We show that \(H^1(G_{\QQ,S_\ell},D)[\ell]=0\). The group \(D[\ell]\) has trivial action and identifies with the additive group of \(k\). Every continuous homomorphism \(G_{\QQ,S_\ell}\to k\) factors through the pro-\(\ell\) part of \(\Gal(\QQ(\mu_{\ell^\infty})/\QQ)\), which is \(1+\ell\ZZ_\ell\). This follows from Kronecker--Weber. In a finite cyclotomic field of modulus \(d\ell^a\), \((d,\ell)=1\), the factors at primes dividing \(d\) are generated by their inertia groups, using total ramification in each prime-power cyclotomic field and unramifiedness of the other factors at that prime. The characters here kill all such inertia and the prime-to-\(\ell\) part of \(\ZZ_\ell^\times\). The space of homomorphisms is therefore one-dimensional over \(k\), by the logarithm on principal units.~\cite{V-III-V-kw24}

Set \(s=1+v_\ell(r)\) and take \(z=\ell^{-(s+1)}+\OO\in J/\OO\). The coboundary \(g\mapsto (\vartheta(g)-1)z\) has values in \(D[\ell]\). All values \(\vartheta(g)-1\) belong to \(\ell^s\ZZ_\ell\), and some have valuation exactly \(s\), by taking \(\kappa_\ell(g)=1+\ell\). Indeed multiplication by \(-r\) on logarithms takes \(\ell\ZZ_\ell\) onto \(\ell^s\ZZ_\ell\). This coboundary represents a nonzero generator of the finite-coefficient homomorphism space. Every \(k\)-multiple is also a coboundary with values in \(D\), by lifting the scalar to \(\OO\). The coefficient sequence for multiplication by \(\ell\) now shows the asserted vanishing of \(H^1(G_{\QQ,S_\ell},D)[\ell]\). Every continuous cocycle with values in the discrete \(\ell\)-primary group \(D\) has finite image, hence is killed by some power of \(\ell\). A nonzero such cohomology group would have an element of order \(\ell\). This proves the first case.

Now assume \(\bar\vartheta\ne1\) globally. There is \(\tau\in G_\QQ\) such that \(\vartheta(\tau)-1\) is a unit. In particular \(D^{G_\QQ}=0\). Every class of \(S_\alpha[\ell^m]\) has a unique finite-coefficient cohomology class above it by the coefficient exact sequence. Restriction of that class to \(K_m\) is injective: on the finite abelian group \(\Gal(K_m/\QQ)\), the cocycle identity with \(\tau\) in the proof of Theorem~\ref{thm:mirror} makes every cocycle with these coefficients a coboundary. After choosing the root of unity, the restricted classes are Kummer classes of character \(\beta_m\), as in that proof.

We check the valuation at every finite place. Away from \(\ell\), on inertia over \(K_m\) one has \(\alpha=1\) and \(\kappa_\ell=1\). A finite representative of a strict class therefore vanishes there: in divisible coefficients it equals a coboundary on the original inertia, and that coboundary restricts to zero. The local Kummer class becomes trivial over the maximal unramified extension. Its valuation is a multiple of \(\ell^m\) because that extension has the same value group. At a place \(w\mid\ell\) the divisible class is locally trivial. For its finite representative there is \(z_w\in D\) with
\[
        b(g)=(\vartheta(g)-1)z_w\qquad(g\in G_{K_{m,w}}).
\]
Here \(b(g)\in D[\ell^m]\), the finite coefficient action over \(K_{m,w}\) is trivial, and \(\alpha(g)=1\). Thus the homomorphism \(b\) vanishes on the kernel of \(\kappa_\ell\). Lemma~\ref{lem:localunit24} shows that each coordinate of its local Kummer class is represented by a unit. This supplies the required valuation at \(w\) even if the finite class there is nonzero.

The ideal-class map in the proof of Theorem~\ref{thm:mirror} therefore applies with \(K_m\). Its kernel consists of unit classes. Its odd part is generated by roots of unity, with character \(\kappa_{\ell,m}\), by the same argument with \(u/\bar u\). On \(\tau\) the difference between \(\beta_m\) and \(\kappa_{\ell,m}\) is a unit because \(\vartheta(\tau)-1\) is a unit. The kernel consequently meets our \(\beta_m\)-classes only in zero. This proves \eqref{eq:all24}.
\end{proof}

\begin{corollary}[The finite test on the entire prime-to-degree range]\leavevmode\label{cor:all24}
On the fields \(E_x\) of Volume~IV, for every odd \(\ell\nmid\operatorname{ord}(\psi)\), every \(1\le r<\lfloor x^{401}\rfloor\) and every character of the complementary sign, the following alternatives give a test. If \(\alpha=\omega^{-r}\) the component vanishes. Otherwise, for any \(m\ge1\) and any positive integer \(b\) coprime to \(f'\ell\), the element
\begin{equation}
 \sum_{\substack{1\le j\le f'\ell^m\\(j,f'\ell)=1}}
   \left\lfloor\frac{bj}{f'\ell^m}\right\rfloor
        (j^{r+1}\alpha(j))^{-1}\pmod{\ell^m}
 \label{eq:test24}
\end{equation}
annihilates \(S_\alpha[\ell^m]\). A nonzero residue of valuation \(v<m\) proves that \(\ell^v S_\alpha=0\). In particular a unit at the first level proves vanishing for every positive twist in the same class modulo \(\ell-1\) with this \(\alpha\).
\end{corollary}
\begin{proof}
The degree of \(E_x\) divides \(4\cdot 2^l\operatorname{ord}(\psi)\). Apply Theorems~\ref{thm:char} and~\ref{thm:all24}. The computation of \((b-\sigma_b)\Theta_n\) and the cyclic-subgroup argument of Theorem~\ref{thm:stick} apply verbatim to the modulus \(f'\ell^m\) and the target in \eqref{eq:all24}. At \(m=1\) the powers \(j^{r+1}\) depend only on \(r\bmod (\ell-1)\); parity also stays fixed in that class. This gives the last assertion.
\end{proof}

For an example with locally trivial residual character take \(E=\QQ(i,\sqrt{13})\), \(\ell=3\), any positive even \(r\), and the quadratic character \(\alpha\) of conductor \(13\). The degree is four and \((3/13)=1\), so \(\bar\vartheta\) is trivial at three. Its restriction to inertia at thirteen is nontrivial. In \eqref{eq:test24} take \(m=1, b=2\), hence \(n=39\). The floor is one on the twelve units between 20 and 38, and zero on the other units. On these twelve the inverse character modulo three has value \(+1\) at \(20,22,25,32\) and value \(-1\) at the other eight, by the squares \(1,3,4,9,10,12\) modulo thirteen. The sum is \(-4\equiv2\pmod3\). The entire component vanishes for every such \(r\).

For the sum of lengths in Volume~I, a unit removes a component outright. Where the first residues vanish, the exponent requires bounds at higher levels, together with bounds for the numbers of cyclic summands. The test covers every twist at primes coprime to the field degree, including primes ramified in the conductor. Section~\ref{sec:norm28} treats the primes dividing the character order.

\subsection{Primes dividing the character order}\label{nav:III-05}
\label{sec:norm28}

We extend the class-group test to the primes dividing the degree of the full field. For this purpose we take the character components over an intermediate field of odd prime-power degree. A determinant over its Galois group will provide an annihilator even though its group order is zero modulo the coefficient prime.

\subsubsection{Descent to the fixed prime-power subfield}\label{nav:III-06}
Let \(E/\QQ\) be finite abelian, \(\ell\) an odd prime, and write \(G=\Gal(E/\QQ)=P\times H\), where \(P\) is its \(\ell\)-Sylow subgroup and \(H\) has order prime to \(\ell\). Let \(K=E^H\), so \(P=\Gal(K/\QQ)\), and write \(d=[K:\QQ]\). The field \(K\) is totally real since complex conjugation has order at most two. Take a finite unramified coefficient ring \(\OO\) over \(\ZZ_\ell\) containing the values of all characters of \(H\), with fraction field \(J\). For \(\alpha\in\widehat H\), let \(\widetilde\alpha\) be its extension to \(G\) that is trivial on \(P\). View it as a Dirichlet character, or equivalently as a character of \(G_\QQ\). Set \(\vartheta=\kappa_\ell^{-r}\widetilde\alpha^{-1}\), \(D=(J/\OO)(\vartheta)\), and use the complementary sign \(\alpha(c)=(-1)^r\). Write \(S_\alpha(K)\) for the classes in \(H^1(K,D)\) trivial on inertia away from \(\ell\) and trivial on all decomposition groups over \(\ell\).

\begin{lemma}\label{lem:des28}
With \(e_\alpha=|H|^{-1}\sum_{h\in H}\alpha(h)^{-1}h\), restriction gives
\[
 S_\alpha(K)\ \simeq\ e_\alpha\big(\Sel_{\ell,r}(E)\otimes_{\ZZ_\ell}\OO\big).
\]
For the fields \(E_x\), the field \(K\) is the subfield of \(K_\psi\) of degree \(\ell^{v_\ell(\operatorname{ord}(\psi))}\); in particular it is fixed as \(x,q,d\) vary.
\end{lemma}
\begin{proof}
Apply inflation--restriction to \(E/K\). All higher cohomology of the finite quotient \(H\) on \(\ell\)-primary modules vanishes by averaging. On \(G_E\) the finite character is trivial, and twisting the quotient action by \(\alpha^{-1}\) selects exactly \(e_\alpha\). At every finite place the corresponding decomposition or inertia quotient has order dividing \(|H|\); its first cohomology vanishes, so the restriction on the relevant local \(H^1\) is injective. This proves the equality of local conditions just as in Theorem~\ref{thm:char}. For \(E_x/K_\psi\) the relative degree is a power of two. Restriction to \(K_\psi\) therefore identifies the \(\ell\)-Sylow subgroups, and the indicated subfield of the cyclic extension \(K_\psi/\QQ\) has precisely that prime-power degree.
\end{proof}

For \(m\ge1\), let \(f_K\) be the conductor of the abelian extension \(K/\QQ\), let \(f_\alpha\) be that of \(\widetilde\alpha\), and put
\[
 \begin{gathered}
 n_m=\operatorname{lcm}(f_K,f_\alpha,\ell^m),\qquad L_m=\QQ(\zeta_{n_m}),\\
 G_m=\Gal(L_m/\QQ),\qquad \Gamma_m=\Gal(L_m/K),\qquad R_m=\OO/\ell^m\OO .
 \end{gathered}
\]
Use \(\widetilde\beta_m=\kappa_{\ell,m}^{r+1}\widetilde\alpha\bmod\ell^m\) on \(G_m\), and write \(\beta_m\) for its restriction to \(\Gamma_m\). For an element \(\sigma\in G_m\) write \([\sigma|_K]\) for the basis vector in the group algebra \(R_m[P]\).

\begin{theorem}[A determinant annihilator]\leavevmode\label{thm:norm28}
Suppose \(\bar\vartheta\) is nontrivial on \(G_K\). For every positive integer \(b\) prime to \(n_m\), form the following element of the group algebra and determinant in \(R_m\):
\begin{align}
 U_{b,m}&=\sum_{\substack{1\le j\le n_m\\(j,n_m)=1}}
   \left\lfloor\frac{bj}{n_m}\right\rfloor
   (j^{r+1}\widetilde\alpha(j))^{-1}[\sigma_j^{-1}|_K],\notag\\
 B_{b,m}&=\det\big(u\mapsto U_{b,m}u\mid R_m[P]\big),
 \qquad \sigma_j(\zeta_{n_m})=\zeta_{n_m}^{\,j}.             \label{eq:norm28}
\end{align}
Then \(B_{b,m}\) annihilates \(S_\alpha(K)[\ell^m]\). If it has valuation \(v<m\) it follows that \(\ell^v S_\alpha(K)=0\). For \(m=1\) the determinant is the \(d\)-th power of the sum of the coefficients of \(U_{b,1}\). Thus one scalar sum modulo \(\ell\) already suffices for the vanishing test, including when \(\ell\mid d\).
\end{theorem}
\begin{proof}
We first prove the finite-level injection
\begin{equation}
 S_\alpha(K)[\ell^m]\ \lhook\joinrel\longrightarrow\
 \big(\Cl(\OO_{L_m})[\ell^m]\otimes_{\ZZ_\ell}\OO\big)[\beta_m],
                                                        \label{eq:inj28}
\end{equation}
where the condition on the right prescribes the action of \(\Gamma_m\) only. Choose \(\tau\in G_K\) with \(\vartheta(\tau)-1\) a unit. This forces \(D^{G_K}=0\), so the finite lift of a given class killed by \(\ell^m\) is unique in cohomology. Its restriction to \(G_{L_m}\) is injective. Indeed the kernel comes from a cocycle \(z\) on the abelian group \(\Gamma_m\); commuting \(\tau\) with \(\sigma\) gives \((\vartheta(\tau)-1)z(\sigma)=(\vartheta(\sigma)-1)z(\tau)\) there, with coefficients in \(R_m\). Hence this cocycle is a coboundary.

Choose a primitive \(\ell^m\)-th root to identify the restricted coefficients with roots of unity. Under Kummer theory the conjugation action on the resulting classes is \(\kappa_{\ell,m}\vartheta^{-1}=\beta_m\) on \(\Gamma_m\). At primes away from \(\ell\) the restricted homomorphism vanishes on inertia. Its Kummer class has valuation divisible by \(\ell^m\), since passing to the maximal unramified extension preserves the value group. At each prime above \(\ell\) the original class is zero in divisible coefficients. Its restricted finite homomorphism is therefore of the form \(g\mapsto(\kappa_\ell(g)^{-r}-1)z\) on \(G_{L_{m,w}}\) for some \(z\in D\), and factors through the cyclotomic image. Lemma~\ref{lem:localunit24} gives unit representatives in that completion in every coefficient coordinate. Consequently division of the principal ideal by the exponent, \((a)=\mathfrak a^{\ell^m}\mapsto [\mathfrak a]\), gives a class-group map. Its kernel consists of unit classes. Complex conjugation belongs to \(\Gamma_m\) since \(K\) is totally real, and \(\beta_m(c)=-1\). As in Theorem~\ref{thm:mirror}, the odd part of the unit quotient is generated by roots of unity, on which \(\tau\) acts by \(\kappa_{\ell,m}(\tau)\). Its difference from \(\beta_m(\tau)\) is a unit. The kernel meets this eigenspace in zero, proving \eqref{eq:inj28}.

Let \(T=(b-\sigma_b)\Theta_{n_m}\) be the integral Stickelberger element of Theorem~\ref{thm:stick}. It annihilates the class group of the full cyclotomic field \(L_m\). Put \(A=R_m[\Gamma_m]\). The group algebra \(R_m[G_m]\) is a free \(A\)-module on any \(d\) coset representatives. Write \(N(T)\in A\) for the determinant of multiplication by the reduction of \(T\) on that module. Its characteristic polynomial, applied to the element \(1\) by Cayley--Hamilton, shows that \(N(T)\) is a multiple of \(T\) in \(R_m[G_m]\). In particular it too annihilates the class group with these coefficients. On the eigenspace in \eqref{eq:inj28}, its action is the scalar obtained by evaluating \(\Gamma_m\) at \(\beta_m\).

After this evaluation there is an explicit isomorphism of \(R_m\)-algebras
\[
 R_m[G_m]\otimes_{A,\beta_m}R_m\ \longrightarrow\ R_m[P],
 \qquad \sigma\otimes1\longmapsto
                    \widetilde\beta_m(\sigma)[\sigma|_K].
\]
It respects multiplication and sends the coset basis to a basis, each vector multiplied by a unit. The image of \(T\) is exactly \(U_{b,m}\) by the floor formula for its coefficients. Formation of the determinant commutes with this base change. The resulting scalar is \(B_{b,m}\), which proves the annihilation. An element of order greater than \(\ell^v\) would give an element of order \(\ell^{v+1}\) already in \(S_\alpha(K)[\ell^m]\). This proves the exponent assertion.

Over the residue field \(k=\OO/\ell\OO\), the augmentation ideal of \(k[P]\) is nilpotent. To see this, generators \(g\) of the finite abelian \(\ell\)-group satisfy \((g-1)^{\operatorname{ord}(g)}=0\) in characteristic \(\ell\). These commuting nilpotents generate the ideal, so any product of sufficiently many generators vanishes. Multiplication by \(U_{b,1}\) is multiplication by its coefficient sum plus a nilpotent endomorphism. Its determinant in dimension \(d\) is the \(d\)-th power of that sum.
\end{proof}

\begin{corollary}[The residual characters left out of the determinant]\label{cor:res28}
On the fields \(E_x\), take the common intervals of Corollary \ref{cor:tower27}, with \(\Pi\) containing the odd primes dividing \(o=\operatorname{ord}(\psi)\). At every odd prime and twist \(1\le r<t(x)\) a complementary component satisfies Theorem~\ref{thm:norm28}, except possibly for the unique character \(\widetilde\alpha=\omega_\ell^{-r}\), where \(\omega_\ell\) is the Teichm\"uller cyclotomic character. At primes \(\ell\nmid o\) this exceptional component vanishes. At the other odd primes the sum of its lengths over the full twist range, weighted by \(\log\ell\), is at most \(\exp(x^{2\eps/3})\) on those same intervals.
\end{corollary}
\begin{proof}
If the residual character is trivial on \(G_K\), the global residual character factors through \(P\), an \(\ell\)-group. Its image in \(k^\times\) has order prime to \(\ell\), hence it is globally trivial. Uniqueness of Teichm\"uller lifts of finite order prime to \(\ell\) gives precisely \(\widetilde\alpha=\omega_\ell^{-r}\), if that character factors through \(E_x\) at all. For \(\ell\nmid o\) we have \(K=\QQ\), and Theorem~\ref{thm:all24} gives vanishing. For the finitely many remaining primes take the intervals sufficiently distant that \(q,d\) differ from each of them. The exceptional character is unramified at \(q,d\). The two quadratic factors of \(E_x/K_\psi C_l\) are generated by the inertia at these two primes: the quadratic field of discriminant \(-q\) ramifies at \(q\), and the other factors are unramified there; the same assertion holds at \(d\). Thus \(\widetilde\alpha\) is trivial on \(H_x^0=\Gal(E_x/K_\psi C_l)\). Its component lies in the \(H_x^0\)-invariants. Corollary \ref{cor:tower27}, bounds their entire weighted sum on the common intervals. Extending scalars to the unramified \(\OO\) preserves normalized length, so it does not change that bound.
\end{proof}

For example, take \(E=\QQ(\zeta_7)^+(i)\), \(\ell=3\), any odd positive \(r\), and the odd character modulo four. Here \(K=\QQ(\zeta_7)^+\) has degree three. Its inertia at three is trivial; the mod-three cyclotomic character on inertia is nontrivial, so the determinant test applies. At \(m=1\), \(n_1=84\), \(b=5\), let \(u\) be the class of \(2\) in \((\ZZ/7)^\times/\{\pm1\}\). Sorting the units \(j\bmod84\) according to the image of \(j^{-1}\) modulo seven gives the following coefficients. Each is the sum of \(\lfloor5j/84\rfloor\) times the sign of the character modulo four, since \(r+1\) is even.
\[
\begin{array}{c|l|r}
 [\sigma_j^{-1}|_K]&j&\text{coefficient before reduction}\\ \hline
 1&1,13,29,41,43,55,71,83&-10\\
 u&11,17,25,31,53,59,67,73&2\\
 u^2&5,19,23,37,47,61,65,79&0
\end{array}
\]
The coefficient sum is \(-8\equiv1\pmod3\), so its cube is a unit. The entire complementary component for this character vanishes, even though three divides the field degree.

For components with a nonzero determinant of valuation \(v<m\), the number of cyclic summands is at most the dimension of the eigenspace in \eqref{eq:inj28} at \(m=1\), by injection on \(\ell\)-torsion. Multiplying this dimension by \(v\) bounds the length. The argument has removed the exclusion of whole primes from the finite test. The remaining numerical problem is to bound the sum of these lengths also where the first residue vanishes and higher determinants give poor bounds.

\subsection{Finite coefficients and Kummer classes on the first cyclotomic layer}\label{nav:III-07}
\label{sec:period30}

We give a test on one finite extension for each odd prime. It applies to the entire field occurring in the count, including at primes dividing its degree. The first step explains why a local condition imposed with divisible coefficients can be compared across twists using finite coefficients.

\subsubsection{A description independent of the integral twist}\label{nav:III-08}
For a number field \(F_0\), an odd prime \(\ell\) and a nonzero integer \(s\), let \(A_s=(\QQ_\ell/\ZZ_\ell)(-s)\), \(D_{s,m}=A_s[\ell^m]\) and \(G=G_{F_0,S_\ell}\). Define \(\Sel_{\ell,s}(F_0)\) using the strict condition at every \(w\mid\ell\) as in Volume~II. Write \(\Gamma=\kappa_\ell(G)\). For the absolute local Galois group at \(w\mid\ell\), write \(H_w=\ker(\kappa_\ell|_{G_{F_{0,w}}})\). Every cohomology group below uses continuous cochains on a profinite group with a discrete coefficient module.

\begin{lemma}[Finite-level comparison]\leavevmode\label{lem:period30}
For \(m\ge1\) there is a natural isomorphism
\begin{equation}
 \Sel_{\ell,s}(F_0)[\ell^m]\simeq
 \frac{\ker\!\left(H^1(G,D_{s,m})\longrightarrow
             \prod_{w\mid\ell} H^1(H_w,D_{s,m})\right)}
      {\operatorname{inf} H^1(\Gamma,D_{s,m})}.              \label{eq:period30}
\end{equation}
In compatible cyclotomic coordinates, two nonzero twists congruent modulo \((\ell-1)\ell^{m-1}\) give the same quotient. The comparison respects the action of field automorphisms when their actions on the finite coefficients agree.
\end{lemma}
\begin{proof}
For either the global or a local cyclotomic image \(J\), its intersection with \(1+\ell\ZZ_\ell\) is an open procyclic group. This follows from finiteness of the number field or local field over \(\QQ\) or \(\QQ_\ell\), respectively, and the cyclotomic tower; globally the tower is unramified away from \(\ell\). If \(u\) is a topological generator of this intersection, \(u^{-s}-1\ne0\). Multiplication by this number is onto on \(\QQ_\ell/\ZZ_\ell\). Thus any continuous cocycle on this subgroup can be matched at the generator by a coboundary. The cocycle law and density of the integral powers give equality everywhere. The remaining quotient has order dividing \(\ell-1\), and its positive-degree cohomology vanishes by averaging. Inflation--restriction gives \(H^1(J,A_s)=0\).

Multiplication by \(\ell^m\) on \(A_s\) is surjective, so every global class killed by \(\ell^m\) lifts to \(H^1(G,D_{s,m})\). On the kernel of the cyclotomic character both modules have trivial action. The map in degree one induced by \(D_{s,m}\subset A_s\) on such a kernel is injective because it is inclusion on homomorphisms. At \(w\mid\ell\) the image in \(H^1(F_{0,w},A_s)\) is zero precisely when the finite lift restricts to zero on \(H_w\). Indeed vanishing there makes its divisible-coefficient class inflated from the local cyclotomic image, whose degree-one cohomology just vanished. The same reasoning globally identifies the kernel of passage from finite to divisible coefficients with inflation from \(H^1(\Gamma,D_{s,m})\): a class mapping to zero kills the global cyclotomic kernel, and every class inflated from \(\Gamma\) maps to zero. Such an inflated class belongs to the numerator. This proves \eqref{eq:period30}. Its data involve \(s\) only through the finite character \(\kappa_\ell^{-s}\bmod\ell^m\). The exponent of \((\ZZ/\ell^m)^\times\) divides \((\ell-1)\ell^{m-1}\). The assertions about congruence and automorphisms follow on cochains.
\end{proof}

\subsubsection{The ideal class and the local generators}\label{nav:III-09}
Let \(L\) be a CM field, let \(c\) be complex conjugation and let a minus superscript mean the image of \((1-c)/2\) on an \(\ell\)-primary group. The group \(\Cl(L)\{\ell\}\) denotes the \(\ell\)-primary part of the ordinary ideal class group. Write \(T_\ell(L)=\Sel_{\ell,-1}(L)^-\); here the coefficient is \(\mu_{\ell^\infty}\).

\begin{proposition}[An injection into the ordinary class group]\label{prop:kummer30}
There is a natural injection \(T_\ell(L)\hookrightarrow\Cl(L)\{\ell\}^-\), compatible with automorphisms commuting with \(c\). At order dividing \(\ell^m\) its image consists of the minus parts of the classes of fractional ideals \(\mathfrak a\) such that \(\mathfrak a^{\ell^m}=(z)\) for a generator satisfying
\begin{equation}
 z\in \mu_{\ell^\infty}(L_w)\,(L_w^\times)^{\ell^m}
                     \qquad(w\mid\ell).                         \label{eq:local30}
\end{equation}
\end{proposition}
\begin{proof}
By the Kummer sequence and Hilbert's theorem 90 a finite lift over \(L\) is represented by \(z\in L^\times/(L^\times)^{\ell^m}\). At a finite place away from \(\ell\) this class is unramified exactly when its valuation is a multiple of \(\ell^m\). Indeed it becomes trivial over the maximal unramified extension precisely when a root exists there. Its value group is still \(\ZZ\), and every unit has an \(\ell^m\)-th root there by taking a root in the algebraic closure of the residue field and applying Hensel's lemma. At \(w\mid\ell\) the kernel of finite-to-divisible local cohomology is the connecting image of \(\mu_{\ell^\infty}(L_w)/\mu_{\ell^\infty}(L_w)^{\ell^m}\). The Kummer cocycle for such a root of unity \(u\) is \(g\mapsto g(v)/v\) with \(v^{\ell^m}=u\), which is exactly that connecting map. Hence strictness at \(w\) is equivalent to \eqref{eq:local30}. This condition implies divisibility of the valuation at \(w\) as well.~\cite{V-III-V-kummer30}

It follows that \((z)=\mathfrak a^{\ell^m}\) for a unique fractional ideal. Changing the representative of \(z\) changes \(\mathfrak a\) by a principal ideal. Changing the finite lift of the same divisible-coefficient class multiplies \(z\) by a global root of unity modulo powers, which leaves its ideal unchanged. These observations give the map to \(\Cl(L)[\ell^m]\) and prove the asserted description of the image after applying the minus projector. If \(\mathfrak a\) is principal, dividing out the power of a generator lets us take \(z\) to be a global unit. For every global unit \(u\), \(u/c(u)\) is a unit of absolute value one at all complex embeddings, hence a root of unity by Kronecker's theorem. Consequently the minus part of the unit classes modulo \(\ell^m\)-th powers is generated by roots of unity, since two is invertible there. Their prime-to-\(\ell\) parts are already powers and their \(\ell\)-parts die in divisible cohomology by the global connecting map. This proves injectivity on the minus part. The inclusion of finite coefficients from level \(m\) to \(m+1\) sends a Kummer representative \(z\) to \(z^\ell\), leaving the ideal class unchanged. Passage to the union proves the full injection. All ideal and cohomology maps used commute with automorphisms.
\end{proof}

\subsubsection{A test for every twist on the full family}\label{nav:III-10}
Let \(E/\QQ\) be abelian and totally imaginary, \(L=E(\mu_\ell)\), \(\Delta=\Gal(L/E)\), and let \(\omega\) denote the Teichm\"uller lift of the mod-\(\ell\) cyclotomic character. For \(r\ge1\), put
\[
 M_{\ell,r}(E)=\Sel_{\ell,r}(E)^{c=(-1)^r},\qquad
 T_{\ell,r}(E)=\{z\in T_\ell(L):\sigma z=\omega(\sigma)^{r+1}z
           \text{ for every }\sigma\in\Delta\}.
\]
Let \(e_{\ell,r}(E)\) be the least nonnegative integer \(e\) such that \(\ell^eT_{\ell,r}(E)=0\). This is finite by Proposition~\ref{prop:kummer30}; the same definition gives \(e=0\) for the zero group. All identifications of finite twists below use a compatible system of bases for the cyclotomic Tate module.

\begin{theorem}[Comparison with one fixed finite group]\label{thm:compare30}
For \(1\le m\le1+v_\ell(r+1)\) there is an isomorphism of finite groups
\begin{equation}
            M_{\ell,r}(E)[\ell^m]\simeq
            T_{\ell,r}(E)[\ell^m].                             \label{eq:compare30}
\end{equation}
In particular, \(M_{\ell,r}(E)=0\) exactly when \(T_{\ell,r}(E)=0\). If \(v_\ell(r+1)\ge e_{\ell,r}(E)\), the entire groups in \eqref{eq:compare30} are isomorphic without truncation. The test uses no restriction on the prime divisors of \([E:\QQ]\).
\end{theorem}
\begin{proof}
Since \([L:E]\mid\ell-1\), restriction identifies the strict group downstairs with the invariants upstairs. To see this globally use inflation--restriction on the absolute groups and average over \(\Delta\). At each place over \(\ell\), restriction on local cohomology is injective since corestriction composed with restriction multiplies by the local degree. Away from \(\ell\) the extension is unramified and the inertia groups agree. These facts give the identification with all the local conditions.

Write any \(\ell\)-adic unit as its Teichm\"uller part times a member of \(1+\ell\ZZ_\ell\). For \(u\in1+\ell\ZZ_\ell\), binomial expansion gives \(u^{r+1}\equiv1\pmod {\ell^{1+v_\ell(r+1)}}\), by raising successively to the \(\ell\)-th power. Thus modulo \(\ell^m\) the ratio of the cyclotomic character to \(\kappa_\ell^{-r}\) is \(\omega^{r+1}\). On \(G_L\) this ratio is one, because \(L\) contains \(\mu_\ell\). In the coordinates just chosen, the numerator and denominator of \eqref{eq:period30} over \(L\) therefore agree for exponents \(-r\) and \(+1\). For an automorphism of \(L/\QQ\) the action on cochains uses conjugation on the argument and the coefficient action on its value. The second action for exponent \(+1\) is \(\omega^{r+1}\) times the one for exponent \(-r\). This factor is independent of the lift of the automorphism since \(\omega|_{G_L}=1\). Consequently the invariant classes under \(\Delta\) become precisely its \(\omega^{r+1}\) eigenspace. The original conjugation eigenvalue \((-1)^r\) becomes \((-1)^r(-1)^{r+1}=-1\). This proves \eqref{eq:compare30}.

Every nonzero \(\ell\)-primary torsion group contains an element of order \(\ell\), so the assertion about vanishing follows by taking \(m=1\). If \(a=v_\ell(r+1)\ge e=e_{\ell,r}(E)\), take \(m=e+1\). The right side is killed by \(\ell^e\), so the left group cannot have an element of order \(\ell^{e+1}\), or of any larger order. This proves the full comparison.
\end{proof}

\begin{corollary}[A simultaneous length bound]\label{cor:family30}
On every sufficiently distant common interval of \hyperref[nav:volume-IV]{Volume~IV}, use its fields \(E=E_x\subset\QQ(\zeta_N)\), and set \(t=\lfloor x^{401}\rfloor\). Let \(\mathcal A_x\) consist of the pairs \((\ell,r)\) with \(\ell\) odd, \(1\le r<t\), and \(v_\ell(r+1)\ge e_{\ell,r}(E)\). For an absolute \(C\),
\begin{equation}
 \sum_{(\ell,r)\in\mathcal A_x}(\log\ell)\len_{\ZZ_\ell}M_{\ell,r}(E)
   \le C[E:\QQ]t^3\log(2Nt)\le \exp(x^{\eps/3}).              \label{eq:family30}
\end{equation}
For the pairs with \(\ell-1\mid r+1\), the same full comparison may also be tested directly in \(T_\ell(E)\) using its exponent.
\end{corollary}
\begin{proof}
Pairs with \(e_{\ell,r}(E)=0\) contribute zero by Theorem~\ref{thm:compare30}. For each contributing pair in the sum one has \(\ell\mid r+1\), hence \(\ell\le t\). Its contribution is at most \(\log|\Cl(E(\mu_\ell))|\). If \(d=[E(\mu_\ell):\QQ]\), the conductor--discriminant formula gives \(D_{E(\mu_\ell)}\le (N\ell)^d\), with \(d\le[E:\QQ]\ell\). Minkowski's bound gives representatives of norm at most \(2^dD_{E(\mu_\ell)}^{1/2}\), and the sum of \((\operatorname{Norm}\mathfrak a)^{-2}\) over integral ideals is at most \(\zeta(2)^d\). Therefore the logarithm of the class number is at most \(Cd\log(2N\ell)\). Sum over fewer than \(t\) twists and at most \(t\) primes, using \(d\le[E:\QQ]t\). On the common intervals \([E:\QQ]\ll_\psi x^{2+8\eps}\) and \(\log(2N)\ll_{\psi,\eps} x^{\eps/4}\). This proves the second inequality for sufficiently large \(x\). If \(\ell-1\mid r+1\) then \(\omega^{r+1}=1\), so the same cochain comparison and Proposition~\ref{prop:kummer30} apply over \(E\) itself.
\end{proof}

\begin{proposition}[A persistent nonzero class on the prime-count fields]\label{prop:example30}
For the primitive quadratic character of discriminant \(-1208\), every field \(E_x\) in the construction has \(M_{3,r}(E_x)[3]\ne0\) for every positive odd \(r\).
\end{proposition}
\begin{proof}
The field cut out by the character is \(F_0=\QQ(\sqrt{-302})\), whose ring of integers is \(\ZZ[\sqrt{-302}]\), since \(302=2\cdot151\) is squarefree. Let \(z=73+9\sqrt{-302}\). Its norm is \(73^2+81\cdot302=29791=31^3\). A common prime divisor of \((z)\) and \((\bar z)\) would lie over 31 and divide \(z+\bar z=146\), which is impossible. As their product is \((31)^3\), it follows that 31 splits and \((z)=\mathfrak p^3\) for a prime of norm 31. This ideal is nonprincipal: \(a^2+302b^2=31\) has no solution in integers. At every place over three, \(z\in1+9\OO_{F_{0,w}}\) is a cube. Indeed for any \(h\) in that valuation ring, the polynomial \(y+3y^2+3y^3-h\) has residue root \(\bar h\) and derivative one modulo the maximal ideal. Its Hensel lift gives \((1+3y)^3=1+9h\).

Thus \(z\) gives a strict Kummer class of order three whose ideal class is nonzero. Conjugation sends the ideal class to its inverse because \(\mathfrak p\bar{\mathfrak p}=(31)\). Its minus projection remains nonzero, so \(T_3(F_0)[3]\ne0\). For odd \(r\), \(r\equiv-1\bmod2\). At \(m=1\) the finite coefficient actions for \(s=r,-1\) agree on the entire absolute Galois group, including conjugation. Lemma~\ref{lem:period30} transfers the nonzero minus class to \(M_{3,r}(F_0)[3]\). Finally \([E_x:F_0]\) is a power of two. Restriction followed by corestriction multiplies the group by this power of two and therefore restriction is injective at three. The local conditions are preserved by restriction, and it commutes with conjugation. This gives the claim on the actual fields.
\end{proof}

The example forces positive lengths in the complementary sign even far along the common intervals. Inequality~\eqref{eq:family30} estimates full lengths on its specified set. At the remaining pairs, \eqref{eq:compare30} determines the first \(1+v_\ell(r+1)\) layers. Uniform control of any higher layers surviving Section~\ref{sec:norm28} is still needed in Volume~I.

\subsection{Large coefficient primes and totally real class groups}\label{nav:III-11}
\label{sec:real79}

On the common intervals of
Section \ref{nav:I-28},
put \(t=\lfloor x^{401}\rfloor\), \(b=\exp(x^\eps)\) and
\(P=b^{1/12}\). Let \(E=E_x\subset\QQ(\zeta_N)\) be the same
abelian CM field as in the theta construction. Enlarge the
lower endpoint so that \(P>2t+3\), \(P>N\), and \(P>[E:\QQ]\).
The last two inequalities follow from
\([E:\QQ]\ll_\psi x^{2+8\eps}\) and
\(\log(2N)\ll_{\psi,\eps}x^{\eps/4}\).
Fix a prime \(\ell>P\) and an unramified coefficient ring
\(\OO\) containing the character values of \(\Gal(E/\QQ)\).
Write \(J\) for its fraction field and \(k\) for its residue field.
For \(1\le r<t\) and \(\alpha(c)=(-1)^r\), set
\[
  \theta=\kappa_\ell^{-r}\alpha^{-1},\qquad
  D=(J/\OO)(\theta),\qquad R_m=\OO/\ell^m\OO,\qquad
  \theta_m=\theta\bmod\ell^m .
\]
Denote the strict group of Theorem~\ref{thm:char} by
\(S_{r,\alpha}\). In this section \(\Cl(F)\{\ell\}\) means
the \(\ell\)-primary part of the ordinary ideal class group.

\begin{theorem}[An exact real class-group formula at every exponent]
\label{thm:real79}
For \(m\ge1\), let \(F_m=(E\,\QQ(\mu_{\ell^m}))^+\) be the
fixed field of complex conjugation and \(G_m=\Gal(F_m/\QQ)\).
There is an isomorphism
\begin{equation}
 S_{r,\alpha}[\ell^m]\ \simeq\
 \operatorname{Hom}_{\OO[G_m]}
 \bigl(\Cl(F_m)\{\ell\}\otimes_{\ZZ_\ell}\OO,\,
                       R_m(\theta_m)\bigr).
 \label{eq:real79}
\end{equation}
The same formula holds with \(F_m\) replaced by the cyclic
totally real subfield cut out by \(\theta_m\) and with its own
Galois group and class group.
\end{theorem}
\begin{proof}
Identify \(D[\ell^m]\) with \(R_m(\theta_m)\) by multiplying
representatives by \(\ell^m\). On inertia at \(\ell\),
\(\bar\theta=\bar\kappa_\ell^{-r}\) is nontrivial: \(E\) is
unramified there, and the cyclotomic image on tame inertia
surjects onto \(\mathbf F_\ell^\times\), with \(0<r<\ell-1\).
Thus \(D\) has zero invariants both globally and on the full
decomposition group at \(\ell\). On inertia at each \(p\ne\ell\)
the character has finite order prime to \(\ell\). Its invariants
are either zero or all of \(D\), which is divisible. The
multiplication-by-\(\ell^m\) sequence now identifies the left
side of \eqref{eq:real79} with finite-coefficient classes
unramified outside \(\ell\) and zero on the decomposition group
at \(\ell\). Indeed the kernel in degree one at any of these
groups is its divisible-coefficient invariants modulo
\(\ell^m\) times those invariants, hence zero.

The finite character \(\theta_m\) factors through \(G_m\),
since it factors through \(E\QQ(\mu_{\ell^m})\) and
\(\theta(c)=1\). Write \(A=R_m(\theta_m)\). Both \(G_m\) and
its decomposition group at \(\ell\) are abelian. Each contains
a cyclic subgroup \(H\) of order prime to \(\ell\), from the
tame part of inertia, on which \(A^H=0\). To see this, take
the prime-to-\(\ell\) component of an inertia element on which
\(\bar\theta\ne1\). Its scalar minus one is a unit. Averaging
over \(H\) gives \(H^i(H,A)=0\) for \(i>0\). The
Hochschild--Serre sequence for this normal subgroup, also
using \(A^H=0\), gives zero cohomology in all degrees on the
two finite groups. In particular global restriction identifies
\(H^1(G_\QQ,A)\) with the \(G_m\)-equivariant continuous
homomorphisms \(G_{F_m}\to A\), by inflation--restriction
in degrees one and two.

For \(p\ne\ell\), the inertia subgroup in \(G_m\) has order
prime to \(\ell\), because the cyclotomic factor is unramified
and \(\ell\nmid[E:\QQ]\). Its first cohomology with values in
\(A\) is zero. Therefore an equivariant homomorphism comes
from an unramified class at \(p\) precisely when it kills
inertia over \(F_m\) there. At \(\ell\), the preceding vanishing
on the finite decomposition group shows that the strict local
condition amounts to killing the full decomposition group over
\(F_m\). Thus class field theory initially gives the ordinary
class group with the primes above \(\ell\) divided out. At
infinite places the difference between narrow and ordinary
class groups is 2-primary and has no effect. The relevant
reciprocity theorem is the Hilbert class field theorem.
\href{https://www.jmilne.org/math/CourseNotes/CFT.pdf}{J.~S. Milne,
\emph{Class Field Theory}.}

The classes of primes above \(\ell\) are killed by every
equivariant homomorphism in \eqref{eq:real79}. Indeed each
prime is fixed by the same tame subgroup \(H\), since \(G_m\)
is abelian. Its image must be \(H\)-invariant in \(A\), and
\(A^H=0\). This removes the quotient by those prime classes.
Since \(A\) is killed by a power of \(\ell\), one may pass to
the \(\ell\)-primary class group and extend scalars by the
tensor--Hom adjunction. The same proof applies to the field
cut out by \(\theta_m\). It is totally real by evenness and
is cyclic because its group is a subgroup of the product of
the cyclic Teichm\"uller group of \(k^\times\) and the cyclic
group of principal rational cyclotomic units modulo
\(\ell^m\). These factors have relatively prime orders.
\end{proof}

\begin{corollary}[All prescribed twists in one class group]
\label{cor:sum79}
Write \(M_{\ell,r}=\tfrac12(1+(-1)^r c)\Sel_{\ell,r}(E)\).
For every \(m\ge1\) one has
\begin{equation}
 \sum_{1\le r<t}\len_{\ZZ_\ell}M_{\ell,r}[\ell^m]
   \le \len_{\ZZ_\ell}\bigl(\Cl(F_m)\{\ell\}/
                                   \ell^m\Cl(F_m)\{\ell\}\bigr).
 \label{eq:allreal79}
\end{equation}
At \(m=1\) the dimension of \(S_{r,\alpha}[\ell]\) equals that of
the \(\bar\theta\)-eigenspace in
\(\Cl(F_1)\otimes_{\ZZ} k\). If that eigenspace vanishes,
then \(S_{r,\alpha}=0\).
\end{corollary}
\begin{proof}
The conductors \(N\) and \(\ell^m\) are coprime. The Chinese
remainder theorem for their cyclotomic fields gives
\(E\cap\QQ(\mu_{\ell^m})=\QQ\). Hence the prime-to-\(\ell\)
subgroup \(\Delta\) of \(G_m\) is
\[
     \bigl(\Gal(E/\QQ)\times\mathbf F_\ell^\times\bigr)/
                                      \langle(c_E,-1)\rangle ,
\]
using Teichm\"uller lifts on the cyclotomic factor. It acts
on \(R_m(\theta_m)\) by the reduction of
\(\delta_{r,\alpha}(g,z)=\alpha(g)^{-1}\widetilde z^{-r}\).
All these characters are distinct. Their restrictions to the
second factor of the product determine \(r\) in
\(1\le r<t<\ell-1\). Restriction to the first then determines
\(\alpha\). Each descends to the displayed quotient by its parity.

For a character \(\delta\) use the idempotent
\(e_\delta=|\Delta|^{-1}\sum_{\sigma\in\Delta}
\delta(\sigma)^{-1}\sigma\) over \(\OO\). The homomorphisms
in \eqref{eq:real79} factor through \(e_{\delta_{r,\alpha}}\)
of the class module modulo \(\ell^m\). The length of the module
of all \(\OO\)-linear homomorphisms from a finite module of
exponent at most \(\ell^m\) to \(R_m\) equals its own length,
as one sees on each cyclic summand. Equivariance can only
decrease this length. The idempotents are pairwise orthogonal,
so the sum is bounded by the length of the full class module
modulo \(\ell^m\). Extension to the unramified ring \(\OO\)
preserves normalized length. Applying Theorem~\ref{thm:char}
and summing over the complementary \(\alpha\)'s for each \(r\)
gives \eqref{eq:allreal79}. For \(m=1\), \(G_1=\Delta\);
there is no further equivariance condition on that eigenspace,
so equality holds in the single-character dimension formula.
A nonzero \(\ell\)-primary torsion group contains an element
of order \(\ell\), which proves the vanishing assertion.
\end{proof}

\begin{theorem}[A uniform sum for finite layers above \(b^{1/12}\)]
\label{thm:cut79}
Put \(U_x=x^{-12}\exp(x^\eps/2)\). For \(\ell>P\) let
\[
  m_x(\ell)=\max\bigl(\{0\}\cup
                     \{m\ge1:\ell^{m+1}\le U_x\}\bigr).
\]
Then on every sufficiently distant common interval
\begin{equation}
 \sum_{\ell>P}(\log\ell)\sum_{1\le r<t}
       \len_{\ZZ_\ell} M_{\ell,r}[\ell^{m_x(\ell)}]
    \ \ll_{\psi,\eps}\ x^{-8}\exp(x^\eps/2).
 \label{eq:cut79}
\end{equation}
Here exponent \(m_x(\ell)=0\) designates the zero subgroup.
This includes the entire first layer at
\(b^{1/12}<\ell\le x^{-6}b^{1/4}\).
\end{theorem}
\begin{proof}
Since \(\ell>b^{1/12}\), one has \(\ell^6>b^{1/2}>U_x\), so
only \(m=1,2,3,4\) can contribute. For each such \(m\) the
field \(F_m\) has degree \(d_m\le[E:\QQ]\ell^m\) and conductor
dividing \(N\ell^m\). The conductor--discriminant formula gives
\(|\operatorname{disc}F_m|\le(N\ell^m)^{d_m}\).
Minkowski's theorem gives an integral representative of norm
at most \(2^{d_m}|\operatorname{disc}F_m|^{1/2}\) in each
ideal class. The number of integral ideals of norm at most \(Y\)
is at most \(Y^2\zeta_{F_m}(2)\le Y^2\zeta(2)^{d_m}\),
the last inequality following prime by prime. Consequently
\[
       \log|\Cl(F_m)|\ll [E:\QQ]\ell^m\log(2N\ell^m)
                           \ll_{\psi,\eps} x^4\ell^m
\]
when \(\ell^{m+1}\le U_x\). We used \([E:\QQ]\ll_\psi x^3\),
\(\log(2N)\ll_{\psi,\eps}x^{\eps/4}\), and
\(\log U_x\le x^\eps/2\).

For primes having \(m_x(\ell)=m\), the term on the left of
\eqref{eq:cut79} is at most \(\log|\Cl(F_m)|\), by
Corollary~\ref{cor:sum79}. For fixed \(m\), the sum of
\(\ell^m\) over these primes is at most
\(U_x\): there are at most \(U_x^{1/(m+1)}\) possible positive
integers and each contributes at most \(U_x^{m/(m+1)}\).
Adding the four bounds gives \(O_{\psi,\eps}(x^4U_x)\), as
claimed. No class-number estimate uniform in an unbounded
cyclotomic exponent is used.
\end{proof}

The theorem bounds actual subgroups of the strict cohomology
of \(E_x\) on the interval used in the theta comparison. The
quotients beyond \(m_x(\ell)\), including every exponent at
larger primes, are still part of the required uniform estimate.

\subsection{A finite class group test at the coefficient prime}\label{nav:III-14}
\label{sec:ray101}

Let \(E/\QQ\) be one of the abelian fields of Volume~I, with
conductor dividing \(N\) and degree \(d\). Fix an odd prime
\(\ell\nmid Nd\). Let \(O\) be the integers of an unramified
extension \(J/\QQ_\ell\) containing the character values of
\(\Gal(E/\QQ)\), and let \(k=O/\ell O\). For \(r\ge1\) and a
character \(\alpha\) of \(\Gal(E/\QQ)\) with
\(\alpha(c)=(-1)^r\), put
\[
 \vartheta=\kappa_\ell^{-r}\alpha^{-1},\qquad
 M=(J/O)(\vartheta),\qquad \tau=\bar\vartheta,\qquad
 h=1+v_\ell(r).
\]
Here \(\kappa_\ell\) is the cyclotomic character and \(c\) is
complex conjugation. By Theorem~\ref{thm:char} the required
character component on \(E\) is the group \(S(\vartheta)\) of
global classes in \(H^1(\QQ,M)\) unramified at \(p\ne\ell\)
and locally zero at \(\ell\). Use continuous cohomology.

Suppose first that \(\tau\ne1\). Let \(L\) be its fixed field,
meaning the fixed field of its kernel, and write \(\Delta=\Gal(L/\QQ)\).
This is a totally real cyclic field in \(E(\mu_\ell)\) because
\(\tau(c)=1\) and \(k^\times\) is cyclic. Its degree is prime to
\(\ell\). Write \(\widetilde\tau\) for the Teichmüller lift, the
unique character of order prime to \(\ell\) reducing to \(\tau\).
For \(1\le b\le h\) the representations of \(G_\QQ\) on
\(A_b=M[\ell^b]\) and on \((O/\ell^b)(\widetilde\tau)\) are
isomorphic. Indeed \(1+\ell\ZZ_\ell\) raised to the \(r\)-th
power lies in \(1+\ell^h\ZZ_\ell\), and the finite values of
\(\alpha\) are their own Teichmüller lifts.

\begin{lemma}[The local condition before changing coefficients]
\label{lem:line101}
Every class in \(S(\vartheta)[\ell^b]\) has a unique global
cohomology class lifting it in \(H^1(\QQ,A_b)\). Its unramified
conditions at \(p\ne\ell\) are the same at both coefficients.
At \(\ell\), if \(\tau|_{G_{\QQ_\ell}}\ne1\), the finite lift is
locally zero. If \(\tau|_{G_{\QQ_\ell}}=1\), its permitted
restriction, under local reciprocity, is any \(O/\ell^b\)-multiple
of the homomorphism
\begin{equation}
 \lambda_b:\QQ_\ell^\times\longrightarrow\ZZ/\ell^b,\qquad
 \lambda_b(\ell^j u)=\ell^{-1}\log\langle u\rangle\bmod\ell^b
       \quad(u\in\ZZ_\ell^\times).                             \label{eq:line101}
\end{equation}
Here \(\langle u\rangle\) is \(u\) divided by its
Teichmüller factor. These conditions on a finite class are also
sufficient.
\end{lemma}
\begin{proof}
Multiplication by \(\ell^b\) on \(M\) is surjective. In cohomology
its exact sequence gives
\[
 0\longrightarrow M^G/\ell^bM^G\longrightarrow H^1(G,A_b)
 \longrightarrow H^1(G,M)[\ell^b]\longrightarrow0.
\]
Globally the term on the left vanishes because some
\(\vartheta(g)-1\) is a unit. At \(p\ne\ell\) apply the same row
to inertia. Its image on \(M\) is finite of order prime to
\(\ell\). The invariants are either all of the divisible module
\(M\) or zero, so again the left term vanishes. The same unit
argument gives the nonsplit assertion at \(\ell\).

In the remaining case \(\alpha\) is unramified at \(\ell\).
Trivial residual action on inertia forces \(\ell-1\mid r\).
The residual unramified value of \(\alpha\) is then one, so its
local character is trivial. Thus the local action is
\(\kappa_\ell^{-r}\) and \(M^{G_{\QQ_\ell}}=\ell^{-h}O/O\).
Indeed the local cyclotomic image is \(\ZZ_\ell^\times\), and
the minimum valuation of its nonzero \(r\)-th power differences
from one is \(1+v_\ell(r)\). This follows by taking the logarithm
on \(1+\ell\ZZ_\ell\). The boundary of \(\ell^{-h}a\), using the
lift \(\ell^{-h-b}a\), is the homomorphism
\[
    g\longmapsto a\,(\kappa_\ell(g)^{-r}-1)/\ell^h
                           \pmod{\ell^b}.
\]
In the exponential of a number in \(\ell^h\ZZ_\ell\), the terms
of degree at least two lie in \(\ell^{2h}\ZZ_\ell\) for odd
\(\ell\). This follows from \(v_\ell(n!)\le(n-1)/2\) for \(n\ge3\)
and from the quadratic term separately. Since \(b\le h\), the
boundary is a unit multiple of
\(a\,\log\langle\kappa_\ell(g)\rangle/\ell\) modulo \(\ell^b\).
On \(\QQ_\ell^\times\), reciprocity for the cyclotomic extension
kills the uniformizer and gives the unit or its inverse according
to convention. Either convention gives the same line. The
logarithm identifies \(1+\ell\ZZ_\ell\) with \(\ell\ZZ_\ell\),
so \(\lambda_b\) is onto. Exactness proves both directions.
\end{proof}

We now define a finite abelian \(\ell\)-group \(B_b(L)\).
If \(\ell\) does not split completely in \(L\), set
\(B_b(L)=\Cl(L)\{\ell\}\). If it splits completely, identify
\(L_w=\QQ_\ell\) canonically at each \(w\mid\ell\) and let
\(B_b(L)\) be the \(\ell\)-primary part of the following idele
class quotient. Divide by \(L^\times\), the full factors at
infinity, the local units at finite places away from \(\ell\),
and \(\ker\lambda_b\) at each \(w\mid\ell\). For \(S=\{w:w\mid\ell\}\)
this gives the exact sequence
\begin{equation}
 \mathcal O_{L,S}^\times\longrightarrow(\ZZ/\ell^b)^S
 \longrightarrow B_b(L)\longrightarrow\Cl(\mathcal O_{L,S})\{\ell\}
 \longrightarrow0,                                             \label{eq:ray101}
\end{equation}
where the first map takes the \(\lambda_b\)-values of the local
components. To see this, first form the quotient with all
\(L_w^\times\) removed at \(w\in S\). Fractional ideals away
from \(S\) identify the result with \(\Cl(\mathcal O_{L,S})\).
Retaining the local quotients at \(S\) gives the indicated kernel,
with precisely the relations of the global \(S\)-units. This
proves finiteness and gives the sequence upon taking primary parts
in the two finite groups.

\begin{theorem}[The full strict group at the first \(h\) levels]
\label{thm:ray101}
Let
\[
 C_b=e_{\widetilde\tau}\bigl(B_b(L)\otimes_{\ZZ_\ell}O\bigr),
 \qquad e_{\widetilde\tau}
 =|\Delta|^{-1}\sum_{\gamma\in\Delta}
                   \widetilde\tau(\gamma)^{-1}\gamma .
\]
For \(1\le b\le h\) there is an isomorphism of \(O\)-modules
\begin{equation}
 S(\vartheta)[\ell^b]\ \simeq\ \operatorname{Hom}_O(C_b,O/\ell^b).
                                                               \label{eq:class101}
\end{equation}
In particular \(S(\vartheta)=0\) exactly when \(C_1=0\).
In the split case this means that the \(\widetilde\tau\)-part
of the \(S\)-class group at \(\ell\) vanishes and that the
projection of some \(S\)-unit to the \(\tau\)-line of \(k^S\)
under \eqref{eq:ray101} is nonzero.
\end{theorem}
\begin{proof}
Restriction to \(G_L\) identifies \(H^1(\QQ,A_b)\) with
\(\Delta\)-equivariant continuous homomorphisms from \(G_L\)
to \(A_b\). Indeed \(G_L\) acts trivially on the
coefficient module, and the outer terms of inflation--restriction
vanish because \(|\Delta|\) is a unit on \(A_b\). The same
averaging argument applies to each quotient of local inertia or
decomposition groups. In particular, away from \(\ell\) the
homomorphism must vanish on inertia. In the nonsplit case it
also vanishes on the decomposition groups at the primes above
\(\ell\). Global class field theory then gives homomorphisms out
of the ordinary ideal class group modulo the classes of those
primes, as in Section~\ref{sec:real79}. Each such prime is fixed
by a decomposition subgroup containing an element on which
\(\tau\ne1\), since \(\tau\) is faithful on \(\Delta\).
Equivariance forces its class to have zero image already.

In the split case each local map at \(w\in S\) is allowed precisely
the line \eqref{eq:line101}. Equivalently it vanishes on
\(\ker\lambda_b\), since this local quotient is cyclic of order
\(\ell^b\). Global reciprocity identifies all these homomorphisms
with those out of the stated idele class quotient. This uses the
usual local and global reciprocity maps with their compatibility
under localization.~\cite{V-III-V-cft101}
In both cases the targets have \(\ell\)-power exponent. Extending
scalars to \(O\) and applying the character idempotent gives
\eqref{eq:class101}. A nonzero \(\ell\)-primary torsion group
contains an element killed by \(\ell\). Nakayama's lemma for the
finite \(O\)-module \(C_1\) proves the vanishing criterion.
In the split case \(\Delta\) permutes \(S\) regularly, so the
\(\tau\)-part of \(k^S\) is one-dimensional. Apply the exact
idempotent functor to \eqref{eq:ray101} to obtain the last
assertion.
\end{proof}

\begin{corollary}[When a finite level already gives every order]
\label{cor:stopped101}
The full \(S(\vartheta)\) is killed by \(\ell^{h-1}\) if and only
if \(C_h/\ell^h C_h\) is killed by \(\ell^{h-1}\). Whenever the
test holds, \eqref{eq:class101} at \(b=h\) computes its entire
order and all its elementary divisors. For a fixed \(\ell,\alpha\),
the support and the dimension of \(S(\vartheta)[\ell]\) are
constant as \(r\) varies in any one class modulo \(\ell-1\)
with \(\tau\ne1\).
\end{corollary}
\begin{proof}
For a finite module over \(O\), duality into \(O/\ell^h\)
preserves the elementary divisors of its quotient modulo
\(\ell^h\). If the full strict group has an element of order
at least \(\ell^h\), one of its multiples has order exactly
\(\ell^h\), and hence occurs in \(S(\vartheta)[\ell^h]\).
This proves the equivalence and the assertion about orders.
At \(b=1\) the field, the character idempotent and the local
line depend only on \(\tau\). Reduction of \(\kappa_\ell\) has
order \(\ell-1\), proving the assertion about varying \(r\).
\end{proof}

\subsection{Vanishing for cubic residual characters}\label{nav:III-15}
\label{sec:cubic105}

In Section~\ref{sec:ray101}, the ordinary class group of a
totally real cyclic field determines whether a complementary
strict component vanishes when the coefficient prime is nonsplit.
For a cubic field we can bound its class number in terms of its
conductor strongly enough to settle the prime-conductor case at
every twist.

\begin{lemma}[Class numbers of cyclic cubic fields]\label{lem:cubic105}
Let \(L/\QQ\) be cyclic of degree three and of conductor \(f\).
Its ordinary class number satisfies
\begin{equation}
 h(L)\ \le\ B(f),\qquad
 B(f)=\frac{f}{4}\left(\frac{\log f}{\log(f/2)}\right)^2<f .
                                                        \label{eq:cubic105}
\end{equation}
\end{lemma}
\begin{proof}
The two nontrivial characters of \(L\) are a primitive Dirichlet
character \(\chi\) and its conjugate, each of conductor \(f\).
They are even because they have order three. The field is totally
real, and the conductor--discriminant formula gives \(D_L=f^2\).
Write \(R\) for the regulator with one row deleted from the
logarithmic unit matrix. The residue of
\(\zeta_L(s)=\zeta(s)L(s,\chi)L(s,\bar\chi)\) at one and the
analytic class number formula give
\[
       |L(1,\chi)|^2=4h(L)R/f .
\]
The factor four is \(2^3/2\), since \(L\) has three real places
and its only roots of unity are \(1,-1\).

Take a unit \(z\ne\pm1\). It generates \(L\), which has prime
degree. Order its three real conjugates \(a,b,c\) by decreasing
absolute value. If \(s=b/a\) and \(t=c/b\), then \(s,t\in[-1,1]\)
and
\[
 |(a-b)(a-c)(b-c)|
    = |a|^2|b|\,|(1-s)(1-t)(1-st)|
    \le 2\,|a|/|c|.
\]
For the inequality use \(|abc|=1\). When \(s,t\ge0\) the three
factors in the second product are at most one. If just one of
\(s,t\) is negative, the product is
\((1-x)(1+y)(1+xy)\) with \(0\le x,y\le1\), and
\((1-x)(1+xy)\le1\). If both are negative the product is
\((1+x)(1+y)(1-xy)\), which is at most \(2(1+xy)(1-xy)\)
because \(x+y\le1+xy\). These estimates give the factor two.
The square of the left hand side is the discriminant of
\(\ZZ[z]\), equal to \(D_L[\OO_L:\ZZ[z]]^2\). Hence
\(\log|a|-\log|c|\ge\log(f/2)\).

The logarithmic unit lattice lies in the plane whose coordinates
sum to zero. By Cauchy--Schwarz each of its nonzero vectors has
length at least \(\log(f/2)/\sqrt2\). Its covolume is
\(\sqrt3 R\), since the cross product of a lattice basis is
parallel to \((1,1,1)\), and each of its coordinates has absolute
value \(R\). If \(\lambda\) is the shortest nonzero length in any
plane lattice, its covolume is at least \(\sqrt3\,\lambda^2/2\).
Indeed a shortest vector is primitive. Complete it to a lattice
basis and subtract an integral multiple of it from the second
vector until the parallel component has length at most
\(\lambda/2\). The second vector still has length at least
\(\lambda\), hence height at least \(\sqrt3\,\lambda/2\).
It follows here that \(4R\ge\log^2(f/2)\).

Ramar\'e's explicit bound gives \(|L(1,\chi)|\le
\tfrac12\log f\) for every nonprincipal primitive even character.
Its hypotheses have been checked above. Substitution proves the
first inequality. A character of order three cannot have conductor
at most four, since each of those unit groups has order at most
two. Thus \(\tfrac12\log f<\log(f/2)\), giving \(B(f)<f\).
~\cite[Cor.~1]{V-III-V-ramare105}
\end{proof}

\begin{theorem}[Two consequences for full primary groups]\label{thm:cubic105}
Let \(\ell\) be an odd prime.
\begin{enumerate}
\item In the setting of Theorem~\ref{thm:ray101}, suppose the
residual character \(\tau\) has order three and \(\ell\) is
nonsplit in its field \(L\). Put \(f=\operatorname{cond}(L)\)
and let \(e\) be the multiplicative order of \(\ell\) modulo three.
If \(\ell^e>B(f)\), the entire strict group \(S(\vartheta)\)
vanishes. In particular \(\ell^e\ge f\) suffices.
\item For every positive even \(r\), if the order of
\(\omega^{-r}\) is at most three, then
\[
                       K_{2r}(\ZZ)\{\ell\}=0 .
\]
Here \(\omega\) is the mod-\(\ell\) cyclotomic character.
The order condition is equivalent to
\(\ell-1\mid 2r\) or \(\ell-1\mid3r\).
\end{enumerate}
\end{theorem}
\begin{proof}
In the first assertion \(\ell\ne3\), since the residue field of
characteristic three has no multiplicative element of order three.
Let \(C=\Cl(\OO_L)\{\ell\}\), and let \(\sigma\) generate
\(\Gal(L/\QQ)\). Extending to \(L\) the ideal norm to \(\QQ\)
acts on \(C\) by \(1+\sigma+\sigma^2\), also at ramified primes.
The class group of \(\QQ\) is trivial, so this operator is zero.
If \(C\ne0\), the nonzero space \(C/\ell C\) is consequently a
module over \(\mathbf F_\ell[X]/(X^2+X+1)\). Every irreducible
factor of this polynomial has degree \(e\), since its roots are
the primitive cube roots of unity and Frobenius sends a root to
its \(\ell\)-th power. In particular \(\ell^e\le|C|\le h(L)\).
The inequality in the hypothesis contradicts
Lemma~\ref{lem:cubic105}. Thus the ordinary class group has
trivial \(\ell\)-part. At a nonsplit prime the first layer of
the strict group is dual to the corresponding ordinary class
group component by Theorem~\ref{thm:ray101}. A nonzero
\(\ell\)-primary torsion group has an element of order \(\ell\),
so vanishing at the first layer proves the assertion at all powers.

For the second assertion, orders one and two are covered by
Proposition \ref{prop:qbound102}.
Suppose the order is three. The unique cubic field \(L\) cut out
by \(\omega^{-r}\) is real of conductor \(\ell\). In the maximal
real subfield \(D=\QQ(\zeta_\ell)^+\), the
\(\omega^{-r}\)-part of \(\Cl(\OO_D)/\ell\) descends to \(L\).
Indeed the subgroup fixing \(L\) has order prime to \(\ell\);
extension followed by norm multiplies by that order, and norm
followed by extension is its group sum on ideal classes.
Dividing by the order gives the projector onto invariants, also
after reduction modulo \(\ell\).
By Lemma~\ref{lem:cubic105}, \(h(L)<\ell\), so this eigenspace
is zero. The odd-prime \(K\)-group comparison of
Proposition \ref{prop:realclass100}
identifies its dimension with that of \(K_{2r}(\ZZ)/\ell\)
for even \(r\). The group is finite, so its \(\ell\)-part
vanishes. Finally the order of \(\omega\) is \(\ell-1\).
Raising \(\omega^{-r}\) to the second or third power gives the
two divisibility conditions stated above.
\end{proof}

On each field \(E_x\) in Section \ref{nav:I-28}, the second part
therefore removes the indicated rational summand for all even
\(1\le r<\lfloor x^{401}\rfloor\) when \(\ell\nmid[E_x:\QQ]\).
The embedding as a direct summand is Corollary \ref{cor:transfer100}.
For other residual character orders, and at a completely split
prime in the ray class calculation, the uniform estimate still
requires additional information.

\subsection{Vanishing on the real quartic field of conductor sixteen}\label{nav:III-17}
\label{sec:quartic115}

Let \(C_2=\QQ(\zeta_{16})^+\), the real quartic layer of the
cyclotomic tower used in the period field. We shall prove that
the strict group on this entire subfield vanishes at every
\(p\equiv3\pmod4\) in one congruence class of twists. These
twists include positive even integers as large as the interval
permits. For the rational character and \(p>3\), the corresponding
cyclotomic class-group eigenspace is also the one in Osburn's
Theorem 1.1.~\cite{V-III-V-osburn115}

Write \(\lambda(u)=(u/p)\) for the Legendre symbol of an odd
prime \(p\equiv3\pmod4\). For a character \(\alpha\) of
\(\operatorname{Gal}(C_2/\QQ)\), use its primitive Dirichlet
representative, of conductor \(f\in\{1,8,16\}\).
The order-four characters are determined by their value \(i\)
or \(-i\) at \(3\bmod16\). The character
\(\beta=\alpha\lambda\) is primitive and odd of conductor \(fp\),
since its two factors are primitive of coprime conductors. Put
\[
 B(\alpha,p)=\frac1{fp}\sum_{a=1}^{fp}a\,\overline{\beta}(a).
\]

\begin{lemma}[A unit at the inert Gaussian prime]\label{lem:bern115}
For the three nontrivial characters of \(C_2\), the number
\(B(\alpha,p)\) lies in \(\ZZ[i]\) and is a unit at the prime
above \(p\). The same holds for the trivial character if \(p>3\).
\end{lemma}
\begin{proof}
For every character here \(B(\alpha,p)=-L(0,\overline\beta)\).
Indeed, expressing the Dirichlet \(L\)-function as a sum of
Hurwitz zeta functions and using \(\zeta(0,z)=1/2-z\) gives
this identity, since the complete character sum is zero. The
Dirichlet functional equation for the primitive odd character
relates \(L(0,\overline\beta)\) to \(L(1,\beta)\) by nonzero
factors, and the latter is nonzero by Dirichlet's nonvanishing
theorem. Thus \(B(\alpha,p)\ne0\).

Suppose first \(f=8\) or \(16\). For \(c\bmod f\), set
\[
 A(c)=\frac1f\sum_{j=0}^{f-1}j\,\overline\alpha(c+j).
\]
Pairing \(j\) and \(f-j\) and using evenness and the vanishing
of the complete sum gives \(A(0)=0\). The terms at \(0,f/2\)
vanish. Reindexing gives \(A(c+1)-A(c)=\overline\alpha(c)\),
so \(A(c)\) is the partial sum of \(\overline\alpha\) through
\(c-1\) for \(0\le c<f\). Every such partial sum has both real
and imaginary parts in \(\{-1,0,1\}\). To check this at
\(f=8\), the successive nonzero values of the character are
\(1,-1,-1,1\). At \(f=16\) they are
\[
       1,\ -i,\ i,\ -1,\ -1,\ i,\ -i,\ 1
\]
or their conjugates, in order at the odd residues \(1,\ldots,15\).

Write \(a=u+pj\), \(1\le u<p,\ 0\le j<f\), dropping multiples
of \(p\), which contribute zero. The terms with coefficient \(u\)
sum to zero at each \(u\), as the \(j\)'s run through all residues
modulo \(f\). Thus
\[
  B(\alpha,p)=\overline\alpha(p)\sum_{u=1}^{p-1}
                  \lambda(u) A(p^{-1}u\bmod f).
\]
Multiplication by \(\overline\alpha(p)\lambda(u)\) preserves the
square \([-1,1]+i[-1,1]\) and the Gaussian integer lattice.
Consequently \(B(\alpha,p)\in\ZZ[i]\) and the absolute value of
each coordinate is at most \(p-1\). As \(p\equiv3\pmod4\),
\(\ZZ[i]/p\ZZ[i]\) is a field. Divisibility by its prime above
\(p\) would make both coordinates multiples of \(p\), forcing
\(B(\alpha,p)=0\).

For \(f=1\) and \(p>3\), Euler's criterion gives
\(\sum_{u=1}^{p-1}u\lambda(u)\equiv\sum_{u=1}^{p-1}
u^{(p+1)/2}=0\bmod p\). The last exponent is strictly between
zero and \(p-1\), so summing over the cyclic multiplicative group
gives zero in \(\mathbf F_p\). Thus \(B(1,p)\) is an integer,
and its absolute value is at most
\(p^{-1}\sum_{u=1}^{p-1}u=(p-1)/2<p\). Its nonvanishing again
proves the assertion.
\end{proof}

\begin{theorem}[Full vanishing through the quartic layer]
\label{thm:quartic115}
Let \(p\equiv3\pmod4\) be prime, and let \(r\ge1\) satisfy
\begin{equation}
       r+1\equiv (p-1)/2\pmod{p-1}.                      \label{eq:leg115}
\end{equation}
Then \(\mathcal S_{p,r}(C_2)=0\).
\end{theorem}
\begin{proof}
The twists in \eqref{eq:leg115} are even. Since
\([C_2:\QQ]=4\), character descent in Theorem~\ref{thm:char}
splits the entire group after extension to an unramified ring
containing \(i\). The four strict characters over \(\QQ\)
are \(\theta=\kappa_p^{-r}\alpha^{-1}\). Their reflections
modulo \(p\) are \(\bar\kappa_p^{\,r+1}(\alpha\bmod p)\).
By Euler's
criterion and \eqref{eq:leg115}, the characteristic-zero lift
of this finite character is \(\beta=\lambda\alpha\).

For clarity, the annihilator of the first layer from
Corollary~\ref{cor:all24} can be evaluated as an ordinary
finite character sum. Set \(n=fp\) and choose a positive
integer \(a_0\) coprime to \(n\). Use the complex conjugate
\(\overline\beta=\beta^{-1}\) on units. Then
\begin{equation}
 \sum_{a=1}^{n}\left\lfloor\frac{a_0a}{n}\right\rfloor
                \overline\beta(a)
       =(a_0-\beta(a_0))B(\alpha,p).                       \label{eq:floor115}
\end{equation}
Indeed, subtract the least nonnegative residue of \(a_0a\)
from \(a_0a\) and divide by \(n\). Multiplication by \(a_0\)
permutes the units, and on replacing \(a_0a\) by a new residue
the character introduces the factor \(\beta(a_0)\). Reducing
the left side of \eqref{eq:floor115} at the prime above \(p\)
gives exactly the annihilator in \eqref{eq:test24} for \(m=1\).
That annihilation follows in Section~\ref{sec:all24} by sending
the finite strict classes to the odd part of the cyclotomic
ideal class group and applying the integral Stickelberger element.

If \(\alpha\ne1\), take \(a_0\equiv1\bmod p\) and choose its
residue modulo \(f\) so that \(\alpha(a_0)\ne1\). These two
conditions are compatible by the Chinese remainder theorem.
The prime-to-\(p\) root \(\alpha(a_0)\) remains nontrivial
modulo \(p\). Globally the residual \(\theta\) is nontrivial
since the nontrivial character at conductor \(f\) cannot be
cancelled by one of conductor \(p\). Both factors on the right
of \eqref{eq:floor115} are units by Lemma~\ref{lem:bern115}.
Thus \(S_\QQ(\theta)[p]=0\).

For \(\alpha=1,\ p>3\), the global residual action is
nontrivial since \(r\equiv(p-3)/2\bmod (p-1)\). Taking \(a_0=2\)
gives \(a_0-\lambda(a_0)\in\{1,3\}\), a unit at \(p\).
Lemma~\ref{lem:bern115} again applies. At \(p=3,\alpha=1\)
use \(O=\ZZ_3\). The residual action is trivial, hence \(F_1=\QQ\).
In \eqref{eq:idele113} the \(3\)-integer class group is trivial
and the single local group has order three. Thus \(|B_1|\le3\).
The global boundary in \eqref{eq:strict113} has order three by
Corollary~\ref{cor:levels113}. Exactness gives \(S_\QQ(\theta)[3]=0\)
here as well.

Every nonzero \(p\)-primary torsion group contains a nonzero
element of order \(p\). All four strict groups therefore vanish
through every power of \(p\). Extension of scalars to the
unramified coefficient ring is faithful, proving the theorem
over \(\ZZ_p\).
\end{proof}

\begin{corollary}[A vanishing summand on every period field]
\label{cor:quarticperiod115}
Let \(E=\widetilde E\) be the period field in
Section \ref{nav:I-28}.
For all sufficiently large \(x\) on its selected intervals,
\(C_2\subset E\). At every prime and twist in
\eqref{eq:leg115} with \(p\equiv3\pmod4\),
\(p\nmid[E:C_2]\) and \(1\le r<\lfloor x^{401}\rfloor\),
one has
\[
       \mathcal S_{p,r}(E)^{\operatorname{Gal}(E/C_2)}=0 .
\]
\end{corollary}
\begin{proof}
The period field contains \(C_{l'}\) with \(l'\ge2\) for large
\(x\). Restriction from \(C_2\) identifies unrestricted
cohomology with the invariants over \(E\), by averaging over
the relative Galois group of order prime to \(p\).
At primes over \(p\), local restriction is injective by
restriction and corestriction. At other finite primes the
coefficients are unramified and an inertia homomorphism that
vanishes upstairs factors through a finite group of order
prime to \(p\), hence vanishes downstairs. This gives the
identification for strict groups and the conclusion follows
from Theorem~\ref{thm:quartic115}.
\end{proof}

The common field, its auxiliary primes and the interval have
all been retained. This removes the indicated invariant summand
also for twists outside the range of the short-subfield bound
in Volume~I. Components on which \(\operatorname{Gal}(E/C_2)\)
acts nontrivially can still contribute to the logarithmic order.

\subsection{Locally trivial twists on the full real cyclotomic tower}\label{nav:III-18}
\label{sec:tower138}

Let \(p\) be an odd prime, \(r>0\) a multiple of \(p-1\), and
\(\theta=\kappa_p^{-r}\). Write \(D=(\QQ_p/\ZZ_p)(-r)\).
For a number field \(K\), let \(S(K)=\mathcal S_{p,r}(K)\)
denote the global classes unramified away from \(p\) whose
restrictions to all completions above \(p\) are zero in
\(H^1(K_w,D)\). Put \(\Sigma=\{w:w\mid p\}\).
We first determine the part of \(S(K)[p]\) detected at these
places when the finite coefficient action is trivial.

\begin{lemma}[Local coefficients of a strict class]\label{lem:slope138}
Suppose \(K/\QQ\) is finite Galois and unramified at \(p\).
There is a natural sequence exact at its first two nonzero terms,
\begin{equation}
  0\longrightarrow
  \operatorname{Hom}\bigl(\Cl(K)/\langle[w]:w\mid p\rangle,
                              \mathbb F_p\bigr)
  \longrightarrow S(K)[p]\longrightarrow
                   \mathbb F_p^\Sigma/\mathbb F_p\mathbf 1 .
  \label{eq:slope138}
\end{equation}
The last map is equivariant for \(\operatorname{Gal}(K/\QQ)\),
acting by permutations on the coordinates on the right.
\end{lemma}
\begin{proof}
At each \(w\mid p\) the cyclotomic character has full image
\(\ZZ_p^\times\), since \(K_w/\QQ_p\) is unramified and the
local extensions generated by \(p\)-power roots of unity are
totally ramified. The global image is therefore full as well.
Put \(b=1+v_p(r)\). The character \(\theta\) has image
\(1+p^b\ZZ_p\) locally and globally. Indeed the torsion units
are killed by \(p-1\mid r\) and the \(p\)-adic logarithm
identifies the \(r\)-th power map on principal units with
multiplication by \(r\) on \(p\ZZ_p\).
Thus
\[
             \lambda(g)=(\theta(g)-1)/p^b\bmod p
\]
is a homomorphism to \(\mathbb F_p\), surjective on every
decomposition group at \(p\). It is invariant under conjugation
by \(G_\QQ\).

By the coefficient sequence for multiplication by \(p\) on
\(D\), a strict class killed by \(p\) has a representative in
\(H^1(G_{K,\Sigma},D[p])\). Identify \(D[p]=p^{-1}\ZZ_p/\ZZ_p\)
with \(\mathbb F_p\) by multiplication by \(p\) on representatives.
The action there is trivial, so this representative is a
homomorphism \(f:G_{K,\Sigma}\to\mathbb F_p\).
On \(G_{K_w}\) it is a coboundary in \(D\). Before identifying
\(D[p]\) with \(\mathbb F_p\), this says
\(p^{-1}f(g)=(\theta(g)-1)y_w\) in \(D\) for some \(y_w\in D\).
Since \(p y_w\in D^{G_{K_w}}=p^{-b}\ZZ_p/\ZZ_p\), we can
write \(y_w=p^{-b-1}z_w\bmod\ZZ_p\). Hence
\(f|_{G_{K_w}}=a_w\lambda\) with unique \(a_w=z_w\bmod p\).
Changing the global representative by a coboundary in \(D\)
changes \(f\) by a common multiple of \(\lambda\), by the same
calculation with \(G_{K,\Sigma}\) in place of \(G_{K_w}\).
Conversely every such multiple is a global coboundary in \(D\).
This gives the displayed map to tuples modulo constants.
Both \(\lambda\) and the residual coefficient identification
are invariant under \(G_\QQ\), proving equivariance.

If all \(a_w\) agree, subtract that common multiple to obtain a
representative trivial on every \(G_{K_w}\). It is unramified
at every other finite place by construction. Global class
field theory identifies precisely these homomorphisms with
those from the ordinary class group killing the prime classes
in \(\Sigma\). Archimedean inertia has order at most two and
has zero image since \(p\) is odd. Conversely each of these
homomorphisms gives a strict class, and a global multiple of
\(\lambda\) trivial on any one \(G_{K_w}\) must be zero. This
proves exactness.
\end{proof}

\begin{theorem}[Descent to a field of degree at most eight]
\label{thm:tower138}
Let \(C_j=\QQ(\zeta_{2^{j+2}})^+\), \(j\ge0\). Suppose that
\(p\not\equiv\pm1\pmod{64}\) and \(p-1\mid r\), and put
\[
        a=\min\{j,v_2(p^2-1)-3\}.
\]
Then \(0\le a\le3\), and restriction is an isomorphism of the
entire strict groups \(S(C_a)\simeq S(C_j)\).
For every \(p\equiv3,5\pmod8\) in this assertion,
\(S(C_j)=0\) at every \(j\).
\end{theorem}
\begin{proof}
By the class-number theorem of Morisawa and Okazaki, \(p\) does
not divide the ordinary class number of \(C_j\) for any \(j\)
whenever \(p\not\equiv\pm1\bmod64\)
\cite[Corollary B]{V-III-V-mo138}. The prime \(p\) is unramified in
\(C_j\), so Lemma~\ref{lem:slope138} embeds \(S(C_j)[p]\) in
the permutation module on the places above \(p\), modulo its
constant line. The field is abelian, so the decomposition group
\(J\) of one place fixes every place above \(p\) and acts
trivially on this module. It therefore acts trivially on
\(S(C_j)[p]\). Its order is a power of two. Averaging over \(J\)
defines an idempotent on every \(p\)-primary group.
If the complementary summand of \(S(C_j)\) were nonzero, some
element in it would have exact order \(p\), by multiplying an
element of finite \(p\)-power order. All classes have finite
\(p\)-power order since a continuous cocycle from a compact
group to \(D\) has finite image. This contradicts triviality
of \(J\) on \(S(C_j)[p]\). Thus the full group is \(J\)-invariant.

For a Galois extension of degree prime to \(p\), restriction
identifies the lower strict group with the invariants of
the upper strict group. Indeed inflation--restriction for
the absolute Galois groups gives an isomorphism on global
\(H^1\) with those invariants, since positive-degree cohomology
of the finite quotient on \(p\)-primary modules vanishes by
averaging. Local restriction at \(p\) is injective by
restriction and corestriction. At a prime away from \(p\) the
coefficients are unramified. An inertia homomorphism vanishing
after restriction factors through a finite quotient of order
prime to \(p\), so vanishes. This proves descent with all the
stated local conditions.

The cyclic group of \(C_j\) is
\((\ZZ/2^{j+2}\ZZ)^\times/\{\pm1\}\), and \(J\) is generated
by \(p\) in this quotient. Choose the sign \(u=\pm p\) with
\(u\equiv1\bmod4\), and put \(e=v_2(u-1)\).
Then \(e=v_2(p^2-1)-1\) because the other of \(p-1,p+1\)
has valuation one. Successive squaring of \(u\) increases
the valuation of its difference from one by exactly one each
time, as \(u^{2^d}+1\equiv2\bmod4\). Its order in the cyclic
quotient is \(2^{\max(0,j+2-e)}\). Hence \(J\) has index
\(2^{\min(j,e-2)}=2^a\), and its fixed field is \(C_a\).
Our condition modulo \(64\) gives \(e\le5\).
This proves the isomorphism. If \(p\equiv3,5\bmod8\), then
\(e=2\), \(a=0\), and \(\Sigma\) consists of a single place
in \(C_j\). The permutation quotient is zero. The injection
above gives \(S(C_j)[p]=0\), so \(S(C_j)=0\).
\end{proof}

\begin{corollary}[The cyclotomic contribution on the period field]
\label{cor:descent138}
Take the actual period field
\(E=K_\psi C_{l'}(\sqrt{-q},\sqrt{-d_0})\) on the intervals of
Section \ref{nav:I-28}. Put \(A=\operatorname{Gal}(E/C_{l'})\),
\(m_E=|A|\), and take \(a\) in Theorem~\ref{thm:tower138} with
\(j=l'\). Under its hypotheses,
\[
      m_E\, S(E)^A\ \subset\ \operatorname{res}_{E/C_a}S(C_a).
\]
If \(p\nmid m_E\), restriction identifies \(S(C_a)\) with
\(S(E)^A\). For \(p\equiv3,5\bmod8\) one has
\(m_E\, S(E)^A=0\) even when \(p\mid m_E\). In those residue
classes and twists the invariant part contributes zero after
multiplication by the depth \(p^{k_{p,r}}\) of Theorem \ref{thm:deep114}, for every sufficiently large \(x\) on the same
intervals.
\end{corollary}
\begin{proof}
Corestriction takes \(S(E)\) to \(S(C_{l'})\). At the primes
over \(p\) this follows from the local restriction--corestriction
formula and the strict zero restrictions upstairs. At each other
finite prime apply the same double-coset formula on inertia.
Every upper restriction there is zero, by unramifiedness of
both the coefficients and the class. Restriction back to \(E\)
is the sum of the \(A\)-translates, hence is multiplication by
\(m_E\) on \(S(E)^A\). Apply Theorem~\ref{thm:tower138} and,
when \(p\nmid m_E\), prime-to-\(p\) descent as in its proof.

For the last claim \(m_E\le 4[K_\psi:\QQ]\) on all the intervals.
Only the primes dividing an integer of this bounded size can
give a nonzero invariant group in these residue classes.
For all such \(p\) the depth \(k_{p,r}\) tends to infinity
uniformly for \(1\le r<\lfloor x^{401}\rfloor\) by Theorem \ref{thm:deep114}. In particular it exceeds \(v_p(m_E)\) for all
sufficiently large \(x\). This kills \(S(E)^A\) at that depth.
\end{proof}

For example all positive even twists at \(p=3\) satisfy the
vanishing on \(C_{l'}\), uniformly up to the terminal twist
of the period construction. The other character components
on \(E\) and coefficient primes outside the given congruences
remain in the full interval estimate.

\section{Special values and integral annihilators}\label{nav:g-III-1}

\subsection{Strict classes and finite even K-groups}\label{nav:III-12}
\label{sec:wild98}

Let \(E\) be a totally imaginary number field of degree \(d\).
For \(r\ge1\) and any prime \(\ell\) write
\[
 M=(\QQ_\ell/\ZZ_\ell)(-r),\qquad
 \Sel_{\ell,r}(E)=
 \ker\!\left(H^1(G_{E,S},M)\longrightarrow
                      \prod_{v\mid\ell}H^1(E_v,M)\right).
\]
Here \(S\) contains the places above \(\ell\) and the archimedean
places, and \(G_{E,S}\) permits ramification only in \(S\).
For this definition take no other finite places in \(S\).
For \(K=E\) or one of the completions \(E_v\) above \(\ell\)
let \(w_{\ell,r}(K)=|H^0(K,M)|\).
Use \(K_j(\OO_E)\) for Quillen's algebraic \(K\)-groups,
\(B\{\ell\}\) for the \(\ell\)-primary subgroup of an abelian
group, and \(B^\vee\) for its Pontryagin dual. The strict group
is the group in Section \ref{nav:I-26}.

\begin{theorem}[The strict kernel in an even \(K\)-group]\label{thm:wild98}
There is a natural exact sequence
\begin{equation}
 \begin{split}
 0\longrightarrow \Sel_{\ell,r}(E)^\vee
 &\longrightarrow K_{2r}(\OO_E)\{\ell\}
 \longrightarrow \bigoplus_{v\mid\ell}H^0(E_v,M)^\vee\\
 &\longrightarrow H^0(E,M)^\vee\longrightarrow 0 .
 \end{split}                                                 \label{eq:wild98}
\end{equation}
In particular each strict group is finite. For fixed \(E,r\)
only finitely many coefficient primes contribute. For all of them,
including two,
\begin{equation}
 |\Sel_{\ell,r}(E)| =
   |K_{2r}(\OO_E)\{\ell\}|\,
       \frac{w_{\ell,r}(E)}{\prod_{v\mid\ell}w_{\ell,r}(E_v)} .
 \label{eq:orders98}
\end{equation}
\end{theorem}
\begin{proof}
For \(j\ge1\) use the finite dual modules
\(A_j=M[\ell^j]\) and \(B_j=\ZZ/\ell^j(r+1)\), paired into
\(\mu_{\ell^j}\). Let \(\operatorname{Sha}^i_S(C)\) in this proof denote the
kernel of restriction in degree \(i\) from \(G_{E,S}\) to the
places in \(S\), with the usual Tate modification at infinity.
The complex places have zero modified groups. Finite-module
Poitou--Tate duality pairs \(\operatorname{Sha}^1_S(A_j)\) perfectly with
\(\operatorname{Sha}^2_S(B_j)\). The last part of its sequence identifies the
cokernel of \(H^2(G_{E,S},B_j)\to\bigoplus_{v\mid\ell}H^2(E_v,B_j)\)
with \(H^0(E,A_j)^\vee\). Local Tate duality identifies these
local degree-two groups with \(H^0(E_v,A_j)^\vee\).
The hypotheses for duality hold since all primes dividing the
orders of the modules are in \(S\).~\cite[I, Cor.~2.3 and Thm.~4.10]{V-III-V-milne98}

The inclusions \(A_j\to A_{j+1}\) correspond under the pairings
to the quotient maps \(B_{j+1}\to B_j\). Continuous cohomology of
discrete modules commutes with this direct limit. The product
of local groups here is finite in the number of factors.
Exactness of the direct limit gives
\(\varinjlim_j\operatorname{Sha}^1_S(A_j)=\Sel_{\ell,r}(E)\).
Take the inverse limit on the dual side. The finite global and
local cohomology groups of the \(B_j\) satisfy the Mittag--Leffler
condition, since their images stabilize at each fixed index.
Inverse limits of the resulting exact systems of finite groups
are exact. In particular the degree-one inverse-limit correction
in continuous degree two vanishes. The global middle term in
\eqref{eq:wild98} at this stage is
\(H^2(G_{E,S},\ZZ_\ell(r+1))\).

We use the number-ring comparison theorem in its even-degree
form. For \(r\ge1\), if \(\ell\) is odd or the field is totally
imaginary, it identifies \(K_{2r}(\OO_E)\{\ell\}\) with
\(H^2_{\rm et}(\OO_E[1/\ell],\ZZ_\ell(r+1))\).
The latter is \(H^2(G_{E,S},\ZZ_\ell(r+1))\) by comparison with
restricted-ramification Galois cohomology for \(\ell\)-primary
coefficients on this ring. In the comparison the residue fields
at \(v\mid\ell\) contribute no \(\ell\)-primary change upon
inverting \(\ell\). Their positive even \(K\)-groups vanish and
their group in degree \(2r-1\) has order \(|\kappa(v)|^r-1\),
prime to \(\ell\). The even global group is finite by the
finite-generation theorem for rings of integers and Borel's rank
formula.~\cite[Ch.~VI, Thms.~8.2 and 8.4]{V-III-V-weibel98}
These results apply here also at two. Substitution gives
\eqref{eq:wild98}. Finiteness of the middle group gives finite
support in \(\ell\), and taking orders gives \eqref{eq:orders98}.
\end{proof}

\begin{lemma}[The local correction]\label{lem:loc98}
Put \(d_v=[E_v:\QQ_\ell]\). Then \(w_{\ell,r}(E_v)\le4rd_v\).
For odd \(\ell>rd_v+1\) it equals one. At an odd prime unramified
in \(E\), the correction
\(\sum_{v\mid\ell}\log w_{\ell,r}(E_v)-\log w_{\ell,r}(E)\)
is exactly
\[
 \begin{cases}
 (g_\ell-1)(1+v_\ell(r))\log\ell,&\ell-1\mid r,\\
 0,&\ell-1\nmid r ,
 \end{cases}
\]
where \(g_\ell\) is the number of places above \(\ell\).
\end{lemma}
\begin{proof}
The cyclotomic image of \(G_{E_v}\) has index \(j_v\) in
\(\ZZ_\ell^\times\) equal to the degree over \(\QQ_\ell\) of
\(E_v\cap\QQ_\ell(\mu_{\ell^\infty})\). This index divides \(d_v\).
The image contains all \(j_v\)-th powers. An element of exact
order \(\ell^a\) in \(M\) is fixed precisely when the \(r\)-th
powers of the image are one modulo \(\ell^a\). Hence the exponent
of \((\ZZ/\ell^a)^\times\) divides \(rj_v\). For odd \(\ell\)
and \(a\ge1\) this exponent is \((\ell-1)\ell^{a-1}\).
It follows that \(\ell\le rd_v+1\) and \(\ell^a\le2rd_v\).
For \(\ell=2,\ a\ge3\) the exponent is \(2^{a-2}\), yielding
\(2^a\le4rd_v\). The smaller exponents obey the same bound.
Since \(M\) has one subgroup of each \(\ell\)-power order,
these bounds apply to the full group of invariants.

If \(v\) is unramified over \(\ell\), the totally ramified
cyclotomic tower intersects \(E_v\) trivially over \(\QQ_\ell\).
Its local image is therefore all of \(\ZZ_\ell^\times\).
The global image contains this image of a decomposition group,
so it too is full. For odd \(\ell\), the largest \(a\) such that
\((\ell-1)\ell^{a-1}\mid r\) is \(1+v_\ell(r)\) if
\(\ell-1\mid r\), and there is no such \(a\ge1\) otherwise.
This gives the asserted difference of logarithms.
\end{proof}

\begin{corollary}[An exact sum and a polynomial error on the common fields]
\label{cor:sum98}
For \(t\ge2\) an integer let \({\cal L}_E(t)\) be the supremum over
finite sets of primes of
\(\sum_{\ell}(\log\ell)\sum_{1\le r<t}\len_{\ZZ_\ell}\Sel_{\ell,r}(E)\)
restricted to the chosen set. Then
\begin{align}
 {\cal L}_E(t)
    &=\sum_{1\le r<t}\log|K_{2r}(\OO_E)|-C_E(t),                \label{eq:ktotal98}\\
 C_E(t)&=\sum_{1\le r<t}\sum_\ell
       \left(\sum_{v\mid\ell}\log w_{\ell,r}(E_v)-
                                      \log w_{\ell,r}(E)\right), \notag\\
 0\le C_E(t)&\ll d^2t^2\log(4dt).                              \notag
\end{align}
For \(E=E_x,\ t=\lfloor x^{401}\rfloor\) from Volume~I, this error
is \(o(\exp(x^\eps/2))\), uniformly as the common intervals recede.
\end{corollary}
\begin{proof}
For a finite \(\ell\)-group \(B\), its length times \(\log\ell\)
is \(\log|B|\). Theorem~\ref{thm:wild98} gives finitely many
nonzero groups for the finitely many \(r\), so the supremum is
attained. Sum its order formula over all primes and twists.
Restriction embeds the global invariants in each local set of
invariants, so each correction is nonnegative. By Lemma~\ref{lem:loc98},
only primes \(\ell\le dr+1\) need be counted, including two. There
are at most \(d\) places at each, and their total logarithm is
at most \(d\log(4dr)\). Summing over \(r<t\) gives the displayed
error bound. The degree \([E_x:\QQ]\) is bounded by a fixed
power of \(x\) on the intervals, whence the last assertion.
\end{proof}

The fixed-scale finiteness assertion was also proved in Section \ref{nav:II-08}. Formula \eqref{eq:ktotal98} expresses the entire strict sum in
finite global groups on \(E_x\) itself. It supplies the correction
to be subtracted when using bounds on their orders. Uniform control
of the sum of those orders' logarithms at the precision of
Volume~I, equation (L), is still needed.

\subsection{Special values and the two conjugation eigenspaces}\label{nav:III-13}
\label{sec:critical99}

Let \(E/\QQ\) be a CM abelian extension of conductor dividing
\(N\), let \(F=E^+\) be its maximal real subfield, and put
\(f=[F:\QQ]\) and \(d=2f\). Suppose that \(E/F\) is unramified
at all finite places. These hypotheses hold for the fields \(E_x\)
of Theorem \ref{thm:norm91}.
Indeed, complex conjugation \(c\) acts nontrivially on the two
imaginary quadratic subfields of coprime discriminants \(-q_0,-d_0\).
At each rational prime inertia is trivial on at least one of them.
Thus \(c\) lies in no finite inertia group.

For \(m\ge2\) let \(B_m^\delta(E)\), with \(\delta=+1\) or \(-1\),
be the \(\delta\)-eigenspace of \(c\) on the subgroup of elements
of odd order in \(K_{2m-2}(\OO_E)\). Write
\(w_m(J)=|H^0(J,(\QQ/\ZZ)(m))|\) for a number field \(J\).
Let \(X\) be the characters of \(\Gal(E/\QQ)\), each represented by
its primitive Dirichlet character. Define
\[
 X_m=\{\chi\in X:\chi(c)=(-1)^m\},
 \qquad P_m=\prod_{\chi\in X_m}L(1-m,\chi).
\]
There are \(f\) characters of each sign. This follows by restriction
to \(\langle c\rangle\) in the character group of a finite abelian
group. Section \ref{nav:II-00} bounded the corresponding sign of the
strict kernels. Here we determine the odd order of the full global
eigenspace before imposing any local condition.

\begin{lemma}[One common denominator for the special values]\label{lem:bern99}
Let \(D_m=mN(m+1)!\). The number \(D_m^fP_m\) is a nonzero
rational integer, and
\[
                    \log |D_m^fP_m|\ll f m\log(2mN).
\]
\end{lemma}
\begin{proof}
Define \(B_j(z)\) by
\(u e^{zu}/(e^u-1)=\sum_{j\ge0}B_j(z)u^j/j!\).
For a primitive character of conductor \(q\mid N\), the exponential
sum at \(y>0\) equals
\(\sum_{a=1}^q\chi(a)e^{-ay}/(1-e^{-qy})\).
Its Mellin integral is \(\Gamma(s)L(s,\chi)\) for \(\real s>1\).
At zero the coefficient of \(y^{m-1}\) is
\(q^{m-1}\sum_a\chi(a)B_m(1-a/q)/m!\).
Subtract the Laurent terms through that degree on \((0,1]\) before
integrating. The integral of the remainder is holomorphic at
\(s=1-m\), as is the exponentially convergent integral on
\([1,\infty)\). Each subtracted monomial \(b y^j\) contributes
\(b/(s+j)\). Divide the resulting residue at \(1-m\) by that of
\(\Gamma(s)\), which is \((-1)^{m-1}/(m-1)!\).
Replacing \(u\) by \(-u\) gives \(B_m(1-z)=(-1)^m B_m(z)\).
Together these calculations give
\[
             L(1-m,\chi)=-\frac{q^{m-1}}{m}
                       \sum_{a=1}^q\chi(a)B_m(a/q).
\]
The formula also applies at \(q=1\) with \(\chi(1)=1\).

Put \(B_j=B_j(0)\). Multiplication of the generating series by
\(e^u-1\) gives \(B_0=1\) and
\((j+1)B_j=-\sum_{i<j}\binom{j+1}{i}B_i\) for \(j\ge1\).
Induction gives \((j+1)!B_j\in\ZZ\). In the expansion
\(B_m(z)=\sum_{j=0}^m\binom mj B_j z^{m-j}\), multiplication by
\(q^{m-1}\) clears the denominator from \(z=a/q\) except possibly
one factor \(q\). Hence \(D_m L(1-m,\chi)\) is an algebraic integer.
The Galois group of the cyclotomic field of character values permutes
the characters in \(X_m\) and preserves their conductors. The product
of the cleared values is therefore a rational integer.

For \(\chi\in X_m\), the gamma factor in the Dirichlet functional
equation is finite and nonzero at \(1-m\), since \(\chi\) has the
parity of \(m\). The Euler product at \(m>1\) is nonzero, so
\(L(1-m,\chi)\ne0\). Finally the Bernoulli generating function is
bounded by an absolute constant on \(|u|=1,\ 0\le z\le1\).
Cauchy's estimate gives \(|B_m(z)|\le C m!\). It follows that
\(|L(1-m,\chi)|\le C(m-1)!N^m\). Multiply over the \(f\)
characters and use \(\log(m!)\le m\log m\).
\end{proof}

\begin{theorem}[The order of the critical eigenspace]\leavevmode\label{thm:critical99}
Set \(J=F\) if \(m\) is even and \(J=E\) if \(m\) is odd.
For every odd prime \(\ell\) one has
\[
 v_\ell |B_m^{(-1)^m}(E)|=v_\ell\!\left(w_m(J)|P_m|\right).
\]
In particular, with an absolute implicit constant,
\begin{align}
 \log |B_m^{(-1)^m}(E)|&\ll d m\log(4dmN),\label{eq:critorder99}\\
 \sum_{r=1}^{t-1}\log |B_{r+1}^{(-1)^{r+1}}(E)|
              &\ll d t^2\log(4dtN) \qquad(t\ge2).\notag
\end{align}
\end{theorem}
\begin{proof}
Fix \(\ell>2\). The comparison used in Theorem~\ref{thm:wild98}
applies to \(E\) and, at this odd prime, also to \(F\). The quadratic
extension is unramified away from the infinite places. On the rings
with \(\ell\) inverted it is an étale Galois cover. Restriction
identifies \(H^2(\OO_F[1/\ell],\ZZ_\ell(m))\) with the invariants
of \(H^2(\OO_E[1/\ell],\ZZ_\ell(m))\), where cohomology is étale.
Indeed all positive cohomology of the covering group on
\(\ell\)-primary coefficients vanishes by averaging. Apply
Hochschild--Serre at finite coefficients and take inverse limits as
in Theorem~\ref{thm:wild98}. Thus the plus group has the odd local
order of \(K_{2m-2}(\OO_F)\), and the minus group has the quotient
of the two local orders for \(E\) and \(F\).

If \(m=2k\), the abelian totally real special-value theorem gives
\[
 v_\ell|\zeta_F(1-m)|
   =v_\ell|K_{4k-2}(\OO_F)|-v_\ell|K_{4k-1}(\OO_F)|.
\]
This is the odd part of the theorem in
\cite[Ch.~VI, Thms.~8.7--8.8]{V-III-V-weibel98}. We verify that
\(v_\ell|K_{4k-1}(\OO_F)|=v_\ell w_m(F)\).
The group \(K_{4k-1}(\OO_F)\) has rank zero
by Borel's theorem. At this prime its completion is
\(H^1(\OO_F[1/\ell],\ZZ_\ell(m))\) by the same number-ring
comparison. The coefficient sequence from \(\ZZ_\ell(m)\) to
\(\QQ_\ell(m)\) identifies the torsion of this group with
\(H^0(F,\QQ_\ell/\ZZ_\ell(m))\). The preceding rational
invariants vanish since the cyclotomic image is open and \(m>0\).
The rational degree-one cohomology is a vector space, so the
connecting map has exactly the stated torsion image.

If \(m\) is odd, then \(m\ge3\). Write \(\xi\) for the quadratic
character of \(E/F\). The higher relative class number formula
gives, at odd \(\ell\),
\[
 v_\ell\!\left(\frac{|K_{2m-2}(\OO_E)|}{|K_{2m-2}(\OO_F)|}\right)
      =v_\ell\!\left(w_m(E)|L(1-m,\xi)|\right).
\]
In this formula the higher class numbers are the orders of
degree-two integral étale cohomology and agree with the indicated
\(K\)-group orders at odd primes. The higher unit index and the
explicit power of two have zero valuation at \(\ell\).
This is the odd-prime conclusion of the CM special-value theorem,
which is unconditional for every CM quadratic extension of a totally
real field.~\cite[Prop.~1.3]{V-III-V-voight99}

The abelian factorization of \(\zeta_E\) gives \(P_m=\zeta_F(1-m)\)
in even weight and \(P_m=L(1-m,\xi)\) in odd weight. Indeed the
characters trivial on \(\Gal(E/F)\) are precisely the even
characters. At a rational prime the local factors are computed on
the inertia invariants of these characters, which gives their
primitive Euler factors, including at the ramified primes.
The identity for \(L(s,\xi)\) follows on the half-plane of absolute
convergence by dividing \(\zeta_E(s)\) by \(\zeta_F(s)\), and then
by continuation of the product of odd characters. This proves the
valuation formula in both parities.

For \(J=E\) or \(F\), the index of its cyclotomic image in
\(\ZZ_p^\times\) divides \([J:\QQ]\le d\). If a twisted invariant
has order \(p^a\), the exponent of \((\ZZ/p^a)^\times\) divides
\(m\) times this index. As in Lemma~\ref{lem:loc98}, one obtains
\(p^a\le 4dm\) and \(p\le dm+1\). Hence
\(\log w_m(J)\ll dm\log(4dm)\).
By Lemma~\ref{lem:bern99}, \(w_m(J)D_m^fP_m\) is a nonzero
integer. At every odd prime its valuation is the group valuation
in the statement plus \(f v_\ell(D_m)\). Thus the odd group order
divides the absolute value of this integer. The lemma bounds its
logarithm and gives \eqref{eq:critorder99}. Summing over
\(m=2,\ldots,t\) gives the second estimate.
\end{proof}

In particular, on the intervals of Volume~I the full range
\(1\le r<\lfloor x^{401}\rfloor\) contributes only \(x^{O(1)}\)
to the logarithms in \eqref{eq:critorder99}. Here \(d\ll_\psi x^3\)
and \(\log(2N)\ll_{\psi,\eps}x\) on the chosen fields. The
eigenspace with sign \((-1)^r\) at twist \(r\) still requires a
bound of its own.

\subsection{Bernoulli sums for characters at every two-power layer}\label{nav:III-19}
\label{sec:bern139}

Let \(p\equiv3\pmod4\), and write \(\lambda(u)=(u/p)\).
For a primitive even character \(\alpha\) of order \(2^j\)
of the real cyclotomic tower, \(j\ge1\), its conductor is
\(f=2^{j+2}\). Put \(D_j=2^{j-1}\),
\(L_j=\QQ(\zeta_{2^j})\), and \(n=fp\). The quantity
\[
    B_\alpha=\frac1n\sum_{v=1}^n v\,\overline{\alpha\lambda}(v)
                 =-L(0,\overline{\alpha\lambda})
\]
is the same Bernoulli sum as in Section~\ref{sec:quartic115}.
The character \(\alpha\lambda\) is primitive and odd, so the
functional equation and nonvanishing at one give \(B_\alpha\ne0\).
Let \(F_j(p)\) denote the residue degree of \(p\) in \(L_j\),
including \(F_1(p)=1\).

\begin{lemma}[Integral coordinates and norms]\label{lem:coord139}
Write \(B_\alpha=\sum_{t=0}^{D_j-1}b_t\zeta_{2^j}^t\).
Then \(b_t\in\ZZ\) and \(|b_t|\le p-1\) for every \(t\).
At least one prime above \(p\) has valuation zero on \(B_\alpha\).
With valuations normalized by \(v_{\mathfrak p}(p)=1\), one has
\begin{align}
 F_j(p)\sum_{\mathfrak p\mid p}v_{\mathfrak p}(B_\alpha)
   &\le D_j\log_p\!\bigl((p-1)\sqrt{D_j}\bigr),\label{eq:ord139}\\
 F_j(p)\sum_{\mathfrak p\mid p}v_{\mathfrak p}(B_\alpha)
   &\le D_j\log_p\!\left(
                 \frac{\sqrt{fp}}{\pi}\bigl(2+\log(fp)\bigr)\right).
                 \label{eq:arch139}
\end{align}
Here \(\log_p y=\log y/\log p\) is a real logarithm.
\end{lemma}
\begin{proof}
For an integer \(c\) modulo \(f\) use the partial sum
\[
 A(c)=\frac1f\sum_{v=0}^{f-1}v\,\bar\alpha(c+v).
\]
Evenness and the vanishing of the complete character sum give
\(A(0)=0\) by pairing \(v\) and \(f-v\). A shift of the summation
index gives \(A(c+1)-A(c)=\bar\alpha(c)\). For \(0\le c<f\)
this is the sum of the first \(c\) character values.

Each coordinate of every such sum in the displayed power basis
is \(0,1\), or \(-1\). To see this, the even character of exact
order \(2^j\) is faithful on
\((\ZZ/f\ZZ)^\times/\{\pm1\}\). Its kernel on the units has
order two, and \(\alpha(1+f/2)=-1\) by primitivity.
The residues contributing to any fixed basis coordinate occur
in a four-element set. Writing its least residue as \(a\) gives
\[
       0<a<f/4,\qquad a,\quad f/2-a,\quad f/2+a,\quad f-a .
\]
Indeed the two operations are negation and translation by
\(f/2\), the latter coinciding on odd residues with multiplication
by \(1+f/2\). The four signs in increasing order are
\((+,-,-,+)\) up to overall sign. Every partial sum of those
signs is \(0,1\) or \(-1\), proving the assertion.

Write a residue prime to \(p\) as \(u+pv\) with \(1\le u<p\),
\(0\le v<f\), and sum first in \(v\). The part weighted by
\(u\) has zero sum since \(\alpha\ne1\). One obtains
\[
      B_\alpha=\bar\alpha(p)\sum_{u=1}^{p-1}
                         \lambda(u)A(p^{-1}u\bmod f).
\]
A sign times a \(2^j\)-th root of unity permutes the power
basis up to signs, since \(\zeta^{D_j}=-1\). This proves the
coordinate bound in the integral basis of \(L_j\). If all
primes above \(p\) divided \(B_\alpha\), then \(B_\alpha\)
would lie in \(p\mathcal O_{L_j}\), since \(p\) is unramified.
Every coordinate would be divisible by \(p\), forcing them all
to be zero. This contradicts nonvanishing.

The average of \(|\sigma B_\alpha|^2\) over the embeddings of
\(L_j\) is \(\sum_t b_t^2\). Indeed the average of
\(\sigma(\zeta)^{t-t'}\) over the odd powers of \(\zeta\)
vanishes for \(0<|t-t'|<D_j\) by a geometric sum.
The geometric mean of \(|\sigma B_\alpha|\) is therefore
at most \((p-1)\sqrt{D_j}\). The left side of \eqref{eq:ord139}
is the \(p\)-adic valuation of the absolute rational norm.
This proves \eqref{eq:ord139}. For a primitive odd character
\(\beta\) modulo \(n\), evaluating its completed functional
equation at zero gives
\(|L(0,\bar\beta)|=\sqrt n\,|L(1,\beta)|/\pi\). The series
for \(L(1,\beta)\) through \(n\) has absolute value at most
\(1+\log n\), and the absolute value of its tail is at most
one by summation by parts with partial sums bounded by \(n\).
Apply this at each embedding to prove \eqref{eq:arch139}.
\end{proof}

\begin{theorem}[Components of the strict group at all layers]
\label{thm:half139}
Suppose \(p>3\), \(p\equiv3\pmod4\), and
\(r+1\equiv(p-1)/2\bmod(p-1)\). For every \(j\ge1\), at least
one full Frobenius orbit of the primitive characters of order
\(2^j\) has zero strict component on \(C_J\) for every \(J\ge j\).
The same orbit works for all the indicated twists.
For \(p\equiv3\bmod8\) and \(J\ge2\), at least \(2^{J-1}+2\)
of the \(2^J\) character labels of \(C_J\) have zero component.

For a fixed \(p\)-adic embedding attached to \(\mathfrak p\mid p\)
put \(b_\alpha=v_{\mathfrak p}(B_\alpha)\). If
\(b_\alpha<1+v_p(r+1)\), the strict component for that embedding
is killed by \(p^{b_\alpha}\).
\end{theorem}
\begin{proof}
Extend coefficients to an unramified ring containing the
character values. Descent from \(C_J\) to \(\QQ\) separates the
characters by Theorem~\ref{thm:char}, since the degree of \(C_J\)
is a power of two. For \(\alpha\ne1\), the residual character
\(\bar\kappa_p^{-r}\bar\alpha^{-1}\) is globally nontrivial
by its conductor at two. The twist \(r\) is even. The reflection
test \eqref{eq:test24} applies with the primitive conductor \(f\),
as in Section~\ref{sec:quartic115}.
Put \(m=1+v_p(r+1)\). At level \(p^m\), every unit \(v\)
satisfies \(v^{r+1}=\lambda(v)\bmod p^m\). In fact \(r+1\)
kills its principal unit part modulo \(p^m\), and on the
Teichm\"uller part its power is exactly the Legendre symbol.

Choose \(a_0\equiv1\bmod p\) with \(\alpha(a_0)\ne1\) at
conductor \(f\) by the Chinese remainder theorem. The scalar
\(a_0-(\alpha\lambda)(a_0)\) is a \(p\)-adic unit since a
nontrivial root of unity of two-power order stays nontrivial
in residue characteristic \(p\). At modulus \(fp^m\) the floor
sum in \eqref{eq:test24} is congruent to
\((a_0-(\alpha\lambda)(a_0))B_\alpha\).
To check the change of modulus, lift each unit modulo \(fp\)
to its \(p^{m-1}\) representatives modulo \(fp^m\). The sum of
their least positive residues divided by \(fp^m\) is the
original residue divided by \(fp\), plus \((p^{m-1}-1)/2\).
The additional term sums to zero against the nontrivial
character. Now permute units by multiplication by \(a_0\),
as in \eqref{eq:floor115}, to obtain the identity for floors.
Thus if \(b_\alpha<m\), the test kills the \(p^m\)-torsion by
\(p^{b_\alpha}\). Any element of greater order would yield
an element of order exactly \(p^{b_\alpha+1}\) by multiplication
by a power of \(p\), a contradiction. This proves the exponent
bound, including vanishing for \(b_\alpha=0\).

All the primitive labels of order \(2^j\) are conjugates of
\(\alpha\). The quadratic character \(\lambda\) is rational,
so \(B_{\sigma\alpha}=\sigma B_\alpha\) under these conjugations.
In a fixed \(p\)-adic closure the conjugate labels run through
all primes of \(L_j\) above \(p\), with one orbit of size
\(F_j(p)\) for each prime. Lemma~\ref{lem:coord139} supplies
at least one unit orbit. For \(p\equiv3\bmod8\) and \(j\ge3\),
\(F_j(p)=2^{j-2}\), since \(v_2(p^2-1)=3\) and successive
squaring adds one to this valuation. The four labels through
order four all vanish by Theorem~\ref{thm:quartic115}.
Adding \(4+\sum_{j=3}^J2^{j-2}\) proves the stated count.
\end{proof}

\begin{corollary}[Uniform exponents at selected twists]
\label{cor:exp139}
Let \(p>3\), \(p\equiv3\bmod8\), \(J\ge2\), and assume the twist
congruence of Theorem~\ref{thm:half139}. If
\[
              v_p(r+1)\ \ge\ 1+\lceil\log_p(2^J)\rceil ,
\]
then the whole group \(S(C_J)\) is killed by
\(p^{\,\lceil2+\log_p(2^J)\rceil}\). If instead
\[
       p> 2^{J+2}\bigl(2+\log(2^{J+2}p)\bigr)^2,\qquad
       p\mid r+1 ,
\]
then \(p S(C_J)=0\).
\end{corollary}
\begin{proof}
The components through order four already vanish. At every
remaining layer \(D_j/F_j(p)=2\). From \eqref{eq:ord139},
\(b_\alpha<2+\log_p D_j\), so the first hypothesis puts
\(b_\alpha\) strictly below \(1+v_p(r+1)\) for every label.
For the second assertion \(\sqrt{fp}(2+\log(fp))/\pi<p\) for
all conductors \(f\le2^{J+2}\). Thus \eqref{eq:arch139}
gives \(\sum_{\mathfrak p\mid p}v_{\mathfrak p}(B_\alpha)<2\).
Every \(b_\alpha\) is a nonnegative integer and is at most one.
Apply Theorem~\ref{thm:half139} and descend the bounds from the
faithfully flat unramified coefficient extension.
\end{proof}

\subsection{An integral annihilator across the twists}\label{nav:III-20}
\label{sec:inter141}

Let \(E/\QQ\) be a finite abelian field, \(n=[E:\QQ]\), of
conductor dividing \(N\). Let \(p\) be odd, \(p\nmid n\), and
choose a character \(\alpha\) of \(\operatorname{Gal}(E/\QQ)\)
with values in an unramified coefficient ring \(O\) over \(\ZZ_p\).
Write \(\alpha=\alpha'\omega^i\) on the units of its conductor
times \(p\), where \(\alpha'\) is primitive of conductor \(f\)
prime to \(p\), and \(\omega\) is the Teichm\"uller character
at \(p\). This factorization follows by the Chinese remainder
theorem. The order of \(\alpha\) divides \(n\), so its local
factor at \(p\) is trivial on the principal units.
Fix \(r_0\bmod(p-1)\) with \(\alpha(-1)=(-1)^{r_0}\), and write
\(h=r_0+1+i\bmod(p-1)\).
For \(r>0\) in this residue class denote the strict component
over \(\QQ\) for \(\kappa_p^{-r}\alpha^{-1}\) by \(S_{\alpha,r}\).
If its residual character is trivial, Theorem~\ref{thm:all24}
already gives \(S_{\alpha,r}=0\). The following calculation
applies to the other components, including those ramified at \(p\).

\begin{theorem}[An interpolating annihilator and a nonzero value]
\label{thm:moment141}
Set \(b=1+fp\), and take the unique \(j\in\{1,\ldots,p-1\}\)
with \(j\equiv-h\bmod(p-1)\). There is a power series
\(\mathcal A(X)=\sum_{d\ge0}c_d X^d\) over \(O\), convergent
on \(\ZZ_p\), with the following properties.
\begin{enumerate}
\item \(v_p(c_d)\ge d-v_p(d!)\). For every \(m\ge1\),
the value \(\mathcal A(r+1)\bmod p^m\) annihilates
\(S_{\alpha,r}[p^m]\).
\item With \(k=j+1\), \(\chi=\bar\alpha'\) and
\(B_{k,\chi}=f^{k-1}\sum_{a=1}^f\chi(a)B_k(a/f)\), one has
\begin{equation}
  \mathcal A(-j)=
     \frac{b^k-1}{k b^{k-1}}\,(1-\chi(p)p^{k-1})\, B_{k,\chi}
      \ne0. \label{eq:moment141}
\end{equation}
Its valuation \(V\) satisfies, with an absolute constant,
\begin{equation}
     0\le V\le 1+C\, n p\,\frac{\log(2Np)}{\log p}.
     \label{eq:V141}
\end{equation}
\end{enumerate}
\end{theorem}
\begin{proof}
Write \(X_f=(\ZZ/f\ZZ)^\times\times\ZZ_p^\times\), with
trivial first group if \(f=1\). For the cylinder of units
represented by \(1\le v\le fp^m\) at modulus \(n_m=fp^m\),
assign the value
\[
       \mu_m(v)=\left\lfloor bv/n_m\right\rfloor-(b-1)/2.
\]
These \(p\)-integral values are compatible under summing over
lifts. Indeed for \(a\) a power of \(p\) the floor sum over
\(v+un_m\), \(0\le u<a\), equals
\(\lfloor bv/n_m\rfloor+(b-1)(a-1)/2\) since multiplication
by \(b\) permutes residues modulo \(a\). Subtracting the
\(a\) copies of \((b-1)/2\) gives \(\mu_m(v)\).
This defines a bounded \(p\)-adic integral on \(X_f\). First integrate
locally constant functions by finite summation. The absolute
value of the result is at most the uniform \(p\)-adic norm
of the function by the ultrametric inequality, so uniform
approximation defines the integral of every continuous function.

For \(v\in X_f\), write \(u=v^{(p)}/\omega(v^{(p)})\),
where \(v^{(p)}\in\ZZ_p^\times\) is its second coordinate.
Thus \(u\in1+p\ZZ_p\). Form the coefficients
\[
       c_d=\int_{X_f}\bar\alpha'(v)\omega(v^{(p)})^{-h}
                        \frac{(-\log u)^d}{d!}\,d\mu(v).
\]
Their asserted valuations follow from \(\log u\in p\ZZ_p\)
and boundedness of the integral. Since \(v_p(d!)\le d/(p-1)\),
the series converges uniformly on \(\ZZ_p\). Summing the
exponential series inside the integral shows that at each
integer \(s\) its integrand is
\(\bar\alpha'(v)\omega(v^{(p)})^{-h}u^{-s}\).
This is odd in \(v\), since \(\alpha'(-1)=(-1)^{h-1}\)
and \(u\) is unchanged by negation. After reduction modulo
\(p^m\) the integrand is constant on the indicated cylinders, and subtracting
the constant \((b-1)/2\) has zero effect modulo \(p^m\) by
pairing \(v\) with \(-v\). At \(s=r+1\) the integrand is
the inverse of \(\alpha(v)v^{r+1}\) on units. Consequently
the resulting floor sum is \eqref{eq:test24}, with prime-to-\(p\)
conductor \(f\). The proof of Corollary~\ref{cor:all24} applies
to the present field by Theorem~\ref{thm:all24}. This proves
annihilation at every level.

At \(s=-j\) the integrand becomes \(\chi(v)v^j\) on integer
representatives, since \(j\equiv-h\bmod(p-1)\).
We compute the limit of the uncorrected floor sums against
this polynomial, which are congruent to the integral modulo
\(p^m\) at level \(m\). For units \(v\bmod n_m\) put
\(q_v=\lfloor bv/n_m\rfloor\) and \(w=bv-n_m q_v\).
The \(w\)'s permute the units and \(\chi(b)=1\) because
\(b\equiv1\bmod f\). The binomial expansion, summed against
\(\chi(v)\), gives
\[
 k b^{k-1}\sum_v\chi(v)v^{k-1}q_v
    =(b^k-1)\frac1{n_m}\sum_v\chi(v)v^k+O(p^m).
\]
Here the error belongs to \(p^m O\), since the terms of degree
at least two in \(n_m q_v\) are divisible by \(n_m^2\) before
division by \(n_m\). The integer \(k\) is fixed as \(m\) tends
to infinity.

To compute the normalized moment, first allow all residues,
extending \(\chi\) by zero at nonunits of \(f\). For a residue
\(a\bmod f\) represented by \(1\le a\le f\), sum the powers of
\(a+f l\), \(0\le l<p^m\), using
\[
   \sum_{l=0}^{M-1}(l+z)^k
      =\frac{B_{k+1}(z+M)-B_{k+1}(z)}{k+1}.
\]
This identity follows from the difference formula for Bernoulli
polynomials. After division by \(M=p^m\) the right side tends
\(p\)-adically to \(B_k(z)\) by the derivative formula. Summing
over \(a\) and dividing by the remaining \(f\) shows that the
moment tends to \(B_{k,\chi}\). Subtracting the terms divisible
by \(p\) subtracts \(\chi(p)p^{k-1}\) times the all-residue
moment at the preceding level. This proves \eqref{eq:moment141}.

The primitive character \(\chi\) has parity \((-1)^k\).
The identity \(B_{k,\chi}=-kL(1-k,\chi)\), proved in
Section~\ref{sec:critical99} by Mellin integration, and the
functional equation show this is nonzero. Indeed the gamma factor at
\(1-k\) is regular for this parity and the Euler product at
\(k\ge2\) is nonzero. The Euler factor in \eqref{eq:moment141}
is a \(p\)-adic unit. Since \(p\nmid f\), the \(p\)-adic
logarithm or the binomial expansion gives
\(v_p(b^k-1)=1+v_p(k)\). Hence \(V=1+v_p(B_{k,\chi})\).

For completeness \(f(k+1)!\,B_{k,\chi}\) is an algebraic integer.
Indeed \((a+1)!B_a\in\ZZ\) follows inductively from the Bernoulli
recurrence. Expanding \(B_k(z)\) in the numbers \(B_a\)
introduces at most one further denominator \(f\) in
\(f^{k-1}B_k(v/f)\). For \(0\le z\le1\) the generating function
\(w e^{zw}/(e^w-1)\) is uniformly bounded on \(|w|=1\),
so \(|B_k(z)|\le C_1k!\). The absolute value of every conjugate
of the cleared integer is thus at most
\(\exp(C_2 k\log(2fk))\). The field of its character values
has degree at most \(n\), as the order of \(\alpha'\) divides
the exponent of \(\operatorname{Gal}(E/\QQ)\). Taking the
absolute rational norm bounds its valuation at any prime
above \(p\), normalized by \(v_p(p)=1\), by
\(C_2 n k\log(2fk)/\log p\). Clearing denominators can only
increase that valuation. The inequalities \(f\le N\) and
\(k\le p\) prove \eqref{eq:V141}.
\end{proof}

\begin{theorem}[A bound for the number of twists of large exponent]
\label{thm:inter141}
Let \(t\ge2\) be an integer. Put \(D=2(V+1)\) and
\(L=V+D\lfloor\log_p t\rfloor\). Among \(1\le r<t\) in the
chosen residue class there are at most \(D\) for which
\(p^L S_{\alpha,r}\ne0\).
\end{theorem}
\begin{proof}
Whenever \(v_p(\mathcal A(r+1))\le L\), Theorem~\ref{thm:moment141}
at precision \(p^{L+1}\) rules out an element of order
\(p^{L+1}\), by applying its scalar annihilator on that torsion.
Any element of greater order would yield one of exactly that
order by taking a multiple. Thus a twist counted in the
assertion must have \(v_p(\mathcal A(r+1))>L\), allowing a zero
value to have infinite valuation.

Suppose there are \(D+1\) such arguments \(s_0,\ldots,s_D\)
in \([2,t]\). Let \(Q(X)=\prod_i(X-s_i)\). For every \(d\),
the remainder \(R_d(X)\) of \(X^d\) modulo this monic polynomial
has integral coefficients and degree at most \(D\).
The series \(\sum_d c_dR_d(X)\) converges coefficientwise to
a polynomial \(R\) of degree at most \(D\), and
\(R(s_i)=\mathcal A(s_i)\). For \(d\le D\) the remainder is
just \(X^d\), and for \(d>D\) the coefficient bound of
Theorem~\ref{thm:moment141} gives \(v_p(c_d)\ge d/2>V\).
Evaluating the integral remainders at \(-j\) shows that
\(R(-j)-\mathcal A(-j)\in p^{V+1}O\). In particular
\(v_p(R(-j))=V\).

On the other hand evaluate the Lagrange interpolation formula
for \(R\) at \(-j\). In the term with node \(s_i\) the numerator
of the basis polynomial is integral, and its denominator has
valuation at most \(D\lfloor\log_p t\rfloor\) since
\(0<|s_i-s_l|<t\) for every \(l\ne i\).
Each term thus has valuation strictly greater than \(V\),
as does their sum. This contradiction proves the theorem.
\end{proof}

\subsection{An annihilator at primes dividing the field degree}\label{nav:III-21}
\label{sec:family143}

Let \(E/\QQ\) be finite abelian of degree \(n\) and conductor
dividing \(N\). Fix any odd prime \(p\). Write
\(\operatorname{Gal}(E/\QQ)=P\times H\), where \(P\) is
the \(p\)-Sylow subgroup. Put \(K=E^H\), \(d=|P|\), and
choose an unramified coefficient ring \(O\) over \(\ZZ_p\)
containing the values of every character of \(H\). As in
Lemma~\ref{lem:des28}, extend a character \(\alpha\) of \(H\)
trivially on \(P\), and write \(S_\alpha(K)\) for the strict
group over \(K\) with divisible coefficient
\(D=(\operatorname{Frac}(O)/O)(\vartheta)\), where
\(\vartheta=\kappa_p^{-r}\alpha^{-1}\). Here \(r\ge1\) and
\(\alpha(-1)=(-1)^r\). The prime-to-\(p\) order of \(\alpha\)
gives \(\alpha=\alpha'\omega^i\) as in Section~\ref{sec:inter141}.
We use \(\overline{\alpha'}\) for the inverse of this finite
character and \(\bar\vartheta\) for the residue-field reduction
of \(\vartheta\).

\begin{lemma}[The two kernels in reflection]\leavevmode\label{lem:loss143}
Use the determinant \(B_{b,m}\) of \eqref{eq:norm28}. For
every \(m\ge1\) and every permissible \(b\),
\[
       p^{2w_r}B_{b,m}\,S_\alpha(K)[p^m]=0,\qquad
       w_r=
       \begin{cases}
       0,&\bar\vartheta|_{G_K}\ne 1,\\
       1+v_p(d)+v_p(r),&\bar\vartheta|_{G_K}=1 .
       \end{cases}
\]
\end{lemma}
\begin{proof}
Choose \(\tau\in G_K\) such that
\(v_p(\vartheta(\tau)-1)=w_r\). In the first case this is the
definition of a nontrivial residual character. In the second,
\(\bar\vartheta\) factors through the \(p\)-group \(P\), so its
image in the multiplicative group of the residue field is
trivial. The prime-to-\(p\) characters \(\alpha\) and
\(\omega^{-r}\) then have the same reduction and are equal.
Choose \(g\in G_\QQ\) with \(\kappa_p(g)=1+p\) and take
\(\tau=g^d\). It fixes \(K\), and
\(\vartheta(\tau)=(1+p)^{-dr}\) has the required difference
from \(1\) by the \(p\)-adic logarithm for odd \(p\).

Take a class \(s\in S_\alpha(K)[p^m]\). The coefficient exact
sequence for multiplication by \(p^m\) lifts \(s\) to some
\(z\in H^1(K,D[p^m])\). Over the cyclotomic field \(L_m\)
of Section~\ref{sec:norm28} the finite coefficient is trivial.
Choose a generator of \(\mu_{p^m}\). Kummer theory turns
the restriction of \(z\) into a Kummer class of character
\(\beta_m=\kappa_{p,m}\vartheta^{-1}\) on
\(\Gamma_m=\operatorname{Gal}(L_m/K)\), with coefficients
extended to \(O/p^mO\). Indeed conjugation of the values
contributes \(\kappa_{p,m}\) and conjugation of the argument
contributes \(\vartheta^{-1}\).

All its finite valuations are divisible by \(p^m\).
On inertia away from \(p\) over \(L_m\) the original strict
class is a coboundary and its divisible coefficient is trivial
there. Thus the restricted finite cocycle vanishes on this
inertia. A \(p^m\)-th root over the maximal unramified extension
then has integral valuation in the same value group.
At \(w\mid p\) the finite cocycle over \(L_{m,w}\), viewed in
\(D\), is \((\kappa_p^{-r}-1)y_w\) for a local \(y_w\in D\).
It factors through the cyclotomic image. Lemma~\ref{lem:localunit24}
supplies unit representatives in every coefficient coordinate.
These arguments use the divisible local conditions on \(s\)
and hold for any choice of its finite lift.

Divide the principal ideal of the Kummer representative by
the exponent, sending \((a)=\mathfrak a^{p^m}\) to the ideal
class of \(\mathfrak a\). The resulting class is in the
\(\beta_m\)-eigenspace. The scalar \(B_{b,m}\) kills that
ideal class. Indeed the integral Stickelberger element kills
the entire class group of \(L_m\). The determinant of its
multiplication on the free group algebra over the subgroup
\(\Gamma_m\) is a multiple of that element by Cayley--Hamilton.
Evaluation on \(\Gamma_m\) at \(\beta_m\) gives \(B_{b,m}\)
by the coset-basis computation of Theorem~\ref{thm:norm28}.
This calculation applies directly to the image of our class.

The Kummer class of \(B_{b,m}z\) over \(L_m\) now comes from
units and still has character \(\beta_m\). This character is odd,
and \(K\) is totally real, since \(d\) is odd. The odd part of
the unit quotient is generated by roots of unity. In fact complex
conjugation \(c\) lies in \(\Gamma_m\), and \(u/c(u)\) is a root
of unity for every unit of the CM cyclotomic field. Multiplication
by two is invertible on the quotient in question. On its roots
of unity \(\tau\) acts as \(\kappa_{p,m}(\tau)\). Subtracting the
two actions shows that \(\vartheta(\tau)-1\) kills this Kummer
class. The assertion includes the case where the odd part is zero.

It follows that \((\vartheta(\tau)-1)B_{b,m}z\) is in the kernel of
restriction to \(L_m\). By inflation--restriction this kernel
comes from \(H^1(\Gamma_m,D[p^m])\). If \(y\) is a cocycle on
this abelian group, commuting \(\tau|_{L_m}\) with \(\sigma\)
in its cocycle identity gives
\[
   (\vartheta(\tau)-1)y(\sigma)
                  =(\vartheta(\sigma)-1)y(\tau|_{L_m}).
\]
The right side is a coboundary. Thus \((\vartheta(\tau)-1)^2B_{b,m}\)
kills \(z\), hence its image \(s\). The square has valuation \(2w_r\).
\end{proof}

Fix \(r\bmod(p-1)\), put \(h=r+1+i\bmod(p-1)\), and let
\(j\in[1,p-1]\) satisfy \(j\equiv-h\bmod(p-1)\). Write \(k=j+1\).
Let \(A\) be the least common multiple of the conductors of \(K\)
and \(\alpha\), \(F=A/p^{v_p(A)}\), \(a=\max(1,v_p(A))\), and
\(M=Fp^a\). In particular \(M\le pN\). We use \(b=1+M\).
For a character \(\xi\) on the units modulo \(M\), set
\[
 B_{k,\xi;M}=M^{k-1}\sum_{\substack{1\le v\le M\\(v,M)=1}}
                            \xi(v) B_k(v/M).
\]

\begin{theorem}[A determinant varying with the twist]\leavevmode\label{thm:family143}
There is a series \(\mathcal B(X)=\sum_{l\ge0}b_l X^l\)
over \(O\) with \(v_p(b_l)\ge l-v_p(l!)\).
For every positive twist in the fixed class and every \(m\ge1\),
\(p^{2w_r}\mathcal B(r+1)\) annihilates \(S_\alpha(K)[p^m]\).
Its value at \(-j\) is nonzero and equals
\begin{equation}
 \mathcal B(-j)=
 \left(\frac{b^k-1}{k b^{k-1}}\right)^d
 \prod_{\eta\in\widehat P} B_{k,\overline{\alpha'}\eta^{-1};M}.
 \label{eq:family143}
\end{equation}
Here \(\eta\) is viewed as the Dirichlet character of \(K\)
through the cyclotomic quotient. The product is evaluated in
a splitting field for \(P\).
\end{theorem}
\begin{proof}
On \((\ZZ/F\ZZ)^\times\times\ZZ_p^\times\), take the measure
whose value on the cylinder \(v\bmod Fp^m\) is
\(\lfloor bv/(Fp^m)\rfloor-(b-1)/2\), with its least positive
representative. Compatibility and boundedness follow from the
floor identity in the proof of Theorem~\ref{thm:moment141}.
Denote by \(\sigma(v)|_K\) the image in \(P\) of the unit
residue of \(v\) modulo \(M\) acting on roots of unity. Write
\(v^{(p)}\) for the second coordinate of \(v\) and
\(\langle v\rangle=v^{(p)}/\omega(v^{(p)})\).
With values in \(O[P]\), integrate the series
\[
 \overline{\alpha'}(v)\,\omega(v^{(p)})^{-h}
      [\sigma(v)^{-1}|_K]\,
              \exp(-X\log\langle v\rangle)
\]
coefficient by coefficient against the measure. Call the result
\(U(X)\). Every coordinate of its \(l\)-th coefficient has
valuation at least \(l-v_p(l!)\) by boundedness of the measure.
Let \(\mathcal B(X)\) be the determinant of multiplication by
\(U(X)\) on the basis \(P\). For degrees \(l_1+\cdots+l_d=l\),
\(\sum_i v_p(l_i!)\le v_p(l!)\) by the integral multinomial
coefficient. The determinant expansion gives the stated bound
and uniform convergence on \(\ZZ_p\).

The integrand is odd in \(v\), since its scalar weight is odd
and complex conjugation acts trivially on \(K\). The constant
correction to the floors therefore cancels. For \(m\ge a\)
the modulus \(n_m\) of \eqref{eq:norm28} equals \(Fp^m\).
At \(X=r+1\) the integrand modulo \(p^m\) is constant on these
cylinders, and its determinant after integration is exactly
\(B_{b,m}\). Lemma~\ref{lem:loss143} gives the asserted
annihilator. For \(m<a\) the same conclusion follows by
including the \(p^m\)-torsion in the \(p^a\)-torsion.

Over a characteristic-zero splitting field the regular
representation of the abelian group \(P\) splits into its
individual characters. At \(X=-j\) evaluation on \(\eta\)
turns the integrand on integer representatives into
\(\xi(v)v^{k-1}\), where \(\xi=\overline{\alpha'}\eta^{-1}\)
on units modulo \(M\). At \(n_m=Mp^{m-a}\), multiplication
by \(b\) permutes the units and fixes \(\xi\), since \(b\equiv1\bmod M\).
The binomial calculation in the proof of Theorem~\ref{thm:moment141}
gives the integral as \((b^k-1)/(k b^{k-1})\) times the limit
of \(n_m^{-1}\sum_v\xi(v)v^k\) over unit residues. Splitting this
sum into progressions \(v=u+Ml\) with \(u\) a unit modulo \(M\)
and applying the Bernoulli difference formula as in that proof
gives exactly \(B_{k,\xi;M}\). Multiplication over \(\eta\)
proves the formula.

To verify nonvanishing, let \(\xi_0\) be the primitive character
inducing \(\xi\), of conductor \(g\mid M\). The Bernoulli
distribution formula gives
\[
 B_{k,\xi;M}=B_{k,\xi_0}
      \prod_{\substack{q\mid M,\ q\ {\rm prime}\\q\nmid g}}
                                 (1-\xi_0(q)q^{k-1}).
\]
Indeed, the sum with the primitive weight extended by zero
at nonunits of \(g\), before imposing the extra coprimality
conditions, has the same normalized value for every multiple
of \(g\). This follows by grouping the terms above each residue
modulo \(g\) and using
\(\sum_{l=0}^{D-1}B_k((z+l)/D)=D^{1-k}B_k(z)\).
Deleting multiples of a squarefree divisor \(u\) of the extra
primes contributes \(\xi_0(u)u^{k-1}\) times that value in
inclusion--exclusion. This proves the product formula.
Every Euler factor is nonzero since \(k\ge2\), and
\(\xi_0(-1)=(-1)^k\) since \(\eta\) is even. The primitive
functional equation at \(1-k\) and the Euler product at \(k\)
give \(B_{k,\xi_0}\ne0\) as in Theorem~\ref{thm:moment141}.
\end{proof}

\begin{corollary}[Counting twists at every odd prime]\label{cor:family143}
Fix an integer \(t\ge2\). For the fixed \(\alpha,h\) put
\(V_{\alpha,h}=v_p(\mathcal B(-j))\) and \(D_0=2(V_{\alpha,h}+1)\).
All but at most \(D_0\) of the twists \(1\le r<t\) in the
chosen class satisfy
\[
     p^{V_{\alpha,h}+D_0\lfloor\log_p t\rfloor+2w_r}
                                                S_\alpha(K)=0 .
\]
For each rational conjugacy orbit \(\mathcal O\) of characters
of \(H\) and each \(h\) with the required parity,
\begin{equation}
 \sum_{\alpha\in\mathcal O}(1+V_{\alpha,h})
       \ll |\mathcal O|\,d\, k\,\frac{\log(2Np)}{\log p}.
 \label{eq:orbit143}
\end{equation}
The implicit constant is absolute.
\end{corollary}
\begin{proof}
The interpolation proof of Theorem~\ref{thm:inter141} uses only
the bound \(v_p(b_l)\ge l/2\), the distinct integral arguments
in \([2,t]\), and the nonzero value at \(-j\) in the unramified
ring \(O\). It bounds the number of arguments with valuation
larger than \(V_{\alpha,h}+D_0\lfloor\log_p t\rfloor\) by \(D_0\).
At each other argument Lemma~\ref{lem:loss143} and
Theorem~\ref{thm:family143} supply the displayed exponent.
Indeed an element of larger order would have a multiple of order
exactly one power higher, where the finite-level annihilator
already applies.

For the orbit estimate \(k\) is fixed by \(h\), and \(M\) is
constant on the orbit. Each
\(M(k+1)! B_{k,\xi;M}\) is an algebraic integer of absolute
value at most \(\exp(C k\log(2Mk))\) at every embedding, by
the Bernoulli polynomial expansion and the estimate on
\([0,1]\) in Theorem~\ref{thm:moment141}. The product of these
integers over \(\widehat P\) lies in the value field of \(\alpha\).
Indeed its prime-to-\(p\) factor \(\alpha'\) takes values there
and every automorphism fixing that field permutes the factors
indexed by \(\eta\). Embeddings of the value field correspond
to \(\mathcal O\). Take the absolute rational norm of the
product. The sum of its valuations over those embeddings is
at most \(C|\mathcal O|d k\log(2Mk)/\log p\). The cleared
integer denominators can only increase that sum. Finally
\(v_p((b^k-1)/(k b^{k-1}))=a\) by \(v_p(b-1)=a\) and the
logarithm. Adding \(|\mathcal O|(1+da)\), using \(p^a\le M\le pN\)
and \(k\le p\), proves \eqref{eq:orbit143}.
\end{proof}

\subsection{Bernoulli units at large two-power conductor}\label{nav:III-22}
\label{sec:fourier145}

Fix an odd prime \(p\) and an embedding of the algebraic character
values in \(\overline{\QQ}_p\). Let \(Q=p^D\) with \(D\ge1\) and
\(Q\equiv1\bmod4\), and put \(e=v_2(Q-1)\). Let \(\phi\) be a
primitive Dirichlet character of odd conductor \(f\) and order
dividing \(Q-1\). Let \(\xi\) be primitive of conductor \(M=2^s\).
We extend characters by zero on nonunits, and suppose
\(\phi(-1)\xi(-1)=-1\). Thus \(\phi\xi\) is a nontrivial odd
character of conductor \(fM\). The following estimate is uniform
at the primes above \(p\), including when \(p\mid f\).

\begin{theorem}[A unit criterion in the two-power tower]
\label{thm:fourier145}
If \(M\ge f4^e\), then
\[
     B_{1,\phi\xi}
       =\frac1{fM}\sum_{a=0}^{fM-1}a\phi(a)\xi(a)
       \quad\text{is a \(p\)-adic unit.}
\]
The conclusion also holds with both characters replaced by
their inverses.
\end{theorem}
\begin{proof}
First group the integers in the sum by their class \(u\bmod M\).
Set
\[
  A(u)=\frac1{fM}\!\!\sum_{\substack{0\le a<fM\\a\equiv u\ (M)}}
                      a\phi(a),\qquad
  G(u)=A(u)-A(0),\qquad S=\sum_{a\bmod f}\phi(a).
\]
The unweighted sum of \(\phi(a)\) in each such class is \(S\),
since \(M\) is prime to \(f\). Write \([u]\in[0,M-1]\) for
the representative of \(u\). For a multiple \(b=ft\le M\),
shifting representatives by \(b\) gives
\begin{equation}
    G(u)-G(u-b)=\frac{tS}{M}
                   -\mathbf1_{[u]<b}\phi([u]).
    \label{eq:jump145}
\end{equation}
Indeed the shift preserves \(\phi\), adds \(b\) to each
representative, and subtracts \(fM\) exactly for representatives
which end in \([0,b-1]\). This interval contains at most one
integer in the given class. Starting at \(G(0)=0\) and using
the shift \(f\), one obtains every \(G(u)\), since \(f\) generates
the additive group modulo \(M\). Formula \eqref{eq:jump145}
therefore makes all \(G(u)\) integral at \(p\), even when
\(f\) is divisible by \(p\). Their reductions lie in
\(\mathbf F_Q\), since the nonzero values of \(\phi\) are
\((Q-1)\)-st roots of unity. We have
\(B_{1,\phi\xi}=\sum_u \xi(u)G(u)\), since
\(\sum_u\xi(u)=0\). This also proves its integrality.

Work now in a finite residue field containing \(\mathbf F_Q\)
and all \(M\)-th roots of unity. Write the reduced \(G,\phi,\xi\)
with the same letters. Put \(m=s-e\), so \(m\ge e\ge2\), and
\[
       C=\{u\in(\ZZ/M\ZZ)^\times:\xi(u)^{2^e}=1\}.
\]
Suppose the Bernoulli sum is zero in this field. For every unit
\(y\bmod M\), changing \(u\) to \(uy\) gives
\(\sum_u\xi(u)G(uy)=0\). This forces
\begin{equation}
               \sum_{u\in C}\xi(u)G(uy)=0
               \qquad \text{for every unit }y .
               \label{eq:filter145}
\end{equation}
To see this, average the equations obtained by applying the
first \(2^m\) powers of the \(Q\)-power Frobenius. They fix
all \(G(uy)\). The powers of \(Q\) modulo \(M\) run through
exactly the residues \(1+2^e j\), \(0\le j<2^m\). Indeed
\(v_2(Q^{2^a}-1)=e+a\) by repeated squaring, since \(Q\equiv1\bmod4\).
This gives order \(2^m\) in the subgroup of that size.
For \(z=\xi(u)\) the average is \(z\) times the average of the
powers of \(z^{2^e}\). The geometric sum is zero unless
\(z^{2^e}=1\). Division by \(2^m\) is permitted in odd
characteristic. This proves \eqref{eq:filter145}.

Every \(u\in C\) satisfies \(u\equiv 1\) or \(-1\bmod 2^m\).
In fact the units modulo \(2^s\) are \(\{\pm1\}\times\langle5\rangle\).
The equality \(v_2(5^{2^a}-1)=2+a\) proves that \(5\) generates
the subgroup congruent to \(1\bmod4\). Primitivity of \(\xi\)
means that it takes the value \(-1\) on \(1+M/2\), the nonidentity
element of the kernel of reduction to \(M/2\). Thus \(\xi(5)\)
has order \(2^{s-2}\), also after reduction at \(p\). Writing
\(u=\pm5^a\), the condition on \(\xi(u)\) forces \(2^{m-2}\mid a\),
as asserted. The two signs are distinct since \(m\ge2\).

Take \(b=f2^e\), so \(1<b\le2^m\), and subtract
\eqref{eq:filter145} at \(y=1-b\) from the same equation at
\(y=1\). For the sign \(+\) one has \(u(1-b)=u-b\bmod M\),
since \(2^m\mid u-1\) and \(2^e\mid b\). In \eqref{eq:jump145}
the indicator is nonzero only at \(u=1\), where \(\xi(1)\phi(1)=1\).
For the sign \(-\) one has \(u(1-b)=u+b\). Use
\eqref{eq:jump145} at \(u+b\) with the opposite sign. Its
indicator is nonzero only if \([u+b]=b-1\), equivalently
\(u=-1\). This time its signed contribution is
\(\xi(-1)\phi(b-1)=\xi(-1)\phi(-1)=-1\).
The constant terms have total zero. Indeed multiplication by
\(1+M/2\) permutes each sign in \(C\) and negates \(\xi(u)\).
The difference of the two zero sums is therefore \(-1-1\),
a contradiction in characteristic \(p\). The inverse characters
satisfy the same hypotheses.
\end{proof}

\begin{corollary}[Residual cocycles unramified after restriction]
\label{cor:coc145}
Let \(\gamma\) be the reduction of \(\phi\xi\) in a finite
field \(k\) of characteristic \(p\), with the hypotheses of
Theorem~\ref{thm:fourier145}. Let \(\delta=\gamma\) or
\(\gamma^{-1}\). Suppose a continuous cocycle
\(z:G_\QQ\to k(\delta)\) vanishes on a subgroup of finite
index prime to \(p\) in \(I_v\cap\ker\delta\) for every finite
inertia group \(I_v\). Then \(z\) is a coboundary.
\end{corollary}
\begin{proof}
The restriction to \(I_v\cap\ker\delta\) is an additive
homomorphism. The index of its kernel divides the stipulated
index, while its image is a finite \(p\)-group. Thus it vanishes
on all of \(I_v\cap\ker\delta\). Over \(L=\QQ(\zeta_{fM})\)
the restriction of \(z\) is an everywhere unramified additive
homomorphism, and \(L\) has no real places. Class field theory
identifies it with a homomorphism \(u:\Cl(\OO_L)\to k\).
The cocycle identity gives
\(u(\sigma c)=\delta(\sigma)u(c)\) for every ideal class \(c\)
and \(\sigma\in\operatorname{Gal}(L/\QQ)\), using the Galois
action on classes under reciprocity.

Choose a positive integer \(b\equiv1\bmod \operatorname{lcm}(p,f)\)
and \(b\equiv1+M/2\bmod M\) by the Chinese remainder theorem.
The integral Stickelberger element \((b-\sigma_b)\Theta_{fM}\)
annihilates the class group by Stickelberger's theorem, used
in Section~\ref{sec:reflect}. Let \(\delta_0\) be the corresponding
character in characteristic zero. Evaluation of the element
at \(\delta_0\) gives
\((b-\delta_0(b))B_{1,\delta_0^{-1}}\) by the defining
sum for \(\Theta_{fM}\). The first factor reduces to \(2\)
and the second is a unit by
Theorem~\ref{thm:fourier145}. The group-ring element thus acts
by an invertible scalar on the image of \(u\), forcing \(u=0\).
The remaining class inflates from the finite abelian group
\(\operatorname{Gal}(L/\QQ)\). Commuting a group element of
character \(-1\) with an arbitrary element in the cocycle
identity shows that multiplication by \(-2\) is a coboundary
on this class. This proves its vanishing.
\end{proof}

The finite Fourier argument above reconstructs
\srcloc{Development/AutomorphicArithmetic/AutomorphicVarieties/LevelCohomology/LemmaX}{78}.
The two translated sums occur in
\srcloc{Development/AutomorphicArithmetic/AutomorphicVarieties/LevelCohomology/LxCore}{139}.
The application to residual cocycles appears in
\srcloc{Development/AutomorphicArithmetic/AutomorphicVarieties/LevelCohomology/ChrInvNeg}{197}.

\subsection{Effective descent in the real cyclotomic tower}\label{nav:III-23}
\label{sec:effect146}

Let \(E/\QQ\) be finite abelian. Use the notation of
Section~\ref{sec:family143} at an odd prime \(p\), so that
\(\operatorname{Gal}(E/\QQ)=P\times H\), \(K=E^H\), and
\(\alpha\) is a character of \(H\) extended trivially on \(P\).
Choose \(r\ge1\) with \(\alpha(-1)=(-1)^r\). Let
\(\beta=\omega^{r+1}\alpha\) denote the characteristic-zero
Teichm\"uller character of the residual reflected action.
It is odd. Let \(M_1=\operatorname{lcm}(p,\operatorname{cond}K,
\operatorname{cond}\alpha)\), and factor the primitive character
inducing \(\beta\) as \(\phi\xi\) of conductor \(f2^s\) with
\(f\) odd and \(\xi\) of two-power conductor \(2^s\).

\begin{proposition}[A vanishing test including primes dividing the degree]
\label{prop:effect146}
Suppose \(D\ge1\), \(Q=p^D\equiv1\bmod4\), the order of
\(\phi\) divides \(Q-1\), and \(2^s\ge f4^{e}\) with
\(e=v_2(Q-1)\). Suppose also that each prime \(q\mid M_1\)
not dividing \(f2^s\) satisfies \(q\not\equiv\pm1\bmod 2^{s-e}\).
Then \(S_\alpha(K)=0\).
\end{proposition}
\begin{proof}
The primitive value \(B_{1,\beta^{-1}}\) is a \(p\)-adic unit by
Theorem~\ref{thm:fourier145}. On enlarging the modulus to \(M_1\)
the Bernoulli number becomes
\[
 B_{1,\beta^{-1};M_1}
   =B_{1,\beta^{-1}}
       \prod_{\substack{q\mid M_1,\ q\ {\rm prime}\\q\nmid f2^s}}
                          (1-\beta^{-1}(q)).
\]
This follows by the Bernoulli distribution formula and
inclusion--exclusion as in Theorem~\ref{thm:family143}, now
at \(k=1\). If any extra factor vanishes modulo \(p\), then
\(\xi(q)\bmod p\) is in \(\mathbf F_Q^\times\). Its order is a
power of two dividing \(2^e\). The description of \(C\) in
the proof of Theorem~\ref{thm:fourier145} would give
\(q\equiv\pm1\bmod 2^{s-e}\). Thus every factor is a unit.

The first-level residual character \(\bar\vartheta\) on \(G_K\)
is nontrivial. Indeed on \(G_\QQ\) it is the reduction of
\(\omega\beta^{-1}\), whose two-power conductor is \(2^s\).
If it were trivial on \(G_K\) it would factor through the
\(p\)-group \(\operatorname{Gal}(K/\QQ)\), which has no
nontrivial character in the multiplicative residue field.
Choose \(b\equiv1\bmod (M_1)_{\rm odd}\) and
\(b\equiv1+2^{s-1}\bmod 2^s\), where \((M_1)_{\rm odd}\)
denotes the odd part of \(M_1\). This makes \(b\) prime to \(M_1\),
\(\beta(b)=-1\), and \(b\equiv1\bmod p\).

The augmentation of the group-algebra element in \eqref{eq:norm28}
at \(m=1\) is the reduction of
\[
 \sum_{v\in(\ZZ/M_1\ZZ)^\times}
            \left\lfloor b[v]/M_1\right\rfloor\beta^{-1}(v)
       =(b-\beta(b))B_{1,\beta^{-1};M_1}.
\]
Here \([v]\) is its least positive representative. To obtain
the equality, permute representatives by \(v\mapsto bv\) and
use \([bv]=b[v]-M_1\lfloor b[v]/M_1\rfloor\). The constant
term in \(B_1(z)=z-1/2\) sums to zero against the odd character.
The augmentation is therefore nonzero. Over the residue field
the determinant of multiplication on its group algebra of \(P\)
is the \(|P|\)-th power of that augmentation, by the nilpotence
calculation in Theorem~\ref{thm:norm28}. It kills \(S_\alpha(K)[p]\)
by that theorem and is a unit. Any nonzero element in the
\(p\)-primary group would have a multiple of order \(p\), so
the entire group vanishes.
\end{proof}

Fix a primitive character \(\psi\) of conductor \(F\), let
\(o=[K_\psi:\QQ]\), and put \(B_j=K_\psi C_j\) with
\(C_j=\QQ(\zeta_{2^{j+2}})^+\). For \(B_j\) or any other
abelian field write \(M_{p,r}=\tfrac12(1+(-1)^r c)\Sel_{p,r}\).
Let \(D=2\varphi(o)\), where \(\varphi(1)=1\). For every odd
prime \(p\) put
\begin{equation}
 e_p=v_2(p^D-1),\qquad
 H_p=4(F+2)(p+1)4^{e_p},\qquad
 a_p=\min\{a\ge0:2^{a+2}\ge H_p\}.
 \label{eq:level146}
\end{equation}

\begin{theorem}[An explicit layer for every complementary twist]
\label{thm:effect146}
For every odd prime \(p\), every \(r\ge1\), and every \(J\ge a_p\),
restriction gives an isomorphism
\[
                      M_{p,r}(B_{a_p})\ \simeq\ M_{p,r}(B_J).
\]
For \(J<a_p\) it gives an injection in the other direction.
The field \(B_{a_p}\) has degree and cyclotomic modulus
\(O_\psi((p+1)4^{e_p})=O_\psi(p^3)\).
\end{theorem}
\begin{proof}
Write a character \(\alpha\) of the prime-to-\(p\) part of
\(\operatorname{Gal}(B_J/\QQ)\) as \(\lambda\zeta\), where
\(\lambda\) comes from \(K_\psi\) and \(\zeta\) from \(C_J\).
Such a factorization holds because the restriction map embeds
the Galois group of the compositum in the product of the two
Galois groups, and a finite abelian subgroup character extends
to the product. To extend across one generator whose coset
has order \(a\), choose an \(a\)-th root of the prescribed
value on its \(a\)-th power. Repeating this proves the claim.
Choose all values as algebraic roots of unity. Since \(\zeta\)
has two-power order and \(\alpha\) has order prime to \(p\),
the pullback of \(\lambda\) to the compositum has order prime
to \(p\). The restriction to \(K_\psi\) is onto, so the order
of \(\lambda\) divides the prime-to-\(p\) part of \(o\).

Let \(2^s\) be the two-part of the conductor of \(\alpha\).
In \(\beta=\omega^{r+1}\alpha=\phi\xi\) the dyadic character
has the same conductor, since \(p\) is odd. The odd conductor
\(f\) of \(\phi\) divides \(pF\), and its order divides
\(\operatorname{lcm}(p-1,o/p^{v_p(o)})\). For any divisor
\(u\) of \(o\), \(\varphi(u)\) divides \(\varphi(o)\), as seen
prime by prime in the formula for Euler's totient. Thus
\(p^D\equiv1\bmod o/p^{v_p(o)}\), and \(D\) also ensures
\(p^D\equiv1\bmod4\). The order condition in
Proposition~\ref{prop:effect146} follows.

The \(p\)-power subfield \(K\) of \(B_J\) is contained in
\(K_\psi\), since \([B_J:K_\psi]\) is a power of two. The odd
prime divisors of \(M_1\) are consequently among the primes
dividing \(pF\). If \(2^s>H_p\), then \(2^s>f4^{e_p}\)
and \(2^{s-e_p}>\max(p,F)+1\). Each odd prime in \(M_1\)
outside the primitive conductor lies strictly between \(1\)
and \(2^{s-e_p}-1\). Proposition~\ref{prop:effect146}
eliminates this component through every exponent.

It remains to interpret the components with \(2^s\le H_p\).
The dyadic conductor of \(\lambda\) divides \(2^{v_2(F)}\),
and that of \(\zeta=\alpha/\lambda\) divides the larger of
\(2^s,2^{v_2(F)}\), hence divides \(2^{a_p+2}\).
As an even character of two-power conductor, \(\zeta\) factors
through \(C_{a_p}\). Therefore \(\alpha\) is trivial on
\(\operatorname{Gal}(B_J/B_{a_p})\) when \(J\ge a_p\).
Over unramified coefficients the integral idempotents for the
prime-to-\(p\) part now show that the whole complementary
group is invariant under this two-group, including its
components with nontrivial \(P\)-action. Descent for a
prime-to-\(p\) extension identifies those invariants with
the strict group downstairs, since averaging gives exact
invariants globally and on each inertia and decomposition
group. This argument respects complex conjugation, and faithful
scalar extension gives the claimed integral result. If \(J<a_p\)
restriction followed by corestriction is multiplication by the
two-power relative degree, so restriction is injective.

Finally
\[
 e_p=v_2(p-1)+v_2(p+1)+v_2(D)-1
    =\max(v_2(p-1),v_2(p+1))+v_2(D).
\]
For the first equality use \(v_2(p^u-1)=v_2(p-1)\) for odd
\(u\) and repeated squaring after the first factorization of
\(p^{2u}-1\). Precisely one of \(p-1,p+1\) has valuation one,
which gives the second equality. Thus \(2^{e_p}\le
2^{v_2(D)}(p+1)\). Both the degree and the modulus are at
most \(F2^{a_p+2}\le 2F H_p\), which gives the bounds.
\end{proof}

\subsection{Bernoulli values shared by different coefficient primes}\label{nav:III-28}
\label{sec:shared201}

Fix a finite abelian extension \(E/\QQ\) of degree \(n\) in
\(\QQ(\zeta_N)\). Let \(c\) denote complex conjugation. For an
odd prime \(p\) and a positive integer \(r\), put
\[
 M_{p,r}(E)=\tfrac12(1+(-1)^r c)\ker\!\left[
 H^1(\mathcal O_E[1/p],(\QQ_p/\ZZ_p)(-r))
 \longrightarrow\prod_{w\mid p} H^1(E_w,(\QQ_p/\ZZ_p)(-r))\right].
\]
We use continuous Galois cohomology, unramified away from \(p\) in
the global term. Suppose first that \(p\nmid nN\). Fix an embedding
of the algebraic character values into \(\overline{\QQ}_p\).
Take \(O\) to be the ring of integers of a finite unramified
extension of \(\QQ_p\) containing the values of all characters of
\(G=\Gal(E/\QQ)\). Such an extension exists since \(p\nmid n\).
Normalize \(v_p(p)=1\). A character of \(G\) will be written as its
primitive Dirichlet character. For \(\alpha(c)=(-1)^r\), let
\(S_{\alpha,r}\) denote the \(\alpha\)-component of \(M_{p,r}(E)\otimes O\).
The idempotents for \(G\) are integral at \(p\). Their sum over this
parity is the identity on the indicated module.

Write \(\chi=\bar\alpha=\alpha^{-1}\) on the units of its conductor
\(f\). Let \(k\ge2\) with \(k+r\equiv0\bmod(p-1)\). Then \(\chi\)
has parity \((-1)^k\). With the convention \(B_1=-1/2\), put
\[
 z_{k,\chi}=\frac{B_{k,\chi}}{k},\qquad
 B_{k,\chi}=f^{k-1}\sum_{a=1}^f\chi(a)B_k(a/f).
\]
For the trivial character take \(f=1\). The identity
\(z_{k,\chi}=-L(1-k,\chi)\) and the functional equation give
\(z_{k,\chi}\ne0\). Indeed the archimedean gamma factors at
\(1-k\) and \(k\) are regular for this parity, and the Euler
product at \(k>1\) is nonzero.

\begin{lemma}[Comparison at a negative integer]\label{lem:shared201}
If \(\alpha=1\) and \(p-1\mid r\), then \(S_{\alpha,r}=0\).
In every other case \(v=v_p(z_{k,\chi})\) is a nonnegative integer.
If \(v<1+v_p(r+k)\), then \(p^v S_{\alpha,r}=0\).
In particular \(v=0\) gives vanishing for every positive twist in
the residue class \(-k\bmod(p-1)\).
\end{lemma}

\begin{proof}
Let \(\omega\) be the Teichm\"uller character at \(p\). As in
Theorem~\ref{thm:all24}, character descent gives the coefficient
character \(\alpha^{-1}\kappa_p^{-r}\) over \(\QQ\), where
\(\kappa_p\) is the cyclotomic character. Its residual character
on \((\ZZ/f\ZZ)^\times\times\mathbf F_p^\times\) is
\(\chi\,\omega^k\bmod p\). Triviality is equivalent to
\(\alpha=1\) and \(p-1\mid r\), by the independence of the two
factors and injectivity of reduction on the roots of unity of
order prime to \(p\). Theorem~\ref{thm:all24} gives the first
assertion in precisely this case.

Otherwise choose a positive integer \(b>1\), prime to \(fp\), at
which \(\chi(b)b^k\not\equiv1\bmod pO\). Such a residue exists
because the residual character is nontrivial, and lifts by the
Chinese remainder theorem. Let \(X_f=(\ZZ/f\ZZ)^\times\times
\ZZ_p^\times\). For a unit residue \(a\) at modulus \(m=fp^l\),
represented in \([1,m]\), give its cylinder the weight
\(\lfloor ba/m\rfloor-(b-1)/2\). These weights define a bounded
\(O\)-valued measure \(\mu_b\). In fact the child residues are
\(a+mi\) for \(0\le i<p\), and \(bi\bmod p\) permutes the
residues modulo \(p\). Summing their floors gives
\(\lfloor ba/m\rfloor+(b-1)(p-1)/2\), proving compatibility
for \(l\ge1\). Each weight is integral since \(p\) is odd.
For the \(p\)-coordinate of \(a\), put
\(\langle a\rangle=a^{(p)}/\omega(a^{(p)})\).
With \(h=r+1\bmod(p-1)\), set
\[
 A_b(s)=\int_{X_f}\chi(a)\omega(a^{(p)})^{-h}
                            \langle a\rangle^{-s}\,d\mu_b(a)
                    \qquad(s\in\ZZ_p).
\]
The bounded integral is the limit of locally constant
approximations. At \(s=r+1\) its integrand modulo \(p^l\)
is \(\alpha(a)^{-1}a^{-r-1}\). Its sign at \(-a\) is negative.
Thus the constant subtracted from the floors contributes zero,
by pairing the residues. The finite floor sum annihilates
\(S_{\alpha,r}[p^l]\) by Corollary~\ref{cor:all24}. This is the
class-group annihilation from the integral Stickelberger element
on the cyclotomic field of conductor \(fp^l\).

At \(s=1-k\) the integrand on integer representatives is
\(\chi(a)a^{k-1}\). To compute its integral let
\(q_a=\lfloor ba/m\rfloor\) and \(w_a=ba-mq_a\). The \(w_a\)
permute the units modulo \(m\) and \(\chi(w_a)=\chi(b)\chi(a)\).
Expanding \(w_a^k\), multiplying by \(\chi(a)\), and summing
over the units gives
\[
 k b^{k-1}\sum_a\chi(a)a^{k-1}q_a
  \equiv(b^k-\chi(b)^{-1})\,m^{-1}\sum_a\chi(a)a^k\pmod{p^lO}.
\]
Every binomial term of degree at least two in \(mq_a\) contains
\(m^2\) before division by \(m\). For the normalized sum on the
right, first allow multiples of \(p\). On each progression modulo
\(f\) the identity
\(\sum_{i=0}^{u-1}(y+i)^k=(B_{k+1}(y+u)-B_{k+1}(y))/(k+1)\)
with \(u=p^l\) gives the limit \(B_{k,\chi}\) after summation
and normalization. Indeed the linear term of the numerator is
\((k+1)B_k(y)u\) and every higher term tends to zero after
division by \(u\). Deleting multiples of \(p\) subtracts
\(\chi(p)p^{k-1}\) times the preceding normalized sum.
The integral of the constant correction again vanishes by oddness.
Letting \(l\) tend to infinity gives
\begin{equation}
 A_b(1-k)=(b^k-\chi(b)^{-1})b^{1-k}
           (1-\chi(p)p^{k-1})\, z_{k,\chi}.
 \label{eq:value201}
\end{equation}
All factors multiplying \(z_{k,\chi}\) are units. Thus its
valuation equals the valuation
of an integral and is nonnegative. Moreover
\(\langle a\rangle^{r+k}\equiv1\bmod p^{1+v_p(r+k)}\) for every
\(a\), by the binomial theorem on \(1+p\ZZ_p\). Boundedness gives
the same congruence between \(A_b(r+1)\) and \(A_b(1-k)\).
Under the strict inequality in the lemma both have valuation \(v\).
The annihilator at \(l=v+1\) rules out a cyclic summand of order
\(p^{v+1}\), hence also an element of any larger order.
\end{proof}

For each \(k\ge2\) take one representative \(\alpha\) from each
rational conjugacy orbit of characters of \(G\) of parity
\((-1)^k\). Write \(F_\alpha=\QQ(\alpha)\) for the field of values.
In its ring of integers factor the principal fractional ideal
of \(z_{k,\bar\alpha}\) into a quotient of relatively prime
integral ideals. Write \(\mathfrak a_{k,\alpha}\) for the numerator
ideal and define the positive rational integer
\begin{equation}
 J_k(E)=\prod_{\alpha\ {\rm representative}}
                \mathrm N_{F_\alpha/\QQ}(\mathfrak a_{k,\alpha}).
 \label{eq:int201}
\end{equation}
An empty product equals one. Ideal norms make this independent of
the representatives. In particular the same integer can be used
at different coefficient primes.

\begin{theorem}[An integer for all primes at a fixed weight]
\label{thm:integer201}
For \(k\ge2\) one has \(\log J_k(E)\ll nk\log(2Nk)\) with an
absolute constant. For every odd \(p\nmid nN\) put \(w=v_p(J_k(E))\).
If \(p-1\mid r+k\) and \(w<1+v_p(r+k)\), then
\[
                         p^w M_{p,r}(E)=0.
\]
In particular \(M_{p,r}(E)=0\) for all positive \(r\equiv-k\bmod(p-1)\)
at an odd prime \(p\nmid nN J_k(E)\).
\end{theorem}

\begin{proof}
For each prime \(p\), the embeddings of \(F_\alpha\) into
\(\overline{\QQ}_p\) over a fixed prime ideal \(\mathfrak p\mid p\)
number \(e_\mathfrak p f_\mathfrak p\), its ramification index
times its residue degree. With \(v_p(p)=1\), their total positive
valuation on \(z=z_{k,\bar\alpha}\) is
\(f_\mathfrak p\max(0,\ord_\mathfrak p z)\). Summing first over
\(\mathfrak p\) and then over the orbits yields
\begin{equation}
 v_p(J_k(E))=\sum_{\alpha(c)=(-1)^k}
                     \max\{0,v_p(z_{k,\bar\alpha})\},
 \label{eq:norm201}
\end{equation}
where every individual character occurs once on the right.

The integer \(a=f k(k+1)!\) clears the denominators of \(z\).
Indeed \((i+1)!B_i\in\ZZ\) follows from the Bernoulli recurrence,
and in \(f^{k-1}B_k(j/f)\) at most one further factor \(f\)
is needed. The numerator ideal of \(z\) divides the integral
ideal \((az)\), so its norm is at most \(|\mathrm N(az)|\).
The generating function \(u e^{yu}/(e^u-1)\) on \(|u|=1\)
gives \(|B_k(y)|\le Ck!\) uniformly for \(0\le y\le1\).
It follows that every complex conjugate of \(az\) has absolute
value at most \(\exp(C'k\log(2Nk))\). The sum of the degrees of
the orbit representatives equals the number of characters of
that parity, at most \(n\). Multiplication proves the height bound.
Equation~\eqref{eq:norm201} bounds the nonnegative valuation at
each nontrivial residual component by \(w\). Apply
Lemma~\ref{lem:shared201} to each component and descend the
annihilation through the faithfully flat extension to \(O\).
\end{proof}

\begin{theorem}[Counting exceptional twists across primes]
\label{thm:shared201}
Let \(t\ge2\) be an integer and \(Y\ge3\) a real number. For
each odd \(p\le Y\) with \(p\nmid nN\) there is a set
\(\mathcal E_p\subseteq\{1,\ldots,t-1\}\) such that
\begin{equation}
 \sum_{\substack{2<p\le Y\\p\ {\rm prime},\ p\nmid nN}}
       |\mathcal E_p|\log p
 \ \le\ 2\sum_{k=2}^{\lfloor Y\rfloor}\log J_k(E)
 \ \ll\ nY^2\log(2NY).
 \label{eq:count201}
\end{equation}
For \(1\le r<t\) outside \(\mathcal E_p\), let \(k\) be the unique integer in
\([2,p]\) congruent to \(-r\bmod(p-1)\). Then
\[
 p^{(1+2\lfloor\log_p t\rfloor)v_p(J_k(E))}M_{p,r}(E)=0.
\]
Here \(\log_p t=(\log t)/(\log p)\).
\end{theorem}

\begin{proof}
Fix a nontrivial residual component and its class of \(r\bmod(p-1)\).
In Lemma~\ref{lem:shared201} we may hold \(b,h,k\) fixed on that
class. Write \(v=v_p(z_{k,\chi})\). At \(v=0\) every group in the
class vanishes. Suppose \(v\ge1\) and put \(D=2v\). In the expansion
\(A_b(s)=\sum_{i\ge0}c_i s^i\), bounded integration of
\((-\log\langle a\rangle)^i/i!\) gives \(v_p(c_i)\ge i-v_p(i!)\ge i/2\).
The series converges uniformly on \(\ZZ_p\).

At most \(D\) integers \(s=r+1\) in the class and in \([2,t]\)
can have \(v_p(A_b(s))>v+D\lfloor\log_p t\rfloor\).
In fact, if there were \(D+1\) distinct such points, reduce each
monomial of the series modulo their monic product
\(\prod_j(X-s_j)\). All remainders are integral, so the sum of the
reduced series is a polynomial \(R\) of degree at most \(D\).
At the nodes it has values \(A_b(s_j)\). Terms of index \(i>D\)
have valuation greater than \(v\), giving
\(R(1-k)\equiv A_b(1-k)\bmod p^{v+1}\). The latter has valuation
\(v\) by \eqref{eq:value201}. On the other hand, evaluating
the Lagrange interpolation formula for \(R\) at \(1-k\) gives
terms of valuation strictly greater than \(v\). Each denominator
is a product of \(D\) nonzero differences of absolute value less
than \(t\), and each numerator factor is integral. This is a
contradiction.

Put those exceptional twists into \(\mathcal E_p\), taking the union
over all characters and classes. At every other twist the
finite-level annihilator bounds the exponent as in the lemma.
The exceptional count for a single character and \(k\) is at
most \(2v\), including when \(v=0\). Components trivial residually
need no exceptions. Summing \eqref{eq:norm201} first at \(p\), and
then interchanging the finite sums over \(p\) and \(2\le k\le p\),
uses each prime factor of \(J_k(E)\) at most to its multiplicity.
This proves \eqref{eq:count201}. The last estimate follows by summing
Theorem~\ref{thm:integer201} over \(k\le Y\).
\end{proof}

\begin{corollary}[Vanishing in shifted residue classes]
\label{cor:shared201}
For every odd prime \(p\) and every positive even \(r\),
\[
 K_{2r}(\ZZ)\{p\}=0
 \quad\text{if }p-1\mid r+k
 \text{ for some }k\in\{2,4,6,8,10,14\}.
\]
In particular \(K_{2p-6}(\ZZ)\{p\}=0\) for every prime \(p\ge5\).
For \(E=\QQ(i)\) the entire group \(M_{p,r}(E)\) vanishes for
every odd prime \(p\) and every positive odd \(r\) with \(p-1\mid r+3\).
\end{corollary}

\begin{proof}
Over \(\QQ\) the Bernoulli recurrence gives, in order of increasing \(k\),
\[
 \frac{1}{12},\quad -\frac{1}{120},\quad \frac{1}{252},\quad
 -\frac{1}{240},\quad \frac{1}{132},\quad \frac{1}{12}
\]
for \(B_k/k\). Thus their numerator ideals all have norm one.
Theorem~\ref{thm:integer201} gives the vanishing of \(M_{p,r}(\QQ)\).
For even \(r\), its Pontryagin dual is \(K_{2r}(\ZZ)\{p\}\).
Indeed the localization sequence and odd-primary Quillen--Lichtenbaum
comparison identify the latter with \(H^2(\ZZ[1/p],\ZZ_p(r+1))\),
and Poitou--Tate identifies the dual of the kernel of localization
there with the strict group. The local \(H^2\) term is dual to
\(H^0(\QQ_p,(\QQ_p/\ZZ_p)(-r))\). Global and local cyclotomic
images are both \(\ZZ_p^\times\), so their invariants agree.
Exactness of the final localization and sum-of-invariants maps
therefore makes all the global \(H^2\) strict. These identifications
are given in Section~\ref{sec:critical99} and in the proof of
Corollary \ref{cor:transfer100}.
For \(E=\QQ(i)\) only \(\chi_{-4}\) has odd parity. Since
\(B_3(u)=u^3-3u^2/2+u/2\), one has
\(4^2(B_3(1/4)-B_3(3/4))/3=1/2\).
Thus \(J_3(E)=1\). Taking \(N=4\) gives \(nN=8\), prime to every odd \(p\).
\end{proof}

\subsection{Character comparison at primes dividing the degree}\label{nav:III-29}
Let \(E/\QQ\) be finite abelian with group \(G\) and degree \(n\).
Fix an odd prime \(p\). Let \(O\) be the integers in a finite
extension \(J/\QQ_p\) containing the values of the characters of
\(G\), and put \(D=J/O\). For a character \(\alpha\) of \(G\)
write \(S_{\alpha,r}\) for the subgroup of
\(H^1(\QQ,D(\alpha^{-1}\kappa_p^{-r}))\) consisting of classes
zero on inertia away from \(p\) and on the decomposition group at
\(p\). Write \(M_{p,r}(E)\) for the sign \((-1)^r\) in the strict
group over \(E\) with coefficients \((\QQ_p/\ZZ_p)(-r)\).
All valuations below have \(v_p(p)=1\).

\begin{lemma}[Integral comparison with all the scalar groups]
\label{lem:fourier231}
There are maps in both directions between \(M_{p,r}(E)\otimes O\)
and \(\bigoplus_{\alpha(c)=(-1)^r} S_{\alpha,r}\), with both
composites multiplication by \(n\). On the summand indexed by
\(\alpha\), the action of a deck transformation \(g\) of \(E/\QQ\)
is \(\alpha(g)\). In particular, for any projector formed from
averages over subgroups of order prime to \(p\), an exponent
\(p^a\) on its scalar images gives exponent \(p^{a+v_p(n)}\)
on its image over \(E\).
\end{lemma}
\begin{proof}
The induced coefficient over \(\QQ\) is
\(U=D[G]\otimes\kappa_p^{-r}\), where \(D[G]\) denotes one copy of
\(D\) for each \(g\in G\). Arithmetic translation on the basis is
\(e_g\mapsto e_{hg}\) for \(h\in G\). The commuting deck action
of \(t\in G\) is \(e_g\mapsto e_{gt^{-1}}\).
Shapiro's lemma gives the degree-one cohomology over \(E\).
Indeed evaluation at the identity coset on restricted cochains
is the Shapiro map. Translation of that coset gives the usual
action on the classes by conjugation and transport of coefficients.
For a decomposition or inertia group \(I\), its orbits on the cosets
give a direct sum of induced modules for the stabilizers
\(I\cap gG_Eg^{-1}\). Evaluation on each orbit identifies the
restriction map with the restrictions at all corresponding places
of \(E\). Applying this at the decomposition group above \(p\)
and at every inertia group away from \(p\) proves the assertion
for the strict local conditions, including at primes ramified in
\(E\). Extension from \(\ZZ_p\) to \(O\) commutes with these kernels
since \(O\) is finite free.

Suppress the common twist. On tuples in the indicated copies of
\(D\) use the two maps
\[
 (x_g)_g\longmapsto
       \left(\sum_g\alpha(g)^{-1}x_g\right)_\alpha ,
 \qquad
 (y_\alpha)_\alpha\longmapsto
       \left(\sum_\alpha\alpha(g)y_\alpha\right)_g .
\]
Their entries belong to \(O\). Arithmetic translation acts on
each output of the first map by \(\alpha(h)^{-1}\), and deck
translation by \(\alpha(t)\). The maps commute with the coefficient
actions in both directions. Character orthogonality gives both
composites \(n\) on the coefficient modules, hence on cohomology
and on its strict kernels. Project to the prescribed sign using
the idempotent for complex conjugation, since \(2\) is a unit.
For the final assertion apply both maps to the projected module.
A scalar bound on the first image kills \(n\) times that module.
The prime-to-\(p\) part of \(n\) acts invertibly, and faithful
flatness descends the resulting bound to \(\ZZ_p\).
\end{proof}

\begin{theorem}[An integral scalar annihilator with ramified character values]
\label{thm:scalar231}
Let \(\chi=\alpha^{-1}\) be primitive of conductor \(f\). Suppose
\(k\ge2\), \(\alpha(c)=(-1)^k\), and
\(\bar\chi\bar\kappa_p^{\,k}\ne1\), where the bar denotes reduction
in the residue field of \(O\). There is a positive integer \(b\)
prime to \(fp\) and a series \(F\in O[[X]]\) such that, on putting
\(X_s=(1+p)^s-1\) and \(A(s)=F(X_s)\), one has
\[
 A(r+1)S_{\alpha,r}=0
       \quad(r\ge1,\ p-1\mid r+k).
\]
For every \(K=k+j(p-1)\), \(j\ge0\), the same series satisfies
\[
 A(1-K)=b^{1-K}(b^K-\chi(b)^{-1})
           (1-\chi(p)p^{K-1})\,{B_{K,\chi}\over K}.
\]
The factors preceding the Bernoulli value are units of \(O\),
with \(\chi(p)=0\) if \(p\mid f\).
\end{theorem}
\begin{proof}
Choose a unit residue \(b\) modulo \(\operatorname{lcm}(f,p)\)
on which \(\bar\chi(b)\bar b^k\ne1\), and take a positive
representative. Fix an admissible \(r\) and put
\(\theta=\chi\kappa_p^{-r}\). Some Galois element has
\(\theta(\tau)-1\in O^\times\), so \(D(\theta)^{G_\QQ}=0\).
For \(l\ge1\) the coefficient sequence for multiplication by \(p^l\)
therefore gives a unique finite-coefficient lift of each strict
class killed by \(p^l\). Set \(m=\operatorname{lcm}(f,p^l)\)
and \(L=\QQ(\zeta_m)\). Restriction of the lift to \(L\) is
injective. A cocycle \(z\) inflated from the abelian quotient
\(\operatorname{Gal}(L/\QQ)\) satisfies
\((\theta(\tau)-1)z(g)=(\theta(g)-1)z(\tau)\) modulo \(p^l\),
and is a coboundary. This proves injectivity by inflation--restriction.

Choose a generator of \(\mu_{p^l}\). Kummer theory identifies the
restricted lift with an element of
\((L^\times/L^{\times p^l})\otimes_{\ZZ/p^l}(O/p^lO)\)
of character \(\beta=\kappa_p\theta^{-1}\bmod p^l\). The factor
\(\kappa_p\) acts on the chosen root and \(\theta^{-1}\) comes from
conjugation on the argument of the restricted homomorphism.
Away from \(p\) the divisible class is a coboundary on inertia.
On the inertia of \(L\) the whole character \(\theta\) is trivial,
so the finite lift restricts to zero there. A Kummer representative
then acquires a \(p^l\)-th root over an unramified extension,
and its valuation is divisible by \(p^l\). At \(w\mid p\) the
finite homomorphism on \(G_{L_w}\), viewed inside \(D\), has the
form \((\kappa_p^{-r}-1)y\) because the divisible class is locally
zero and \(L\) kills \(\chi\). It factors through the cyclotomic
image. Lemma~\ref{lem:localunit24} gives a unit representative in
each coordinate of \(O/p^lO\), a free \(\ZZ/p^l\)-module.
Thus division of the principal ideal by its exponent defines the
equivariant map \((a)=\mathfrak a^{p^l}\mapsto[\mathfrak a]\).
Its kernel consists of global unit classes. On the odd part those
are generated by roots of unity, since \(u/\bar u\) is a root of
unity for each global unit and \(2\) is invertible. The roots have
character \(\kappa_p\), whereas \(\beta\) differs from it by a
unit at \(\tau\). As \(\beta(c)=-1\), this proves injection on the
required classes.

The integral Stickelberger element \((b-\sigma_b)\Theta_m\) kills
the class group of \(L\), as in Theorem~\ref{thm:stick}.
Its action on this eigengroup is the sum over units \(1\le a\le m\)
of \(\lfloor ba/m\rfloor\chi(a)a^{-r-1}\bmod p^l\). This argument
uses an eigengroup defined by equations and requires no division
by a group order.

Write \(f_0\) for the prime-to-\(p\) part of \(f\). On
\((\ZZ/f_0\ZZ)^\times\times\ZZ_p^\times\) use the bounded measure
which assigns \(\lfloor ba/m\rfloor-(b-1)/2\) to the unit cylinder
\(a\bmod m\), for \(m=f_0p^v\), \(v\ge\max(1,v_p(f))\).
Summing the floors over the \(p\) lifts adds
\((p-1)(b-1)/2\) to the old floor, since multiplication by \(b\)
permutes residues modulo \(p\). This gives compatibility and
the measures have \(p\)-integral values. Integrate
\[
 \chi(a)\omega(a_p)^{k-1}\langle a_p\rangle^{-s},
 \qquad \langle a_p\rangle=a_p/\omega(a_p).
\]
The integrand is odd, so the constant correction cancels in
each relevant finite sum. At \(s=r+1\) the integral reduces to
the Stickelberger sum modulo \(p^l\). This proves annihilation
on all elements of the scalar group. Write
\(\langle a_p\rangle=(1+p)^u\), \(u\in\ZZ_p\), and expand
\((1+X_s)^{-u}\) in integral binomial coefficients. Uniform
convergence for \(X_s\in p\ZZ_p\) permits termwise integration
and gives \(F\in O[[X]]\).

At \(s=1-K\) the integrand at integer representatives is
\(\chi(a)a^{K-1}\). Put \(w_a=ba-m\lfloor ba/m\rfloor\).
These representatives permute the units. Expanding \(w_a^K\)
and dividing by \(m\) gives, up to a term tending to zero,
\[
 Kb^{K-1}\sum_a\chi(a)a^{K-1}\lfloor ba/m\rfloor
 = (b^K-\chi(b)^{-1})\,m^{-1}\sum_a\chi(a)a^K+o_p(1).
\]
Here \(o_p(1)\) denotes a quantity tending to zero \(p\)-adically
as \(v\to\infty\), with \(K\) fixed. Summing the powers first on
the progressions modulo \(f\) gives the limit \(B_{K,\chi}\)
when all residues are allowed, by the Bernoulli polynomial
difference and derivative formulas in Theorem~\ref{thm:moment141}.
If \(p\nmid f\), deletion of the multiples of \(p\) subtracts
\(\chi(p)p^{K-1}B_{K,\chi}\). If \(p\mid f\), such terms already
have character value zero. This proves the formula. The Euler
factor is a unit since \(K\ge2\), and the choice of \(b\) works
for all these \(K\) because they agree modulo \(p-1\).
\end{proof}

\begin{proposition}[Higher finite differences and the full scalar exponent]
\label{prop:digits231}
In Theorem~\ref{thm:scalar231} put \(u_j=A(1-k-j(p-1))\) and
\(y=(r+k)/(p-1)\). Let \(\Delta u_j=u_{j+1}-u_j\). For every
\(h\ge1\) one has
\[
 A(r+1)\equiv\sum_{j=0}^{h-1}\binom{-y}{j}\Delta^ju_0
                          \pmod {p^hO}.
\]
If this residue is nonzero then \(p^hS_{\alpha,r}=0\).
In particular suppose \(u_0\in pO\), and denote the residues of
\(u_0/p\) and \((u_1-u_0)/p\) in \(O/pO\) by \(a_0,a_1\).
If \(a_0\notin\mathbf F_p a_1\), then
\(p^2S_{\alpha,r}=0\) for every positive \(r\equiv-k\bmod(p-1)\).
When \(O/\ZZ_p\) is unramified the sharper bound
\(pS_{\alpha,r}=0\) holds on this entire class of twists.
\end{proposition}
\begin{proof}
In the measure just constructed, \(\Delta^ju_0\) is the integral
of \(\chi(a)a_p^{k-1}(a_p^{p-1}-1)^j\), hence belongs to \(p^jO\).
Multiply \(\chi(a)a_p^{k-1}\) by the binomial expansion of
\((1+(a_p^{p-1}-1))^{-y}\). This recovers the integrand at
\(r+1\), and its terms from index \(h\) onward lie uniformly in
\(p^hO\). If the residue is nonzero, the valuation of the
annihilator is less than \(h\), so it divides \(p^h\) in \(O\).
For the last assertion the residue of \(A(r+1)/p\) is
\(a_0-ya_1\). Only elements of the prime field occur as residues
of the ordinary integer \(y\). Under the stated condition the
annihilator belongs to \(pO\) and lies outside \(p^2O\). In an
unramified valuation ring this gives valuation exactly one.
\end{proof}

\begin{corollary}[A quartic character at eleven]
\label{cor:quartic231}
Let \(\chi\) be the character modulo \(41\) with \(\chi(6)=i\).
For \(p=11\) the full scalar groups \(S_{\chi^{-1},r}\) are
killed by \(11\) for every positive \(r\equiv2\bmod10\).
\end{corollary}
\begin{proof}
Modulo \(41\), the powers \(6^{20}=-1\), \(6^8=10\) show that
\(6\) has order forty. Take \(k=8\), \(O=\ZZ_{11}[i]\) and \(b=12\).
The ring is unramified of degree two since \(-1\) is a nonsquare
modulo eleven. As \(12=6^{27}\bmod41\) its character value is
\(-i\), giving the required smoothing unit. The following four
columns are enough to reproduce the two moment reductions.
For row \(e\) and column \(j\), sum
\(\lfloor12t/4961\rfloor t^e\) modulo \(121\) over
\(1\le t\le4961=41\cdot121\), \(11\nmid t\), and
\(t\bmod41\in\{6^{4a+j}:0\le a<10\}\).
\[
\begin{array}{c|rrrr}
 e\backslash j&0&1&2&3\\ \hline
 7 &15&29&59&95\\
 17&70&117&114&84
\end{array}
\]
For an alternative direct evaluation the Bernoulli polynomial
formula gives
\(B_{8,\chi}/8=1705510378+4864413026i\).
Combining the columns with the weights \(1,i,-1,-i\) gives
\(u_0=77+55i\) and \(u_1=77+33i\) modulo \(121\).
Thus \(a_0=7+5i\), \(a_1=9i\) in \(\mathbf F_{121}\).
Every \(a_0-ya_1\) for \(y\in\mathbf F_{11}\) has real coordinate
seven. Proposition~\ref{prop:digits231} proves the assertion.
\end{proof}

\subsection{Omitted Euler factors at the nonpositive integers}\label{nav:III-43}
A long vertical theta string can put the argument of a Hecke
function at zero in a horizontal doubling integral. The finite
places removed from that integral then matter. We record the
calculation for the CM fields and angular characters in the
manuscript.

\begin{proposition}[A split Euler factor and the functional equation]\leavevmode
\label{prop:splitvalue257}
Let \(E/\QQ\) be abelian CM of degree \(2f\). At a choice of one
embedding \(\tau\) from each conjugate pair, let a unitary Hecke
character \(\rho\) have infinity factors
\((z/|z|)^{2a_\tau}\), with \(a_\tau\in\ZZ\) and some
\(a_\tau\ne0\). If \(w\) lies over a completely split rational
prime \(p\) and \(\rho\) is unramified at \(w\), put
\(z_w=\rho_w(\varpi_w)\), \(J=E(z_w)\) and \(D=[J:\QQ]\), using a
fixed inclusion \(E\subset\CC\). Then \(J\) is a number field and
\[
 |1-z_w|\ \ge\
        2^{1-D/2}p^{-[J:E]\sum_\tau|a_\tau|/2}>0 .
\]
In fact \(z_w\) has infinite order. Let \(S_E\) be a finite set of
primes over completely split rational primes, and put
\(A=|D_E|N\mathfrak f_\rho\). For \(u\ge0\) integral with
\(M_\tau=|a_\tau|>u\) for every \(\tau\), one has
\[
\begin{aligned}
 |L_E^{S_E}(-u,\rho)|
   ={}& A^{u+1/2}(2\pi)^{-f(2u+1)}
       \prod_\tau\prod_{b=-u}^{u}(M_\tau+b)\,
       |L_E^{S_E}(1+u,\bar\rho)|\\
    &{}\times
       \prod_{\substack{w\in S_E\\ \rho_w\ {\rm unramified}}}
       \frac{|1-z_w(Nw)^u|}
            {|1-\bar z_w(Nw)^{-1-u}|}>0 .
\end{aligned}
\]
The first bound applies at each omitted place when \(u=0\).
\end{proposition}
\begin{proof}
Let \(H\) be the order of the class of \(w\) in the ray class
group modulo the conductor of \(\rho\). We can write
\(w^H=(\alpha)\) with \(\alpha\equiv1\bmod\mathfrak f_\rho\).
Evaluate the character on the principal idele \(\alpha\).
At the conductor primes its residue is trivial and at the
other primes only \(w\) contributes. Thus
\[
 z_w^H=\prod_\tau
      \left(\frac{\tau(\alpha)}{\bar\tau(\alpha)}\right)^{-a_\tau}.
\]
All embeddings on the right may be viewed as automorphisms of
\(E\). The ideals \(\tau w,\bar\tau w\) as \(\tau\) varies are
pairwise distinct since the decomposition group of \(w\) is
trivial. The valuations of the right side at these two primes
are \(-Ha_\tau,Ha_\tau\). This proves algebraicity and gives
a nonzero valuation at some prime, excluding every root of
unity. At each embedding of \(E\) the numerator and denominator
in each ratio are complex conjugates, since conjugation is
central in \(\operatorname{Gal}(E/\QQ)\). All conjugates of
\(z_w\) in \(\CC\) therefore have absolute value one.

At a prime of \(J\) over \(\tau w\) the valuation of \(z_w\) is
\(-a_\tau\) times its ramification index. Over \(\bar\tau w\)
the sign is reversed. There are no other nonzero valuations.
Using the local degree formula, the denominator ideal of
\(z_w\) in \(J\) has absolute norm
\(p^{[J:E]\sum_\tau|a_\tau|}\). The denominator of \(1-z_w\)
divides this ideal. Its numerator ideal has norm at least one.
In the field norm of \(1-z_w\) the chosen embedding of \(J\)
and its complex conjugate contribute \(|1-z_w|^2\). Each of
the other \(D-2\) absolute factors is at most two. Comparing
the two estimates for the absolute norm gives the asserted bound.

The completed Hecke function has conductor factor \(A^{z/2}\)
and gamma factors \(\Gamma_{\CC}(z+M_\tau)\), where
\(\Gamma_{\CC}(z)=2(2\pi)^{-z}\Gamma(z)\). Take absolute values
in the functional equation at \(z=-u\). The root number has
modulus one. Each quotient of gamma factors contributes
\((2\pi)^{-(2u+1)}\prod_{b=-u}^u(M_\tau+b)\) by the gamma
recurrence, since \(M_\tau-u>0\). Restoring the omitted Euler
factors at \(1+u\) and removing the ones at \(-u\) gives the
second formula. The value on the positive side is nonzero by
absolute convergence for \(u>0\), and by Lemma~\ref{lem:hecke171}
for \(u=0\). For \(u>0\) the numerator of each correction
factor is at least \((Nw)^u-1\); the denominators are positive
for every \(u\ge0\). The preceding split-prime argument handles
the remaining numerator at \(u=0\).
\href{https://sites.math.rutgers.edu/~alexk/2023S572/Tate1950.pdf}
{Tate, global functional equation for Hecke characters}.
\end{proof}

\section{Norms and cyclic character fields}\label{nav:g-III-2}

\subsection{Norms of ideal classes and the exponents of strict classes}\label{nav:III-24}
\label{sec:norm191}

Fix one of the fields \(E\) of Section~\ref{sec:ray101} and an odd
prime \(p\). Write \(G=\Gal(E/\QQ)=P\times H\), where \(P\) is the
\(p\)-Sylow subgroup, and put \(K=E^H\), \(a=\log_p|P|\).
Choose an unramified coefficient ring \(O\) with fraction field \(J\)
and residue field \(\mathbf f\) containing the required roots of
unity of order prime to \(p\) (in \(O\) we use their Teichmüller
lifts). For \(\alpha\in\widehat H\), extend \(\alpha\) trivially
over \(P\). At a complementary twist \(r\ge1\) we have
\(\alpha(c)=(-1)^r\). On \(G_K\) write
\[
 \theta=\alpha^{-1}\kappa_p^{-r}
       =\tau\langle\kappa_p\rangle^{-r},\qquad
 \tau=\alpha^{-1}\omega^{-r},\qquad
 \kappa_p=\omega\langle\kappa_p\rangle .
\]
Here \(\omega\) is the Teichmüller character and
\(\langle\kappa_p\rangle\) takes values in \(1+p\ZZ_p\).
Let \(S_\alpha\) be the strict group with coefficients
\((J/O)(\theta)\) over \(K\). By Lemma~\ref{lem:des28} it is the
\(\alpha\)-summand of the original strict group on \(E\) after
extension to \(O\). This descent includes primes dividing \([E:\QQ]\).

Let \(\QQ_\infty/\QQ\) be the cyclotomic \(\ZZ_p\)-extension,
with layer \(\QQ_i\) of degree \(p^i\). Set
\[
 L=(E(\mu_p))^+,\qquad
 [K\cap\QQ_\infty:\QQ]=p^d,\qquad
 L_i=L\QQ_{d+i},\qquad L_\infty=L\QQ_\infty .
\]
The field \(K\) is totally real, \(L/K\) has degree prime to \(p\),
and the even character \(\tau\) factors through
\(\Delta=\Gal(L/K)\). We may equally use the smaller field cut out
by \(\ker\tau\) over \(K\) in place of \(L\). In either case
\(L\cap\QQ_\infty=K\cap\QQ_\infty\), by the prime-to-\(p\)
degree. Thus \(\Gal(L_\infty/K)=\Delta\times\Gamma\), where
\(\Gamma=\Gal(L_\infty/L)\). Choose the generator \(\gamma\) mapping
to \((1+p)^{p^d}\) under the principal cyclotomic character. Put
\[
       u=(1+p)^{-rp^d},\qquad h=v_p(u-1)=1+d+v_p(r),
       \qquad n_0=a-d .
\]
The valuation follows by taking the \(p\)-adic logarithm for odd \(p\).
Every place over \(p\) is totally ramified in \(L_m/L_n\) for
\(m\ge n\ge n_0\). Indeed inertia at \(p\) surjects onto
\(\Gal(\QQ_\infty/\QQ)\). Passing to a completion \(L_w\) replaces
the absolute inertia subgroup by one of index \(e(L_w/\QQ_p)\).
The index of its image in the cyclotomic group therefore divides
this ramification index and has \(p\)-part at most \(p^a\).
It follows that its image contains \(\Gamma^{p^{a-d}}\).

Write \(A_i=\Cl(\mathcal O_{L_i}[1/p])\{p\}\). These are the
ordinary \(p\)-primary ideal class groups modulo the classes of
the primes above \(p\). Define
\[
 C_i=e_\tau(A_i\otimes_{\ZZ_p}O),\qquad
 e_\tau=|\Delta|^{-1}\sum_{z\in\Delta}\tau(z)^{-1}z .
\]
We use the unramified ring \(O\) large enough for the characters
under consideration. At fixed \(E,p\), the fields \(L_i\) in the
larger choice above work for every \(\alpha,r\).

\begin{lemma}[Strict classes over the cyclotomic tower]\label{lem:direct191}
With transition maps induced by ideal norms, there is an identification
\begin{equation}
 S_\alpha=\varinjlim_{i\ge n_0}
     \operatorname{Hom}_O\!\big(C_i/(\gamma-u)C_i,J/O\big).
 \label{eq:direct191}
\end{equation}
The maps on the right are injective. The formula also holds when
\(\tau\) is trivial on a decomposition group above \(p\).
\end{lemma}
\begin{proof}
Over \(L_\infty\) the action on \(D=(J/O)(\theta)\) is trivial.
Restriction identifies \(H^1(K,D)\) with the equivariant continuous
homomorphisms from \(G_{L_\infty}\) to \(D\). For this assertion
use inflation--restriction. The quotient \(\Delta\times\Gamma\)
has vanishing \(H^2\) on \(D\), because averaging removes the
prime-to-\(p\) factor and \(\ZZ_p\) has \(p\)-cohomological dimension
one. Its \(H^1\) also vanishes. Indeed \(D^\Delta\) is zero or
all of \(D\), and multiplication by \(u-1\ne0\) is surjective on
the divisible module \(D\). This computes the procyclic \(H^1\).
At a place over \(p\) the same argument gives injectivity on
restriction in \(H^1\). Its decomposition quotient is a product
of a prime-to-\(p\) group and an open subgroup of \(\Gamma\).
The corresponding power of \(u\) still differs from one. On inertia
away from \(p\) the quotient in \(L_\infty/K\) has order prime to
\(p\), so restriction there is injective as well.

Represent a strict class by a continuous cocycle \(f\). Its image
is finite by compactness. At each of the finitely many places of
\(K\) over \(p\) it is a coboundary \((\theta(g)-1)b_w\).
Choose \(i\) large enough that \(\theta(G_{L_i})\) acts trivially
on the image of \(f\) and on all \(b_w\). Then \(f\) restricts to a
homomorphism on \(G_{L_i}\), zero on the decomposition groups
over \(p\). For conjugate places the cobounding elements differ
from translates of the chosen elements by values of \(f\), so the
same \(i\) works. Away from \(p\) the cocycle is a coboundary on
inertia and \(\theta\) is trivial on inertia over \(L_i\).
The restricted homomorphism is consequently unramified there.
Global reciprocity at the finite layer now factors it through
\(A_i\).~\cite{V-III-V-cft101}

Let \(H'_i\) be the class field of \(A_i\). It is the maximal
abelian unramified \(p\)-extension of \(L_i\) split at all places
over \(p\). The archimedean distinction has order a power of two.
For \(i\ge n_0\), total ramification gives
\(H'_i\cap L_\infty=L_i\). Thus \(G_{L_\infty}\) surjects onto
\(A_i\). Conjugation of the restricted cocycle on this group
selects the character \(\tau\) of \(\Delta\) and imposes
\(f(\gamma z)=u f(z)\). These are exactly the relations on the
right side of \eqref{eq:direct191}. Conversely, a homomorphism
there gives an equivariant homomorphism on \(G_{L_\infty}\) and
hence a unique global class. It satisfies the two local conditions
by the local injectivity just proved. Base change sends
\(H'_i L_m\) into \(H'_m\) for \(m\ge i\), and restriction of class
field automorphisms is the ideal norm. Surjectivity on \(A_i\)
follows from \(H'_i\cap L_m=L_i\). This proves the formula and
the injectivity of its maps.
\end{proof}

\begin{lemma}[Relations at two successive layers]\label{lem:norm191}
For \(m\ge n\ge n_0\) let \(Y_{m,n}=\ker(A_m\to A_n)\).
Write \(g=\gamma^{p^n}\) acting on \(A_m\), and
\(\nu_n=1+g+\cdots+g^{p-1}\). For \(m\ge n+1\),
\begin{equation}
       (g-1)A_m\ \subseteq\ Y_{m,n},\qquad
       Y_{m,n+1}=\nu_n Y_{m,n}.
 \label{eq:norm191}
\end{equation}
\end{lemma}
\begin{proof}
The field \(H'_m\) is Galois over \(L_n\) by the defining maximal
property of the class field. A decomposition group of a prime over
\(p\) maps isomorphically onto \(\Gal(L_m/L_n)\), since the prime
splits in \(H'_m/L_m\) and is totally ramified in \(L_m/L_n\).
Choose one such group and lift \(g\) to its generator \(s\).
The full group is the semidirect product of \(A_m\) by
\(\langle s\rangle\). For each prime \(w\mid p\) of \(L_n\),
choose one prime above it in \(H'_m\). The generator of its
decomposition group mapping to \(g\) has the form \(z_w s\)
with \(z_w\in A_m\), written additively. For the reference
prime \(z_w=0\).

Abelianizing imposes \((g-1)A_m=0\). Killing the decomposition
groups then imposes \(s=1\) and \(z_w=0\) for every \(w\).
The resulting extension of \(L_n\) is unramified outside \(p\)
and split at \(p\). It is exactly \(H'_n\), since \(H'_nL_m\)
is contained in \(H'_m\). Hence
\[
      Y_{m,n}=(g-1)A_m+\sum_{w\mid p}\ZZ_p z_w .
\]
There is a unique prime of \(L_{n+1}\) above each \(w\).
Its decomposition generator for the next base is
\((z_ws)^p=(\nu_n z_w)s^p\). The same presentation, now with
\(g^p-1=\nu_n(g-1)\), gives \eqref{eq:norm191}.
\end{proof}

Let \(J_n=\ker(C_{n+1}\to C_n)\). The following integer counts
generators of the relative class group after reducing both the
coefficients and the cyclotomic action. It is given by
\[
       \rho_n(\tau)=
       \dim_{\mathbf f} J_n/\big(pJ_n+(\gamma-1)J_n\big)
            \qquad(n\ge n_0).
\]

\begin{theorem}[Bounds for high elementary divisors]\label{thm:tower191}
The sequence \(\rho_n(\tau)\) is nonincreasing. At every
\(n\ge n_0\) the number of elementary divisors of \(S_\alpha\)
strictly larger than \(p^{h+n}\), measured over \(O\), is at
most \(\rho_n(\tau)\). If \(C_{n+1}\to C_n\) is an isomorphism,
then all \(C_m\to C_n\) for \(m\ge n\) are isomorphisms and
\begin{equation}
     S_\alpha\simeq\operatorname{Hom}_O
               \big(C_n/(\gamma-u)C_n,J/O\big),
     \qquad p^{h+n} S_\alpha=0 .
 \label{eq:stable191}
\end{equation}
This one equality of class-group orders applies to every twist
with the same \(\tau\).
\end{theorem}
\begin{proof}
Put \(Y^\tau_{m,n}=e_\tau(Y_{m,n}\otimes O)\), and write
\(\mathfrak m=(p,\gamma-1)\) for the indicated operators.
Surjectivity of norms and Lemma~\ref{lem:norm191} identify
\(J_n=Y^\tau_{m,n}/\nu_nY^\tau_{m,n}\) whenever \(m\ge n+1\).
Since \(\nu_n\equiv p\bmod(\gamma-1)\), it follows that
\[
       Y^\tau_{m,n}/\mathfrak mY^\tau_{m,n}
                      \simeq J_n/\mathfrak m J_n .
\]
For \(m\ge n+2\), multiplication by \(\nu_n\) surjects from
\(Y^\tau_{m,n}\) onto \(Y^\tau_{m,n+1}\). Reduction modulo
\(\mathfrak m\) proves monotonicity. If the displayed dimension is
zero, \(\gamma-1\) acts surjectively on \(Y^\tau_{m,n}/p\).
It also acts nilpotently there since \(\gamma\) has \(p\)-power
order. Thus \(Y^\tau_{m,n}/p=0\), which forces the finite group to vanish.
In particular \(J_n=0\) gives the stabilization in the statement.

Put \(B_m=C_m/(\gamma-u)C_m\). In this module
\(\gamma^{p^n}-1\) acts by \(u^{p^n}-1\), of valuation \(h+n\).
By \eqref{eq:norm191}, \(p^{h+n}B_m\) is contained in the image
\(I\) of \(Y^\tau_{m,n}\) in \(B_m\). On \(I/pI\) both \(p\)
and \(\gamma-1\) vanish. Hence \(\dim_{\mathbf f}I/pI\le\rho_n(\tau)\).
The subgroup \((p^{h+n}B_m)[p]\) injects into \(I[p]\), which
has the same dimension as \(I/pI\) by the elementary divisor
theorem over \(O\). This bounds the number of long summands of
\(B_m\), and likewise of its \(O\)-dual into \(J/O\).
The strict group is finite by the number-ring comparison of
Section~\ref{nav:III-12} and finiteness of the even \(K\)-group there.
Lemma~\ref{lem:direct191} therefore attains the whole group at
some finite \(m\), which may be increased beyond \(n\).
This proves the bound. Under stabilization all the norms are
equivariant isomorphisms, so the limit is the dual at \(n\).
Here \(\gamma^{p^n}=1\) on \(C_n\), proving the stated exponent.
\end{proof}

\subsubsection{An exact nonzero component in every matching twist}\label{nav:III-25}
Write \(\chi_D\) for the primitive quadratic character of
fundamental discriminant \(D\).

\begin{proposition}\label{prop:67data191}
For \(k=\QQ(\sqrt{67})\), \(k_i=k\QQ_i\) in the cyclotomic
\(\ZZ_3\)-extension, one has
\[
           |\Cl(\mathcal O_{k_i}[1/3])\{3\}|=3
                  \qquad(i\ge3).
\]
The nonidentity element of \(\Gal(k/\QQ)\) acts by inversion
on these groups.
\end{proposition}
\begin{proof}
We use Theorem 2 of Fukuda and Komatsu for a real quadratic field
in which \(p\) splits. Let \(\mathfrak p^{\,h_k}=(z)\), where \(h_k\)
is its ordinary class number, and let \(\varepsilon\) be a fundamental
unit. In the completion at the conjugate prime write
\(n_1=v_p(z^{p-1}-1)\), \(n_2=v_p(\varepsilon^{p-1}-1)\).
In the range \(2\le n_1<n_2\) the theorem applies if
\(p\nmid h_k\) and the \(p\)-class group of \(k(\mu_p)\) is
elementary abelian. For \(i\ge n_1+n_2-2\) it gives
\((A:D)=p^{n_1-1}\), where \(A=\Cl(k_i)\{p\}\) and \(D\)
is its subgroup represented by powers of the prime above
\(\mathfrak p\).~\cite[Thm.~2, pp.~95--96]{V-III-V-fk191}

Here \(p=3\) splits, since \(67\equiv1\bmod3\).
The ring of integers is \(\ZZ[\sqrt{67}]\), of discriminant
\(268\). Minkowski's bound \(\sqrt{268}/2<9\) shows that
prime ideals of norm at most eight generate its class group.
The prime 5 is inert. The elements \(221+27\sqrt{67}\),
\(8+\sqrt{67}\) and \(90+11\sqrt{67}\) have norms
\(-2,-3,-7\). Together with their conjugates they generate the
prime ideals needed over \(2,3,7\), so \(h_k=1\). Take
\(z=8+\sqrt{67}\). The continued fraction
\[
 \sqrt{67}=[8;\overline{5,2,1,1,7,1,1,2,5,16}]
\]
gives \(\varepsilon=48842+5967\sqrt{67}\) from the first full
period by the continued-fraction theorem for Pell's equation.
Norm \(-1\) is impossible by reduction modulo four.
In the completion where \(z\) is a unit we have
\(\sqrt{67}\equiv38\bmod81\), since this is a root with residue
\(-1\bmod3\) and the derivative is a unit. Therefore
\(z\equiv46\), \(\varepsilon\equiv26\bmod81\).
Squaring and subtracting one gives \(n_1=2\) and \(n_2=3\).

The biquadratic field \(k(\mu_3)\) has the other quadratic
subfields \(\QQ(\sqrt{-3})\) and \(\QQ(\sqrt{-201})\).
The former has class number one. The latter has discriminant
\(-804\) and class number twelve. For a direct check the leading
coefficients of its reduced positive primitive forms, with
multiplicities, are
\[
        1,\ 2,\ 3,\ 5,\ 5,\ 6,\ 7,\ 7,\ 10,\ 10,\ 14,\ 14 .
\]
Indeed in a form \((A,B,C)\) take \(A\le\sqrt{804/3}<17\),
\(|B|\le A\), \(C=(B^2+804)/(4A)\) integral and at least \(A\),
and \(\gcd(A,B,C)=1\). At \(|B|=A\) or \(A=C\) take \(B\ge0\).
These give respectively \(B=0,2,0,\pm4,6,\pm6,\pm6,\pm6\)
for the eight distinct leading coefficients above. The
classification of reduced binary quadratic forms gives the count.
On the \(3\)-class group of the biquadratic field the four character
idempotents are integral at 3. The invariant subgroup for each
quadratic involution descends to its fixed field because extension and
norm compose in the two directions to multiplication by two
and to the sum of the involution and the identity. The common
rational summand is trivial. Hence the whole \(3\)-part is the
one from \(\QQ(\sqrt{-201})\), of order three and elementary
abelian. All hypotheses of the quoted theorem now hold.

The unique prime over 3 of \(\QQ_i\) is principal. A generator
is the relative norm of \(1-\zeta_{3^{i+1}}\) from the full
cyclotomic field to \(\QQ_i\), since the residue degree is one.
Its extension to \(k_i\) is the product of the two primes above
3, each to the first power because 3 splits in \(k\).
Their classes are therefore inverse, and \(A/D\) is precisely
the \(3\)-class group with both primes divided out. The published
index formula gives \(3^{2-1}=3\) for \(i\ge3\).

Finally the corresponding group over \(\QQ_i\) is zero. Apply
the presentation in Lemma~\ref{lem:norm191} to
\(\QQ_i/\QQ\). There is just one prime over 3, so every \(z_w\)
in that presentation is zero. The bottom group is trivial and
the upper group equals its image under \(\gamma-1\), which
forces it to vanish by reduction modulo 3 and nilpotence.
The quadratic norm and extension argument applies also after
inverting 3. It kills the invariant part on \(k_i\). An
involution on a group of order three with this property acts
by inversion.
\end{proof}

\begin{theorem}[An exact component on the period fields]\label{thm:67strict191}
Take \(p=3\). If the primitive character specifying the period
field is \(\psi=\chi_{268}\), let \(r\) run over the positive
even integers. If \(\psi=\chi_{-804}\), let \(r\) run over the
positive odd integers. For every field
\(E=K_\psi C_J(\sqrt{-q},\sqrt{-d_0})\) in either construction,
the \(\psi\)-component of the complementary strict group is
cyclic of order three at every indicated twist. Its total
contribution for \(1\le r<t\) is
\(\lfloor(t-1)/2\rfloor\log3\) in the first case and
\(\lfloor t/2\rfloor\log3\) in the second.
\end{theorem}
\begin{proof}
The Galois group of \(E/\QQ\) has order a power of two, so here
\(K=\QQ\) and descent is an isomorphism on every character
component over \(\ZZ_3\). The Teichmüller character at 3 is
\(\chi_{-3}\), and \(\chi_{-804}=\chi_{268}\chi_{-3}\).
In both cases \(\tau=\psi^{-1}\omega^{-r}=\chi_{268}\).
Use its field \(k\) in Lemma~\ref{lem:direct191}.
Proposition~\ref{prop:67data191} says that for \(i\ge3\)
the entire group belongs to this character and has order three.
Every norm on this part is surjective, hence an isomorphism.
The pro-\(3\) group \(\Gamma\) acts trivially on a cyclic group
of order three, whose automorphism group has order two. Also
\(u=4^{-r}\equiv1\bmod3\). Consequently its quotient by
\(\gamma-u\) and its dual into \(\QQ_3/\ZZ_3\) both have order
three. This computes the full group in \eqref{eq:direct191}.
Counting the even or odd integers in the required range gives
the two totals.
\end{proof}

\subsection{A single relative class group for all twists at a prime}\label{nav:III-26}
\label{sec:joint193}

Retain \(K,L_i,\gamma,d,n_0\) and the coefficient ring \(O\)
from Section~\ref{sec:norm191}, using \(L=(E(\mu_p))^+\).
Fix a character \(\tau\) of \(\Delta=\Gal(L/K)\).
Consider finitely many complementary pairs \((\alpha_i,r_i)\)
with residual character \(\tau\), taking each pair once.
Their positive twists \(r_i\) are distinct. Indeed at fixed \(r\)
the restriction of \(\alpha\) to \(G_K\) determines \(\alpha\),
since \(G_K\) surjects onto \(\Gal(E/K)=H\). Put
\[
   u_i=(1+p)^{-r_i p^d},\qquad h_i=1+d+v_p(r_i),\qquad
   D_i=\sum_{\ell\ne i}(1+d+v_p(r_i-r_\ell)).
\]
An empty sum is zero. The numbers \(D_i\) measure the
denominators for interpolation at the different \(u_i\).

\begin{theorem}[One norm bound for several specializations]\label{thm:joint193}
Let \(n\ge n_0\), \(j\ge1\), and let \(c_i\) be integers with
\(c_i\ge n+h_i+D_i\). Write \(S_i\) for the full strict group
of the indicated pair, and \(C_n=e_\tau(A_n\otimes O)\).
Then
\begin{equation}
     \sum_i \log\big|(p^{c_i}S_i)[p^j]\big|
       \ \le\ \log \big|\ker(C_{n+j}\longrightarrow C_n)\big| .
       \label{eq:joint193}
\end{equation}
If the right side is less than \(j\log|\mathbf f|\), then
the same bound holds for \(\sum_i\log|p^{c_i}S_i|\).
\end{theorem}
\begin{proof}
Take \(m\ge n+j\) large enough for the direct limit in
Lemma~\ref{lem:direct191} to attain each of the finitely many
finite groups \(S_i\). Thus \(S_i\) is the \(O\)-dual into
\(J/O\) of \(B_i=C_m/(\gamma-u_i)C_m\). Write \(Y\) for
\(\ker(C_m\to C_n)\), and \(I\) for its image under the
simultaneous quotient map into \(\bigoplus_i B_i\).
For a given \(i\) and \(b_i\in B_i\), choose \(z_i\in C_m\)
mapping to \(b_i\). Apply to it the operator
\[
    (\gamma^{p^n}-1)\prod_{\ell\ne i}(\gamma-u_\ell).
\]
The result belongs to \(Y\) because \(\gamma^{p^n}\) acts
trivially on \(C_n\). Its image in every \(B_\ell\) for
\(\ell\ne i\) is zero. In \(B_i\) its image is
\((u_i^{p^n}-1)\prod_{\ell\ne i}(u_i-u_\ell)b_i\).
The first factor has valuation \(n+h_i\). The difference
\(u_i-u_\ell\) has valuation \(1+d+v_p(r_i-r_\ell)\) by the
principal-unit logarithm. The scalar is therefore \(p^{n+h_i+D_i}\)
times a unit of \(O\). Varying \(b_i\) independently gives
\[
                 \bigoplus_i p^{c_i}B_i\ \subseteq I .
\]

Set \(g=\gamma^{p^n}\) and \(v=1+g+\cdots+g^{p^j-1}\).
On \(B_i\) this operator acts by
\((u_i^{p^{n+j}}-1)/(u_i^{p^n}-1)\), of valuation \(j\).
It follows that \(\bigoplus_i(p^{c_i}B_i)[p^j]\) is contained
in the kernel of \(v\) on \(I\). For an endomorphism of a
finite group the kernel and cokernel have equal orders. Also
\(Y/vY\) maps onto \(I/vI\). Iterating Lemma~\ref{lem:norm191}
gives \(vY=\ker(C_m\to C_{n+j})\), since the \(j\) successive
geometric sums multiply to \(v\). Surjectivity of the norms
identifies \(Y/vY\) with the kernel in \eqref{eq:joint193}.
Dual finite \(O\)-modules have the same elementary divisors,
which proves the inequality for the \(S_i\).
If a group \(p^{c_i}S_i\) had exponent at least \(p^j\), its
subgroup killed by \(p^j\) would contain \(O/p^jO\). The strict
inequality in the final assertion rules this out for every \(i\),
so the displayed truncated orders already give the full orders.
\end{proof}

\begin{corollary}[Summing over the characters of the period field]
\label{cor:joint193}
On the common fields and intervals of Volume~I, put
\(\bar k_{p,r}=\max(1,k_{p,r})\), where \(k_{p,r}\) is the
depth in its equation \eqref{eq:tail114}. For every sufficiently large \(x\),
every odd prime \(p\), and every integer \(j\ge1\), one has
\begin{equation}
 \sum_{1\le r<t}\log\big|(p^{\bar k_{p,r}}M_{p,r})[p^j]\big|
    \ \le\ \log\frac{|A_{n_0+j}|}{|A_{n_0}|}
    \ \ll_\psi\ [E:\QQ]p^{j+1}\log(2Np^{a+j+1}).
 \label{eq:alljoint193}
\end{equation}
Here \(M_{p,r}\) is the complementary strict group on \(E\),
\(N\) is its prescribed cyclotomic modulus, and \(a=v_p([E:\QQ])\).
\end{corollary}
\begin{proof}
For \(1\le r_i<t\), distinct differences on either side of \(r_i\)
are distinct positive integers after taking absolute values.
Consequently
\[
 D_i\le (t-2)(1+d)
       +v_p\big((r_i-1)!(t-1-r_i)!\big)
       \le (t-2)(1+d+1/(p-1)).
\]
For \(p\le P=\exp(x^\eps/12)\), the common lower bound for
\(k_{p,r}\) in Theorem \ref{thm:deep114} is
\(\lfloor x^{-420}\allowbreak\exp(x^\eps/3)\rfloor\). It exceeds
\(n_0+h_i+D_i\) uniformly for large \(x\), since \(t\le x^{401}\)
and \(p^a\le\operatorname{ord}(\psi)\).
For \(p>P\) we have \(p>t+1\), \(p\nmid N[E:\QQ]\) and
\(K=\QQ\). Inertia at \(p\) distinguishes the residues \(-r\)
modulo \(p-1\), and hence distinguishes all the twists in this
range. Each character \(\tau\) occurs at most once, giving
\(D_i=0\), \(h_i=1\) and \(n_0=0\).

Apply Theorem~\ref{thm:joint193} with \(n=n_0\) separately for
each \(\tau\), then sum. Each summand of \(M_{p,r}\otimes O\)
occurs once. On the right the integral character idempotents
partition the relative class group after extension to \(O\).
Divide the logarithms by \([J:\QQ_p]\) to undo this extension
on both sides. Norm surjectivity gives the first inequality.

For the size estimate put \(T=L_{n_0+j}\).
It lies in the cyclotomic field of modulus \(Np^{a+j+1}\)
and has degree at most \([E:\QQ](p-1)p^{a-d+j}\).
The factor \(p^{a-d}\) is bounded by \(\operatorname{ord}(\psi)\).
The conductor--discriminant formula bounds \(\log|\operatorname{disc}T|\)
by \([T:\QQ]\log(Np^{a+j+1})\). Minkowski's theorem gives an
integral representative of each class of norm at most
\(2^{[T:\QQ]}|\operatorname{disc}T|^{1/2}\).
The number of ideals of norm at most \(B\ge1\) is at most
\(B^2\zeta_T(2)\le B^2\zeta(2)^{[T:\QQ]}\).
These estimates give the asserted upper bound even for the
ordinary class number of \(T\).
\end{proof}

At a prime where \(k_{p,r}=0\), the first layer of \(M_{p,r}\)
enters the original sum separately. Its logarithmic order is
exactly \(\log|M_{p,r}|-\log|pM_{p,r}|\). The estimate above
controls all the twists jointly beyond the stated depths. The
elementary class-number bound grows with \(j\), leaving a
uniform bound for the full surviving orders to be established.

\subsection{Real cyclic quartic and sextic fields}\label{nav:III-27}
\label{sec:units195}

Let \(F/\QQ\) be totally real and cyclic of degree \(m=4\) or
\(6\), with conductor \(f\). Write \(k_j\) for its subfield of
degree \(j\), where \(j=2\) in both cases and \(j=3\) also occurs
when \(m=6\). Write \(f_j\) for the corresponding conductor.
Let \(h_L\) denote the ordinary ideal class number of a field \(L\)
and \(R_L\) its regulator, with one row deleted in the logarithmic unit
matrix. Use the following two quotients.
\[
 (H,R_*)=
 \begin{cases}
 (h_F/h_{k_2},\ R_F/R_{k_2})&m=4,\\
 (h_F/(h_{k_2}h_{k_3}),\ R_F/(R_{k_2}R_{k_3}))&m=6.
 \end{cases}
\]
We will bound \(H\), then use the bound to exclude entire
primary components from the strict groups in Section~\ref{sec:ray101}.

\begin{lemma}[The two faithful characters]\label{lem:relative195}
The number \(H\) is the order of the joint kernel of the ideal class
norms to the maximal proper subfields of \(F\). For a faithful
primitive Dirichlet character \(\chi\) of \(F\),
\[
              H R_* = \frac f4 |L(1,\chi)|^2,
      \qquad H \le \frac{f(\log f)^2}{16 R_*}.
\]
If \(p\nmid m\) and a faithful constituent occurs in \(\Cl(F)/p\Cl(F)\),
then \(p^{\operatorname{ord}_m(p)}\le H\).
\end{lemma}
\begin{proof}
There are two faithful characters, \(\chi\) and its complex conjugate.
Both have conductor \(f\), since both cut out \(F\). The remaining
characters come from the indicated proper subfields. The
conductor--discriminant formula gives \(D_F=f^2 f_2\) for \(m=4\)
and \(D_F=f^2 f_2 f_3^2\) for \(m=6\), with \(D_{k_2}=f_2\) and
\(D_{k_3}=f_3^2\). Divide the factorization of \(\zeta_F\) by
\(\zeta_{k_2}\) in the quartic case. In the sextic case take
\(\zeta_F\zeta_{\QQ}/(\zeta_{k_2}\zeta_{k_3})\). The result in
either case is \(L(s,\chi)L(s,\bar\chi)\). At \(s=1\), the
residue of the zeta function of a totally real field of degree
\(n\) is \(2^{n-1}hR/\sqrt D\). The displayed identity follows.
Ramar\'e's inequality \(|L(1,\chi)|\le(\log f)/2\) applies since
\(\chi\) is even, primitive and nonprincipal.~\cite[Cor.~1]{V-III-V-ramare105}

Each extension \(F/k_j\) has prime degree and ramifies at a finite
place. Indeed \(D_F/D_{k_j}^{[F:k_j]}>1\) by the discriminants
above, since \(f_j\mid f\) and each proper subfield listed here has
conductor greater than one. Thus \(F\) meets the ordinary Hilbert
class field of \(k_j\) in \(k_j\). Its compositum with \(F\) is
unramified abelian over \(F\) and split at real places. Restriction
onto its Galois group from the Hilbert class field of \(F\) is
surjective. By functoriality of the Artin map this is surjectivity
of the class norm to \(k_j\).~\cite{V-III-V-cft101}

For \(m=6\) the joint norm is also onto. An extended ideal class
from \(k_2\) has norms \((u^3,1)\), and one from \(k_3\) has
norms \((1,v^2)\). The cross norm is principal because the
conjugates from one subfield run through its full Galois group
over \(\QQ\), giving the extension of its rational ideal norm.
This identity can be checked on ideals after extension to \(F\),
where extension of fractional ideals is injective. On each primary
part either cubing or squaring is invertible. The image therefore
contains one direct factor, and its projection onto the other is
onto. This proves the assertion about \(H\).

Finally the action on each norm target factors through a proper
quotient of \(\Gal(F/\QQ)\). On the \(p\)-parts, the projector to
the faithful characters is defined integrally over \(\ZZ_p\) by
summing their character idempotents in a splitting ring. The sum
is invariant under Galois conjugation and \(m\) is a unit.
Its image on \(\Cl(F)\{p\}\) belongs to the joint kernel and
reduces onto the faithful part modulo \(p\). Each faithful
irreducible over \(\mathbf F_p\) has dimension
\(\operatorname{ord}_m(p)\), the length of the Frobenius orbit of a
primitive \(m\)-th root of unity. This proves the lower bound.
\end{proof}

\begin{lemma}[Relative units and discriminants]\label{lem:unit195}
There is a unit \(z\) generating \(F\) as a field with the following
properties. Let \(\sigma\) generate \(\Gal(F/\QQ)\), and set
\(s=\log|z|\) and \(t=\log|\sigma z|\) at one fixed real embedding.
Then \(R_*\ge Q\), where
\[
 Q=\begin{cases}s^2+t^2&m=4,\\s^2-st+t^2&m=6.\end{cases}
\]
In the quartic case,
\begin{equation}
       \log(f\sqrt{f_2})\le\log16+\sqrt{10}\sqrt Q .
       \label{eq:quart195}
\end{equation}
In the sextic case,
\begin{align}
 \log(f\sqrt{f_2} f_3)&\le\log216+\sqrt{208/3}\sqrt Q,\notag\\
 \log(f f_3/f_2)&\le\log4+4\sqrt Q,\notag\\
 \log(f\sqrt{f_2}/f_3)&\le\log8+(4/\sqrt3)\sqrt Q .
       \label{eq:sext195}
\end{align}
\end{lemma}
\begin{proof}
Put \(U_L=\OO_L^\times/\{\pm1\}\) and let \(W\) be the joint
kernel of the norms from \(U_F\) to the indicated \(U_{k_j}\).
Dirichlet's unit theorem realizes \(U_F\) as a full lattice
in the real hyperplane of coordinate sum zero. The embedded
subfield lattices map by the norms onto subgroups of finite
index in the product of the subfield lattices. On their own
factor these maps are multiplication by \([F:k_j]\) and the
cross norms are zero. Thus \(W\) has rank two. Its logarithmic
vectors, in cyclic order, have the form
\[
  (s,t,-s,-t)\quad\text{or}\quad(s,t,t-s,-s,-t,s-t),
\]
respectively. For degree six the norm to \(k_3\) makes opposite
coordinates negatives of each other, while the norm to \(k_2\)
makes the sum of the even coordinates zero. Thus \(\sigma\)
on \(W\) satisfies \(\Phi_m(\sigma)=0\). The irreducible polynomial
\(\Phi_m\) has degree two, so \(W\otimes\QQ\) has dimension one
over \(\QQ(\zeta_m)\). Hence \(W\) is a fractional ideal up to
scaling in \(\ZZ[\zeta_m]\). This ring is Euclidean in both cases
(the Gaussian or Eisenstein integers). Choose a generator and
represent it by \(z\in\OO_F^\times\). Then \(z,\sigma z\) give a
basis of \(W\). If \(z\) belonged to a maximal proper subfield,
its norm there would make a positive power of \(z\) equal to
\(\pm1\), contradicting its nonzero logarithm. Hence \(\QQ(z)=F\).

The logarithmic covolume of \(U_L\) is \(\sqrt{[L:\QQ]}R_L\)
for a totally real field. Indeed all its signed maximal minors
agree up to sign and the square of the covolume is their sum of
squares. The subfield spaces and the relative space are mutually
orthogonal by summing logarithms over the norm fibers. In degree
four the embedded quadratic lattice has length \(2R_{k_2}\).
The two basis vectors of \(W\) are orthogonal of squared length
\(2Q\). The index of the sum of the subfield and relative lattices
is at most two, by the norm modulo \(2U_{k_2}\). Its kernel modulo
the subfield lattice is exactly \(W\), since subtracting a
subfield unit with that norm leaves norm zero. It follows that
\(R_*\ge (2R_{k_2})(2Q)/(2\sqrt4 R_{k_2})=Q\).

In degree six the embedded subfield covolumes are
\(\sqrt6R_{k_2}\) and \(2\sqrt3 R_{k_3}\), since their logarithmic
coordinates are repeated three and two times respectively. The
two relative basis vectors have squared length \(4Q\) and inner
product \(2Q\), so their area is \(2\sqrt3 Q\). The quotient by
all three lattices injects by the norms into
\(U_{k_2}/3U_{k_2}\times U_{k_3}/2U_{k_3}\). Indeed if the two
norms are divisible by three and two, subtracting the corresponding
subfield units leaves an element of \(W\). The index is at most
\(3\cdot2^2=12\). Multiplying the three covolumes and dividing
by \(12\sqrt6R_{k_2}R_{k_3}\) yields \(R_*\ge Q\).

For the discriminants we record a determinant estimate. List the
conjugates \(x_i\) of \(z\) with logarithms of their absolute
values \(\lambda_1\le\cdots\le\lambda_n\). Put
\(S=\sum_i(i-1)\lambda_i\). Then
\[
        \prod_{i<j}|x_i-x_j|\le n^{n/2}e^S .
\]
To see this, put \(b_1=0\), \(b_{j+1}-b_j=\lambda_j\) and
\(a_i=(i-1)\lambda_i-b_i\). The exponent
\((j-1)\lambda_i-b_j\) is maximized at \(j=i\), as its successive
differences are \(\lambda_i-\lambda_j\). Multiplying the rows and
columns of the Vandermonde matrix by \(e^{-a_i}\) and \(e^{-b_j}\)
gives rows of norm at most \(\sqrt n\). Hadamard's inequality
gives the claim since \(\sum_i(a_i+b_i)=S\).
The square of the Vandermonde product is the discriminant of
\(\ZZ[z]\), a multiple of \(D_F\) by the index formula. In degree four, if \(A\ge B\) are
\(|s|,|t|\) in decreasing order, \(S=3A+B\le\sqrt{10Q}\).
This gives \eqref{eq:quart195}.

In degree six the three nonnegative logarithm magnitudes, one
for each opposite pair, have the form \(A+B,A,B\) with
\(A\ge B\ge0\), since they are \(|s|,|t|,|s-t|\). Here
\(Q=A^2+AB+B^2\). Sorting all six gives \(S=8A+6B\), and
\((208/3)Q-S^2=(4/3)(2A-5B)^2\ge0\). This proves the first
line of \eqref{eq:sext195}.

For an order \(\OO_{k_j}[z]\), the relative discriminant is
a multiple of \(\mathfrak d_{F/k_j}\), by the index formula
locally at every prime. Thus the product of absolute
Vandermonde determinants over the embeddings of \(k_j\) is
at least \((D_F/D_{k_j}^{[F:k_j]})^{1/2}\). Over each embedding
of \(k_2\) the three conjugates have absolute product one.
Their logarithms have squared length \(2Q\), hence their range
is at most \(2\sqrt Q\) by Cauchy--Schwarz. The three-variable
inequality in the proof of Lemma~\ref{lem:cubic105} bounds each
Vandermonde by twice the exponential of that range. Over \(k_3\)
each pair has opposite logarithms \(\lambda,-\lambda\), so its
difference is at most \(2e^{|\lambda|}\). The sum of the three
magnitudes is \(2(A+B)\le4\sqrt{Q/3}\), since
\(4Q-3(A+B)^2=(A-B)^2\). These two relative estimates give the
remaining lines of \eqref{eq:sext195}.
\end{proof}

\begin{proposition}[An upper bound for the joint norm kernel]
\label{prop:conduct195}
For arbitrary conductor, \(H<f\) if \(m=4\) and \(\log f\ge16\),
or if \(m=6\) and \(\log f\ge90\). For prime conductor \(f=p\),
\(H<p\) if \(m=4\) and \(p\ge64\), or if \(m=6\) and \(p\ge1024\).
\end{proposition}
\begin{proof}
It suffices to show \(\sqrt Q>(\log f)/4\), by the preceding
lemmas. For a quartic field, \eqref{eq:quart195} gives this
when \(\log f\ge16\), since \(\sqrt{10}/4<4/5\) and \(\log16<3\).
If \(f=p\) then \(f_2=p\). For \(p\ge64\) we have
\((3/2)\log p-\log16\ge (5/6)\log p>(\sqrt{10}/4)\log p\).
In degree six, add the three inequalities in \eqref{eq:sext195}
with weights \(1,2,3\). Both proper conductors cancel. This gives
\[
 6\log f\le \log(216\cdot4^2\cdot8^3)
               +(\sqrt{208/3}+8+4\sqrt3)\sqrt Q .
\]
The constant is less than \(15\) and the coefficient is less
than \(70/3\). If \(\log f\ge90\), then \(15\le(\log f)/6\).
The desired strict inequality follows.

Suppose now \(m=6\) and \(f=p\) is prime. The unique prime
above \(p\) is totally ramified with residue field \(\mathbf F_p\).
All conjugates of \(z\) therefore have the same residue \(a\).
Its norms to \(k_2,k_3\) are signs, so \(a^6=a^4=1\). Change
the sign of \(z\) to arrange \(a=1\). Both norms are then
\(+1\), as \(p\) is odd. Each pair over \(k_3\) consists of
actual reciprocals. The conjugates with indices \(0,2,4\) are
representatives of the three reciprocal pairs and their product
is one. At most two representatives are negative.

The nonzero integer \(|N_{F/\QQ}(1-z)|\) is divisible by \(p\).
For a reciprocal pair \(x,x^{-1}\), put \(L=|\log|x||\).
Its contribution to this integer is \(e^L(1-e^{-L})^2\) for
\(x>0\), and \(e^L(1+e^{-L})^2\) for \(x<0\). Multiplication
over the three pairs gives
\[
    p\le16e^{2(A+B)}\le16e^{4\sqrt{Q/3}}.
\]
If \(p\ge2^{10}\), then \(\log p-\log16\ge(3/5)\log p\)
and \(3\sqrt3/5>1\). This gives \(\sqrt Q>(\log p)/4\).
\end{proof}

\begin{theorem}[Prime conductor at every layer]\label{thm:prime195}
Let \(p\) be an odd prime and let \(F\) be any subfield of
\(\QQ(\zeta_p)^+\) of degree \(m\in\{1,2,3,4,6\}\). Let
\(\QQ_n\) be the field of degree \(p^n\) in
\(\QQ(\zeta_{p^{n+1}})\). Then \(p\nmid h(F\QQ_n)\) for every
\(n\ge0\).
\end{theorem}
\begin{proof}
Begin at \(n=0\). The case \(m=1\) is immediate. For the
quadratic field its discriminant is \(p\) and each ideal class
has an integral representative of norm at most \(Y=\sqrt p/2\)
by Minkowski's theorem. At most \(\sum_{ab\le Y}1\le Y^2<p\)
such ideals exist, because at a rational prime power \(q^j\)
there are at most \(j+1\) choices. In the cubic case the claim
follows from Lemma~\ref{lem:cubic105} with \(f=p\). In degrees
four and six the proper subfields therefore have class numbers
prime to \(p\). Proposition~\ref{prop:conduct195} proves the
claim above its prime-conductor thresholds.

Here is an exact check below the thresholds. Such a real
subfield exists precisely when \(2m\mid p-1\).
Let \(L=\QQ(\zeta_p)\) and let \(\omega\) be its mod-\(p\)
cyclotomic character. The two faithful characters of \(F\) on
\(\mathbf F_p\) lift to \(\omega^k\) and \(\omega^{p-1-k}\) on
\(\Gal(L/\QQ)\), where \(k=(p-1)/m\). If \(\Cl(F)/p\Cl(F)\)
were nonzero, a faithful character would occur. Indeed every
nonfaithful summand is invariant under the subgroup defining a
proper subfield. Its group sum is extension of the ideal norm,
and acts on invariants by the subgroup order, a unit modulo \(p\).
The preceding proper-subfield conclusion rules this out.
The same norm and extension identities show that a nonzero
faithful summand would persist in \(\Cl(L)/p\Cl(L)\).

For either exponent \(u\in\{k,p-1-k\}\), suppose such a summand
exists. The group \(\Gal(L/\QQ)\) has order \(p-1\), so there is a
nonzero equivariant homomorphism from \(\Cl(L)\) to
\(\mathbf F_p(\omega^u)\). It gives an unramified homomorphism
on \(G_L\) by class field theory. Its restriction at the prime
above \(p\) is zero since that prime is generated by \(1-\zeta_p\).
After identifying its values with \(\mu_p\), Kummer theory gives
a nonzero class of character \(\beta=\omega^{1-u}\) in
\(L^\times/L^{\times p}\). Conjugation on values contributes
\(\omega\) and precomposition with inverse conjugation contributes
\(\omega^{-u}\). Its valuations are multiples of \(p\).
Indeed the class is trivial at \(p\), and at each other finite
place it becomes trivial over the maximal unramified extension,
which preserves the value group. Thus, writing \((y)=\mathfrak a^p\),
it maps to \([\mathfrak a]\in\Cl(L)[p]\). The kernel consists of
unit classes. Their odd part is generated by roots of unity. Indeed,
for any unit \(v\), the quotient \(v/\bar v\) has absolute value
one under every embedding and is a root of unity. Since \(p\) is
odd, the operator \(1-c\) spans the odd part, where \(c\) is
complex conjugation. The character on its roots of unity modulo
\(p\)-th powers is \(\omega\). But \(u\) is even and
\(u\not\equiv0\pmod{p-1}\), so \(\beta\) is odd and differs from
\(\omega\). Our image in \(\Cl(L)[p]\) is therefore nonzero.

For \(p\nmid b\), the integral Stickelberger element
\((b-\sigma_b)\Theta_p\) annihilates this image. Its action
on the \(\beta\)-component is multiplication by
\begin{equation}
 T_b(u)=\sum_{j=1}^{p-1}\left\lfloor\frac{bj}{p}\right\rfloor
                        j^{u-1}\pmod p ,
 \label{eq:res195}
\end{equation}
since the coefficient of \(\sigma_j^{-1}\) in this element is
\(\lfloor bj/p\rfloor\) and \(\beta(j)^{-1}=j^{u-1}\).
This is also the group-ring computation in
Theorem~\ref{thm:stick}.~\cite[Ch.~6]{V-III-V-stick}
It suffices to exhibit a nonzero value for each \(u\).

For degree four the two rows \((p,b,T_b(k),T_b(p-1-k))\) are
\((17,2,13,7)\) and \((41,2,5,29)\). For degree six the table
gives the two residues in that order. A superscript star uses
\(b=3\), and every other entry uses \(b=2\).
\begin{center}
\small
\begin{tabular}{r rr@{\hspace{2em}}r rr@{\hspace{2em}}r rr}
\hline
\(p\)&\(T_b(k)\)&\(T_b(p-1-k)\)&
\(p\)&\(T_b(k)\)&\(T_b(p-1-k)\)&
\(p\)&\(T_b(k)\)&\(T_b(p-1-k)\)\\
\hline
13&5&10&337&12&128&673&545&222\\
37&24&36&349&169&240&709&449&512\\
61&34&27&373&190&214&733&495&393\\
73&40&63&397&184&275&757&539&455\\
97&94&56&409&367&200&769&465&734\\
109&10&90&421&192&301&829&109&85\\
157&86&18&\(433^*\)&82&387&853&779&180\\
181&133&104&\(457^*\)&367&312&877&449&270\\
193&80&54&541&184&274&937&467&284\\
229&78&108&577&323&364&997&587&991\\
241&134&196&\(601^*\)&242&65&1009&610&546\\
277&122&26&613&229&139&1021&388&291\\
313&155&92&661&615&221&&&\\
\hline
\end{tabular}
\end{center}
For reproducibility, enumerate \(p=1+2ma\) below the relevant
threshold, starting with \(a=1\). Discard candidates divisible
by a prime at most \(\sqrt p\). To verify each table entry using
only integer arithmetic, initialize \(T=0\) and for
\(j=1,\ldots,p-1\) compute \(j^{u-1}\bmod p\) by repeated
squaring. Add its product with the integer quotient of \(bj\)
by \(p\) to \(T\), reducing modulo \(p\) at each step. For
example at \(m=6,p=13,k=2,b=2\) the first residue is
\(7+8+9+10+11+12\equiv5\pmod{13}\).
All residues are nonzero, proving the claim at the base.

It remains to pass up the tower. At each step
\(L_1=F\QQ_{n+1}\) over \(L_0=F\QQ_n\), the extension is
cyclic of degree \(p\). Only the unique prime above \(p\)
ramifies, and it is totally ramified. These facts follow by
restriction from the prime-power cyclotomic field, since
\([F:\QQ]\) is prime to \(p\). Suppose \(p\nmid h_{L_0}\)
and let \(B\) be the maximal unramified elementary abelian
\(p\)-extension of \(L_1\) split at its real places.
It is Galois over \(L_0\) by uniqueness, and
\(G=\Gal(B/L_0)\) is a finite \(p\)-group.
The inertia subgroup above \(p\) has order \(p\), because
\(B/L_1\) is unramified. Its normal closure is all of \(G\), since
the quotient defines an everywhere unramified extension of
\(L_0\), split at the real places, and a nontrivial finite
\(p\)-group quotient would have a cyclic quotient of order \(p\).
Class field theory excludes this by \(p\nmid h_{L_0}\).
Every maximal subgroup of a finite \(p\)-group is normal.
The inertia group therefore cannot be proper, since it would
be contained in a maximal subgroup. Thus \(|G|=p\), giving
\(B=L_1\) and \(p\nmid h_{L_1}\). Induction proves the theorem.
\end{proof}

\begin{corollary}[Strict groups of residual order four or six]
\label{cor:strict195}
In Theorem~\ref{thm:ray101}, let the even residual character
\(\tau\) cut out a field \(F\) of degree \(m=4\) or \(6\).
Suppose the coefficient prime \(p\) does not split completely
in \(F\). If \(p^{\operatorname{ord}_m(p)}>H\), then
\(S(\vartheta)=0\). Vanishing holds whenever \(f=p\). For
arbitrary conductor it also holds if \(p^{\operatorname{ord}_m(p)}\ge f\) and \(\log f\ge16\)
for \(m=4\) or \(\log f\ge90\) for \(m=6\).
\end{corollary}
\begin{proof}
The order \(m\) of the residual character is prime to \(p\).
Since it is faithful on \(\Gal(F/\QQ)\), the nonsplitting
assumption makes it nontrivial on a decomposition group at
\(p\). By Theorem~\ref{thm:ray101}, the elements of
\(S(\vartheta)\) killed by \(p\) are dual to the
\(\tau\)-part of the ordinary class group modulo \(p\) after
extension of residue scalars. Indeed the classes of the primes
above \(p\) have zero image on this character, since their
stabilizers act nontrivially on it. A nonzero such part requires
a faithful irreducible constituent over \(\mathbf F_p\).
Lemma~\ref{lem:relative195} and Proposition~\ref{prop:conduct195}
give the stated tests for arbitrary conductor.
Theorem~\ref{thm:prime195} treats every prime conductor equal to
the coefficient prime. Finally a nonzero \(p\)-primary torsion
group contains an element of order \(p\), so the first-layer
vanishing gives vanishing at every exponent.
\end{proof}

\subsection{Norms of ideal classes in the remaining cohomology estimate}\label{nav:III-31}
Fix \(x\) in the same interval as in \eqref{eq:tail114} and keep
the field \(E=K_\psi C_J(\sqrt{-q},\sqrt{-d_0})\) used for the
theta comparison. For an odd prime \(p\), let \(K\) be the
maximal subfield of \(E\) of \(p\)-power degree, \(p^a=[K:\QQ]\).
Let \(\QQ_i\) be the degree-\(p^i\) layer of the cyclotomic
\(\ZZ_p\)-extension \(\QQ_\infty/\QQ\), and set
\[
 L=(E(\mu_p))^+,\quad p^d=[K\cap\QQ_\infty:\QQ],\quad
 L_i=L\QQ_{d+i},\quad A_i=\operatorname{Cl}(\mathcal O_{L_i}[1/p])\{p\}.
\]
For a character \(\alpha\) of the prime-to-\(p\) part of
\(\operatorname{Gal}(E/\QQ)\), extended trivially over the
\(p\)-part, consider its complementary twists
\(\alpha(c)=(-1)^r\). Write \(\tau=\alpha^{-1}\omega^{-r}\)
on \(G_K\), with \(\omega\) the Teichmüller character. With
coefficients in a splitting unramified ring \(O\) of residue
field \(\mathbf f\), put \(C_i=e_\tau(A_i\otimes O)\), and
let \(S_\alpha=e_\alpha(M_{p,r}\otimes O)\). One tower of
fields works for all these characters and twists at this \(p\).

\begin{theorem}[Class-group growth forced by a surviving strict class]
\label{thm:norm192}
Write \(k=k_{p,r}\), \(h=1+d+v_p(r)\), \(n_0=a-d\), and
\(D=\dim_{\mathbf f}(p^kS_\alpha)[p]\). When \(k\ge h+n_0\),
for every \(n_0\le n\le k-h\) one has
\[
 \dim_{\mathbf f}
 \frac{\ker(C_{n+1}\to C_n)}
      {(p,\gamma-1)\ker(C_{n+1}\to C_n)}\ \ge D ,
\]
where \(\gamma\) generates \(\operatorname{Gal}(L\QQ_\infty/L)\)
and the maps are norms. In particular
\begin{equation}
 \log\frac{|C_{k-h+1}|}{|C_{n_0}|}
           \ \ge\ (k-h-n_0+1)D\log|\mathbf f| .
 \label{eq:norm192}
\end{equation}
If one of these transitions has equal class-group orders, this
entire character contributes zero beyond depth \(k\) in
\eqref{eq:tail114}. The assertion applies at primes dividing
the degree of the period field as well.
\end{theorem}
\begin{proof}
Restriction to \(E/K\) and averaging over its prime-to-\(p\)
Galois group identify \(S_\alpha\) with the strict cohomology
of \(\alpha^{-1}\kappa_p^{-r}\) over \(K\). Over \(L\QQ_\infty\)
this character becomes trivial. Its strict classes are dual to
the quotients of the \(C_m\) by
\(\gamma-(1+p)^{-rp^d}\) as \(m\) grows. Local restriction in
this descent is injective even at split primes above \(p\)
because the principal cyclotomic image is open and multiplication by a
nonzero power of its generator minus one is surjective on the
divisible coefficient group. The complete argument is in
Lemma \ref{lem:direct191}.

For \(m>n\ge n_0\), total ramification implies that the kernel
of \(C_m\to C_{n+1}\) is the image of the kernel of \(C_m\to C_n\)
under \(1+\gamma^{p^n}+\cdots+\gamma^{(p-1)p^n}\).
This operator belongs to \((p,\gamma-1)\). In the specialization
for the twist, \(\gamma^{p^n}-1\) is a unit times \(p^{h+n}\).
Consequently the rank on the left bounds the number of elementary
divisors of \(S_\alpha\) with exponent greater than \(h+n\), as
proved in Lemmas \ref{lem:direct191}--\ref{lem:norm191} and Theorem \ref{thm:tower191}. At \(h+n\le k\) this includes
all \(D\) summands of exponent greater than \(k\). The norms are
surjective, so summing the logarithms of their kernel orders gives
\eqref{eq:norm192}. Equal orders at any indicated step force
\(D=0\), and then \(p^k S_\alpha=0\) by finiteness.
This also kills its contribution after the projector used in
\eqref{eq:tail114}.
\end{proof}

For \(p\le b^{1/12}\), the lower bound on \(k_{p,r}\) in
Theorem~\ref{thm:deep114} tends to infinity uniformly faster
than \(\log(2x)\). Here \(a\le v_p(\operatorname{ord}\psi)\)
and \(v_p(r)\le\log_p t\). Hence every surviving component in
this range forces growth at every one of the indicated consecutive
layers. At the larger primes the same test applies whenever
\(k\ge h+n_0\).

\begin{corollary}[Exact lengths for two primitive characters]
\label{cor:67strict192}
On every period field for the primitive quadratic character of
discriminant \(268\), its own complementary component at \(3\)
has order three for every positive even twist. For the character
of discriminant \(-804\) the corresponding statement holds for
every positive odd twist.
\end{corollary}
\begin{proof}
In both cases \(\alpha=\psi\) and
\(\psi^{-1}\omega^{-r}=\chi_{268}\). On the cyclotomic
\(\ZZ_3\)-tower of \(\QQ(\sqrt{67})\) the class group modulo
the primes above 3 has \(3\)-part of order three from layer
three onward. Fukuda and Komatsu's quadratic-field theorem
gives this index. The verification of its numerical hypotheses,
the character action and descent to the original period field is in
Proposition \ref{prop:67data191} and Theorem \ref{thm:67strict191}. The pro-\(3\) action on the limiting cyclic group
and the twisting character modulo three are both trivial. Its
twisted dual therefore still has order three. In particular the
contribution of this component to the original sum is exactly
\(\log 3\) at each indicated twist.
\end{proof}

To obtain the cohomological bound in Section \ref{nav:I-28} one still needs
an upper estimate for the total orders of the surviving components
in \eqref{eq:tail114}. Equation~\eqref{eq:norm192} gives a
necessary growth condition for them. The required lower arithmetic
comparison for the integral theta period at its original finite
level also remains open in this manuscript.

\subsection{A bound for all twists from two fields at each prime}\label{nav:III-32}
Use the cyclotomic tower and norm groups of
Theorem~\ref{thm:norm192}. The residual character determines
which twists must be treated together. Polynomial interpolation
in the cyclotomic generator gives the following simultaneous
bound.

\begin{theorem}[A simultaneous estimate for the remaining classes]
\label{thm:joint194}
Put \(\bar k_{p,r}=\max(1,k_{p,r})\) and \(n_0=a-d\).
For all sufficiently large \(x\) on the same intervals,
for each odd prime \(p\), and for every \(j\ge1\), one has
\begin{equation}
  \sum_{1\le r<t}\log\big|(p^{\bar k_{p,r}}M_{p,r})[p^j]\big|
     \ \le\ \log|\ker(A_{n_0+j}\longrightarrow A_{n_0})| .
     \label{eq:joint194}
\end{equation}
The groups on the right are \(p\)-parts of ideal class groups
with all primes above \(p\) divided out. In particular the
number of cyclic summands surviving these depths, summed over
every twist, is at most the base-\(p\) logarithm of the first
norm kernel (\(j=1\)). If the right side for any \(j\) is
strictly less than \(j\log p\), it bounds the sum of the
full orders in logarithmic scale at this prime.
\end{theorem}
\begin{proof}
For a fixed residual character write \(u_r=(1+p)^{-rp^d}\).
At a common finite level \(m\) all the strict groups in this
character are dual to the quotients by \(\gamma-u_r\) of the
same class-group component. Multiplying a lift of a class in
the \(r\)-th quotient by
\((\gamma^{p^{n_0}}-1)\prod_{s\ne r}(\gamma-u_s)\) kills its
other coordinates and puts it in the norm kernel down to
\(n_0\). The product ranges over the other twists in that
residual character. Its loss, including the first factor,
is at most \(\bar k_{p,r}\). For \(p\le\exp(x^\eps/12)\) this
follows from the factorial bound for the valuations of the
distinct \(r-s\) and Theorem~\ref{thm:deep114}. For larger
\(p\) the twists are distinct on inertia and the loss is one.

In these independent coordinates the sum of the first \(p^j\)
powers of \(\gamma^{p^{n_0}}\) acts by a scalar of valuation
\(j\) on each coordinate. Its kernel therefore contains the
pieces killed by \(p^j\) in \eqref{eq:joint194}. Its cokernel
on the joint image is a quotient of the norm kernel from
\(n_0+j\) to \(n_0\). Kernel and cokernel have equal sizes on
the finite image. Summing over the disjoint residual characters
and descending the unramified coefficients gives the inequality.
The joint image, the exact norm relation and the full valuation
bound are proved in Theorem \ref{thm:joint193} and Corollary \ref{cor:joint193}. Finally an elementary divisor
of a surviving group of order at least \(p^j\) would contribute
at least \(j\log p\) to the left side. This proves the last assertion.
\end{proof}

When \(k_{p,r}=0\), the term of \eqref{eq:tail114} also uses
the first layer of \(Q_{p,r}\). Its exact contribution at
that layer is \(\log|Q_{p,r}[p]|\). At every other depth
\eqref{eq:joint194} applies directly to the tails in
\eqref{eq:tail114} because \(Q_{p,r}\) is a summand of
\(M_{p,r}\). The estimate required for Dirichlet GRH concerns
full orders summed over the coefficient primes. The explicit
class-number bound in Corollary \ref{cor:joint193} increases with
\(j\) and does not supply that estimate.

In the accompanying source, \srcloc{Development/AutomorphicArithmetic/AutomorphicVarieties/HeightAndVanishing/Counts}{53}
states one of the bounds needed for full orders. For its first
case it sums over characters and over any finite set of allowed
coefficient primes, at fixed twist \(r<t\), with upper bound
\((Nt)^C\). The proof in
\srcloc{Development/AutomorphicArithmetic/AutomorphicVarieties/ModularityArgument/ThresholdArguments/DenseThresholdCases/LinkedUnlE}{49}
uses a cohomology estimate, rigidity and automorphy inputs.
Equation~\eqref{eq:joint194} proves a finite-layer estimate by
class field theory. A uniform bound for the full global sum and
the lower arithmetic comparison for the designated integral
theta period remain the two outstanding steps of Section \ref{nav:I-28}.

\subsection{Vanishing at quartic and sextic residual characters}\label{nav:III-33}
We first consider the classes descending from \(\QQ\). In the
rational summand of the cohomological estimate the same twist \(r\)
must be treated at many coefficient primes \(p\). Two additional
families can be removed at every exponent.

\begin{theorem}[The rational summand for characters of orders four and six]
\label{thm:units196}
For every positive even integer \(r\) and every odd prime \(p\),
\[
         p-1\mid4r\quad\text{or}\quad p-1\mid6r
        \qquad\Longrightarrow\qquad K_{2r}(\ZZ)\{p\}=0 .
\]
Consequently the supremum in \eqref{eq:qbound102} is unchanged
when its finite prime sets are restricted to primes \(p\ge5\)
satisfying both \(p-1\nmid4r\) and \(p-1\nmid6r\).
\end{theorem}
\begin{proof}
Put \(D=\QQ(\zeta_p)^+\). By Proposition~\ref{prop:realclass100},
the dimension of \(K_{2r}(\ZZ)/p\) is the dimension of the
\(\omega^{-r}\)-part of \(\operatorname{Cl}(D)/p\), where
\(\omega\) describes the action on \(\mu_p\). The character is
even. Its order is \(m=(p-1)/\gcd(p-1,r)\), so our two divisibility
conditions give precisely \(m\in\{1,2,3,4,6\}\). Let \(F\) be
the field cut out by this character. Extension and norm of ideal
classes identify \(\operatorname{Cl}(F)/p\) with the invariants in
\(\operatorname{Cl}(D)/p\) under \(\operatorname{Gal}(D/F)\).
Indeed their two composites are multiplication by \([D:F]\)
and the subgroup sum, respectively, and the degree is a unit
modulo \(p\). The desired eigenspace belongs to these invariants.

The class number of every such \(F\) is prime to \(p\). For degrees
four and six one divides the analytic class number formula for
\(F\) by the formulas for its proper subfields. Only two primitive
Dirichlet values remain. Relative units bound the quotient of
the regulators from below and give a class number quotient less
than \(p\) above explicit thresholds. For the smaller primes
the two faithful characters are excluded by exact nonzero
Stickelberger residues. The calculation, including all
coefficients in the regulator bounds and the small primes, is
proved in Lemmas \ref{lem:relative195}--\ref{lem:unit195}, Proposition \ref{prop:conduct195} and Theorem \ref{thm:prime195}.
The latter theorem also treats degrees one, two and three.
Thus the eigenspace is zero and finiteness of \(K_{2r}(\ZZ)\)
gives vanishing at every exponent. Proposition~\ref{prop:qbound102}
identifies the orders in its supremum with the corresponding primary
orders, so deleting the zero terms preserves the equality.
\end{proof}

\begin{corollary}[Character components on the theta fields]\label{cor:units196}
Fix an admissible \(x\), its field \(E\) in \eqref{eq:tail114},
an odd prime \(p\nmid[E:\QQ]\), and \(1\le r<t(x)\). Extend
coefficients to an unramified ring containing the values of
the characters of \(\operatorname{Gal}(E/\QQ)\). The trivial
character summand of \(M_{p,r}\) is zero for the even twists
and primes of Theorem~\ref{thm:units196}.

More generally, let \(\alpha(c)=(-1)^r\) and suppose the even
residual character \(\tau=\overline{\kappa_p^{-r}\alpha^{-1}}\)
has order \(m=4\) or \(6\). Let \(F\) be its field and \(f\) its
conductor. Assume that \(p\) does not split completely in \(F\).
The \(\alpha\)-component of \(M_{p,r}\) vanishes if \(f=p\). It
also vanishes if \(p^{\operatorname{ord}_m(p)}\ge f\) and
\(\log f\ge16\) in degree four or \(\log f\ge90\) in degree six.
\end{corollary}
\begin{proof}
Averaging over the prime-to-\(p\) group
\(\operatorname{Gal}(E/\QQ)\) identifies the character summand
with the strict group over \(\QQ\) for the coefficient
\((J/O)(\kappa_p^{-r}\alpha^{-1})\), where \(O\) is the chosen
coefficient ring and \(J\) is its fraction field. The same
averaging on inertia away from \(p\) and on decomposition groups
at \(p\) preserves the strict local conditions, as proved in
Theorem \ref{thm:char}.
For \(\alpha=1\) and \(r\) even, its Pontryagin dual over
\(\ZZ_p\) before extending scalars is \(K_{2r}(\ZZ)\{p\}\),
by Proposition~\ref{prop:realclass100} and its proof.

In the second assertion, the finite coefficient lift of an
order-\(p\) class is locally zero at \(p\), because some local
element acts by \(\tau(g)\ne1\) modulo \(p\). Restricting to
\(F\) makes the residual coefficient trivial. Class field theory
therefore identifies these classes with homomorphisms on the
faithful component of the ordinary class group modulo \(p\).
Classes of primes over \(p\) are killed automatically on this
component by their stabilizers. The norm-kernel and prime-conductor
bounds give its vanishing under the indicated conditions, by
Corollary \ref{cor:strict195}.
A nonzero primary torsion group has an element of order \(p\).
This proves vanishing of the entire component on the original
field for each indicated \(x\).
\end{proof}

In the supplied source,
\srcloc{Development/AutomorphicArithmetic/AutomorphicVarieties/HeightAndVanishing/Counts}{49}
uses a rational vanishing statement for \(p-1\mid2r\).
The two results above extend the range of such vanishing.
The other residual characters still occur in the sum. At a prime \(p>6t+1\), all positive even twists
\(r<t\) already have order greater than six on the rational
summand, since \((p-1)/\gcd(p-1,r)\ge(p-1)/r>6\).
The uniform total in \eqref{eq:tail114} and the lower arithmetic
estimate for the original finite-level integral theta period
are still needed on the common intervals.

\subsection{A bound in every twist on a cyclic sextic subfield}\label{nav:I-62}
\begin{theorem}[A uniform positive contribution when three divides the degree]\leavevmode
\label{thm:sextic278}
Let \(B=\QQ(\zeta_7)^+\), \(K=\QQ(\sqrt{257})\) and \(F=BK\).
For each positive even \(r\), the order of the actual strict
group \(\mathcal S_{3,r}(F)\) is either \(3\) or \(9\).
It depends only on \(r\bmod6\). For any character cutting out
\(F\), put \(h_x=\lfloor(t(x)-1)/2\rfloor\).
On its original interval fields one has
\[
 h_x\log3\ \le\
   \sum_{1\le r<t(x)}
       \log\left|
       \bigl(\mathcal S_{3,r}(E_x)^{c=(-1)^r}\bigr)^{\Gal(E_x/F)}
       \right|
             \ \le\ h_x\log9 .
\]
\end{theorem}
\begin{proof}
The prime \(3\) is inert in \(B\), since \(3\) generates
\((\ZZ/7)^\times\). Theorem~\ref{thm:oneprime278} gives
\(\mathcal S_{3,r}(B)=0\) at every even twist. The other summand
under the involution of \(F/B\) is \(S(B,A)\) with
\(A=(\QQ_3/\ZZ_3)(\chi\kappa_3^{-r})\) for the quadratic
character \(\chi\) of conductor \(257\). Quadratic descent
applies integrally at \(3\). Here \(\chi(3)=\chi(7)=-1\).
For example the \(128\)-th powers of \(3\) and \(7\) are both
\(-1\bmod257\). Thus Theorem~\ref{thm:cyclic278} applies with
\(p=3\), \(a=1\), \(b=7\). It identifies the lengths of the two
character quotients with \(d_0,d_1\) and bounds the length on
the cubic field between their maximum and their sum.
Theorem~\ref{thm:quadratic278} gives \(d_0=1\).

For \(d_1\), take the cubic character \(\lambda\) modulo \(7\)
with \(\lambda(3)=\zeta_3\). Put \(\pi=\zeta_3-1\), so that
\(\pi^2=-3(1+\pi)\). In \(\ZZ_3[\zeta_3]/3\) one has
\(\zeta_3^j=1+j\pi\), and \(\pi\ne0\). The reduction argument
in Theorem~\ref{thm:cyclic278} gives \(d_1\ge1\). For the
reverse bound evaluate the class at \(s=43177=8\cdot5397+1\),
where \(5397=257\cdot7\cdot3\). Trial division through \(207\)
proves primality. Take \(g=5\) and \(w=g^8=2032\bmod s\).
For \(d=2,3,7,257\), the residues \(g^{(s-1)/d}\) are
\(43176,15074,18105,21688\). This checks that \(g\) is a generator.
For \(1\le a<5397\) prime to \(5397\), write
\(a=3^{j(a)}\bmod7\), with \(j(a)\) taken modulo three after
computing the logarithm modulo six. Form the three sums
\[
 C_j=\sum_{j(a)=j}\chi(a)\operatorname{ind}_g(1-w^a)\bmod3,
                   \qquad (C_0,C_1,C_2)=(0,2,1).
\]
To reproduce them, advance \(y\) from \(1\) by multiplying by
\(w\bmod s\) for each integer \(a=1,\ldots,5396\). Retain the
terms prime to \(5397\). The index modulo three is the unique
\(i\in\{0,1,2\}\) such that
\((1-y)^{(s-1)/3}=w^{1799i}\bmod s\). Give it sign \(+1\) or
\(-1\) according as \(a^{128}\) is \(1\) or \(-1\bmod257\).
The successive residues of \(3^j\bmod7\) for
\(j=0,\ldots,5\) are \(1,3,2,6,4,5\), which assigns the entry.

For every positive even \(r\), the twist factor \(a^r\) is one
modulo three. Proposition~\ref{prop:residue269} therefore
evaluates the corestricted class, up to a unit, as
\(\sum_j C_j(1+j\pi)=\pi\bmod\pi^2\). Thus \(d_1=1\).
This proves both bounds on the actual group over \(F\).
As its exponent divides nine, it equals its nine-torsion.
Lemma~\ref{lem:period30} identifies the nine-torsion at any two
positive twists congruent modulo six. Equivalently, at these
twists the finite coefficient actions agree modulo nine and
the local cyclotomic kernels in that lemma agree as well.
This proves the asserted periodicity. Finally \([E_x:F]\) is
a power of two on the fields of Theorem~\ref{thm:conductor232}.
Restriction gives the same invariant group there. The sign on
the real field \(F\) excludes the odd twists and counting the
even twists gives the two inequalities.
\end{proof}

These results evaluate part of the cohomology sum on the fields
used for the theta period. The sum over the other components and
the upper and lower estimates for the normalized integral period
in Section~\ref{sec:missing} remain open.

\subsection{The full strict groups over the cyclic sextic field}\label{nav:I-63}
\begin{theorem}[Exact cyclic groups at all positive even twists]\leavevmode
\label{thm:orders279}
Let \(F=\QQ(\sqrt{257},\zeta_7+\zeta_7^{-1})\). Then
\[
  \mathcal S_{3,r}(F)\ \simeq\
    \begin{cases}
       \ZZ/9\ZZ,& r\equiv0\pmod6,\ r>0,\\
       \ZZ/3\ZZ,& r\equiv2,4\pmod6,\ r>0.
    \end{cases}
\]
For any Dirichlet character cutting out \(F\), put
\(t=\lfloor x^{401}\rfloor\). At every point on the common
intervals the invariant contribution on its original fields \(E_x\) is
\[
 \sum_{r=1}^{t-1}
       \log\left|
       \bigl(\mathcal S_{3,r}(E_x)^{c=(-1)^r}\bigr)^{\Gal(E_x/F)}
       \right|
       =\left(\left\lfloor\frac{t-1}{2}\right\rfloor+
              \left\lfloor\frac{t-1}{6}\right\rfloor\right)\log3 .
\]
\end{theorem}
\begin{proof}
Keep the quadratic character \(\chi\), the cubic character
\(\lambda\) and the open set \(U\) of Proposition~\ref{prop:lines279}.
Both integral character lattices have degree-two length one.
All the local invariant terms vanish in their quadratic
character sector at \(3,7,257\), as proved in
Theorem~\ref{thm:cyclic278}. The sector for the trivial
quadratic character contributes zero by
Theorem~\ref{thm:oneprime278} for \(\QQ(\zeta_7)^+\).
Thus the integral degree-two group for
\[
       T_R=\ZZ_3(\chi\kappa_3^{r+1})\otimes R,\qquad
       R=\ZZ_3[X]/(X^3-1),
\]
is dual to the entire strict group over \(F\), by Shapiro's
lemma and the duality sequence of Theorem~\ref{thm:common279}.
For the integral group ring the dual lattice agrees with the
regular lattice, since the basis of group elements is orthonormal
for its invariant pairing. By Theorem~\ref{thm:intersection279}
the length is \(1+1-1+\dim(L_0\cap L_1)\). The proposition
therefore gives length two when \(6\mid r\) and one for the
other positive even twists.

To determine the abelian group, put \(M=H^2(U,T_R)\) and
\(\Phi=X^2+X+1\). One has \(M/\Phi M=H^2(U,T_R/\Phi T_R)\).
Indeed degree two is the last cohomological degree. Applying
the long exact sequence first to the image of multiplication
by \(\Phi\), and then to its inclusion, shows that its image
on cohomology is \(\Phi M\) and that the displayed quotient
map is surjective. The lattice quotient is the one for
\(\chi\lambda\), whose degree-two group has order three.
Consequently, for \(\bar M=M/3M\) and \(Y=X-1\), one has
\(\dim_{\mathbb F_3}\bar M/Y^2\bar M=1\), since
\(\Phi=Y^2\bmod3\).
Now \(Y^3=0\) on \(\bar M\). Its nonzero quotient modulo
\(Y^2\) therefore has nonzero quotient modulo \(Y\).
Both quotients have dimension one, so \(Y\bar M=Y^2\bar M\),
and nilpotence gives \(Y\bar M=0\). It follows that
\(\bar M\) has dimension one. The finite abelian \(3\)-group
\(M\), and hence its dual, is cyclic. This proves the asserted
isomorphism in each weight.

Finally \([E_x:F]\) is a power of two and restriction
identifies the strict group with the invariant part as in
Theorem~\ref{thm:sextic278}. The complementary part at an odd
weight vanishes on the real field \(F\). Each even weight
contributes \(\log3\), with a second \(\log3\) at each positive
multiple of six. Counting them gives the equality.
\end{proof}

The full interval estimate also includes the other characters
of \(E_x\) and the other coefficient primes. Their total and the
two estimates for the prescribed integral theta period in
Section~\ref{sec:missing} remain to be proved.

\subsection{Two infinite families of character fields with exact strict orders}\label{nav:I-64}
Keep \(K=\QQ(\sqrt{257})\), \(F,L\) and \(\gamma\) of
Theorem~\ref{thm:control280}. Let \(\tau\) be either generator
of \(\Gal(F/K)\), extended to fix \(Q_\infty\). For
\(m\ge2\) and \(e\in\{1,-1\}\) define \(D_{m,e}\) as the
fixed field in \(L\) of the closed subgroup generated by
\(\delta=\tau^e\gamma^{3^{m-1}}\).
Let \(D_{m,e}^+=D_{m,e}\), let \(B_{m,e}\) be its degree-\(3^m\)
subfield, and put \(D_{m,e}^-=B_{m,e}\QQ(\sqrt{-771})\).

\begin{theorem}[Exact orders for character fields of unbounded conductor]\leavevmode
\label{thm:fields280}
Each \(D_{m,e}^{\pm}\) is cyclic of degree \(2\cdot3^m\)
over \(\QQ\) and has conductor \(257\cdot7\cdot3^{m+1}\).
At every positive twist the complementary strict group is
given by
\[
 \begin{array}{c|cc}
 \mathcal S_{3,r}(D)^{c=(-1)^r}&r\ {\rm even}&r\ {\rm odd}\\ \hline
 D=D_{m,e}^+&\ZZ/3\ZZ&0\\
 D=D_{m,e}^-&0&\ZZ/3\ZZ
 \end{array}
\]
In particular, for any primitive Dirichlet character cutting
out one of these fields \(D\), at every point of the original
common intervals,
\[
     \sum_{1\le r<t(x)}
       \log\left|
       \bigl(\mathcal S_{3,r}(E_x)^{c=(-1)^r}\bigr)^{\Gal(E_x/D)}
       \right|
   =
   \begin{cases}
   \lfloor(t(x)-1)/2\rfloor\log3,&D=D_{m,e}^+,\\
   \lfloor t(x)/2\rfloor\log3,&D=D_{m,e}^- .
   \end{cases}
\]
\end{theorem}
\begin{proof}
The Galois group of \(L/\QQ\) is the product of the quadratic
group of \(K\), the cubic group generated by \(\tau\), and
\(\Gamma\) generated topologically by \(\gamma\). In the last
two factors the quotient by \(\delta\) is generated by the
image of \(\gamma\), since \(\tau^e=\gamma^{-3^{m-1}}\) there.
It has exact order \(3^m\), because the intersection of the
closed subgroup generated by \(\delta\) with \(\Gamma\) is
generated by \(\delta^3=\gamma^{3^m}\). The quadratic factor
is independent, so \(D_{m,e}\) has the asserted cyclic Galois
group. The two signs give distinct subgroups, since they give
different cubic components above \(\gamma^{3^{m-1}}\).

For the conductor, choose the character on \(\Gamma\) with
value \(\zeta_{3^m}\) on \(\gamma\), taken trivial on the
Teichmüller factor. Its conductor is \(3^{m+1}\), since the
order of \(4\bmod 3^{m+1}\) is \(3^m\) and the preceding
modulus gives order \(3^{m-1}\). Multiply this character by
the cubic character taking \(\tau\) to \(\zeta_3^{-e}\) and
the quadratic character of conductor \(257\), choosing
\(\zeta_{3^m}^{\,3^{m-1}}=\zeta_3\). The product kills
\(\delta\) and has order \(2\cdot3^m\), so it cuts out \(D_{m,e}\).
The three primitive factors have pairwise coprime conductors.
Each retains its conductor on its own inertia group, giving
the product conductor. Multiplication by the quadratic character
\(\omega\) of \(\QQ(\sqrt{-3})\) replaces the quadratic factor
by that of \(\QQ(\sqrt{-771})\). The \(3\)-part of the product
character's conductor is still \(3^{m+1}\), since \(\omega\) is trivial on principal
units and the character on \(\Gamma\) is primitive there.
This proves the statements on the degree and conductor for
the imaginary fields as well.

The extension \(L/D_{m,e}\) has group \(\ZZ_3\) generated by
\(\delta\), as follows by projection to \(\Gamma\).
Outside \(3\) the only possible ramification of \(L/\QQ\)
is at \(7\) and \(257\), whose inertia groups are finite and
meet this torsion-free subgroup trivially. Thus the twisted
quotient in Theorem~\ref{thm:control280} applies.
By Theorem~\ref{thm:tower280}, \(X\) is cyclic of order nine,
\(\gamma\) acts trivially, and \(\tau\) acts by a unit of
order three in \((\ZZ/9)^\times\). Such a unit is \(4\) or
\(7\), so subtracting one gives a multiple of three of
valuation exactly one. Since \(m\ge2\),
\(u(\delta)^{-r}=4^{-3^{m-1}r}\equiv1\bmod9\).
The twisted quotient \(S(D_{m,e}^+,A_r^{\rm pr})\) is
consequently cyclic of order three for either \(e\) and
every \(r\ge1\). The quadratic involution \(\sigma\) acts on
this group by \(-1\) by Theorem~\ref{thm:tower280}. Indeed
it fixes \(Q_\infty\), hence acts trivially on the coefficient
in the equivariant homomorphisms of Theorem~\ref{thm:control280}.

Write \(B=B_{m,e}\). Quadratic descent from \(D_{m,e}^+=BK\)
therefore gives \(S(B,A_r^{\rm pr})=0\) and
\(S(B,A_r^{\rm pr}\otimes\chi_{257})\simeq\ZZ/3\ZZ\),
where \(\chi_{257}\) is the character of \(K\). Here the
two summands arise by applying the two integral quadratic
projectors to the induced coefficient. Restriction detects
unramifiedness on inertia and the zero local class at three
since all relative local degrees divide two.
On the real field \(D_{m,e}^+\) the even twist uses
\(A_r^{\rm pr}\), and the odd complementary sign is absent.
On \(D_{m,e}^-/B\) complex conjugation is the nontrivial
quadratic automorphism with character \(\chi_{257}\omega\).
At an even twist its plus summand is \(S(B,A_r^{\rm pr})\).
At an odd twist its minus summand has coefficient character
\(\kappa_3^{-r}\chi_{257}\omega
   =\langle\kappa_3\rangle^{-r}\chi_{257}\).
This gives both rows of the table.

Finally the relative degree \([E_x:D]\) for these characters
is a power of two. Restriction and averaging preserve the
global and local conditions, as in Theorem~\ref{thm:oneprime278}.
Counting the positive even or odd twists less than \(t(x)\)
now gives the two formulas.
\end{proof}

The calculation follows the indicated invariant part as the
conductor grows. The other parts of the strict groups on \(E_x\)
and the required integral theta-period estimates in Section~\ref{sec:missing}
remain open.


\subsection{Exact strict groups for cyclic fields with several ramified primes}
\label{sec:cyclicray657}

Let \(p\) be an odd prime and \(b_1,\ldots,b_m\) distinct primes
congruent to \(1\bmod p\), where \(m\ge1\). Choose surjections
\(\nu_i:(\ZZ/b_i\ZZ)^\times\to(\mathbb F_p,+)\). They define a
Dirichlet character \(\psi=\zeta_p^{\,\sum_i\nu_i}\). Its field
\(K\) is cyclic of degree \(p\), is totally real and has conductor
\(\prod_i b_i\), since each local component has order \(p\).
Define a matrix over \(\mathbb F_p\) by
\[
       D_{ij}=\nu_j(b_i)\quad(i\ne j),\qquad
       D_{ii}=-\sum_{j\ne i}\nu_j(b_i).
\]
We assume that the rational primes satisfy
\[
          \operatorname{rank}D=m-1,\qquad \sum_i\nu_i(p)\ne0.
\]
The second says that \(p\) is inert in \(K\). The empty minor
has determinant one when \(m=1\).

\begin{lemma}[The strict character groups over the rational field]
\leavevmode\label{lem:scalarcyclic657}
Suppose \(r>0\) is a multiple of \(p-1\). For every
\(s\in\mathbb F_p\), the strong group \(S(\kappa_p^{-r}\psi^{-s})\)
over \(\QQ\) with coefficients in \(\QQ_p(\zeta_p)/\ZZ_p[\zeta_p]\)
vanishes. Here strictness is imposed at \(p,b_1,\ldots,b_m\)
and classes are unramified elsewhere.
\end{lemma}
\begin{proof}
Use \(\pi=\zeta_p-1\), so \(e=p-1\), and write
\(\theta_s=\kappa_p^{-r}\psi^{-s}\). Its reduction is trivial.
At precision one the global finite characters unramified away
from the indicated primes have basis \(\nu_0,\nu_1,\ldots,\nu_m\).
Here \(\nu_0\) generates the order-\(p\) character group at
\(p\) of conductor \(p^2\). Indeed Kronecker--Weber identifies
these global additive characters with characters of the product of
the corresponding cyclotomic factors. At each \(b_i\) the higher
units are pro-\(b_i\) and the residue multiplicative group is cyclic.
At \(p\), the \(p\)-primary unit factor \(1+p\ZZ_p\) is procyclic
and reduction modulo \(p^2\) gives its quotient of order \(p\).
There is no infinite-place contribution. Write the coordinates
of a character in this basis as \(X_0,\ldots,X_m\).

First take \(s\ne0\). At \(p\), \(\psi\) is unramified with
nontrivial value at Frobenius. The cyclotomic part
\(\kappa_p^{-r}\) takes values in \(1+p\ZZ_p\), hence is
one modulo \(\pi^2\). It follows that \(c_p=1\) in the notation
of Theorem~\ref{thm:raylength656}. The local correcting line
at \(p\) is unramified and nonzero on Frobenius. Membership
in that line gives \(X_0=0\).

At \(b_i\) choose an inertia element of \(\nu_i\)-value one,
and a Frobenius lift with \(\nu_i\)-value zero. This lift exists
by adjustment by inertia and has \(\nu_j\)-value \(\nu_j(b_i)\)
for \(j\ne i\). Again \(c_{b_i}=1\). On these two elements
the correcting character has values
\(-s\) and \(-s\sum_{j\ne i}\nu_j(b_i)\), respectively, since
\((\zeta_p^z-1)/\pi\equiv z\bmod\pi\). A finite character
with \(X_0=0\) has values \(X_i\) and
\(\sum_{j\ne i}\nu_j(b_i)X_j\). In the abelian quotient of
exponent \(p\) a uniformizer and the cyclic tame unit quotient
generate the local group. Proportionality on these two elements
therefore gives the entire local condition. For every \(i\) it reads
\[
             \sum_{j\ne i}\nu_j(b_i)(X_j-X_i)=0.
\]
By the rank assumption only the constant vectors survive.
These constitute exactly the global correcting line of
\(\theta_s\), whose \(m\) entries are all \(-s\).
Theorem~\ref{thm:raylength656} gives zero \(\pi\)-torsion.

If \(s=0\), the local correcting characters at the \(b_i\)
are unramified, so every \(X_i\), \(i\ge1\), vanishes. The
global correcting character generates the remaining line
\(\mathbb F_p\nu_0\). Indeed the cyclotomic image over \(\QQ\)
and its inertia image at \(p\) are both \(\ZZ_p^\times\).
Thus \((\kappa_p^{-r}-1)/\pi^{c_{\QQ}}\) is a unit on some
element of this inertia group. This again gives zero
\(\pi\)-torsion. Each divisible-coefficient cohomology group
under consideration is \(\pi\)-primary torsion, so vanishing
at precision one proves the lemma.
\end{proof}

\begin{lemma}[The \(p\)-class group detected by the ramified primes]
\leavevmode\label{lem:genuscyclic657}
The \(p\)-primary subgroup of \(\mathrm{Cl}(K)\) is elementary
abelian of rank \(m-1\). It is generated by the prime classes above
the \(b_i\).
\end{lemma}
\begin{proof}
Let \(K'\) be the compositum of the degree-\(p\) fields cut out
by the individual \(\nu_i\). Their disjoint conductors identify
\(\operatorname{Gal}(K'/\QQ)\) with \(\mathbb F_p^m\), where
restriction to \(K\) is the sum of the coordinates. Each inertia
line in \(K'/\QQ\) maps injectively to this quotient. Thus
\(K'/K\) is unramified, including at the real places since
all fields are totally real.

Let \(H_p\) be the Hilbert \(p\)-class field of \(K\) and
\(C=\operatorname{Gal}(H_p/K)\). Global class field theory
identifies \(C\) equivariantly with \(\mathrm{Cl}(K)\{p\}\).
The extension \(H_p/\QQ\) is Galois by the defining maximality.
Write \(\sigma\) for a generator of \(\operatorname{Gal}(K/\QQ)\).
The commutator subgroup over \(\QQ\) is \((\sigma-1)C\), since
the group is generated by the abelian normal subgroup \(C\)
and one lift of \(\sigma\), and commuting that lift with
each member of \(C\) gives exactly the claimed subgroup.
The maximal subextension abelian over \(\QQ\) contains \(K'\).
It is unramified outside the \(b_i\), with inertia of exponent
at most \(p\) there since \(H_p/K\) is unramified. By
Kronecker--Weber its characters are products of local characters
at the \(b_i\) of order dividing \(p\), all of which factor
through the cyclic residue quotients defining \(K'\).
It therefore equals \(K'\), and
\[
 C/(\sigma-1)C \simeq
   \operatorname{Gal}(K'/K)=
   \ker\!\left(\mathbb F_p^m\xrightarrow{\ \sum\ }\mathbb F_p\right).
\]

Let \(\mathfrak b_i\) be the unique prime above \(b_i\) in \(K\).
Its class is killed by \(p\), since \((b_i)=\mathfrak b_i^p\).
Its Frobenius in \(K'/K\) has coordinates \(\nu_j(b_i)\) for
\(j\ne i\), since its residue degree over \(\QQ\) is one and those
factors are unramified there. Its remaining coordinate is forced
by membership in the kernel of the sum. The images of these
classes are exactly the rows of \(D\) under arithmetic reciprocity.
By rank they span \(C/(\sigma-1)C\). If \(B\) is the subgroup
they generate, the finite \(p\)-group \(C/B\) consequently
satisfies \(C/B=(\sigma-1)(C/B)\). The operator \(\sigma-1\)
is nilpotent on this quotient because its \(p\)-th power
is zero modulo \(p\). Iteration increases divisibility by \(p\) until
the group is killed. Thus \(C=B\). Every generating prime class
is fixed by \(\sigma\), so \(C=C/(\sigma-1)C\), as asserted.
\end{proof}

\begin{theorem}[The complete strict group in the cyclic subfield]
\leavevmode\label{thm:strictcyclic657}
Under the two residue hypotheses, for every positive multiple
\(r\) of \(p-1\) one has
\[
 \mathcal S_{p,r}(K)\simeq(\ZZ/p\ZZ)^{m-1},\qquad
 \mathcal S_{p,r}(E_x)^{\operatorname{Gal}(E_x/K)}
                 \simeq\mathcal S_{p,r}(K)
\]
on every original period field of \(\psi\).
\end{theorem}
\begin{proof}
Apply the integral Fourier maps in the proof of
Theorem~\ref{thm:integralblock656} to \(K/\QQ\), with strong
conditions at \(p\) and its conductor. Over \(K\), unramified
classes away from \(p\) are locally zero in the divisible
coefficient by the Frobenius calculation in that proof.
All the scalar groups vanish by Lemma~\ref{lem:scalarcyclic657},
so multiplication by \(p\) kills the original group after
extension to \(\ZZ_p[\zeta_p]\). This extension is faithfully
flat, hence \(p\) kills \(\mathcal S_{p,r}(K)\) itself.

Compute its \(p\)-torsion using \(\QQ_p/\ZZ_p\) as coefficient.
The mod-\(p\) twist is trivial. There is a unique \(w\mid p\),
and \(K_w/\QQ_p\) is unramified. The \(p\)-power cyclotomic
extension of \(\QQ_p\) is totally ramified, so the local
cyclotomic image is \(\ZZ_p^\times\), as is the global image
since it contains this local image. In the coefficient sequence
for multiplication by \(p\) the global correcting line thus
restricts isomorphically to the local correcting line at \(w\).
In detail, if \(c=\min v_p(\kappa_p^{-r}-1)\), the minimum is
the same globally and locally, and the same generator
\((\kappa_p^{-r}-1)/p^c\bmod p\) is nonzero at \(w\).
Every strict class therefore has a unique finite-character
representative vanishing there, by subtracting its unique
global correcting multiple. Away from \(p\) the finite
character must be unramified, since inertia acts trivially
on the divisible coefficient and its connecting image is zero.
Conversely, each such everywhere-unramified character split
at \(w\) represents a strict class. Class field theory gives
\[
    \mathcal S_{p,r}(K)[p]\simeq
       \operatorname{Hom}\bigl(\mathrm{Cl}(K)/\langle[w]\rangle,
                                     \mathbb F_p\bigr).
\]
The prime \(w=(p)\) is principal by inertness. The class group
lemma and the exponent bound now give the first isomorphism.
For the second, the extension \(E_x/K\) has degree a power of
two. Restriction and averaging identify global classes with
the invariants upstairs and detect each local kernel by local
restriction, whose degree also divides that power of two.
\end{proof}

For example, take \(p=3\) and \((b_1,b_2,b_3)=(7,13,37)\).
Take \(\nu_i\) to be discrete logarithm modulo three to bases
\((3,2,2)\), respectively. Their powers of exponent
\((b_i-1)/3\) are \(2,3,26\), each of order three in its
respective residue field. For a residue \(z\), comparison of
\(z^{(b_i-1)/3}\) with the three powers of this element gives
\(\nu_i(z)\). This gives \((\nu_i(3))=(1,1,2)\) and
\[
       D=\begin{pmatrix}2&2&2\\0&1&2\\2&1&0\end{pmatrix}
                          \quad\text{over }\mathbb F_3 .
\]
The upper left minor is \(2\), the row sums vanish and the
three values at \(3\) sum to \(1\bmod3\). Thus on the cyclic
cubic field of this character of conductor \(3367\), the entire
strict group has order nine and exponent three at every positive
even weight. On its invariant summand in \(E_x\) the exact
logarithmic contribution for \(1\le r<t(x)\) divisible by
\(p-1\) is \(\lfloor(t(x)-1)/(p-1)\rfloor(m-1)\log p\)
in the general theorem. Other character summands on the same
fields still enter the estimate in Section~\ref{sec:missing}.


\subsection{Tame inertia and the strict group in every dyadic layer}
\label{sec:tameray662}

Keep the notation \(K,P,H,\Delta_H,\chi,M,R_h\) of
Theorem~\ref{thm:wholeblock661}. Assume \(\chi(D_p)\ne1\).
Only finitely many primes ramify in the cyclic \(p\)-extension
\(K/\QQ\). Among them put
\[
 T=\{b\ne p:b\text{ ramifies in }K,\ \chi(D_b)=1\},\qquad
 u_b=v_p(e_b(K/\QQ)),\qquad \mathfrak m=\prod_{b\in T}b\mathcal O_H .
\]
Every \(b\in T\) splits completely in \(H\), and \(u_b\ge1\).
Use the ordinary ray class group \(\Cl_{\mathfrak m}(H)\).
For \(b\in T\) let \(I_b^{\mathrm{ray}}\) be the image in this
group of \(\prod_{w\mid b}(\mathcal O_H/w)^\times\) under the
residue-unit map in the ray class sequence. In additive notation put
\[
 C(K,\chi)=\Cl_{\mathfrak m}(H)\Big/
                      \sum_{b\in T}p^{u_b}I_b^{\mathrm{ray}} .
\]
This finite group and its \(\Delta_H\)-action require only the
residual field and the ramification indices of \(K\).

\begin{theorem}[The ray class group at the split ramified primes]
\leavevmode\label{thm:tameray662}
For \(1\le h\le1+v_p(r)\) there is an isomorphism
\[
 (M[p^h])^P \simeq
 \operatorname{Hom}_{\Delta_H}
                \big(C(K,\chi), R_h(\widetilde\chi)\big).
\]
In particular, \(M=0\) if and only if
\(\operatorname{Hom}_{\Delta_H}(\Cl_{\mathfrak m}(H),k(\chi))=0\).
\end{theorem}
\begin{proof}
At the indicated precision the action on \(R_h\) is
\(\widetilde\chi\). By Theorem~\ref{thm:wholeblock661} and its
proof, \(M[p^h]\) is the group over \(K\) with this coefficient,
unramified away from \(p\) and locally zero above \(p\).
Some value of \(\widetilde\chi-1\) at a place over \(p\) is a
unit, so the coefficient invariants over \(K\) vanish.
Inflation--restriction identifies the \(P\)-fixed global
cohomology with first cohomology over \(\QQ\). We determine the
local conditions on that cohomology.

At \(p\), restriction is injective because the same local
invariants vanish. At a prime \(v\ne p\) unramified in \(K\) the
inertia group is unchanged. If \(v\ne p\) ramifies in \(K\) and
\(\chi(D_v)\ne1\), its image stays nontrivial on \(G_{K_w}\)
by the \(p\)-power local index. The unramified local cohomology
over each of the two fields is zero. Indeed, either inertia
has a value whose difference from one is a unit, or the
coefficient is unramified and Frobenius minus one is a unit.
Local restriction is injective by vanishing of the invariants
over \(K_w\). Thus at these primes unramifiedness descends too.

Now take \(b\in T\). The coefficient is trivial on \(D_b\),
so local classes are additive characters. Their restrictions
to wild inertia vanish since \(b\ne p\). The maximal pro-\(p\)
quotient of tame inertia is procyclic. In this quotient the
image of inertia over \(K_w\) has index \(p^{u_b}\), since
the image of inertia in \(\operatorname{Gal}(K_w/\QQ_b)\) has
exactly that order. Hence a character is unramified over \(K_w\)
precisely when its values on inertia are killed by \(p^{u_b}\).
It can take any unramified value allowed by global reciprocity.
Local reciprocity identifies the tame inertia character with a
character of \(\mathbb F_b^\times\).

Restriction from \(\QQ\) to \(H\) now gives the
\(\widetilde\chi\)-equivariant characters by averaging over
\(\Delta_H\), whose order is prime to \(p\). Away from \(T\),
restriction detects unramifiedness, since the corresponding
inertia indices are also prime to \(p\). At \(b\in T\) the
completions of \(H\) equal \(\QQ_b\). Thus finite class field
theory gives the ray modulus \(\mathfrak m\) and the indicated
quotient by the multiples of the inertia images. At a prime of
\(H\) over \(p\), a stabilizer element acts on the target by a root
whose difference from one is a unit. The ray class of this prime
is fixed by that element. Its value must vanish, so every such
ray character already satisfies local splitting at \(p\).
Real places contribute no odd-order characters. This proves the
isomorphism in both directions.

At \(h=1\) the extra multiples \(p^{u_b}I_b^{\mathrm{ray}}\)
are automatically killed. The space \(M[p]\) is finite by
Theorem~\ref{thm:wholeblock661}. The operator \(\sigma-1\) for
a generator of \(P\) is nilpotent on it, since
\((\sigma-1)^{|P|}=0\) in characteristic \(p\).
Every nonzero such space has nonzero invariants. Every nonzero
\(p\)-primary torsion group has a nonzero element of order \(p\).
This proves the vanishing equivalence.
\end{proof}

\begin{corollary}[A unit test in the residue fields]
\leavevmode\label{cor:unitrows662}
Suppose \(\operatorname{Hom}_{\Delta_H}(\Cl(H),k(\chi))=0\).
Put \(U=\mathcal O_H^\times\), \(V=(\mathcal O_H/\mathfrak m)^\times\),
and \(e_\chi=|\Delta_H|^{-1}\sum_{\delta\in\Delta_H}
\chi(\delta)^{-1}\delta\in k[\Delta_H]\). Then
\[
 \dim_k(M[p])^P
  = |T|-\dim_k e_\chi\operatorname{im}
       \big(U\otimes_{\ZZ}k\longrightarrow V\otimes_{\ZZ}k\big).
\]
In particular, the residues of global units decide whether the
entire \(P\)-block vanishes.
\end{corollary}
\begin{proof}
In the ray class sequence the kernel of the surjection onto
\(\Cl(H)\) is \(V/\operatorname{im}U\). First take the finite
\(p\)-parts, extend scalars to \(O\), and apply
\(e_{\widetilde\chi}\), defined by the same projector with
Teichm\"uller values. These operations preserve exactness.
The projected ordinary class group has zero reduction modulo \(p\)
by the character hypothesis, and is zero by finiteness.
Consequently the projected ray group modulo \(p\) is the quotient
of \(e_\chi(V\otimes k)\) by the displayed unit image. Its dual
has the required dimension by Theorem~\ref{thm:tameray662}.

For each \(b\in T\) one has \(p\mid b-1\). Indeed Frobenius
raises a tame inertia generator to its \(b\)-th power in the
abelian local quotient of \(K/\QQ\), where that generator has
order \(p^{u_b}\). The primes over \(b\) in \(H\) are permuted
freely and transitively by \(\Delta_H\), and their residue
fields are all canonically \(\mathbb F_b\). Thus this part of
\(V\otimes k\) is one copy of the regular representation of
\(\Delta_H\), contributing exactly one dimension to \(e_\chi\).
Summing over \(b\) proves the formula.
\end{proof}

\begin{theorem}[The cubic field of conductor \(3367\) in the real dyadic tower]
\leavevmode\label{thm:cubictower662}
Let \(K\) be the cyclic cubic field with the three characters
at \(7,13,37\) specified in Section~\ref{sec:cyclicray657}.
For every \(J\ge0\) and every positive even \(r\), restriction gives
\[
 \mathcal S_{3,r}(K)\ \xrightarrow{\ \sim\ }\ 
       \mathcal S_{3,r}(K C_J)\simeq(\ZZ/3\ZZ)^2 .
\]
On the original fields \(E_x\) this group is exactly the subgroup
fixed by \(\operatorname{Gal}(E_x/K C_J)\).
\end{theorem}
\begin{proof}
Split the action of \(\operatorname{Gal}(C_J/\QQ)\) over an
unramified coefficient extension containing its character values.
The preceding descent proofs apply to \(K C_J\) directly, since
the complementary factor to its cubic Sylow group is the dyadic
group. For a character \(\beta\) of exact order \(2^j>1\) the
residual character is \(\bar\beta^{-1}\), as \(r\) is even.
Its field is \(C_j\). Its ordinary class number is prime to
three by Corollary~1.1 of Fukuda and Komatsu~\cite{V-III-weberFK524}.

For \(v=3,13,37\) one has \(v_2(v^2-1)-3=0\), and for \(v=7\)
the value is one. The Frobenius order calculation of
Corollary~\ref{cor:dyadiccyclic661} shows that all four
decomposition groups have nontrivial character image if \(j\ge2\).
Theorem~\ref{thm:cyclicdescent661} kills these blocks. At \(j=1\)
only \(7\) splits. In \(H=C_1=\QQ(\sqrt2)\) the two residues
of \(\sqrt2\) at \(7\) are \(3,-3\). The unit \(1+\sqrt2\)
has residues \(4,5\). Their discrete logarithms to base three
in \(\mathbb F_7^\times\) are \(4,5\), distinct modulo three.
Conjugation exchanges the two places, so their logarithms have
nonzero projection to the one-dimensional nontrivial character
space. Corollary~\ref{cor:unitrows662} kills this last
nontrivial dyadic block as well.

The trivial dyadic block descends to \(K\) by restriction and
averaging, globally and on the local kernels, since \(2^J\) is
a unit at three. The calculation over \(K\) in
Theorem~\ref{thm:strictcyclic657} gives the claimed group.
The coefficient extension is faithfully flat, so restriction is
already an isomorphism over \(\ZZ_3\). The further extension
\(E_x/K C_J\) has degree a power of two and the same restriction
argument identifies its invariants.
\end{proof}

Thus on this whole dyadic subfield the even weights
\(1\le r<t\) contribute exactly
\(2\lfloor(t-1)/2\rfloor\log 3\) to the logarithmic sum of
strict orders. The terms on which a moving quadratic acts
nontrivially still enter the sum on the full period field.


\subsection{Cyclic descent at every precision and the split coefficient prime}
\label{sec:cycpres663}

Let \(E,K,P,\alpha,\theta,\chi,H,O\) and \(M\) be as in
Section~\ref{sec:beforecyclic661}, and assume \(\chi\ne1\).
The prime \(p\) may split completely in \(H\). Write \(Q_n\)
for the degree \(p^n\) layer of the cyclotomic \(\ZZ_p\)-extension
of \(\QQ\), with \(Q_0=\QQ\). Its completion at \(p\) is \(Q_{p,n}\),
and \(Q_{p,\infty}\) denotes their union. For an integer \(h\ge1\)
and a subgroup \(U\le P\), write \(K_0=K^U\) and
\(a=v_p([K:\QQ])\). Put
\[
 n=\max\{0,h-1-v_p(r)\},\qquad m=\max\{n,a\},\qquad
 F=K_0 H Q_n .
\]
On \(G_{K_0}\) the character \(\theta\bmod p^h\) factors through
\(G=\operatorname{Gal}(F/K_0)\). Indeed the principal-unit part
of the cyclotomic character raised to the power \(-r\) is
trivial modulo \(p^h\) on \(G_{Q_n}\). This follows from
\(v_p((1+p)^{r p^n}-1)=1+v_p(r)+n\) for odd \(p\).
The prime-to-\(p\) part is \(\widetilde\chi\).

Define a finite quotient \(C_h(F;K)\) of the idele class group
of \(F\) by the following local subgroups. All completions used
together are embedded in one algebraic closure of the rational
local field. The extensions in question are abelian, so the
subgroups are independent of the choices and are permuted by \(G\).
\begin{enumerate}
\item At \(w\mid p\), take the norm subgroup from
\(K_v H_{v'} Q_{p,m+h}\) to \(F_w\), where \(v,v'\) are the
compatible places of \(K,H\).
\item At \(w\mid b\ne p\) where \(K/K_0\) ramifies and
\(\chi(D_b)=1\), let \(e_b\) be the ramification index of
\(K/K_0\) there. Take the group generated by the principal
units and the \(e_b\)-th powers of the local units.
\item At the remaining finite places take the whole unit group.
At every infinite place take the whole multiplicative group.
\end{enumerate}
The local subgroups at the finitely many exceptional places
contain open subgroups of the units. The quotient is consequently
a quotient of a finite ray class group. The norm in the first
rule is a local field norm.

\begin{theorem}[All finite torsion over each intermediate field]
\leavevmode\label{thm:everycyclic663}
With \(R_h=O/p^hO\), there is an isomorphism
\[
       (M[p^h])^U \simeq
        \operatorname{Hom}_{G}\big(C_h(F;K),R_h(\theta)\big).
\]
In particular \(U=1\) computes the whole group killed by \(p^h\).
For \(U=P\) the global field used in the computation is \(H Q_n\).
\end{theorem}
\begin{proof}
Descent through the prime-to-\(p\) part of \(E/\QQ\) identifies
\(M\) with the strict group for \(A(\theta)\) over \(K\), where
\(A=O[1/p]/O\), as in Theorem~\ref{thm:wholeblock661}.
The image of the nontrivial residual character on \(G_K\) is
its full image, since \([K:\QQ]\) is a power of \(p\).
Thus \(A(\theta)^{G_K}=0\). Multiplication by \(p^h\) on
\(A(\theta)\), with kernel identified with \(R_h(\theta)\)
by division by \(p^h\), gives unique finite-coefficient lifts
of \(M[p^h]\). They are unramified away from \(p\). Indeed,
the inertia action there has order prime to \(p\), so its
divisible-coefficient invariants are either \(A\) or zero.
The map on inertia cohomology from finite to divisible
coefficients is therefore injective. An unramified class maps
to zero in divisible local cohomology at such a place, since
on nonzero inertia invariants Frobenius acts by a root of unity
times the inverse \(r\)-th power of the residue cardinality,
a scalar different from one. At \(p\) the local condition on
the finite lift is that its divisible-coefficient class vanish.

The finite coefficient also has zero invariants over \(K\).
Inflation--restriction from \(K_0\) to \(K\) identifies
\(U\)-invariant global classes with global classes over \(K_0\).
Restrict these further to \(F\). The group \(G\) is the product
of \(\Delta_H\) and the group of \(K_0 Q_n/K_0\), since
\(K_0 Q_n\) has \(p\)-power degree. The cohomology of
\(\Delta_H\) with coefficients \(R_h(\theta)\) vanishes in
every degree, including degree zero. Indeed its order is a unit
and some value of \(\widetilde\chi-1\) is a unit. The group
extension sequence then gives \(H^1(G,R_h(\theta))=
H^2(G,R_h(\theta))=0\). Restriction identifies global classes
over \(K_0\) with \(G\)-equivariant additive characters over
\(F\).

Consider a finite place over \(b\ne p\). If \(K/K_0\) is
unramified there, inertia is unchanged in this first descent.
At a ramified prime for which \(\chi(D_b)\ne1\), its image
remains nontrivial over both local \(p\)-extensions of
\(\QQ_b\). The unramified finite subgroup is zero on both
local fields, and restriction is injective by vanishing of
local invariants. These assertions follow by testing inertia
first and then Frobenius on its invariants, as in
Theorem~\ref{thm:tameray662}. In passing from \(K_0\) to \(F\),
the inertia index is prime to \(p\), since \(Q_n\) is unramified
at \(b\). Restriction thus detects unramifiedness in both cases.

At the exceptional tame primes \(H\) splits and \(F/K_0\) is
unramified. The finite coefficient can have nontrivial
Frobenius but is trivial on inertia. Its restriction to inertia
is an additive character on the procyclic \(p\)-part of tame
inertia. Inertia over \(K\) maps to the subgroup of index \(e_b\)
in that quotient. Thus the class is unramified over \(K\) exactly
when \(e_b\) kills its inertia image. Local reciprocity over \(F_w\)
gives the second rule.

Finally fix a place over \(p\) and write
\(L_\infty=K_v H_{v'}Q_{p,\infty}\). The coefficient \(A(\theta)\)
is trivial over \(L_\infty\). The group
\(\operatorname{Gal}(L_\infty/K_v)\) is a product of a finite
group of order prime to \(p\) and an open subgroup of the
cyclotomic \(\ZZ_p\). Indeed it embeds with finite index in
\(\operatorname{Gal}(H_{v'}Q_{p,\infty}/\QQ_p)\), which has
this form by disjointness of the two factors. A generator of
its procyclic factor acts on \(A\) by a scalar different from
one since \(r>0\). Subtracting a coboundary can therefore kill
the value of a cocycle on this generator, as multiplication
by that scalar minus one is surjective on \(A\). Density of
integral powers and averaging over the finite factor give
\(H^1(\operatorname{Gal}(L_\infty/K_v),A(\theta))=0\).
By inflation--restriction the local divisible class is zero
over \(K_v\) exactly when it is zero on \(G_{L_\infty}\).
There the action is trivial, so this is precisely the vanishing
of the finite character on \(G_{L_\infty}\).

The restriction of that character to \(G_{F_w}\) must therefore
factor through \(L_\infty/F_w\). To make the condition finite put
\(L_i=K_v H_{v'}Q_{p,i}\). The intersection of \(K_v H_{v'}\)
with \(Q_{p,\infty}\) has degree at most \(p^a\), since
\([H_{v'}:\QQ_p]\) is prime to \(p\). For \(i\ge m\) restriction
therefore identifies \(\operatorname{Gal}(L_\infty/L_i)\) with
\(\operatorname{Gal}(Q_{p,\infty}/Q_{p,i})\). In particular
\(\operatorname{Gal}(L_\infty/L_{m+h})\) is the subgroup
of \(p^h\)-th powers in \(\operatorname{Gal}(L_\infty/L_m)\).
The latter group fixes \(F_w\). Every character of exponent
dividing \(p^h\) which factors through \(L_\infty/F_w\) thus
factors through \(L_{m+h}/F_w\). The converse follows by
restriction. Local reciprocity gives exactly the first norm
subgroup. Global reciprocity with these local rules now gives
the required isomorphism. Odd-order characters impose no
further condition at real places.
\end{proof}

\begin{corollary}[A ray of modulus \(p^2\) at a split coefficient prime]
\leavevmode\label{cor:splitray663}
Suppose \(p\) splits completely in \(H\). Put
\[
 T=\{b\ne p:b\text{ ramifies in }K,\ \chi(D_b)=1\},\qquad
 \mathfrak m=p^2\mathcal O_H\prod_{b\in T}b\mathcal O_H .
\]
Let \(K_v\) be the completion at \(p\). Start with
\(C=\Cl_{\mathfrak m}(H)\). If \(K_v\subset Q_{p,\infty}\),
quotient \(C\) by the classes of the ideles equal to \(p\)
at one place over \(p\) and one everywhere else, for all such
places. Use \(C\) itself when \(K_v\not\subset Q_{p,\infty}\).
Writing \(C'\) for the resulting group, one has
\[
             \dim_k (M[p])^P
            =\dim_k\operatorname{Hom}_{\Delta_H}(C',k(\chi)).
\]
In particular \(M=0\) exactly when the right side is zero.
\end{corollary}
\begin{proof}
Take \(h=1\) and \(U=P\), so \(F=H\) and each \(F_w\) at \(p\)
is \(\QQ_p\). By local reciprocity the maximal abelian pro-\(p\)
quotient over \(\QQ_p\) is the product of two copies of
\(\ZZ_p\), from the uniformizer and the principal units. In
the cyclotomic \(\ZZ_p\)-extension the uniformizer acts trivially
and the principal units act faithfully. In the compositum with
\(K_v\) the uniformizer therefore generates a cyclic group of
order \([K_v:K_v\cap Q_{p,\infty}]\) in the kernel of restriction
to \(Q_{p,\infty}\). To see it generates that kernel, decompose
each local reciprocity element into a uniformizer power and a
unit. A trivial image in the cyclotomic quotient forces the
principal-unit part to be trivial. The principal-unit direction
maps isomorphically to the cyclotomic quotient and has trivial
intersection with this finite kernel. At exponent \(p\), the
local character is consequently arbitrary on
\((1+p\ZZ_p)/(1+p^2\ZZ_p)\). Its value at \(p\) is free when
the finite kernel is nontrivial and is zero otherwise. This gives
the stated modulus and uniformizer classes. At each \(b\in T\)
the tame index is divisible by \(p\), so its condition at this
precision is simply the squarefree ray condition. Apply the
theorem. The equivalence for the whole block follows from
nilpotence of \(\sigma-1\) on \(M[p]\), as in
Theorem~\ref{thm:tameray662}.
\end{proof}

\begin{corollary}[Simultaneous tests for the moving quadratic characters]
\leavevmode\label{cor:allmoving663}
Fix an odd \(p\) dividing \(\operatorname{ord}\psi\) and take a
sufficiently distant original interval. For every \(\alpha\)
nontrivial on at least one of the two moving quadratic factors,
one ray group over \(H\) decides the vanishing of \(M\)
simultaneously at all positive weights in any one residue class
modulo \(p-1\) of the prescribed sign. Use
Corollary~\ref{cor:splitray663} when \(p\) splits completely
in \(H\), and Theorem~\ref{thm:tameray662} otherwise.
The defining field \(H\) has conductor dividing
\(\operatorname{lcm}(p,N_*)\). The finite ray modulus divides
\(p^2\prod_{b\ne p,\ b\text{ ramifies in }K} b\mathcal O_H\),
where the factor \(p^2\) also denotes its extended ideal.
\end{corollary}
\begin{proof}
The two auxiliary primes are distinct and avoid \(2p\) and the
fixed conductor for all sufficiently large intervals. At the
prime of a nontrivial moving factor, its quadratic inertia
value survives reduction modulo odd \(p\). The other factors
and \(\bar\kappa_p\) are unramified there, so \(\chi\ne1\).
It is cut out by a character of conductor dividing the displayed
least common multiple. At precision one its value depends only
on \(\alpha\) and \(r\bmod(p-1)\). The local ramification indices,
the split-prime uniformizer condition and the residual field
are consequently the same throughout that weight class. The two
ray formulas give the conclusion, including the asserted modulus.
\end{proof}

Vanishing of one such character space removes its whole block
from the sum in Section~\ref{sec:missing}. If it is nonzero, the
test at precision one determines the number of indecomposable
summands of \(M[p]\). Their lengths at higher precision still
enter the required uniform bound.


\subsection{Tame ramification in the moving quadratic components}
\label{sec:movinggenus665}

Use the fields and character notation of
Section~\ref{sec:cycpres663}, with \(\chi\ne1\). Write
\(S_{\QQ}(\theta)\) for the subgroup of \(H^1(\QQ,A(\theta))\)
locally zero at \(p\) and unramified away from \(p\), with
\(A=O[1/p]/O\). At a finite prime away from \(p\) its unramified
local subgroup is zero, by the Frobenius calculation in
Theorem~\ref{thm:wholeblock661}. Put \(s=1+v_p(r)\), and
use the set \(T\) and tame exponents \(u_b\) of
Theorem~\ref{thm:tameray662}, allowing either splitting behavior
at \(p\). Let \(w\) be a place of \(K\) over \(p\), and put
\[
       p^\ell=[K_w:K_w\cap Q_{p,\infty}],\qquad
       v=\begin{cases}
             \min(\ell,s),&\chi(D_p)=1,\\
             0,&\chi(D_p)\ne1.
          \end{cases}
\]

\begin{theorem}[The length added by cyclic descent]
\leavevmode\label{thm:descentcost665}
Restriction gives an exact sequence at its first three terms
\[
 0\longrightarrow S_{\QQ}(\theta)\longrightarrow M^P
   \longrightarrow A[p^v]\oplus\bigoplus_{b\in T}A[p^{u_b}].
\]
In particular the difference of the \(O\)-lengths of \(M^P\)
and \(S_{\QQ}(\theta)\) lies between zero and \(v+\sum_{b\in T}u_b\).
The same bounds apply to the difference of the lengths of their
images under multiplication by \(p^j\), for every \(j\ge0\).
\end{theorem}
\begin{proof}
Since \(\chi\ne1\) has prime-to-\(p\) order, \(A(\theta)^{G_K}=0\).
Global inflation--restriction identifies \(P\)-invariant
cohomology over \(K\) with cohomology over \(\QQ\).
The local conditions for \(M\) are those over \(K\) with
coefficient \(A(\theta)\). Away from \(p\), unramifiedness
descends at every prime unramified in \(K\), because inertia
is unchanged. At a ramified prime with \(\chi(D_b)\ne1\),
the local invariants over \(K_w\) vanish, so restriction on
local \(H^1\) is injective. The unramified local subgroup over
\(K_w\) is zero. Therefore a descended strict class is locally
zero at such a prime. The same injectivity holds at \(p\) if
\(\chi(D_p)\ne1\).

At \(b\in T\) the coefficient is unramified. Restriction of a
local class to tame inertia determines it, since the unramified
subgroup is zero. Wild inertia gives no \(p\)-primary character.
Choose a generator of the pro-\(p\) tame inertia quotient.
The condition of becoming unramified over \(K_w\) forces its
value to be killed by \(e_b(K/\QQ)=p^{u_b}\). Evaluation
therefore embeds the possible local images into \(A[p^{u_b}]\).

If \(\chi(D_p)=1\), then over \(\QQ_p\) the full character
\(\theta\) is \(\langle\kappa_p\rangle^{-r}\), with the brackets
denoting projection to \(1+p\ZZ_p\). Its action is trivial over
\(Q_{p,\infty}\). Restriction to that field is injective on local
cohomology. Indeed a generator of the cyclotomic \(\ZZ_p\) acts
on \(A\) by a scalar whose difference from one is surjective,
so the first cohomology of that quotient vanishes. A class
becoming zero over \(K_w\) now gives a homomorphism on
\(\operatorname{Gal}(K_wQ_{p,\infty}/Q_{p,\infty})\).
This group is cyclic of order \(p^\ell\). The compositum is
abelian over \(\QQ_p\), and equivariance forces the image into
the fixed subgroup of \(A(\theta)\) under the cyclotomic
quotient. This is \(A[p^s]\), by
\(v_p((1+p)^{-r}-1)=s\). Evaluation at a generator embeds
the possible local images into \(A[p^v]\).

The kernel of all these local evaluations is precisely
\(S_{\QQ}(\theta)\). The assertion concerns finite groups
by Proposition~\ref{prop:finite15} and the injection just
proved. Each \(A[p^u]\) has \(O\)-length \(u\). Finally
multiplication by \(p^j\) induces a surjection from the
original quotient of the two groups to the quotient of their
images, which proves the assertion at every cutoff.
\end{proof}

For the fixed odd divisors \(p\) of \(\operatorname{ord}\psi\),
the sum of these discrepancies over all moving components and
\(1\le r<t(x)\), with their weights \(\log p\), is
\(O_\psi([E_x:\QQ]\,t(x))\). Indeed \(K\) depends only on
\(p,\psi\), so \(\ell\), the ramified primes and their indices
are bounded independently of \(x\). There are at most
\([E_x:\QQ]\) characters at each such prime. Every moving
character has \(\chi\ne1\) by Corollary~\ref{cor:allmoving663}.
This bound compares the rational strict groups with the
\(P\)-fixed subgroups. The lengths of the rational groups
themselves still occur in the estimate.

\begin{theorem}[Lower bounds from the ramified primes]
\leavevmode\label{thm:genusranks665}
Let \(h\) satisfy \(1\le h\le s\), and let \(N_h\) count
the primes \(b\in T\) with \(u_b\ge h\). The number of
elementary divisors of \(M^P\) of length at least \(h\) over
\(O\) is at least
\[
        \max\{0,N_h-1+
                   {\bf1}_{\{\chi(D_p)=1,\ \ell\ge h\}}\}.
\]
\end{theorem}
\begin{proof}
For a finite \(O\)-module \(B\), denote the number in question by
\(m_h(B)\). The elementary divisor theorem identifies it both
with \(\dim_k (p^{h-1}B)[p]\) and with
\(\dim_k(p^{h-1}B/p^hB)\). Thus a submodule cannot have a
larger \(m_h\) than the ambient group. Quotienting by a cyclic
\(O\)-submodule decreases \(m_h\) by at most one. Indeed in
the second description the kernel is the image modulo \(p^h B\)
of the intersection of that cyclic submodule with
\(p^{h-1}B\), so it has dimension at most one.

At this precision the coefficient on the residual field is
trivial. Theorem~\ref{thm:tameray662} applies if \(\chi(D_p)\ne1\).
Use the ray class group there with the specified tame relations.
If \(\chi(D_p)=1\), use Theorem~\ref{thm:everycyclic663} with
\(U=P\), so \(F=H\). The local calculation of
Corollary~\ref{cor:splitray663} works through precision \(h\)
as follows. At each \(w\mid p\) the principal-unit direction
has quotient \((1+p\ZZ_p)/(1+p^{h+1}\ZZ_p)\), cyclic of order
\(p^h\). The value on the uniformizer idele is killed by
\(p^\ell\), since its image generates the cyclic kernel of
order \(p^\ell\) in the local compositum. Thus in the split
case one starts with the additional modulus
\(p^{h+1}\mathcal O_H\) and the local uniformizer relations.
If \(h\le\ell\), the latter are automatically satisfied by
characters into \(R_h\).

In either computation take the \(p\)-primary part, extend
scalars to \(O\) and apply \(e_{\widetilde\chi}\).
This is an exact operation since \(|\Delta_H|\) is a unit.
The residue-unit term in the ray class sequence has one
cyclic summand on this component from each \(b\in T\), as
in Corollary~\ref{cor:unitrows662}. Its length after imposing
the tame multiples is \(u_b\), since
\(p^{u_b}\mid b-1\). In the split case there is also one
summand of length \(h\) from the principal units above \(p\).
Indeed those places are freely permuted by \(\Delta_H\).
The order of the residue field's multiplicative group at \(p\)
is prime to \(p\).

The projected global unit group is cyclic over \(O\).
To see this, Dirichlet's unit theorem identifies the real
representation on the logarithms of the units of the totally
real \(H\) with the hyperplane of coordinate sum zero in the
permutation representation on its embeddings. This is the regular
representation less its trivial character. Its rational character
therefore contains \(\widetilde\chi\) once after extending
scalars to the splitting field. The torsion in the unit group has order at most two. Hence the
completed unit group on this component is a free rank-one
\(O\)-module. In the ray sequence
the image of the residue units, with the tame multiples imposed,
is thus a submodule of the ray quotient obtained from the
listed summands by one cyclic relation. More explicitly, its
kernel in that residue-unit presentation is the image of
the global units. This lowers \(m_h\) by at most one.

In the split case with \(\ell<h\), imposing the multiples
\(p^\ell\) of all uniformizer classes costs at most one further
cyclic relation on the projected ray group. The classes are an
orbit indexed by the freely permuted places over \(p\), so on
this component they are generated by the projection of any
one of them. This potential loss cancels the extra local
principal-unit summand in the lower bound. For \(h\le\ell\)
there is no loss at this step. Finally the equivariant
characters into \(R_h\) are the \(R_h\)-dual of the projected
group through precision \(h\). A cyclic summand contributes
to \(m_h\) of this dual exactly when its original length is
at least \(h\). The ray formulas identify that dual with
\((M^P)[p^h]\), which has the same \(m_h\) as \(M^P\).
\end{proof}

\begin{corollary}[A nonzero moving component at every odd weight]
\leavevmode\label{cor:movingcubic665}
Use the cubic character of conductor \(3367\) from
Section~\ref{sec:cyclicray657}, and take \(p=3\).
On every sufficiently distant interval of
Theorem~\ref{thm:conductor232} choose
\(\alpha=\eta_{-d_0}\), the quadratic character of the second
auxiliary field. At every positive odd weight in the required
range the \(P\)-fixed subgroup of this moving component has
at least two elementary divisors. Its \(\ZZ_3\)-length is at
least that of \(S_{\QQ}(\theta)\) and
the difference is at most three. Moreover
\(\Cl(K\QQ(\sqrt{3d_0}))\) has \(3\)-rank at least four.
\end{corollary}
\begin{proof}
At \(p=3\) the residual cyclotomic character is the quadratic
character of \(\QQ(\sqrt{-3})\). At odd \(r\) its product with
\(\eta_{-d_0}\) cuts out \(H=\QQ(\sqrt{3d_0})\), and \(3\)
ramifies in \(H\). The only primes ramified in \(K\) are
\(7,13,37\), all totally ramified of index three. They are
\(1\bmod 3\), so they split in \(\QQ(\sqrt{-3})\).
The choice in Theorem~\ref{thm:long232} gives
\(d_0\equiv-1\) at every odd prime up to a bound tending to
infinity along the chosen intervals. Hence each of these
three fixed primes also splits in \(\QQ(\sqrt{-d_0})\), and
therefore in \(H\). Thus \(N_1=3\) and \(v=0\).
The lower bound follows from Theorem~\ref{thm:genusranks665}
at \(h=1\), and the length comparison from
Theorem~\ref{thm:descentcost665}. The character takes values
in \(\ZZ_3\), so its projector is already defined there.
Extension to the unramified splitting ring preserves length
and the number of elementary divisors.

By Theorem~\ref{thm:wholeblock661}, the nontrivial quadratic
component of the \(3\)-class group of \(KH\) has rank at least
two, by duality at precision one. On the trivial component
extension of ideals embeds the \(3\)-class group of \(K\), since
the norm of an extended class is its square, and squaring is an
automorphism on a finite \(3\)-group. This supplies two further
independent directions by Lemma~\ref{lem:genuscyclic657}.
The two components are disjoint since two is a unit at three.
\end{proof}

This example occurs on the period fields themselves, with the
selected moving prime. Vanishing of all moving \(3\)-components
at all weights would be too strong. The estimate in
Section~\ref{sec:missing} concerns their weighted lengths
beyond the prescribed truncation.


\subsection{All weights on a family of cyclic sextic character fields}
\label{sec:sexticstrict669}

Write \(\omega\) for the order-two character of conductor three,
viewed also as the Teichm\"uller part of the \(3\)-adic cyclotomic
character \(\kappa_3\). The principal-unit part is
\(\langle\kappa_3\rangle=\kappa_3/\omega\). We construct primitive
characters for which the complementary-sign strict group over
the entire fixed character field and every real dyadic layer
can be computed at every positive weight. Their moving quadratic
components on the same intervals also give lower bounds of
arbitrarily large rank.

\begin{lemma}[Realizing the cubic residue conditions]
\leavevmode\label{lem:cubicconductors669}
For each \(m\ge2\) there are infinitely many squarefree products
of \(m\) distinct primes \(b_i\equiv13\bmod24\) with cubic
characters \(\nu_i:(\ZZ/b_i\ZZ)^\times\to\mathbb F_3\) such that
the matrix of Section~\ref{sec:cyclicray657} has rank \(m-1\)
and \(\sum_i\nu_i(3)\ne0\).
\end{lemma}
\begin{proof}
Start with \(b_1=13\) and discrete logarithm to base \(2\)
modulo three. The element \(2\) has order \(12\) modulo \(13\)
and \(2^4=3\bmod13\). Thus \(\nu_1(3)=1\), and the one-row
matrix is zero. Suppose \(b_1,\dots,b_j\) and their characters
have been chosen. Let \(K_i\) be the cubic field of \(\nu_i\).
Consider the finite Galois extensions
\[
 M=\QQ(\zeta_3,\sqrt[3]{3},\sqrt[3]{b_1},\ldots,\sqrt[3]{b_j}),
 \qquad B=\QQ(\zeta_3,\zeta_8)K_1\cdots K_j ,
\]
using real cube roots. The group over \(\QQ(\zeta_3)\) in \(M\)
embeds in a product of groups of cube roots of unity, by its
multipliers on the indicated radicals. Complex conjugation
inverts this group. Its commutators with conjugation exhaust
the group since squaring is invertible on a group of exponent
three. The maximal abelian subfield of \(M/\QQ\) is therefore
\(\QQ(\zeta_3)\), so \(M\cap B=\QQ(\zeta_3)\).

The cubic factors of \(B\) are independent. A product of their
characters trivial on every inertia group has exponent zero at
each factor, since the \(i\)-th character alone is ramified at
\(b_i\). Their compositum has degree \(3^j\) and is disjoint from
\(\QQ(\zeta_3,\zeta_8)\), which has degree a power of two.
Choose an automorphism of \(B\) fixing \(\zeta_3\), acting
by fifth power on \(\zeta_8\), and having cubic coordinates
\((0,\ldots,0,1)\). It agrees on \(M\cap B\) with the identity
of \(M\), so the two extend together to \(MB\). Chebotarev's
theorem for this fixed finite extension gives infinitely many
unramified rational primes with this Frobenius conjugacy class.
Every representative has the specified restrictions, since \(B\)
is abelian and the restriction to \(M\) is the identity.
See \href{https://www.jmilne.org/math/CourseNotes/ANT.pdf}{Milne,
\emph{Algebraic Number Theory}, Chapter~8} for the density theorem.

Take any such new prime \(b\) outside the previous list. We have
\(b\equiv1\bmod3,\ b\equiv5\bmod8\). Splitting in \(M\) makes
\(3,b_1,\ldots,b_j\) cubes modulo \(b\), so any new surjection
\(\nu_{j+1}\) vanishes on them. The new last column is zero in
the old rows and their diagonals do not change. The new row has
entry one in column \(j\), entry minus one in column \(j+1\),
and its other entries are zero. Thus its addition raises the
rank by one while the sum of the character values at three is
unchanged. Iteration, with infinitely many choices at the last
step, proves the assertion.
\end{proof}

Let \(\varphi=\zeta_3^{\,\sum_i\nu_i}\) have cyclic cubic field
\(K\), and use the primitive Dirichlet character
\(\psi=\omega\varphi\) of conductor \(3\prod_i b_i\). The local
conductors are coprime. The kernel of the product is the
intersection of the two kernels because their character orders
are coprime. Its field is \(K\QQ(\sqrt{-3})\). For \(J\ge0\) put \(L_J=K\QQ(\sqrt{-3})C_J\).
Let \(c\) be complex conjugation in this abelian field.

\begin{theorem}[The complementary sign at every weight]
\leavevmode\label{thm:sexticweights669}
For every \(J\ge0\) and every integer \(r\ge1\),
\[
  \mathcal S_{3,r}(L_J)^{\,c=(-1)^r}\simeq(\ZZ/3\ZZ)^{m-1}.
\]
The real dyadic Galois group acts trivially on this part.
Restriction identifies it with the same sign in the subgroup
of \(\mathcal S_{3,r}(E_x)\) fixed by
\(\operatorname{Gal}(E_x/L_J)\), for \(J=j_x\) on the common
fields of \(\psi\).
\end{theorem}
\begin{proof}
The prime-to-three group over \(K\) is the product of the real
dyadic group and \(\operatorname{Gal}(\QQ(\sqrt{-3})/\QQ)\).
After an unramified scalar extension containing its character
values, the characters of the prescribed sign are
\(\alpha=\omega^r\beta\) with \(\beta\) a real dyadic character.
Twisting by \(\alpha^{-1}\) and descending to \(K\) by averaging
gives coefficient character
\(\langle\kappa_3\rangle^{-r}\beta^{-1}\).
For \(\beta\ne1\) its residual field is \(C_i\) if \(\beta\)
has order \(2^i\). The primes \(3,b_1,\ldots,b_m\) generate the
Frobenius group of each \(C_i\). Indeed their residues are
\(3\) or \(5\bmod8\), so their images in the cyclic dyadic
group already generate its first quotient of order two.
The class number of every \(C_i\) is prime to three by
Fukuda and Komatsu~\cite[Corollary~1.1]{V-III-weberFK524}.
Theorem~\ref{thm:cyclicdescent661} now gives zero on this
entire summand.

On the remaining summand let
\(\vartheta=\langle\kappa_3\rangle^{-r}\). We first show that
multiplication by three kills the strict group over \(K\)
with this coefficient. In the integral character comparison of
Theorem~\ref{thm:integralblock656} the rational scalar characters
are \(\xi=\vartheta\varphi^{-s}\), \(s=0,1,2\), over
\(O=\ZZ_3[\zeta_3]\). Their strong groups vanish. Here are the
details for odd as well as even \(r\). Put \(\pi=\zeta_3-1\).
Each \(\xi\) is one modulo \(\pi\). In the coefficient sequence
for multiplication by \(\pi\), the line to quotient out of
finite-character cohomology at a global or local group \(G\)
is generated by \((\xi(g)-1)/\pi^{a_G}\bmod\pi\), with
\(a_G=\min_{g\in G}v_\pi(\xi(g)-1)\). If \(\xi|_G=1\) the
connecting image is zero, by divisibility of the trivial module.
For nontrivial restriction the formula follows by lifting
a generator of the invariants of \(O[1/3]/O\) through
multiplication by \(\pi\) and taking its coboundary.
Since \(\xi(g)^2-1=(\xi(g)-1)(\xi(g)+1)\) and
\(\xi(g)+1\equiv2\bmod\pi\), the minimum and the line are the
same for \(\xi^2\). The residual character is trivial in both
cases. Thus the same finite characters pass the local tests
and the same line is divided out globally. This includes the
inertia tests away from the prescribed primes. It identifies the
\(\pi\)-torsion of the two strong groups. Now
\(\xi^2=\kappa_3^{-2r}\varphi^{-2s}\), whose strong group is
zero by Lemma~\ref{lem:scalarcyclic657} at the positive even weight
\(2r\). A nonzero \(\pi\)-primary group has nonzero \(\pi\)-torsion,
so the assertion follows for \(\xi\). The forward and reverse
integral character maps in Theorem~\ref{thm:integralblock656}
compose to multiplication by \([K:\QQ]=3\). Factoring this
composite through the zero scalar groups proves the exponent
bound, also over \(\ZZ_3\) by faithful flatness.

To compute the resulting group take coefficients over \(\ZZ_3\).
Its residual character is trivial. The prime \(w\) above three
is unique and unramified of degree three by the condition
\(\sum\nu_i(3)\ne0\). The local and global images of the
principal-unit part of the cyclotomic character are both \(1+3\ZZ_3\), since the
local cyclotomic tower is totally ramified and the local
extension at \(w\) is unramified. The global connecting line
therefore restricts isomorphically to the line at \(w\).
Subtracting its unique required multiple gives each strict
class a unique finite character zero at \(w\). At other places
the coefficient is unramified, so the finite character is
unramified on inertia. Conversely an unramified finite character
zero at \(w\) satisfies these conditions in the divisible
coefficient. By class field theory this group is the dual
modulo three of \(\Cl(K)/\langle[w]\rangle\). The ideal \(w=(3)\)
is principal. Lemma~\ref{lem:genuscyclic657} gives its dimension
\(m-1\). This proves the formula and the dyadic action assertion.
Finally the extension from \(L_J\) to the full field has degree
a power of two. Restriction and averaging preserve the global
classes and detect the local kernels, giving the last assertion.
\end{proof}

\begin{corollary}[A moving quadratic component at the same weights]
\leavevmode\label{cor:sexticmoving669}
For these characters, take
\(\alpha=\omega^{r-1}\eta_{-d_0}\) on \(E_x\) on the intervals of
Theorem~\ref{thm:conductor232}. For every sufficiently distant
interval and all \(1\le r<t(x)\), the subgroup fixed by the
cubic Sylow group \(P\) in this component has at least \(m-1\)
elementary divisors. Its length exceeds the rational strict
length by an integer in \([0,m]\). The \(3\)-class rank of
\(K\QQ(\sqrt{3d_0})\) is at least \(2(m-1)\).
\end{corollary}
\begin{proof}
This \(\alpha\) has the required sign. At every weight the
residual character of \(\kappa_3^{-r}\alpha^{-1}\) equals
\(\omega\eta_{-d_0}\), cutting out \(H=\QQ(\sqrt{3d_0})\).
It is ramified at three. The \(b_i\) are all \(1\bmod3\);
they split in \(\QQ(\sqrt{-3})\) and also in \(\QQ(\sqrt{-d_0})\)
since \(d_0\equiv-1\bmod b_i\) on the selected intervals by
Theorem~\ref{thm:long232}. Thus all \(m\) primes split in \(H\)
and have ramification index three in \(K\). Apply
Theorems~\ref{thm:genusranks665} and~\ref{thm:descentcost665}
with \(h=1\) and zero loss at three. Theorem~\ref{thm:wholeblock661}
at precision one puts at least \(m-1\) directions on the
nontrivial quadratic part of the class group of \(KH\) by
duality. Extension of ideals from \(K\) injects its \(3\)-part
into the trivial quadratic part since norm after extension is
squaring. Lemma~\ref{lem:genuscyclic657} supplies \(m-1\) further
directions there.
\end{proof}

Thus the contribution from the complementary sign fixed by
the moving factors in the theorem is exactly
\((t-1)(m-1)\log3\) for an integer cutoff \(t\ge1\). Its image
under multiplication by three is zero. The length of the
moving component in the corollary still depends on the rational
strict group and the part outside \(P\)-invariants.

\section{Cyclotomic indices at finite primes}\label{nav:g-III-3}

\subsection{Cyclotomic indices and strict orders}\label{nav:III-46}
Let \(\ell\) be an odd prime, \(\alpha\) a primitive Dirichlet
character of conductor \(f\), and \(r\ge1\). Assume that
\(\ell\nmid f\varphi(f)\) and \(\alpha(-1)=(-1)^r\).
Take \(O=\ZZ_\ell[\alpha]\) in a fixed algebraic closure,
\(J=\operatorname{Frac}(O)\), and \(k=r+1\).
This ring is an unramified discrete valuation ring. Write
\[
 T=O(\alpha\kappa_\ell^k),\qquad
 A=(J/O)(\alpha^{-1}\kappa_\ell^{-r}).
\]
Let \(S(A)\) be the strict group over \(\QQ\) of
Theorem~\ref{thm:char}. Put \(U=\operatorname{Spec}\ZZ[1/\ell]\).
For \(H^i(U,T)\) use the direct image from the generic point,
so that the stalk at a finite prime of \(U\) consists of the
inertia invariants. Cohomology of the integral lattice is
understood by passage to the inverse limit of finite coefficients.

Choose compatible primitive roots \(\zeta_m\) and a generator
\(t_\alpha\) of \(O(\alpha)\). For \(h\ge1\) take the Kummer
class of \(1-\zeta_{f\ell^h}\) over \(\QQ(\zeta_{f\ell^h})\)
at exponent \(\ell^h\), and tensor it with
\(\zeta_{\ell^h}^{\otimes r}\otimes t_\alpha\). The result has
coefficients in \(T/\ell^hT\). Corestrict it to \(\QQ\).
These are the reductions of the Soul\'e--Deligne class
\(c_{\alpha,k}\in H^1(U,T)\). Indeed the norm of the unit from
level \(h+1\) to \(h\) is \(1-\zeta_{f\ell^h}\), by multiplying
\(1-X\) over the roots of \(X^\ell-\zeta_{f\ell^h}\).
The tensor factors at precision \(\ell^h\) are already defined
over level \(h\). This gives compatibility by the projection
formula. The unit is integral outside \(\ell\), including
when \(f=1\). This is the construction in
\href{https://arxiv.org/pdf/math/0101071\#page=22}{Huber and
Kings, Definitions 3.1.1--3.1.3}.

\begin{theorem}[Length of the strict group in a cyclotomic lattice]\leavevmode
\label{thm:circindex267}
The module \(H^1(U,T)\) is free of rank one over \(O\). Write
\(c_{\alpha,k}=a u\) for an \(O\)-basis \(u\). Then \(a\ne0\)
and
\[
 v_\ell(a)=\operatorname{length}_O S(A)
       +\operatorname{length}_O H^0(\QQ_\ell,A)
       -\operatorname{length}_O H^0(\QQ,A).
\]
In particular \(v_\ell(a)\) bounds the full length from above.
For \(\ell>2\) and \(0<r<\ell-1\), the two invariant groups
vanish and equality holds with \(\operatorname{length}_O S(A)\).
\end{theorem}
\begin{proof}
Let \(D=\QQ(\zeta_f)\) and \(G=\operatorname{Gal}(D/\QQ)\).
Over \(D\) the coefficient \(T\) is the ordinary Tate lattice
with a chosen basis for its finite character. Cohomology on
\(U\) is the \(G\)-invariant part of the cohomology of
\(\mathcal O_D[1/\ell]\) with this character included in the
action. In fact under the finite map of spectra the pushforward
at a geometric point has one summand for each point above it
and has no higher direct images, by proper base change at
finite coefficients. The invariants of that permutation
sheaf tensored with \(T\) have generic fiber \(T\) and stalk
\(T^{I_v}\) at \(v\). Averaging by \(|G|^{-1}\) is integral
and exact. Taking inverse limits gives the assertion for \(T\).
The same argument identifies the finite-character descent
globally and locally in the duality sequence of
Theorem~\ref{thm:wild98}. More explicitly, dualizing sends
the isotypic factor to its inverse character, and summing over
the places over \(\ell\) gives the invariants for one
decomposition group. The trace pairing identifies the
\(O\)-dual into \(J/O\) with the usual Pontryagin dual here.
Thus \(H^2(U,T)\) is finite and its length is the right side
of the asserted formula. The finiteness also follows from
its inclusion as a summand of the \(\ell\)-adic even \(K\)-group.
All statements apply to \(D=\QQ\) at an odd prime.

Complex conjugation acts on \(T\) by \(-1\).
Consequently \(H^0(U,T/\ell T)=0\), and the multiplication
sequence shows that \(H^1(U,T)\) has no \(\ell\)-torsion.
Its rank is one. Indeed the global Euler characteristic over
\(\QQ\) is minus the dimension of the odd subspace, and both
\(H^0(U,T)\) and \(H^2(U,T)\) have rank zero. Equivalently,
one may calculate on the open set also removing \(f\).
At each prime removed there the primitive character is
nontrivial on inertia of order prime to \(\ell\) in its image.
The inertia invariants of \(T\otimes J\) and its \((-1)\)-twist
are zero, so the unramified-extension stalk and the degree-one
inertia term make no rational contribution. The finite generation
in the global Euler characteristic theorem now gives the stated
free rank.

We use the integral main conjecture for Dirichlet characters in
the form proved by Huber and Kings. For \(\ell\) odd and
\(\alpha(-1)=(-1)^{k-1}\), their cyclotomic class generates
the inverse determinant of the global Iwasawa cohomology
complex. Specialization at the trivial character of the
cyclotomic \(\ZZ_\ell\)-extension gives the inverse determinant
of \(R\Gamma(U,T)\).
\href{https://arxiv.org/pdf/math/0101071\#page=36}{Huber and Kings,
Theorem 4.4.1 and the specialization argument of Theorem 5.1.1}.

Here is why specialization gives the actual class in degree
one. Derived specialization is the triangle for multiplication
by \(\gamma-1\) on the inverse-limit complex, where \(\gamma\)
generates the cyclotomic tower. This follows by Shapiro's lemma
from the sequence for multiplication by \(\gamma-1\) on the
completed group ring. The cokernel in degree two is the finite
\(H^2(U,T)\), since degree three vanishes at odd \(\ell\).
After localizing the Iwasawa algebra at \((\gamma-1)\),
Nakayama's lemma therefore kills degree two. Degree one has
no \((\gamma-1)\)-torsion there, by the same triangle and the
vanishing in degree zero downstairs. The determinant consequently
specializes on the rational cohomology line by taking the base
layer of the norm-compatible classes. In Definition 4.1.7 of
the cited paper this layer is a corestriction from \(\QQ(\zeta_\ell)\).
At precision \(\ell^h\), taking the norm from \(f\ell^{h+1}\)
to \(f\ell^h\) first gives exactly the class defined above.
The alternative normalization removing the Euler factor at
\(\ell\) changes the lattice generator by a unit since \(k>1\).

Over a discrete valuation ring the inverse determinant of
this complex, inside \(H^1(U,T)\otimes J\), is
\(\ell^{\operatorname{length}_O H^2(U,T)}H^1(U,T)\).
To see the sign, a cyclic summand \(O/\ell^j\) in degree two
is represented by \(O\xrightarrow{\ell^j}O\) in degrees one
and two. The inverse determinant basis maps to \(\ell^j\).
Multiplying over the elementary divisors proves the formula
and shows \(a\ne0\). Global invariants inject into local ones.
If \(0<r<\ell-1\), tame cyclotomic inertia at \(\ell\) contains
a scalar whose \((-r)\)-th power differs from one modulo \(\ell\),
and \(\alpha\) is unramified there. This kills both invariant
groups.
\end{proof}

\subsection{The distribution of the cyclotomic valuations}\label{nav:III-47}
Fix the interval, list of triples and precisions of
Theorem~\ref{thm:commonlength267}. In particular the original
period fields, their splitting conditions and the power-residue
conditions are included in the field \(K\) of that theorem.
Write \(j=(\ell,r,[\alpha])\) for a triple, with one
representative of its residual Frobenius orbit, and put \(\ell_j=\ell\). Write
\(h_j=h_\ell\), \(b_j=[\mathbb F_\ell(\bar\alpha):\mathbb F_\ell]\)
and \(d_j=\operatorname{length}_{O_j}S(A_j)\), where
\(O_j=\mathbb Z_\ell[\alpha]\). At a prime \(s\) split in \(K\) set
\[
 W_j(s)=w_{h_j}(D_{s,h_j}(\alpha,r)),\qquad
 m_j=\min(d_j,h_j).
\]
The moment, its valuation and the convention \(w_h(0)=h\)
are those of Proposition~\ref{prop:moment267}. Probability
below means relative natural density among the rational
primes split in \(K\).

\begin{theorem}[Joint distribution at every precision]\leavevmode
\label{thm:laws268}
For any integers \(0\le k_j\le h_j\),
\[
 \Pr\bigl(W_j(s)\ge k_j\ \text{for all }j\bigr)
       =\prod_j\ell_j^{-b_j\max(0,k_j-m_j)} .
\]
Thus the valuations from the different triples are independent,
including different twists at the same coefficient prime. The
whole law of the excess \(W_j-m_j\) is that of the valuation
of a uniform element modulo \(\ell^{h_j-m_j}\), with truncation
at \(h_j-m_j\).
\end{theorem}
\begin{proof}
Choose one integral basis class \(u_j\) as in
Theorem~\ref{thm:circindex267} for each triple. The character
\(\alpha\kappa_\ell^{r+1}\) is trivial on \(G_K\) at
precision \(\ell^{h_j}\). A cocycle for \(u_j\), restricted
to \(G_K\) and reduced at that precision, is consequently an
additive homomorphism into \(O_j/\ell^{h_j}\). Its value on
this subgroup is independent of the cocycle representative.

The joint map of these homomorphisms is onto. To see this,
fix \(\ell\) and write \(M=\prod_{j:\,\ell_j=\ell}
O_j/\ell^{h_j}\), where the product is over this prime.
The argument of Theorem~\ref{thm:commonlength267} shows that
the image modulo \(\ell\) is \(M/\ell M\). It applies to
this \(K\), including its prescribed radicals. In detail, a
primitive residual class stays nonzero over \(K\). Its
coefficient module is simple over \(\mathbb F_\ell\), and its
conjugates span that simple module. Distinct listed twists
have different characters on inertia at \(\ell\). Within one
twist the simple modules are distinct precisely after the
Frobenius conjugates have been identified. The idempotents for
these residual characters therefore separate the factors.
If \(I\) is the image at the stated precisions, this proves
\(M=I+\ell M\). The abelian group \(M/I\) is killed by a
power of \(\ell\) and equals \(\ell(M/I)\), so it is zero.
The joint image across all \(\ell\) is the product too. In a
subgroup of a finite abelian product of primary groups one
isolates each primary coordinate by integer multiplication,
using the Chinese remainder theorem.

The kernel of the joint map on \(G_K\) is open and normal in
\(G_{\mathbb Q}\). Indeed conjugation by a global element acts
on each evaluation by the character of its coefficient module.
The term from changing a cocycle representative vanishes on
\(G_K\). The kernel thus defines a finite Galois field over
\(\mathbb Q\). Its group over \(K\) is the full additive product
just calculated. Every condition on the valuations alone
defines a conjugacy-stable subset. Chebotarev's theorem
therefore gives its fraction in this product as its relative
density among split primes. This argument is only applied to
valuation conditions, so choices of generators of residue
fields have no bearing on the density.

At the large primes in the list the invariant correction in
Theorem~\ref{thm:circindex267} is zero. The cyclotomic class is
therefore \(\ell^{d_j}\) times a unit times \(u_j\).
Proposition~\ref{prop:moment267} identifies its local valuation
with \(W_j(s)\). For a uniform element of \(O_j/\ell^{h_j}\)
the condition of belonging to \(\ell^aO_j/\ell^{h_j}\) has
probability \(\ell^{-b_j a}\) for \(0\le a\le h_j\).
Multiplication by the class coordinate and truncation give the
claimed formula for each factor. Multiplication of these
fractions gives the joint law.
\end{proof}

\subsection{The cyclotomic index at an arbitrary odd prime}\label{nav:III-48}
Let \(p\) be odd, \(f\mid N\), and \(\alpha\) a primitive
Dirichlet character of conductor \(f\), viewed \(p\)-adically.
Take \(r\ge1\) with \(\alpha(-1)=(-1)^r\), and set \(k=r+1\).
Let \(O=\mathbb Z_p[\alpha]\), \(J=\operatorname{Frac}(O)\),
and fix a uniformizer \(\pi\). This is the integer ring of a
local cyclotomic field. Write \(b=[O/\pi:\mathbb F_p]\),
\(e=v_\pi(p)\), \(T=O(\alpha\kappa_p^k)\), and
\(U_N=\operatorname{Spec}\mathbb Z[1/(Np)]\).

\begin{theorem}[The index on the common open set]\leavevmode
\label{thm:openindex269}
The module \(H^1(U_N,T)\) is free of rank one. The
Soul\'e--Deligne class \(c=c_{\alpha,k}\) restricts there to a
nonzero multiple of a basis. If \(d=d_{p,r,\alpha}(N)\)
is the \(\pi\)-valuation of its coordinate, then
\[
 \operatorname{length}_O H^2(U_N,T)
  =d+\sum_{\substack{q\mid N,\ q\ {\rm prime}\\q\ne p,\ q\nmid f}}
                      v_\pi(1-\alpha(q)q^r).
\]
The formula includes \(p\) dividing the conductor, the order
of \(\alpha\), or the degree of its full cyclotomic field.
\end{theorem}
\begin{proof}
On \(U_N\) the lattice is unramified and its cohomology is
restricted-ramification Galois cohomology. In rational degree
two it is a summand after passing to the field cut out by
\(\alpha\). Over that field the coefficient is a Tate twist,
whose rational degree two vanishes by the even \(K\)-group
comparison of Theorem~\ref{thm:wild98}. At odd \(p\) this
comparison applies also over real fields. Inverting the
additional primes preserves rational vanishing, since the
positive \(K\)-groups of their finite residue fields are finite.
Global finite generation now gives finite integral degree two.
Degree zero vanishes even modulo \(\pi\), since complex
conjugation acts by \(-1\). The multiplication sequence for
\(\pi\) gives torsion freeness in degree one. The rational
Euler characteristic over \(\mathbb Q\) is minus the dimension
of the odd subspace. It gives rank one.

Put \(U=\operatorname{Spec}\mathbb Z[1/p]\) and extend \(T\)
to \(U\) by the ordinary generic-point direct image. We use
the integral main conjecture of Huber and Kings in its full
Dirichlet-character form. For this parity and \(p\ne2\),
the cyclotomic class generates the inverse determinant of the
global cyclotomic Iwasawa complex. Their theorem includes the
primes dividing \(\varphi(f)\). Their coefficient field can
be enlarged beyond \(J\), since an equality of lattices over
a discrete valuation ring descends under finite faithfully
flat extension. The defining class at precision \(p^h\) is
the corestriction from \(\mathbb Q(\zeta_{fp^h})\) of the
Kummer class of \(1-\zeta_{fp^h}\) tensored with
\(\zeta_{p^h}^{\otimes r}\) and the character basis. This
description also holds when \(p\mid f\), by transitivity of
corestriction. Removing a redundant \(p\)-level uses the norm
of \(1-\zeta_{fp^{h+1}}\) to \(1-\zeta_{fp^h}\).
\href{https://arxiv.org/pdf/math/0101071\#page=22}{Huber and
Kings, Definitions 3.1.1--3.1.3, 4.1.7 and Theorem 4.4.1}.

Here specialization gives the actual class on the rational
cohomology line. Indeed the tower coefficient modulo
\(\gamma-1\) is the base coefficient for a topological
generator \(\gamma\). Shapiro's lemma and the multiplication
triangle give derived specialization to \(R\Gamma(U,T)\).
The tower is unramified away from \(p\), so the inertia
invariants in the direct image on \(U\) commute with adjoining
this free tower coefficient. Rationally the specialized
complex is concentrated in degree one. To see this last point
one may open the primes \(q\mid N\) as below. Their rational
local complexes are acyclic, since \(r>0\). At the
height-one prime \((\gamma-1)\), cancel unit entries in a
bounded free representative of the Iwasawa complex.
The remaining differentials reduce to zero. Its rational
special fiber has dimension one in degree one, leaving a
single free generator in that degree. The determinant
generator therefore specializes on the rational line by
taking the base layer of the norm-compatible classes. This
is \(c\) by the corestriction formula. In particular \(c\ne0\).

We compute the effect of removing \(q\mid N,\ q\ne p\).
The cone from the direct-image complex on \(U\) to the
complex on \(U_N\) has at \(q\) the term
\(R\Gamma(\mathbb F_q,H^1(I_q,T))[-1]\). This follows
from the stalks of the derived direct image and the
\(p\)-cohomological dimension one of inertia \(I_q\).
It holds first for the finite coefficient systems and then
in their inverse limit, by stabilization of the finite
images. When \(q\nmid f\), the inertia term is \(T(-1)\).
Hence the cone has zero degree one and a finite degree two of
length \(v_\pi(1-\alpha(q)q^r)\). When \(q\mid f\), the
inertia term is finite. Indeed first take invariants of wild
inertia and of the prime-to-\(p\) part of tame inertia.
These are exact operations. If their invariants survive, one
computes degree one by the cokernel of \(\tau-1\) on the
lattice for a generator of the tame pro-\(p\) quotient.
A nontrivial ramified character remaining here has
\(\alpha(\tau)\ne1\), giving a finite cokernel. On this finite
module the kernel and cokernel of Frobenius minus one have
the same order. Thus its cone has inverse determinant of
valuation zero.

Under the rational identification of the two degree-one lines,
the inverse determinant for the open complex is consequently
generated by \(c\) times the product of the unramified local
factors, up to a unit. For \(U_N\) the same inverse determinant
is \(\pi^{\operatorname{length}_O H^2(U_N,T)}H^1(U_N,T)\).
Indeed restricted Galois cohomology over \(\mathbb Q\) has degree at
most two at odd \(p\), and a cyclic group in degree two
contributes its elementary divisor to that lattice. Comparing
valuations proves the formula.
\end{proof}

\subsection{Cyclotomic classes that reduce to rational Kummer classes}\label{nav:III-50}
Use the rank-one group and its cyclotomic class \(c\) in
Theorem~\ref{thm:openindex269}. For each odd \(p\), write \(O\)
for the character integer ring, \(\pi\) for its uniformizer,
\(e=v_\pi(p)\), \(b=[O/\pi:\mathbb F_p]\) and
\(d=d_{p,r,\alpha}(N)\). Set
\[
       u=\min_{g\in G_{\mathbb Q}}v_\pi
                          \bigl(\alpha(g)\kappa_p(g)^r-1\bigr).
\]
The value is finite since \(r\ge1\) and the cyclotomic character
has infinite image on the kernel of the finite character.
Here the conductor \(f\) of \(\alpha\) divides \(N\).

\begin{theorem}[The rational Kummer contributions at all odd primes]\leavevmode
\label{thm:kummer273}
If \(u\ge1\) and \(f\) is a power of \(p\), including \(f=1\),
then \(d=0\). If \(u\ge1\) and some prime other than \(p\) divides
\(f\), then \(d\ge u\). On a field \(E_x\) of
Theorem~\ref{thm:conductor232} the total contribution of these
known initial \(u\) layers in the second case, with their weights
\(b\log p\), is at most
\[
   [E_x:\mathbb Q]\,t(x)
       \sum_{\substack{p\mid\operatorname{ord}(\psi)\\p\ {\rm odd}}}
                          \frac{\log p}{p-1}.
\]
In particular the trivial-character index vanishes for every odd
\(p\) and every positive multiple \(r\) of \(p-1\).
\end{theorem}
\begin{proof}
For \(1\le v\le u\) put \(h=\lceil v/e\rceil\). The line with
character \(\alpha\kappa_p^r\) is trivial modulo \(\pi^v\).
Choose its fixed basis given by the tensor of the character
basis and the chosen Tate basis. The defining formula for \(c\)
at precision \(p^h\) is the corestriction of the Kummer class
of \(1-\zeta_{fp^h}\) tensored with that basis. Reduction to
\(O/\pi^v\) makes the basis a constant in the corestriction.
Thus the reduced global class is the Kummer class of
\[
      N_{\mathbb Q(\zeta_{fp^h})/\mathbb Q}(1-\zeta_{fp^h})
                            =\Phi_{fp^h}(1)
\]
with the indicated coefficient extension. Indeed corestriction
of the Kummer connecting homomorphism is the connecting
homomorphism of the field norm. This follows by multiplying
the conjugates of a chosen \(p^h\)-th root, whose \(p^h\)-th
power is that norm. The finite-level class and twist convention
are those in Theorem~\ref{thm:openindex269}.
\href{https://arxiv.org/abs/math/0101071}{A.~Huber and G.~Kings,
Sections 3.1 and 4.1}.

For \(n>1\) one has \(\prod_{j\mid n,\ j>1}\Phi_j(1)=n\) by
dividing \(X^n-1\) by \(X-1\) before evaluating at one.
The prime-power factors each have value their underlying
prime by their geometric sums. Induction on \(n\) gives value
one at every \(n\) divisible by two distinct primes.
In the prime-power case take \(v=1\). The Kummer class of \(p\)
over \(\mathbb Q\) modulo \(p\) is nonzero by its valuation at
\(p\), and extension to the residue field \(O/\pi\) is injective
on this class by a choice of an \(\mathbb F_p\)-linear retraction.
In the other case the Kummer class of one vanishes for all
\(v\le u\). Inflation in degree one from the restricted
Galois group to the absolute group is injective for these
coefficients. The multiplication sequence of \(O(\alpha
\kappa_p^{r+1})\) shows that the kernel of reduction modulo
\(\pi^v\) on \(H^1\) is exactly \(\pi^vH^1\). Since \(H^1\)
is free on one generator, the two assertions about \(d\)
follow.

In the second case take a prime \(q\mid f\), \(q\ne p\).
The cyclotomic character at \(p\) is trivial on \(q\)-inertia.
The image of \(\alpha\) there is nontrivial by primitivity.
If \(u\ge1\), it reduces to one and hence has \(p\)-power
order. Its image contains a primitive \(p\)-th root \(\zeta\).
Then \(u\le v_\pi(\zeta-1)=e/(p-1)\), by the factorization of
\(\Phi_p(1)=p\) and the equality of the valuations of the
factors. In particular \(p\mid [E_x:\mathbb Q]\). The degree
over the field of the fixed character \(\psi\) is a power of
two on these intervals, so \(p\mid\operatorname{ord}(\psi)\).
For each \(p,r\) the sum of \(be\) over the character orbits
is at most \([E_x:\mathbb Q]\) by character decomposition in
characteristic zero. This proves the uniform estimate. For
\(\alpha=1\) and \(p-1\mid r\), the residual cyclotomic
character to the power \(r\) is trivial, giving the last assertion.
\end{proof}

\subsection{A vanishing result for fields ramified at one prime}\label{nav:III-51}
\begin{theorem}[The strict group with an inert coefficient prime]\leavevmode
\label{thm:oneprime278}
Let \(p\) be odd and \(a\ge1\). Let \(B\) be the degree-\(p^a\)
subfield of \(\QQ(\zeta_q)\) for a prime \(q\ne p\), and assume
that \(p\) is inert in \(B\). For every positive multiple \(r\)
of \(p-1\),
\[
                       \mathcal S_{p,r}(B)=0 .
\]
For each \(p,a\) there are infinitely many such \(q\).
On the common intervals for any character cutting out \(B\),
the \(\Gal(E_x/B)\)-invariant strict group has the same vanishing.
For \(p=3\) its complementary-sign part vanishes at every twist.
\end{theorem}
\begin{proof}
We first show that \(p\) does not divide the class number of \(B\).
Let \(H\) be its Hilbert \(p\)-class field. This is Galois over
\(\QQ\), since its defining property is preserved by every
automorphism of \(B\). Write \(G=\Gal(H/\QQ)\), a finite
\(p\)-group. Only \(q\) can ramify. Its inertia group \(I\)
maps isomorphically onto \(\Gal(B/\QQ)\), by total ramification
in \(B/\QQ\) and absence of ramification in \(H/B\).
If \(I\) were proper in \(G\), it would be contained in a
maximal subgroup. Every maximal subgroup of a finite \(p\)-group
is normal of index \(p\). The corresponding cyclic degree-\(p\)
field would be unramified at \(q\), because that subgroup contains
all conjugates of \(I\). It would be unramified at every other
finite prime, and also at infinity since \(p\) is odd. The
class number of \(\QQ\) excludes this by class field theory.
Thus \(I=G\), giving \(H=B\).

Put \(A=(\QQ_p/\ZZ_p)(-r)\). The cyclotomic image over \(B\)
and over its unique completion \(B_w\) at \(p\) is
\(\ZZ_p^\times\). Indeed \(\QQ_p(\mu_{p^n})/\QQ_p\) is totally
ramified, while \(B_w/\QQ_p\) is unramified. The global image
contains this local image. Thus \(A^{G_B}=A^{G_{B_w}}\)
inside \(A\). For \(G=G_{B,S_p}\) the exact sequence for
multiplication by \(p\) gives
\[
  0\longrightarrow A^G/pA^G\longrightarrow H^1(G,A[p])
      \longrightarrow H^1(G,A)[p]\longrightarrow0 .
\]
The local row is the same, and restriction identifies the two
left terms because there is just one prime above \(p\).
It follows that the localization kernels in the middle and
on the right are isomorphic. Since \(p-1\mid r\), \(A[p]\)
is the constant module \(\mathbb F_p\). The middle kernel
therefore consists of homomorphisms from the ideal class group
of \(B\) to \(\mathbb F_p\) that kill the class of \(w\), by
the Hilbert class field theorem. It is zero. Any nonzero
\(p\)-primary torsion group has an element of order \(p\),
which proves the claimed vanishing.

For existence let \(C=\QQ(\mu_{p^a})\).
The real degree-\(p\) field \(\QQ(p^{1/p})\) is not normal,
since it omits the nonreal roots of \(X^p-p\). It cannot be
a subfield of the abelian field \(C\).
Since \(\mu_p\subset C\), Kummer theory gives a cyclic
degree-\(p\) extension \(C(p^{1/p})/C\), Galois over \(\QQ\)
as a splitting field after adjoining \(C\).
Chebotarev applied to the conjugacy class of a nonidentity
element over \(C\) yields infinitely many unramified rational
primes \(q\) split in \(C\) for which \(p\) is a non-\(p\)-th
power in \(\mathbb F_q^\times\). In particular \(p^a\mid q-1\).
The image of \(p\) generates the quotient of order \(p^a\) of
the cyclic group \(\mathbb F_q^\times\), since its exponent
on a generator is prime to \(p\). This gives the asserted
inertness.

In Theorem~\ref{thm:conductor232} the relative degree
\([E_x:B]\) is a power of two for these characters.
Restriction identifies the strict group with the invariant
part upstairs. Indeed averaging gives the isomorphism in
global cohomology, and restriction is injective in local
cohomology and on inertia by the prime-to-\(p\) degrees.
This preserves the local conditions in both directions.
For \(p=3\) the even twists are the multiples of \(p-1\).
The field \(B\) is real, so complex conjugation acts trivially
on its group, and the complementary sign is negative for
each odd twist.
\end{proof}

\subsection{Two Frobenius evaluations at conductor 1799}\label{nav:III-52}
Use \(\chi=(\,\cdot\,/257)\), the cubic character
\(\lambda\bmod7\) with \(\lambda(3)=\zeta_3\), and
\(U=\operatorname{Spec}\ZZ[1/(3\cdot7\cdot257)]\) as in
Theorem~\ref{thm:sextic278}. For \(r\) positive and even, let
\(L_0,L_1\) be the two lines of
Theorem~\ref{thm:intersection279} for \(\chi,\chi\lambda\).
Their common residual coefficient is
\(\bar T=\mathbb F_3(\chi\omega)\), where
\(\omega=\kappa_3\bmod3\).

\begin{proposition}[The two residue lines in each even weight]\leavevmode
\label{prop:lines279}
One has \(L_0=L_1\) exactly when \(6\mid r\).
In each of the other two even residue classes the lines
span \(H^1(U,\bar T)\).
\end{proposition}
\begin{proof}
Write \(c_0,c_1\) for the two cyclotomic classes and
\(\pi=\zeta_3-1\). Their coordinate valuations are both
one by Theorem~\ref{thm:sextic278}. Thus the reductions of
\(c_0/3\) and \(c_1/\pi\) generate \(L_0,L_1\).
The cohomology group \(H^1(U,\bar T)\) has dimension two.
Indeed the multiplication-by-three sequence for the first
lattice gives one dimension from its free rank-one group in
degree one and one from its degree-two group of order three.
This calculation applies to a single residual module for all
even \(r\), since \(\omega^r=1\).

Use the two rational primes and generators \((s,g)\) in the
table. Both \(s-1\) are multiples of \(9\cdot7\cdot257\).
At each prime choose the place in the common cyclotomic
field where a compatible primitive root \(\zeta_n\) reduces to
\(g^{(s-1)/n}\) for each \(n\mid9\cdot7\cdot257\).
These assignments come from the single root of the largest
order. Evaluating a residual cocycle at the corresponding
arithmetic Frobenius gives a linear form, since that Frobenius
acts trivially on \(\bar T\) and coboundaries give zero.
Choose the tensor bases in all even weights by identifying
\(\zeta_3^{\otimes(r+1)}\) with the fixed residual basis
and taking the finite character basis for \(\lambda\) to reduce to one.
These local choices define the same pair of linear forms when
comparing the two classes and all three weights.

For the first class compute the six entries \(B_i\) of
Theorem~\ref{thm:quadratic278}, in the order \(1,2,4,5,7,8\).
For the second compute the three \(C_j\) of
Theorem~\ref{thm:sextic278}. Use the indicated roots \(w_0,w_1\)
of orders \(2313,5397\) for the respective computations.
\[
\begin{array}{r|r|r|r|c|c}
 s&g&w_0&w_1&(B_i)\bmod9&(C_j)\bmod3\\ \hline
 129529&37&63673&59753&(7,6,2,2,6,7)&(2,1,0)\\
 161911&3&131000&88463&(5,6,4,4,6,5)&(2,1,0)
\end{array}
\]
Trial division to \(359\) and \(402\), respectively, tests
primality. The residues \(g^{(s-1)/d}\) at the distinct prime
divisors \(d\) of \(s-1\), taken in increasing order, are
\((129528,120445,60611,35657)\) for \(d=2,3,7,257\) and
\((161910,47683,63009,55156,24972)\) for \(d=2,3,5,7,257\).
For an alternative direct verification of the two last columns,
list \(g^j\bmod s\) for \(0\le j<s-1\), labeling each residue
by \(j\). For each defining sum multiply \(1-w^a\) over the
indices belonging to its specified residue class and having
\(\chi(a)=1\), then divide by the same product for \(\chi(a)=-1\).
Read the label of the result modulo nine or three, respectively.
This gives the entries by additivity of the logarithm.

In the chosen bases the first Frobenius coordinate before
division is \(\sum_iB_i i^r\bmod9\). For the second it is
\(\sum_jC_j(1+j\pi)\bmod\pi^2\). To check the normalizations,
restriction of the corestriction is a sum over \(a\) prime to
the conductor times the coefficient modulus. Conjugation
sends \(1-\zeta_n\) to \(1-\zeta_n^a\) and acts on the extra
tensor factor by \(\chi(a)a^r\), with an additional
\(\lambda(a)\) for the second class. The Kummer cocycle of a
residue \(z\) at Frobenius has value \(z^{(s-1)/p^h}\) for
modulus \(p^h\). Its coordinate in the reduction of
\(\zeta_{p^h}\) is exactly \(\operatorname{ind}_g z\bmod p^h\).
This gives both formulas with compatible bases. The two
Frobenius elements already act trivially at the precision of
each numerator, so that division gives the evaluations of
the integral quotients in residual cohomology.

In the order of the table the coordinate vector for \(c_1/\pi\)
is \((1,1)\) at every even twist. For \(c_0/3\) the vectors at
\(r\equiv0,2,4\bmod6\) are, respectively,
\((1,1),(0,2),(2,0)\) in \(\mathbb F_3^2\).
In residue class two the two vectors are independent, so the
evaluation map from the common two-dimensional cohomology
group is injective. Equality and independence of the lines
in the other residue classes can therefore be read from the
same two coordinates, which proves the result.
\end{proof}

\subsection{Class groups and strict classes throughout one cyclotomic extension}\label{nav:III-53}
Put \(F=\QQ(\sqrt{257},\zeta_7+\zeta_7^{-1})\),
\(Q_m=\QQ(\zeta_{3^{m+1}})^+\), \(F_m=FQ_m\) for \(m\ge0\),
and \(L=\bigcup_{m\ge0}F_m\). Thus \(Q_0=\QQ\).
The prime \(3\) is inert of degree six in \(F\). Its Frobenius
has order two on \(\QQ(\sqrt{257})\) and order three on
\(\QQ(\zeta_7)^+\). The extension \(Q_m/\QQ\) is totally
ramified at \(3\), so \(F\cap Q_m=\QQ\) and \(F_m\) has just
one prime over \(3\). Let \(\Gamma=\Gal(L/F)\) and take
\(\gamma\) with image \(4\) in \(1+3\ZZ_3\).
Here the projection of the cyclotomic character to principal
units identifies \(\Gamma\) with \(1+3\ZZ_3\). Indeed
\(v_3(4^{3^i}-1)=1+i\) by induction and binomial expansion.

Let \(C_m\) be the \(3\)-part of the ideal class group of
\(\mathcal O_{F_m}[1/3]\), and \(X=\varprojlim C_m\) under
the ideal norms. The group \(\Gal(L/\QQ)\) acts on \(X\).
For \(\delta\) in this group, write \(u(\delta)\in1+3\ZZ_3\)
for its principal cyclotomic character.

\begin{theorem}[Untwisted and twisted quotients of the same class group limit]\leavevmode
\label{thm:control280}
For every \(m\), \(C_m=\Cl(F_m)\{3\}\) and
\(C_m\simeq X/(\gamma^{3^m}-1)X\).
Let \(D\subset L\) be a number field such that \(L/D\) is
unramified away from \(3\) and \(\Gal(L/D)\simeq\ZZ_3\)
has generator \(\delta\) with \(u(\delta)\ne1\).
For \(r\ge1\) put
\(A_r^{\rm pr}=(\QQ_3/\ZZ_3)(\langle\kappa_3\rangle^{-r})\),
where the brackets remove the sign of the \(3\)-adic unit.
With the same strict local conditions one has
\[
         S(D,A_r^{\rm pr})^\vee
                  \ \simeq\ X/(\delta-u(\delta)^{-r})X .
\]
At even \(r\) this is the dual of \(\mathcal S_{3,r}(D)\).
The identification uses a fixed basis of the coefficient on \(G_L\).
\end{theorem}
\begin{proof}
The prime over \(3\) in \(Q_m\) is generated by
\((1-\zeta_{3^{m+1}})(1-\zeta_{3^{m+1}}^{-1})\).
Its norm to \(\QQ\) is three, by the cyclotomic polynomial at
one. Over \(Q_m\) the extension \(F_m\) is unramified at this
prime with a single prime above it, so the same element
generates that prime in \(F_m\). Inverting it preserves the
class group.

Let \(H_j/F_j\) be the class field of \(C_j\), so that the
prime over \(3\) splits completely there. For \(j\ge m\),
\(H_m\cap F_j=F_m\), by total ramification in \(F_j/F_m\).
Base change gives \(H_mF_j\subset H_j\). Under reciprocity
restriction of automorphisms gives the norm \(C_j\to C_m\),
which is therefore surjective. In the extension \(H_j/F_m\)
the decomposition group at the single prime over \(3\) maps
isomorphically onto \(\Gal(F_j/F_m)\). Thus
\(\Gal(H_j/F_m)\) is a semidirect product of \(C_j\) with
this cyclic decomposition group. Its largest abelian quotient
in which the prime splits is obtained by killing a generator
of the latter group and its conjugates. It is
\(C_j/(\gamma^{3^m}-1)C_j\). The corresponding field is
unramified everywhere, split at \(3\), and contains \(H_m\),
so is exactly \(H_m\). This proves the finite quotient formula.
Passing to \(X\) preserves the kernel description. Indeed all
projections from \(X\) are surjective. Given compatible elements
in the kernels, the conditions on a preimage under
\(\gamma^{3^m}-1\) at finitely many levels can be met at their
largest level. Compactness supplies a preimage at all levels.

For the twisted formula put \(A=A_r^{\rm pr}\). Its action
factors through the principal-unit cyclotomic quotient and
is therefore trivial on \(G_L\). At even \(r\) the sign of a
\(3\)-adic unit disappears on raising it to the \(r\)-th power.
The procyclic quotient \(\Gal(L/D)\) has zero degree two on
\(A\) and degree one \(A/(u(\delta)^{-r}-1)A=0\), by
divisibility. Thus inflation--restriction identifies degree
one over \(D\) with the equivariant homomorphisms on \(G_L\).
The same local restriction at \(3\) is injective. Its
decomposition quotient is all of \(\Gal(L/D)\), since each
finite layer has just one prime over \(3\). Away from \(3\)
inertia does not change in \(L/D\).

To see which homomorphisms occur, take a strict cocycle over
\(D\). Its image is finite by continuity and compactness.
On the local decomposition group at \(3\) it is the coboundary
of some element of \(A\). Choose \(j\) with \(F_j\supset D\)
so large that \(G_{F_j}\) acts trivially on all these finitely
many elements. The restriction of the cocycle to \(G_{F_j}\)
is then a homomorphism unramified everywhere and zero at
the prime over \(3\). It factors through \(C_j\). Conversely
the class fields \(H_jL/L\) are nested, with groups \(C_j\)
and transition maps the norms. They realize \(X\) as a quotient
of \(G_L\). A continuous equivariant homomorphism from \(X\)
to \(A\) is unramified away from \(3\) and zero at \(3\) on
\(G_L\), since it factors through some \(C_j\). Its unique
global lift is strict by the local injectivity and inertia
comparison above. Equivariance means that it kills
\((\delta-u(\delta)^{-r})X\). This image is compact and closed.
Pontryagin duality now gives the claimed quotient.
\end{proof}

\subsection{An exact strict order at conductor 257}\label{nav:II-58}
\begin{theorem}[Every positive even twist of the quadratic field]\leavevmode
\label{thm:quadratic278}
Let \(K=\QQ(\sqrt{257})\). For every positive even \(r\),
\(\mathcal S_{3,r}(K)\) is cyclic of order three.
For the primitive quadratic character of conductor \(257\),
the full logarithmic contribution at \(3\) of the complementary
groups fixed by \(\Gal(E_x/K)\), summed over \(1\le r<t(x)\),
is exactly \(\lfloor(t(x)-1)/2\rfloor\log3\) on every common interval.
\end{theorem}
\begin{proof}
The discriminant of \(K\) is \(257\). The real quadratic
Minkowski bound \(\sqrt{257}/2<9\) shows that the primes of
norm at most eight generate the class group. The rational primes
\(3,5,7\) are inert, by the quadratic residue test. Only the
two primes over \(2\) remain, since \(257\equiv1\bmod8\).
Their classes are inverses. The element
\((15+\sqrt{257})/2\) has norm \(-8\) and odd trace,
so only one of these two primes divides it. Its ideal is the
cube of that prime. Neither prime is principal. Indeed if \(z\)
had absolute norm two, multiply it by a power of the unit
\(\epsilon=16+\sqrt{257}\), of norm \(-1\), so that in the
positive embedding its absolute value \(u\) lies between
\(\sqrt{2/\epsilon}\) and \(\sqrt{2\epsilon}\).
Then \(u+2/u\le\sqrt{2\epsilon}+\sqrt{2/\epsilon}<\sqrt{257}\).
The difference of the two embeddings of an algebraic integer
is an integral multiple of \(\sqrt{257}\). The bound forces
the difference to be zero, in which case its norm is an integer
square. This proves that \(\Cl(K)\) has order three.

There is one prime over \(3\), and its class is principal.
Since it is unramified, the global and local cyclotomic
images are both \(\ZZ_3^\times\). The finite-coefficient
localization diagram in the proof of Theorem~\ref{thm:oneprime278}
therefore identifies \(\mathcal S_{3,r}(K)[3]\) with
\(\operatorname{Hom}(\Cl(K),\mathbb F_3)\) for every even \(r\).
The same diagram over \(\QQ\) gives zero. By quadratic descent
the full strict group over \(K\) is thus its nontrivial
character summand, namely the strict group over \(\QQ\) for
\((\QQ_3/\ZZ_3)(\chi\kappa_3^{-r})\), where
\(\chi(a)=(a/257)\) is the Legendre character. In particular
this summand has an element of order three.

It remains to bound its full order. Since \(\chi(3)=-1\), the
residual character is nontrivial on the decomposition group at
\(3\). Theorem~\ref{thm:circindex267} gives its length as the
\(3\)-adic valuation of the coordinate of the cyclotomic class.
We evaluate this class modulo \(9\) at the single prime
\(s=23131=10\cdot257\cdot9+1\).
Take \(g=3\), \(w=g^{10}=12787\bmod s\).
Primality follows by trial division through \(152\).
For the distinct prime divisors \(2,3,5,257\) of \(s-1\)
the residues \(g^{(s-1)/d}\) are respectively
\(23130,19250,3529,15632\). Thus \(g\) generates
\(\mathbb F_s^\times\).

For \(i\in(\ZZ/9)^\times\) put
\[
 B_i=\sum_{\substack{1\le a<2313,\ (a,2313)=1\\a\equiv i\bmod9}}
             \chi(a)\operatorname{ind}_g(1-w^a)\pmod9 .
\]
Here and below \(\operatorname{ind}_g\) is the discrete logarithm.
The entries, in the displayed order, are
\[
 \begin{array}{c|rrrrrr}
 i&1&2&4&5&7&8\\ \hline
 B_i&4&1&7&7&1&4
 \end{array}
\]
The following recurrence reproduces the table using only
modular multiplication. Initialize the six entries to zero
and \(y=1\). For \(a=1,\ldots,2312\), replace \(y\) by
\(yw\bmod s\) and skip the entry if \(3\mid a\) or \(257\mid a\).
Read \(j\) from the nine-entry table
\((w^{257j}\bmod s,j)\), \(0\le j<9\), at the key
\((1-y)^{(s-1)/9}\bmod s\). Add \(j\) to \(B_{a\bmod9}\)
when \(a^{128}=1\bmod257\), and subtract it otherwise.
By the Kummer corestriction formula of
Proposition~\ref{prop:residue269}, the Frobenius coordinate is
\(\sum_iB_i i^r\bmod9\), up to a unit. The powers depend only
on \(r\bmod6\). For \(r\equiv0,2,4\bmod6\), all three sums
are \(6\bmod9\). A class divisible by nine would evaluate to
zero. Its coordinate consequently has valuation at most one,
and the element of order three found above gives equality.

The fields for \(\chi\) have degree a power of two over \(K\).
Restriction and averaging identify the invariant strict group
with the one over \(K\), also at all the local conditions.
The odd twists have zero complementary part on the real
field \(K\). Counting even integers \(1\le r<t(x)\) proves
the formula.
\end{proof}

\subsection{Class numbers in every layer of the cyclotomic tower at three}\label{nav:II-60}
\begin{theorem}[An exact class group and strict group in every layer]\leavevmode
\label{thm:tower280}
For \(F_m,L,\gamma,X\) of Theorem~\ref{thm:control280} one has
\[
     X\simeq\ZZ/9\ZZ,\qquad \gamma|_X=1,\qquad
     \Cl(F_m)\{3\}\simeq\ZZ/9\ZZ\quad(m\ge0).
\]
All the norm maps on these primary class groups are isomorphisms.
For \(m\ge1\), \(\mathcal S_{3,r}(F_m)\simeq\ZZ/9\ZZ\) at
every positive even twist. The cyclic group fixing
\(\QQ(\sqrt{257})\) in \(F\), acting on \(L\) by fixing the
cyclotomic tower \(Q_\infty=\bigcup_m Q_m\), acts faithfully on \(X\).
The involution \(\sigma\) of \(F/\QQ(\zeta_7)^+\), also
extended to fix \(Q_\infty\), acts on \(X\) by inversion.
\end{theorem}
\begin{proof}
Give \(X\) its \(\ZZ_3[[T]]\)-module structure via \(T=\gamma-1\).
The action exists on each finite norm quotient since \(\gamma\)
has \(3\)-power order there and \(T\) is nilpotent modulo three.
The twisted quotient for \(D=F\) evaluates this module at
\(z(r)=4^{-r}-1\). By Theorem~\ref{thm:orders279}, it has order
three for \(r=2\) and nine for \(r=6,12,18\). The valuation
of \(z(2)\) is one. Since \(4^6\equiv19\bmod27\), the values
\(z(6),z(12),z(18)\) are \(9,18,0\bmod27\), respectively.
Theorem~\ref{thm:specializations280} with \(b=2\) proves the
assertions about \(X\) from these four twists.

The untwisted quotient in Theorem~\ref{thm:control280} is
therefore \(X\) for every \(m\). These identifications use the
natural projections, so identify the norms as well. In the
twisted quotient the generator is \(\gamma^{3^m}\). The scalar
\(4^{-3^m r}\) is \(1\bmod9\) when \(m\ge1\). This proves the
assertion about strict groups.

For the last assertion take \(r=6\), and let \(\tau\) generate
the cubic group in the statement. The quotient in
Theorem~\ref{thm:control280} for \(D=F\) is all of \(X\).
It identifies the action on the dual strict group with that
on \(X\), because \(\tau\) can be lifted while fixing all
\(3\)-power roots of unity. Indeed \(F\) is unramified at
three and is disjoint from the full \(3\)-power cyclotomic
extension. In the notation of Theorem~\ref{thm:orders279} the
same dual group is \(M=H^2(U,T_R)\).
Right exactness in degree two gives
\[
            M/(X_0-1)M=
              H^2\bigl(U,\ZZ_3(\chi\kappa_3^{7})\bigr),
\]
where \(X_0\) denotes the generator in the finite cubic group
ring \(R\). This follows by factoring multiplication by
\(X_0-1\) through its image as in that theorem.
The group on the right has order three, while \(M\) has order
nine. Under Shapiro's identification \(X_0\) acts by the cubic
generator or its inverse, according to the choice of regular
action. The class-field identification and duality commute with
field automorphisms, so \(\tau\) acts nontrivially on \(X\).
As it has order three it acts faithfully. In the same twist
the \(\sigma\)-invariant strict group over \(F\) is the strict
group over \(\QQ(\zeta_7)^+\), by quadratic descent, and the
latter vanishes by Theorem~\ref{thm:oneprime278}. The lift of
\(\sigma\) also fixes the cyclotomic coefficients. Thus the
invariants on \(X\) vanish, using the integral plus projector
in the dual. Since two is invertible, \(\sigma\) acts by
\(-1\) on all of \(X\).
\end{proof}

\section{Positive twists and explicit strict orders}\label{nav:g-III-4}

\subsection{Local conditions for the positive cyclotomic twist}\label{nav:II-71}
Let \(E=E_x\) be the abelian period field of
Theorem~\ref{thm:conductor232}, with cyclotomic modulus \(N\),
degree \(g\), and \(t=\lfloor x^{401}\rfloor\).
Take a prime \(\ell>2t+3\) with \(\ell\nmid Ng\).
Fix \(1\le r<t\) and a character \(\alpha\) of
\(\operatorname{Gal}(E/\QQ)\) with \(\alpha(c)=(-1)^r\).
Let \(O=\ZZ_\ell[\alpha]\) denote the ring generated by its
values on the Galois group. This is an unramified integer ring,
with fraction field \(J\) and residue degree \(b_\alpha\).
Write \(\kappa\) for the cyclotomic character and
\(\Psi=\alpha\kappa^r\).

Use classes over \(\QQ\) unramified away from \(\ell\), meaning
that they restrict to zero in inertia cohomology there.
Define \(P\) for \((J/O)(\Psi)\) and \(C\) for its opposite
Tate twist \((J/O)(\alpha^{-1}\kappa^{1-r})\) by the following
conditions at \(\ell\). For \(r\ge2\) allow all local classes
in \(P\) and require zero in \(C\). For \(r=1\) use Kummer
classes of units in \(P\) and unramified classes in \(C\).
Here the Kummer condition is read after the unramified extension
cut out locally by \(\alpha\), or equivalently by requiring
zero valuation on inertia. Use the same conditions with finite
coefficients \(O/\ell^j\), and denote the resulting groups by
\(P_j,C_j\).

\begin{theorem}[The positive twist and its finite dual orders]\leavevmode
\label{thm:positive306}
For every \(j\ge1\),
\[
       |P[\ell^j]|=|P_j|=|C_j|=|C[\ell^j]|.
\]
Both modules are finite and have the same elementary divisors
over \(O\). For \(r\ge2\), \(C\) is the \(\alpha\)-component
of the strict group on \(E\) at twist \(r-1\), after extending
scalars to \(O\). For \(r=1\) it is the corresponding
component of
\(\operatorname{Hom}(\mathrm{Cl}(E),J/O)\), with its natural
contragredient Galois action.
\end{theorem}
\begin{proof}
The pairing of \(O/\ell^j\) with itself induced by the
unramified trace and division by \(\ell^j\) is perfect into
\(\QQ_\ell/\ZZ_\ell\). Thus the second finite module is the
Cartier dual of \(M=(O/\ell^j)(\Psi)\). Away from \(\ell\)
unramified local conditions are orthogonal complements. Their
cup product vanishes by inflation from the procyclic unramified
quotient. Each has size equal to that of its local invariants,
by the kernel and cokernel orders of Frobenius minus one on the
inertia invariants. Local duality and the local Euler
characteristic formula give the complementary orders.

At \(\ell\) and \(r\ge2\) both invariant groups vanish, since
the inertia exponents \(r,1-r\) are nonzero modulo \(\ell-1\).
The whole local \(H^1(M)\) consequently has order \(|M|\) and
is orthogonal to zero. For \(r=1\), pass to the unramified
extension \(L/\QQ_\ell\) killing \(\alpha\), of degree prime
to \(\ell\). Restriction identifies local classes with the
appropriate twisted invariants. The unit condition pairs
trivially with unramified characters by local reciprocity.
The logarithm on principal units identifies their pro-\(\ell\)
completion with \(\ell\mathcal O_L\). A normal basis of the
residue field lifts to an integral normal basis of
\(\mathcal O_L\), since its determinant is a unit. Taking
the twisted invariants of that regular module tensored with
\(O/\ell^j\) gives one copy of \(O/\ell^j\). The local
invariants of \(M\) vanish. The local Euler formula now gives
the full orthogonal complement the size of the invariants of
the dual, precisely the size of its unramified classes.

Apply finite Poitou--Tate duality and the global Euler
characteristic formula in their Selmer order form. With
orthogonal local conditions \(L_v,L_v^\perp\), this form gives
the ratio of the two Selmer orders as the ratio of the global
invariant orders times
\(\prod_v |L_v|/|H^0(\QQ_v,M)|\), with zero local condition at
the real place. It follows by taking orders in the finite
localization sequence and substituting the global Euler
characteristic for the alternating product of global orders.
\cite[Chapter I]{V-III-II-duality13}
The global invariants vanish on both sides. For the dual at
\(r=1\) this uses \(\alpha(c)=-1\). The finite-place factors
are one away from \(\ell\), and the factor at \(\ell\) is
\(|M|\). Complex conjugation acts trivially on \(M\), producing
the factor \(1/|M|\). This proves the middle equality.

The coefficient sequence identifies global finite classes with
the \(\ell^j\)-torsion of the divisible-coefficient cohomology,
since its degree-zero invariants vanish. At inertia away from
\(\ell\) the invariants of the divisible module are either all
of it or zero, as the inertial character has order prime to
\(\ell\). They are divisible, so the finite classes inject
locally on inertia. For the strict condition at \(\ell\) in
twists \(r\ge2\) the local invariants are zero. At \(r=1\)
the unramified condition is detected on inertia, where the
dual coefficient is trivial and divisible. On the Kummer side
the valuation on the finite module embeds in the valuation on
the direct limit, by naturality under
\(O/\ell^j\hookrightarrow J/O\). These facts give both outer
equalities.

Restriction to \(E\) and the character projection identify
\(C\) as stated, since \(g\) is a unit in \(O\). In particular
restriction is injective on local cohomology when tested over
the corresponding local extension or inertia subgroup, by
restriction followed by corestriction.
For \(r=1\) the coefficient over \(E\) is constant and the
everywhere unramified characters factor through the ideal
class group. For \(r\ge2\) use Theorems~\ref{thm:char}
and~\ref{thm:critical13}; the parity at twist \(r-1\) is
\((-1)^r\). In both cases \(C\) is finite. The finite-level
equalities bound the increasing union forming \(P\), so it is
finite as well. The successive differences of the lengths of
the \(\ell^j\)-torsion count the elementary divisors of length
at least \(j\). This proves their equality.
\end{proof}

\subsection{Primitive special values for the positive twist}\label{nav:II-72}
Keep the characters, coefficient prime, and modules \(P,C\)
of Theorem~\ref{thm:positive306}. Write \(\alpha^*\) for
the associated primitive Dirichlet character, of conductor
\(f_\alpha\), using arithmetic Frobenius. Evaluate its
algebraic special values at the chosen \(\ell\)-adic embedding
and normalize \(v_\ell(\ell)=1\).

\begin{theorem}[The order in each character component]\leavevmode
\label{thm:primitive307}
For every \(1\le r<t\) under those hypotheses,
\[
    \operatorname{length}_O P=\operatorname{length}_O C
                     =v_\ell\big(L(1-r,\alpha^*)\big).
\]
In particular the positive twist has a nonzero class precisely when
the primitive value has positive valuation, including at \(r=1\)
when \(\alpha(\ell)=1\).
\end{theorem}
\begin{proof}
Let \(K\) be the fixed field of the kernel of \(\alpha\).
Its degree is prime to \(\ell\). Character descent in
Theorem~\ref{thm:positive306} applies equally to this field,
since its proof uses only prime-to-degree restriction. The
character is faithful on \(G=\operatorname{Gal}(K/\QQ)\), so
it is nontrivial on inertia at every ramified prime of \(K\).
The ring \(O\) is the factor of \(\ZZ_\ell[G]\) obtained by
evaluating at this character. We use \(\alpha^{-1}\) when
passing to a contragredient module.

First suppose \(r\ge2\). At every \(w\mid\ell\) of \(K\),
the mod-\(\ell\) cyclotomic character has full inertia image
because \(K_w\) is unramified. The invariants of
\((\QQ_\ell/\ZZ_\ell)(1-r)\) therefore vanish, as
\(0<r-1<\ell-1\). By local duality
\(H^2(K_w,\ZZ_\ell(r))=0\). Poitou--Tate duality as in
Lemma~\ref{lem:fit13} identifies the dual of \(C\) with the
entire \(\alpha^{-1}\)-component of
\(H^2(\mathcal O_K[1/\ell],\ZZ_\ell(r))\) after taking
the coefficient factor. Inverting the primes over
\(f_\alpha\) does not change this component. The adjacent
terms of the localization sequence are sums of residue-field
cohomology with twist \(r-1\) at these primes. Their inertia
groups act trivially on the residue field and on the twist,
whereas \(\alpha^{-1}\) is ramified. Averaging over inertia
kills their character components.

Apply the equivariant Coates--Sinnott equality of Greither
and Popescu, in the form stated in Lemma~\ref{lem:fit13},
with the real base \(\QQ\) and the infinite and ramified
places. The required cyclotomic \(\mu\)-invariant vanishes
by Ferrero--Washington. The group of invariants of
\((\QQ_\ell/\ZZ_\ell)(r)\) over \(K\) vanishes by the same
inertia test, now using \(0<r<\ell-1\). The torsion factor in
the equality thus has unit annihilator. The parity factor
is one on \(\alpha^{-1}\). In that factor the zeroth Fitting
ideal of the group just found is generated by
\(L(1-r,\alpha^*)\). Each omitted ramified Euler factor
already equals one on a faithful character. If the prime
\(\ell\) is also omitted, its factor is an \(\ell\)-adic unit.
The valuation of the Fitting generator over a discrete
valuation ring is the length. Duality preserves this length.
\cite[Theorem 6.11]{V-III-II-gp13}

For \(r=1\) the character is odd, so \(K\) is CM. Now \(C\)
is the \(\alpha\)-component of the dual of the ordinary class
group. Choose two distinct auxiliary primes away from
\(f_\alpha\ell\) congruent to one modulo \(f_\alpha\) and to
two modulo \(\ell\). They exist by Dirichlet's theorem and
split completely in \(K\). Use their set \(T\) as smoothing
primes in the theorem of Dasgupta and Kakde on the Fitting
ideal of the dual ray class group. Two residue characteristics
satisfy the smoothing hypothesis. The ray-to-ordinary kernel
is a quotient of the product of the residue unit groups at
these primes, of orders prime to \(\ell\). The \(\ell\)-primary
modules therefore agree.
\href{https://arxiv.org/pdf/2010.00657}{Dasgupta and Kakde,
\emph{On the Brumer--Stark Conjecture}, Theorem~1.4.}

For clarity, at a ramified prime \(q\) the local ideal in
that formula is generated by the sum of the elements of the
inertia group \(I_q\) and by \(1-\sigma_q e_q\), where
\(e_q\) is their average and \(\sigma_q\) lifts Frobenius.
Its image under \(\alpha\) is the unit ideal, since the average
has image zero. The involution on the Stickelberger factor of
the dual class group makes its \(\alpha\)-value
\(L(0,\alpha^*)\prod_{q\in T}(1-q)\). The smoothing factors
are units. Taking the length again proves the formula.
Theorem~\ref{thm:positive306} gives the same length for \(P\).

For example the odd quadratic character modulo \(167\) has
\(L(0,\alpha^*)=11\). The sum of the least positive quadratic
residues modulo \(167\), obtained by reducing the squares of
\(1,\ldots,83\), is \(6012\), and the sum of the nonresidues is
\(7849\). Their difference in the order nonresidues minus
residues is \(167\cdot11\), giving the value by the first
Bernoulli polynomial. Also \(11\) is a quadratic residue.
Thus for this field, \(r=1\), and \(\ell=11\), the two groups
have order eleven even with the unit condition at the split
prime. This example uses the local hypotheses with \(t=2\).
\end{proof}

\subsection{Every strict twist at eleven in the conductor forty-one tower}\label{nav:II-76}
Let \(D\) be any finite subfield of
\(\QQ(\zeta_{41},\mu_{2^\infty})\). Write \(c\) for complex
conjugation and use the original strict local conditions on
\(\mathcal S_{11,r}(D)\) of Section~\ref{sec:coho12}. Thus the
coefficient over \(D\) is \((\QQ_{11}/\ZZ_{11})(-r)\).

\begin{theorem}[Vanishing on the full dyadic tower over conductor forty-one]\leavevmode
\label{thm:all41tower317}
For every such \(D\) and every \(r\ge1\),
\[
                  \mathcal S_{11,r}(D)^{c=(-1)^r}=0 .
\]
On any of the original intervals take
\(D_x=E_x\cap\QQ(\zeta_{41},\mu_{2^\infty})\).
If \(11\nmid[E_x:D_x]\), the \(\operatorname{Gal}(E_x/D_x)\)-
invariant part of the complementary strict group on \(E_x\)
is zero for every twist on that interval. In particular for
\(\psi\) of conductor dividing \(41\cdot2^a\) this applies to
all the classes descending to \(K_\psi C_J\subset E_x\), with
the actual dyadic layer \(C_J\) of the interval.
\end{theorem}
\begin{proof}
Every finite cyclotomic field used here has degree prime to
eleven. The character descent of Theorem~\ref{thm:char}
reduces the question to \(S_{\alpha,r}\) over \(\QQ\) for the
primitive characters inducing its characters \(\alpha\) with
\(\alpha(-1)=(-1)^r\). Its coefficient is
\((J/O)(\alpha^{-1}\kappa_{11}^{-r})\), with \(O=\ZZ_{11}[\alpha]\)
and \(J\) its fraction field. Exactness of the character
decomposition holds over the unramified coefficient integers.

First let the dyadic conductor \(F\) of \(\alpha\) be at least
\(32768\). In Proposition~\ref{prop:effect146} the maximal
\(11\)-subfield is \(\QQ\). The reflected character
\(\omega^{r+1}\alpha\) has odd conductor \(f_o\mid 11\cdot41\).
Its odd-conductor factor has order dividing \(40\), since
the orders of the two factors divide \(10\) and \(40\).
Take \(Q=121\) in that proposition. Then \(v_2(Q-1)=3\),
\(40\mid Q-1\) and \(F\ge451\cdot 4^3\). Any extra odd prime
in its Euler factors is \(11\) or \(41\), and neither is
congruent to either sign of one modulo \(F/8\). These verify
all the hypotheses of the vanishing test.

For \(F\le16384\) apply Theorem~\ref{thm:towerres317}.
It remains to treat its seven labels, each at \(r\equiv0\bmod10\)
with \(\alpha\ne1\). Take \(k=10\) and \(\chi=\alpha^{-1}\).
The character \(\bar\chi\bar\omega^{10}=\bar\chi\) is nontrivial.
If the primitive value \(B_{10,\chi}/10\) is a unit, the
interpolation of Theorem~\ref{thm:scalar231} supplies a series
\(F_1\in O[[X]]\) taking a unit value at
\(X_{1-10}\in11\ZZ_{11}\). Its constant coefficient is a
unit, hence so is \(F_1(X_{1+r})\) for these twists. The
annihilator conclusion there gives \(S_{\alpha,r}=0\).

Here are the residues of \(B_{10,\chi}/10\) in the field
\(k_0=\mathbb F_{11}[i]\) of Theorem~\ref{thm:towerres317}.
\[
 \begin{array}{r|r|r|l}
 F&l&e&B_{10,\chi}/10\bmod11\\ \hline
 8&20&1&0\\
 16&10&1&10+9i\\
 16&30&3&10+2i\\
 32&5&5&9+8i\\
 32&15&7&9+3i\\
 32&25&1&3+5i\\
 32&35&3&3+6i
 \end{array}
\]
For a check using only small integers, put \(M=41F\).
The generalized Bernoulli formula gives the table by summing
over \(1\le a<M\) prime to \(M\) the expression
\[
 \chi(a)\sum_{j=0}^{9}
           \tfrac1{10}\binom{10}{j} B_j M^{j-1} a^{10-j}.
\]
The omitted constant term sums to zero since \(\chi\ne1\).
The Bernoulli numbers in this sum are integral at \(11\),
so each summand can be reduced there. The nonzero \(B_j\)
in index order are \(1,-1/2,1/6,-1/30,1/42,-1/30\),
at \(j=0,1,2,4,6,8\). For
\(a=6^u\bmod41,\ a=(-1)^b5^v\bmod F\), use
\(\chi(a)=(-1)^{bl}C^{-lu}(\beta^{120e/d})^{-v}\) in \(k_0\),
where \(d=F/4\). This proves vanishing in six cases. In the
first case the character is the even quadratic character of
conductor eight times the quadratic character of conductor
forty-one. It cuts out \(\QQ(\sqrt{82})\) of discriminant
\(328\). Theorem~\ref{thm:splitunits317} proves vanishing for
every positive multiple of ten in this last case as well.

This proves the assertion first for full cyclotomic layers.
Restriction to a larger layer followed by corestriction is
multiplication by its relative degree, prime to eleven. It
follows that the assertion holds for their subfields.
For \(E_x/D_x\) of degree prime to eleven, averaging or
inflation--restriction gives the isomorphism from the strict
group over \(D_x\) onto the invariant subgroup over \(E_x\).
The same argument locally preserves the unramified and zero
local conditions in both directions. Conjugation commutes
with the isomorphism. The final assertion follows since
\(K_\psi C_J\) lies in the tower and the remaining extensions
have degree prime to eleven for this \(\psi\).
\end{proof}

As a consequence one can also read the full relaxed order
on the quartic branch of Corollary~\ref{cor:quartic231}.
Use its \(\chi(6)=i,\ r\equiv2\bmod10\) and
\(\alpha=\chi^{-1}\). The reductions of the two moments in
that corollary show that the value of its interpolation
series has valuation one for every such \(r\), since the
real coordinate of its quotient by eleven is seven modulo
eleven. Theorem~\ref{thm:localvalue309} gives relaxed length
one over \(\ZZ_{11}[i]\). Its entire local image is therefore
isomorphic to \(\ZZ_{11}[i]/11\), of order \(121\), since
the strict kernel has just been shown to vanish.

\subsection{Testing the order of a strict group at finite level}\label{nav:III-70}
Keep the actual field and the large prime of
Theorem~\ref{thm:simultaneous307}. Write \(S^-\) for the
whole strict group with coefficient \((J/O)(\alpha^{-1}\kappa^{-r})\),
and put \(\beta=\alpha^{-1}\). For \(j\ge1\) define the
positive weight and the integer
\[
        s_j=(\ell-1)\ell^{j-1}-r,\qquad
        v_j=v_\ell\big(L(1-s_j,\beta^*)\big).
\]
Let \(K_0\) be the cyclic field cut out by the residual
character \(\delta=\bar\alpha^{-1}\bar\kappa^{-r}\) over
\(\QQ\). Let \(h_0\) denote its ordinary class number and put
\(b=[O/\ell:\mathbb F_\ell]\).

\begin{theorem}[Finite levels of the full complementary group]\leavevmode
\label{thm:layers308}
For every \(j\ge1\),
\[
 \operatorname{length}_O S^-[\ell^j]\le v_j,\qquad
 \dim_{O/\ell} S^-[\ell]\le
       \min\left(v_1,\frac{\log h_0}{b\log\ell}\right).
\]
If \(v_j<j\) at one level, the full group already has
length at most \(v_j\). If its elementary divisors are written
as powers of \(\ell\), the number of them with exponent at
least \(j\) is at most
\(\min_{1\le i\le j}\lfloor v_i/i\rfloor\).
\end{theorem}
\begin{proof}
The finite classes for \(O/\ell^j(\delta_j)\), with
\(\delta_j=\alpha^{-1}\kappa^{-r}\bmod\ell^j\), recover
\(S^-[\ell^j]\) with exactly the unramified conditions away
from \(\ell\) and the zero condition at \(\ell\). Indeed
the inertia character at \(\ell\) has nontrivial reduction, so
both the global and local divisible-coefficient invariants
vanish. At every other prime the inertia action has order
prime to \(\ell\). Averaging over its image exhibits the
invariants as a direct summand of a divisible module.
They are therefore divisible as well. These are
the injection tests in the proof of Theorem~\ref{thm:mirror}.

Every cyclotomic unit modulo \(\ell^j\) has exponent dividing
\((\ell-1)\ell^{j-1}\). Thus this finite coefficient equals
the reduction of \(\beta\kappa^{s_j}\), and the strict classes
inject into its positive-twist Selmer group at level \(j\).
The latter uses all local classes at \(\ell\).
We spell out the permitted weights in applying
Theorems~\ref{thm:positive306} and~\ref{thm:primitive307}
here. For weights \(s\ge2\), their finite dual-order and
Fitting arguments need both \(s\) and \(1-s\) nonzero
modulo \(\ell-1\). This condition kills both local invariant
modules, the global torsion factor in the Coates--Sinnott
formula, and the local degree-two term. Prime-to-degree
character descent and the residue-field terms at ramified
primes are independent of \(s\). For \(s=s_j\) the two residues
are \(-r,1+r\), both nonzero since \(\ell>2t+3\).
Also \(s_j\ge2\) and \(\beta(c)=(-1)^{s_j}\). Thus the
positive group over divisible coefficients has length exactly
\(v_j\), giving the asserted bound for the finite subgroup.
A cyclic summand of exponent at least \(j\) contributes
\(j\) to the length at level \(j\). If \(v_j<j\), there is
no such summand, and \(S^-=S^-[\ell^j]\). Each summand
of exponent at least \(j\) contributes
\(i\) at level \(i\le j\), giving the last assertion.

For the other rank bound, restriction to \(K_0\) is injective
at level one since its degree is prime to \(\ell\). Over this
field our classes are additive homomorphisms unramified
everywhere and zero at primes above \(\ell\). They therefore
lie in the appropriate character space of the dual class group
tensored with \(O/\ell\). The values of \(\delta\) generate
\(O/\ell\) over \(\mathbb F_\ell\), since the residual
cyclotomic factor takes its values in \(\mathbb F_\ell^\times\).
Consequently the residual character has \(b\) distinct
Frobenius conjugates. They occur with equal multiplicity in
a module defined over \(\mathbb F_\ell\), and are disjoint by
prime-to-degree semisimplicity. Thus \(b\) times the dimension
in question is at most the \(\mathbb F_\ell\)-dimension of
\(\mathrm{Cl}(K_0)/\ell\), which is at most
\(\log h_0/\log\ell\).
\end{proof}

For two examples, the upper length of
Theorem~\ref{thm:primitive307} is one for the trivial character
at \((\ell,r)=(691,12)\) and for the odd quadratic character
of conductor \(167\) at \((11,1)\). In the first case
\(L(-11,\mathbf 1)=691/32760\), and the second value is
eleven as computed there. In both cases the whole
complementary strict group is zero. The last column gives
the test at \(j=1\).
\[
\begin{array}{c|r|r|r|r}
 (\ell,r)&f_\beta&s_1&B_{s_1}\bmod\ell&
                    L(1-s_1,\beta^*)\bmod\ell\\ \hline
 (691,12)&1&678&48&110\\
 (11,1)&167&9&0&3
\end{array}
\]
To reproduce the entries put \(u_0=1\) and compute
\(u_m=-\sum_{i<m}\binom{m+1}{i}u_i/(m+1)\) in the residue
field up to \(m=s_1\). All divisors are units since
\(s_1<\ell-1\). With \(f=f_\beta\), multiply
\(\sum_{a=1}^f\beta^*(a)\sum_{i=0}^{s_1}\binom{s_1}{i}
u_i(a/f)^{s_1-i}\) by \(-f^{s_1-1}/s_1\).
For \(f=167\), \(\beta^*(a)\) is the quadratic residue sign,
with value zero at \(167\). The nonzero last entries give
\(v_1=0\) independently of the upper length.

\subsection{A local correction of unbounded length at 37}\label{nav:III-71}
\begin{proposition}[The rational character at \(37\)]\leavevmode
\label{prop:rational37exact309}
For every positive even \(r\) the strict group over \(\QQ\)
for \((\QQ_{37}/\ZZ_{37})(-r)\) vanishes. For
\(r\equiv4\bmod36\), let \(\xi\in\ZZ_{37}\) be the zero in
Proposition~\ref{prop:example235}. When \(r\ne\xi\) the group
with the unrestricted condition at \(37\) is cyclic of order
\(37^{1+v_{37}(r-\xi)}\). Every positive length occurs.
If \(r=\xi\) is a positive integer on this branch, the group
with that local condition is \(\QQ_{37}/\ZZ_{37}\).
\end{proposition}
\begin{proof}
The strict vanishing also follows from Theorem~\ref{thm:towerzero267}.
For a smaller calculation in this rational case use \(s=149\).
The element \(2\) generates \(\mathbb F_{149}^\times\), since
\(2^{74}=-1\bmod149\) and \(2^4\ne1\bmod149\). Put \(w=2^4=16\).
For \(1\le a\le36\) take \(j_a\bmod37\) with
\(2^{j'}=1-w^a\bmod149\) and \(j_a=j'\bmod37\).
The following are the moments
\(D_r=\sum_{a=1}^{36}a^r j_a\bmod37\).
\[
\begin{array}{c|rrrrrrrrr}
 r&0&2&4&6&8&10&12&14&16\\ \hline
 D_r&35&3&26&35&25&23&26&3&6\\[2pt]
 r&18&20&22&24&26&28&30&32&34\\ \hline
 D_r&29&16&18&13&16&32&35&33&29
\end{array}
\]
For a direct reproduction, enumerate \(2^i\bmod149\) for
\(0\le i<148\), recording \(i\bmod37\) at the corresponding
residue. Start at \(z=1\). For each \(a=1,\ldots,36\) replace
\(z\) by \(16z\bmod149\), read the recorded entry at \(1-z\),
and add it times \(a^r\) modulo \(37\). This is exactly the
split-prime Kummer formula of Proposition~\ref{prop:moment267}
at precision one and conductor one. Every entry is nonzero,
also for the residue of positive multiples of \(36\). That
proposition and Theorem~\ref{thm:circindex267} give vanishing
of the strict group. The moments have period \(36\) in \(r\).

On the indicated irregular branch take \(B=2\) in
Theorem~\ref{thm:localvalue309}. Proposition~\ref{prop:example235}
calculates \(v_{37}(F(X_{1+r}))=1+v_{37}(r-\xi)\).
With \(S=0\) the entire relaxed group is its image in the local
line, giving the two assertions about that group. For each
\(j\ge0\) choose a residue modulo \(37^{j+1}\) at distance
\(37^{-j}\) from \(\xi\). The Chinese remainder theorem supplies
a positive \(r\) in that residue with \(r=4\bmod36\).
\end{proof}

Whenever \(37\nmid[E_x:\QQ]\), character descent in
Theorem~\ref{thm:char} identifies the same vanishing group with
the trivial-character summand over the period field.

\subsection{Positive twists with a ramified or exceptional local character}\label{nav:III-72}
Let \(p\) be odd and let \(\gamma\) be a primitive Dirichlet
character of order prime to \(p\), of conductor \(f\). The prime
\(p\) may divide \(f\). Write \(O=\ZZ_p[\gamma]\),
\(J=\operatorname{Frac}O\), \(\kappa=\kappa_p\), and let \(\omega\)
be its Teichm\"uller part. For a discrete coefficient with this
finite character part, write \(R\) for cohomology classes of degree
one over \(\QQ\) unramified at every finite prime away from \(p\),
with all local classes allowed at \(p\).

\begin{theorem}[The relaxed positive twist at every tame local type]\leavevmode
\label{thm:criticalall310}
For \(k\ge2\), \(\gamma(-1)=(-1)^k\), and
\(\bar\gamma\bar\omega^k\ne1\) as a global character,
\[
    \operatorname{length}_O R\big((J/O)(\gamma\kappa^k)\big)
                  =v_p(L(1-k,\gamma)).
\]
If \(\bar\gamma\bar\omega^m=1\) globally for a nonzero integer
\(m\), then \(R((J/O)(\gamma\kappa^m))=0\). This includes
negative \(m\).
\end{theorem}
\begin{proof}
Write \(T_k=O(\gamma\kappa^k)\), \(V_k=T_k\otimes_O J\) and
\(U=\operatorname{Spec}\ZZ[1/fp]\). At every \(q\mid f\)
other than \(p\), inertia acts nontrivially, of order prime to
\(p\), on both the coefficient modulo \(p^j\) and its Cartier
dual. Their invariants vanish. The local Euler characteristic
and local duality at \(q\ne p\) give vanishing of both groups
in degree one. Taking the direct limit shows that \(R=H^1(U,V_k/T_k)\).

Take \(L\) to be the compositum of the field cut out by \(\gamma\) with
\(\QQ(\mu_p)\). This abelian CM extension has degree prime to \(p\)
and is unramified outside \(fp\). Over \(L\) the finite character
becomes trivial. The rational degree-two group of the positive
twist is zero by the comparison with the finite even \(K\)-group
of the ring of integers. Inverting additional primes preserves
this vanishing since the residue-field terms are finite for
\(k\ge2\). Averaging descends the vanishing to \(H^2(U,V_k)\).
Also \(H^0(U,V_k)=0\) since the cyclotomic power has infinite
image. The global Euler characteristic of \(V_k\) over \(\QQ\)
is minus the dimension of its odd subspace, here zero.
Consequently \(H^1(U,V_k)=0\), and the coefficient sequence
identifies \(R\) with \(H^2(U,T_k)\).

Averaging identifies \(H^2(U,T_k)\) with the
\(\gamma^{-1}\)-factor of degree two over \(L\) with coefficient
\(\ZZ_p(k)\), extended to \(O\). Use the refined Coates--Sinnott
equality of Greither and Popescu with the infinite places and
the primes dividing \(fp\). The cyclotomic \(\mu\)-invariant
vanishes by Ferrero--Washington since \(L/\QQ\) is abelian.
In the selected factor the torsion in integral degree one is
\(H^0(\QQ,V_k/T_k)\), by the coefficient sequence and the same
descent. It vanishes because the global residual character is
nontrivial. Its annihilator is therefore the unit ideal.
The parity idempotent acts as the identity. The resulting
zeroth Fitting ideal of \(H^2(U,T_k)\) is generated by the
primitive value \(L(1-k,\gamma)\) up to a unit. At \(q\mid f\)
the primitive Euler factor is already one. At \(p\nmid f\) the
removed factor \(1-\gamma(p)p^{k-1}\) is a unit. The valuation
of a zeroth Fitting generator over \(O\) is the length.
\href{https://arxiv.org/pdf/1103.3069}{Greither and Popescu,
\emph{An equivariant main conjecture in Iwasawa theory and applications},
Theorem 6.11}.

In the last assertion reduction on roots of unity of order prime
to \(p\) is injective. Thus \(\gamma=\omega^{-m}\), \(O=\ZZ_p\),
and \(A=(\QQ_p/\ZZ_p)(\langle\kappa\rangle^m)\) is unramified
outside \(p\). Its global invariants have order \(p^{1+v_p(m)}\)
because \(1+p\) generates the principal units topologically.
The global multiplication sequence injects \(A^{G_\QQ}/pA^{G_\QQ}\)
into \(H^1(\ZZ[1/p],A[p])\), with cokernel \(R(A)[p]\).
Here \(A[p]\) is trivial, and the middle group is one-dimensional
over \(\mathbb F_p\). Indeed its characters factor through the
maximal abelian extension of \(\QQ\) unramified away from \(p\),
whose pro-\(p\) quotient is \(1+p\ZZ_p\) by Kronecker--Weber.
The injection from the group of order \(p\) fills it. Every
nonzero \(p\)-primary group has an element of order \(p\), so
\(R(A)=0\).
\end{proof}

\subsection{Two residue computations at eleven for conductor forty-one}\label{nav:III-73}
Work in \(k=\mathbb F_{11}[i]\), with \(i^2=-1\).
The powers of \(\beta=1+4i\) with exponents \(24,40,60\)
are \(3,5+3i,-1\), so \(\beta\) generates \(k^\times\).
Write \(C=\beta^3\), of order forty. The element \(6\)
generates \((\ZZ/41)^\times\), since \(6^8=10\) and
\(6^{20}=-1\bmod41\). For a primitive character \(\alpha\)
with odd conductor \(f=1\) or \(41\), use \(0\le l<40\)
for the exponent such that its odd-conductor factor at \(6\)
reduces to \(C^l\). Thus \(f=1\) precisely for \(l=0\).
Let \(F\) denote its dyadic conductor, and put
\(r_0=r\bmod10\) in \(\{0,\ldots,9\}\). We always impose
\(\alpha(-1)=(-1)^r\).

If \(F\ge8\), put \(d=F/4,\ m=\max(1,d/8)\) and \(D=d/m\).
The dyadic character at \(5\) has reduction \(Y\) of order
\(d\). Specify it up to Frobenius over \(k\) by an odd
\(e\in[1,D)\) with \(Y^m=\beta^{120e/D}\). For \(F=1,4\)
put \(d=m=D=1,\ e=0,\ Y=1\). At \(F=1\) admissibility
requires \(l+r_0\) even, and at \(F=4\) it requires \(l+r_0\)
odd. At \(F\ge8\) the value of the dyadic character at \(-1\)
is fixed by this same parity and imposes no other condition on \(l,r_0\).

\begin{theorem}[Finite-field tests through dyadic conductor \(16384\)]\leavevmode
\label{thm:towerres317}
For \(F\le 16384\) and every \(r\ge1\) of the prescribed parity,
the strict group \(S_{\alpha,r}\) at \(11\) with coefficient
\((J/O)(\alpha^{-1}\kappa_{11}^{-r})\) over \(\QQ\) vanishes
outside the labels in this table at \(r_0=0\).
\[
 \begin{array}{c|l}
 F &(l,e)\\ \hline
 8 &(20,1)\\
 16 &(10,1),\ (30,3)\\
 32 &(5,5),\ (15,7),\ (25,1),\ (35,3).
 \end{array}
\]
\end{theorem}
\begin{proof}
Use the cyclotomic moment of Proposition~\ref{prop:moment267}
at precision one. Its hypotheses hold since
\(11\nmid fF\varphi(fF)\). A nonzero value proves vanishing
of the full strict group for every positive \(r\) in that
residue class. We give the finite calculations explicitly.
The two auxiliary primes \(s\), with a generator \(g\) of
\(\mathbb F_s^\times\), are
\[
          (140394497,3),\qquad(177340417,5).
\]
They have \(s-1=19\cdot451\cdot 2^{14}\) and
\(s-1=24\cdot451\cdot 2^{14}\), respectively.
Primality is checked by trial division through
\(\lfloor\sqrt s\rfloor\). The powers \(g^{(s-1)/q}\) for
the distinct prime divisors \(q\mid s-1\), in increasing
order, are \(140394496,84793338,132557353,109523072\) in
the first case and \(177340416,168997293,28977015,148801780\)
in the second.

For \(F\ge8\), the odd residue classes modulo \(F\) are
\((-1)^b5^v\) for \(b=0,1\) and \(0\le v<d\).
Indeed \(v_2(5^{2^j}-1)=j+2\) by successive squaring, so
the subgroup generated by \(5\) consists of the classes equal
to one modulo four. A primitive character has order \(d\)
at \(5\), since \(5^{d/2}=1+F/2\bmod F\). If \(d>8\),
the order of \(121\) modulo \(d\) is \(d/8\), since
\(v_2(121^{2^a}-1)=3+a\) for all \(a\ge0\). Thus the
minimal polynomial over \(k\) of each possible \(Y\) is
\(Z^m-\eta\), where \(\eta=\beta^{120e/D}\).
Indeed it divides this polynomial and has degree \(m\).
For \(d\le8\) the same description with \(m=1\) follows
from \(d\mid120\).

Put \(w=g^{(s-1)/(11fF)}\bmod s\). For
\(1\le a<11fF\) prime to \(11fF\), define \(L(a)\) as the
unique integer \(0\le L<11\) for which
\[
       (1-w^a)^{(s-1)/11}=g^{L(s-1)/11}\pmod{s}.
\]
Write \(a=(-1)^b5^v\bmod F\) and \(a=6^j\bmod41\) for the
respective coordinates. Set \(j=0\) at \(f=1\). At \(F=4\)
use \(v=0\), and at \(F=1\) use \(b=v=0\).
The moment at \(Y\) is the value of the polynomial whose
coefficient of \(Z^v\) is the sum of
\((-1)^{b(l+r_0)} C^{lj}a^{r_0}L(a)\) over these \(a\).
After reduction modulo \(Z^m-\eta\), a nonzero coefficient
proves nonvanishing at every root in that orbit.

For reproducibility the following are the full counts of
labels left after testing the coefficients with indices
\(0\le k_0<\min(m,3)\). The last column is cumulative. The
last row applies to each of the nine separate conductors
\(64,128,\ldots,16384\).
\[
 \begin{array}{r|r|r|r}
 F&\text{admissible labels}&s=140394497&\text{both primes}\\ \hline
 1&200&5&0\\
 4&200&8&0\\
 8&400&13&1\\
 16&800&8&2\\
 32&1600&16&4\\
 64,\ldots,16384&1600&0&0
 \end{array}
\]
Here are direct recurrences to reproduce the counts and surviving
labels with integer arithmetic. Start an accumulator in \(k\)
at zero for each \(s,F,l,e,r_0,k_0\). Loop over \(0\le q<D\)
and set \(v=k_0+mq\). For each admissible sign \(b\), each
\(0\le j<40\) (just \(j=0\) at \(f=1\)), and \(1\le t\le10\),
take by the Chinese remainder theorem
\[
 a\equiv(-1)^b5^v\bmod F,\quad
 a\equiv6^j\bmod f,\quad a\equiv t\bmod11 .
\]
Use just \(b=0\) for \(F=1\) and both signs for \(F\ge4\).
Add \(\eta^q(-1)^{b(l+r_0)}C^{lj}t^{r_0}L(a)\).
In this recurrence multiplication of coordinate pairs is
\((a,b)(c,d)=(ac-bd,ad+bc)\bmod11\). Discard a label as
soon as one accumulator is nonzero. The labels remaining in
the last column are exactly those displayed in the statement.
This computes selected coefficients sufficient for every
vanishing claimed here.
\end{proof}

\subsection{The local correction to the interpolated order}\label{nav:IV-59}
Let \(E/\QQ\) be one of the abelian period fields, of degree \(n\)
and cyclotomic modulus \(N\). Let \(p\) be an odd prime with
\(p\nmid nN\), and let \(r\ge1\) satisfy
\(r\not\equiv0,-1\bmod(p-1)\). Take a character \(\alpha\) of
\(\operatorname{Gal}(E/\QQ)\) with \(\alpha(-1)=(-1)^r\).
Use its primitive character of conductor \(f\), and write
\(\beta=\alpha^{-1}\), \(O=\ZZ_p[\alpha]\), \(J=\operatorname{Frac}O\).
Write \(\kappa\) for the cyclotomic character. The two coefficients are
\[
 A=(J/O)(\beta\kappa^{-r}),\qquad T=O(\alpha\kappa^{r+1}).
\]
In \(H^1(\QQ,A)\) let \(R\) consist of classes unramified at every
finite prime other than \(p\), and let \(S\) be the kernel of
localization at \(p\) on \(R\). Write \(U_T\) for the compact group
of classes for \(T\) unramified away from \(p\).

Choose \(B>1\) prime to \(fp\) with \(\beta(B)B^{-r}\ne1\bmod p\).
Such a residue exists by varying \(B\bmod p\) with \(B=1\bmod f\).
Put \(X_z=(1+p)^z-1\) for \(z\in\ZZ_p\) and \(h=r+1\bmod(p-1)\).
We use the integral series \(F\) of Lemma~\ref{lem:shared201},
normalized by \(F(X_z)=A_B(z)\) there. More explicitly, on
\((\ZZ/f)^\times\times\ZZ_p^\times\) integrate
\(\beta(a)\omega(a_p)^{-h}(1+X)^{-\log\langle a_p\rangle/\log(1+p)}\)
against its bounded measure with cylinder weights
\(\lfloor Ba/(fp^j)\rfloor-(B-1)/2\).
Here \(\omega\) is the Teichm\"uller character and
\(\langle a_p\rangle=a_p/\omega(a_p)\). The binomial coefficients
are integral, giving \(F\in O[[X]]\). Formula~\eqref{eq:value201}
gives, for \(k\ge2\) with \(k\equiv-r\bmod(p-1)\),
\[
 v_p(F(X_{1-k}))=v_p(L(1-k,\beta)).
\]
We normalize \(v_p(p)=1\) and put \(v_p(0)=+\infty\).

\begin{theorem}[Order of the relaxed group and the strict local correction]\leavevmode
\label{thm:localvalue309}
Both \(U_T\) and \(H^1(\QQ_p,T)\) are free \(O\)-modules of rank
one. Write localization between bases as multiplication by
\(\rho\in O\), and put \(\lambda=\operatorname{length}_O S\).
Then
\[
  (F(X_{r+1}))=(p^\lambda\rho)\quad\hbox{as ideals of }O.
\]
If the value is nonzero, \(R\) is finite of length
\(v_p(F(X_{r+1}))\) and \(R/S\) is cyclic of length
\(v_p(\rho)\). If the value is zero, then \(\rho=0\) and
\(R/S\simeq J/O\). When \(p\nmid f\varphi(f)\), the same ideal is
generated by the local coordinate of the cyclotomic class
\(c_{\alpha,r+1}\) of Theorem~\ref{thm:circindex267}.
\end{theorem}
\begin{proof}
For large \(j\) put \(k_j=(p-1)p^{j-1}-r\ge2\).
The characters \(\beta\kappa^{k_j}\) and \(\beta\kappa^{-r}\)
agree modulo \(p^j\). With either divisible coefficient, its global
invariants vanish by the inertia action at \(p\). At other inertia
groups the invariants are divisible since averaging over the finite
image of order prime to \(p\) projects onto them. The multiplication
sequence for \(p^j\) therefore recovers the same group \(R[p^j]\)
from the finite coefficient with unramified conditions away from
\(p\) and unrestricted condition at \(p\).

For the positive weight \(k_j\) this finite group is the \(p^j\)-torsion
of a finite module of length \(v_j=v_p(L(1-k_j,\beta))\).
Indeed the proof of Theorem~\ref{thm:layers308} applies with
\(k_j,1-k_j\) nonzero modulo \(p-1\). The field cut out by \(\beta\)
has degree prime to \(p\) and is unramified at \(p\). These are
the descent and local hypotheses of Theorems~\ref{thm:positive306}
and~\ref{thm:primitive307} in weights at least two. Now
\(X_{1-k_j}\equiv X_{1+r}\bmod p^j\), by the binomial theorem.
If \(w=v_p(F(X_{1+r}))<\infty\), we have \(v_j=w<j\) for all
sufficiently large \(j\). A module of length \(w\) is then killed
by \(p^j\). Thus \(R[p^j]\) has length \(w\) for all these \(j\).
The nested groups exhaust \(R\), proving its order. If \(F(X_{1+r})=0\),
then \(v_j\ge j\). The \(p^j\)-torsion of a finite module of length
at least \(j\) has length at least \(j\), by its elementary divisors.
Hence \(R\) is infinite in this case.

At \(q\mid f\) inertia acts nontrivially through a group of order
prime to \(p\) on \(T/p^j\) and its Cartier dual. Both groups of
invariants vanish. Local duality and the Euler characteristic at
\(q\ne p\) then give zero in degree one for both coefficients.
Consequently \(U_T=H^1(\operatorname{Spec}\ZZ[1/fp],T)\).
Its rank is one by Theorem~\ref{thm:openindex269}. In detail,
its rational degree two is zero by descent to the finite even
\(K\)-group, its degree zero vanishes, and its rational Euler
characteristic is minus one since \(T\) is odd. The absence of
invariants modulo \(p\) gives torsion freeness. For the local
group the two residual inertia characters \(\omega^{-r}\) and
\(\omega^{r+1}\) are nontrivial. The multiplication sequence,
local duality and the Euler characteristic over \(\QQ_p\) give a
torsion-free module of rank one in degree one. Its dual
\(H^1(\QQ_p,A)\) identifies with \(J/O\) using the unramified
trace pairing.

The annihilator of the image of \(R\) in \(H^1(\QQ_p,A)\) is
the localization of \(U_T\). To see this first use finite
Poitou--Tate duality for \(A[p^j]\) and \(T/p^j\), with the
unramified conditions away from \(p\), which are exact orthogonal
complements by the local calculation in Theorem~\ref{thm:positive306}.
The image at \(p\) for one coefficient is the orthogonal complement
of the image for the other. If a local integral class annihilates
\(R\), at every \(j\) its reduction therefore has a global lift
with the required conditions. Those sets of lifts are finite.
Reduction carries higher lifts to lower lifts, so every finite
list of levels admits compatible lifts by starting at the largest.
Compactness gives a compatible lift at all levels, which belongs
to \(U_T\). The other inclusion follows from the same finite
pairings. This proves the annihilator assertion. See
\href{https://www.jmilne.org/math/Books/ADTnot.pdf}{Milne,
\emph{Arithmetic Duality Theorems}, Chapter I, Sections 2 and 4}.

The image in \(J/O\) is therefore \((J/O)[\rho]\). The kernel
\(S\) is finite by Theorems~\ref{thm:wild98} and~\ref{thm:char}.
If \(R\) is finite, taking lengths proves \(w=\lambda+v_p(\rho)\).
If \(R\) is infinite, its image is infinite and \(\rho=0\).
This proves the ideal equality in both cases. Finally
Theorem~\ref{thm:circindex267} writes the cyclotomic class as
an \(O\)-basis of \(U_T\) times a scalar of valuation \(\lambda\),
since \(r\not\equiv0\bmod(p-1)\) kills both invariant terms
in its formula. Localization gives the last assertion.
\end{proof}

\subsection{Removing the factor forced by local torsion}\label{nav:IV-60}
Keep \(r\ge1\), \(A=(J/O)(\alpha^{-1}\kappa^{-r})\),
\(T=O(\alpha\kappa^{r+1})\), and \(\alpha(-1)=(-1)^r\).
Assume only that \(p\) is odd and the order of \(\alpha\) is
prime to \(p\). Write \(\alpha=\alpha_0\omega^a\) with
\(\alpha_0\) of conductor prime to \(p\). This is the character
decomposition on the two factors of the modulus. The factor at
\(p\) is tame since its order is prime to \(p\).
If the global residual character of \(A\) is trivial,
Theorem~\ref{thm:criticalall310} already gives \(R=S=0\).
Otherwise take the integral series \(F\) of Theorem~\ref{thm:scalar231}
for the branch \(k\equiv-r\bmod(p-1)\), using the primitive
character \(\alpha^{-1}\). Its unit smoothing factor and
interpolation hold also when \(p\) divides the conductor.
Put
\[
  e=\begin{cases}
  1 & a+r+1\equiv0\bmod(p-1)\ \hbox{and }\alpha_0(p)=1,\\
  0 & \hbox{otherwise.}
  \end{cases}
\]

\begin{theorem}[An integral factor and the full strict order]\leavevmode
\label{thm:forced310}
The series \(F^\circ(X)=F(X)/X^e\) belongs to \(O[[X]]\).
The compact global group \(U_T\) is free of rank one. Let
\(\rho\) represent its localization in the free quotient of
\(H^1(\QQ_p,T)\), in integral bases. For
\(\lambda=\operatorname{length}_O S\),
\[
    (F^\circ(X_{1+r}))=(p^\lambda\rho).
\]
At a positive sample \(k\ge2\) on the branch one has
\[
 v_p(F^\circ(X_{1-k}))
      =v_p(L(1-k,\alpha^{-1}))-e(1+v_p(k-1))\ \ge0 .
\]
If this number is zero at one sample, the strict group vanishes
for every positive twist on the branch. The relaxed group at
each of these twists is cyclic of length \(e(1+v_p(r+1))\).
\end{theorem}
\begin{proof}
Take \(k_j=(p-1)p^{j-1}-r\) for sufficiently large \(j\).
The residual character at this positive weight is the same
nontrivial global character. The finite level argument in
Theorem~\ref{thm:localvalue309} identifies \(R[p^j]\) with
the \(p^j\)-torsion of the positive relaxed group. It uses
global residual nontriviality and divisibility of invariants at
inertia away from \(p\), with no condition at \(p\).
Theorem~\ref{thm:criticalall310} gives the positive group length
\(v_p(F(X_{1-k_j}))\). Since \(X_{1-k_j}=X_{1+r}\bmod p^j\),
the same nested-group argument proves that \(R\) has length
\(v_p(F(X_{1+r}))\) if this is finite, and is infinite otherwise.

The global line and the local orthogonality calculation in
Theorem~\ref{thm:localvalue309} still apply. To check the latter
when local invariants occur, use the finite local pairings and
pass to the compact inverse limit on the \(T\) side.
An integral class pairs with finite classes through its reduction,
irrespective of the kernel of the map to local divisible
coefficients. The finite lifting and compactness argument there
gives the annihilator of the image of \(R\) exactly as before.
Thus, writing \(H=H^1(\QQ_p,T)\), the dual of \(R/S\) is
\(H/\operatorname{loc} U_T\). Finiteness of \(S\) follows by
restriction to the compositum of the field cut out by \(\alpha\)
and \(\QQ(\mu_p)\). Its degree is prime to \(p\), so restriction is
injective by corestriction. Over this field finiteness follows
from the even \(K\)-group comparison of Theorem~\ref{thm:wild98}.

Rationally \(H\) has rank one by the local Euler characteristic,
since both nonzero cyclotomic powers have zero rational invariants.
Its torsion is \(H^0(\QQ_p,(T\otimes_O J)/T)\).
The inertia character there is \(\omega^{a+r+1}
\langle\kappa\rangle^{r+1}\). A Frobenius lift fixing the local
cyclotomic extension acts by \(\alpha_0(p)\). Such a lift exists
since that extension is totally ramified. If either residual
factor is nontrivial, the invariants are zero. Otherwise their
order as an \(O\)-module is determined by
\(v_p((1+p)^{r+1}-1)=1+v_p(r+1)\). Thus the torsion is cyclic
of length \(\tau=e(1+v_p(r+1))\). When \(\rho\ne0\) the torsion
injects into \(H/\operatorname{loc} U_T\) with quotient
\(O/\rho O\). When \(\rho=0\) that quotient of \(H\) has positive
rank. Together with finiteness of \(S\) this proves
\[
          (F(X_{1+r}))=(p^{\lambda+\tau}\rho).
\]

If \(e=1\), choose positive integers \(r_i\) on the same branch
with \(r_i+1=0\bmod p^i\), using the Chinese remainder theorem.
The preceding equality gives \(F(X_{1+r_i})\in p^{i+1}O\),
while \(X_{1+r_i}\to0\). Continuity gives \(F(0)=0\).
Division by \(X\) is then integral. At every nonzero integer
\(z\), \(v_p(X_z)=1+v_p(z)\), proving both formulas. An integral
series taking a unit value at one point of \(pO\) has unit constant
term and takes unit values everywhere on \(pO\). In that case
\(\lambda=0\) and \(\rho\) is a unit, so \(H/\operatorname{loc} U_T\)
is its cyclic torsion subgroup up to isomorphism. This proves
the last assertion.
\end{proof}

For example take \(p=3\) and the primitive character
\(\alpha=\chi_{-8}\) of conductor eight. It is odd and
\(\alpha(3)=1\), so \(e=1\) on every odd twist. The residual
character is nontrivial on inertia at two. At \(k=3\) the values
of \(\alpha\) on \(1,3,5,7\) are \(1,1,-1,-1\).
With \(B_3(u)=u^3-3u^2/2+u/2\), the expression
\(-8^2\sum_b\alpha(b)B_3(b/8)/3\) gives \(L(-2,\alpha)=-3\).
The corrected valuation is \(1-(1+v_3(2))=0\).
For every positive odd \(r\) the strict group vanishes and the
relaxed group has order exactly \(3^{1+v_3(r+1)}\).

\subsection{A split uniformizer test for strict quadratic classes}\label{nav:IV-61}
Let \(D/\QQ\) be real quadratic, \(\alpha\) its quadratic
character, and \(p\) an odd split prime with primes \(w,\bar w\)
in \(D\). Identify both completions with \(\QQ_p\). For
\(z\in\QQ_p^\times\) set
\[
        l_p(z)=\frac{(zp^{-v_p(z)})^{p-1}-1}{p}\pmod p.
\]
This additive-valued homomorphism is zero on \(p\) and generates
the one-dimensional space of \(\mathbb F_p\)-valued characters
of the units. Indeed it is nonzero at \(1+p\), trivial on the
prime-to-\(p\) roots of unity, and zero on \(1+p^2\ZZ_p\).

\begin{theorem}[Vanishing from a split principal-ideal relation]\leavevmode
\label{thm:splitunits317}
Assume that the quotient of the class group of \(D\) by the
classes of \(w,\bar w\) has order prime to \(p\). Suppose
there is an element \(z\in\OO_D[1/p]^\times\) with
\(l_p(z_w)\ne l_p(z_{\bar w})\).
For every positive \(r\equiv0\bmod(p-1)\) the entire strict
group over \(\QQ\) with coefficient
\(A=(\QQ_p/\ZZ_p)(\alpha\kappa_p^{-r})\) is zero.
In particular this holds for \(D=\QQ(\sqrt{82})\) and \(p=11\).
\end{theorem}
\begin{proof}
The residual character of \(A\) is \(\alpha\), which is
nontrivial. Hence \(A^{G_\QQ}=0\) and the multiplication-by-\(p\)
sequence identifies a global class killed by \(p\) with a
unique class with coefficients in \(A[p]\). On restriction
to \(D\) this becomes an additive character into \(\mathbb F_p\).
Restriction is injective since \([D:\QQ]=2\) is prime to \(p\).
The character changes sign under the nonidentity element of
\(\operatorname{Gal}(D/\QQ)\). It is unramified away from \(p\).
Indeed there the original strict class is unramified over \(D\),
and inertia acts trivially on the divisible module \(A\) over
\(D\). On inertia the map in degree one from \(A[p]\) to \(A\)
is injective by the multiplication sequence.

At a place of \(D\) above \(p\) the original class is zero.
Its finite-coefficient character is therefore a boundary of
a local invariant \(a\in A^{G_{\QQ_p}}\). On choosing \(y\)
with \(py=a\), its cocycle is
\((\kappa_p(g)^{-r}-1)y\). Local reciprocity sends \(p\)
to an element acting trivially on all \(p\)-power roots of
unity. Indeed \(p\) is the norm of \(1-\zeta_{p^j}\) in
each of the local cyclotomic extensions. The local character
thus kills \(p\). On identifying the two local fields, their
characters are \(t\,l_p\) and \(-t\,l_p\) for a single
\(t\in\mathbb F_p\), by the sign under quadratic conjugation.
Global reciprocity on the principal idele \(z\) gives
\[
           t\big(l_p(z_w)-l_p(z_{\bar w})\big)=0 .
\]
The other finite components are units, and the archimedean
components contribute zero for odd \(p\). The hypothesis
forces \(t=0\). The character is now everywhere unramified
and trivial on both primes over \(p\). It factors through
the class-group quotient in the statement and so vanishes.
A nonzero element of a \(p\)-primary torsion group would have
a multiple of order \(p\). This proves the assertion at all orders.

For the example, \(\OO_D=\ZZ[\sqrt{82}]\) has discriminant
\(328\), and the two roots of \(X^2-82\) modulo \(121\)
are \(18,-18\). Minkowski's bound \(\sqrt{328}/2<10\)
shows that prime ideals of norm less than ten generate the
class group. The primes \(5,7\) are inert, since \(82\) is a
nonsquare modulo each, so only primes over \(2,3\) are needed.
The prime \(\mathfrak q\) above \(2\) is ramified,
and \(3\) splits. The element \(8+\sqrt{82}\) has norm
\(-18\) and is not divisible by the ideal \((3)\). Its ideal
is therefore \(\mathfrak q\mathfrak r^2\) for one prime
\(\mathfrak r\) over \(3\). The other prime over \(3\) has
inverse ideal class. Since \([\mathfrak q]^2=1\), this proves
that the class group has exponent dividing four.

Take \(z=203+18\sqrt{82}\), of norm \(14641=11^4\).
At the embedding with \(\sqrt{82}\equiv18\bmod121\),
\(z_w\equiv43\bmod121\) and \(43^{10}\equiv45\bmod121\).
Thus \(l_{11}(z_w)=4\). Since
\(z_{\bar w}/11^4=1/z_w\), the value at the other prime is
\(-4\). Their difference is nonzero modulo eleven.
\end{proof}

\subsection{Local cyclotomic degrees and ordinary class groups}\label{nav:IV-62}
Let \(F/\QQ\) be abelian and \(p\) odd. At each \(w\mid p\)
local reciprocity followed by the principal cyclotomic
character gives a homomorphism
\(\lambda_w:F_w^\times\to\Gamma_w\subseteq\Gamma_F\).
Use additive coordinates on the procyclic groups. The image
is all of \(\Gamma_w\). Indeed local reciprocity has dense
image, and the image of a sufficiently deep principal unit
group contains an open subgroup. This last assertion follows
also from the trace formula for the logarithm of the local
norm, since the trace of a full lattice contains a nonzero
ideal of \(\ZZ_p\). Put
\[
 Z=\bigoplus_{w\mid p}\Gamma_w,\qquad
 I=\sum_{w\mid p}\Gamma_w\subseteq\Gamma_F,\qquad
 Z^0=\ker(Z\xrightarrow{\ \sum\ } I).
\]
Write \(C=\Cl(\mathcal O_F[1/p])\{p\}\) for the \(p\)-part
of the ordinary class group after inverting \(p\), and
\(\lambda=(\lambda_w)_{w\mid p}\) on the \(p\)-units.

\begin{theorem}[The integral contribution from the local degrees]\leavevmode
\label{thm:logunits319}
There is a surjection \(C\to\Gamma_F/I\) and an exact sequence
\[
 0\longrightarrow
 Z^0/\overline{\lambda(\mathcal O_F[1/p]^\times)}
 \longrightarrow \widetilde C(F)
 \longrightarrow \ker(C\to\Gamma_F/I)\longrightarrow0.
\]
Let \(s\) be the number of places over \(p\). Take \(s-1\)
independent columns of logarithms of \(p\)-units in a
\(\ZZ_p\)-basis of \(Z^0\), forming a square matrix \(B\).
Such a list exists. Put \(a_B=v_p(\det B)\) and let \(e_B\)
be the largest exponent among its Smith elementary divisors.
For \(s=1\) take both numbers zero. Then
\[
 \log_p|\widetilde C(F)|\le
       a_B+\log_p|C|-\log_p[\Gamma_F:I],
 \qquad
 \log_p\exp\widetilde C(F)\le e_B+\log_p\exp C .
\]
Here \(\exp\) of a finite group denotes its exponent, with
value one for the trivial group.
\end{theorem}
\begin{proof}
At depth \(n\), quotient the ideles by the principal ideles,
the infinite factors, the units away from \(p\), and the
inverse images of \(p^n\Gamma_w\) under \(\lambda_w\) at
the places over \(p\). The result is a finite abelian group.
Indeed quotienting further by \(Z/p^nZ\) gives the ordinary
class group after inverting \(p\). Its local kernel is the
image of \(Z/p^n Z\), and the only relations there are
the logarithms of the \(p\)-units. To verify the last
statement, a principal idele representing a tuple supported
over \(p\) modulo the other local factors has valuation
zero at every prime away from \(p\). Its generator is a
\(p\)-unit. Conversely such a generator gives the stated
relation. Local surjectivity gives all of \(Z/p^n Z\)
before imposing principal relations.

Take the \(p\)-primary parts and pass to the inverse limit
in \(n\). Finite abelian primary decomposition is exact,
and the transition maps on the local quotients are onto.
Global reciprocity identifies the middle limit with \(X(F)\),
since every finite abelian \(p\)-quotient with these local
conditions factors through one of the displayed depths.
On \(Z\) the kernel in the limit is the closure of the
subgroup generated by the \(p\)-units, by reduction modulo
each \(p^n\). The map \(X(F)\to\Gamma_F\) is surjective
and acts on the local factors by their sum. Principal
relations lie in \(Z^0\) by reciprocity. This induces the
asserted surjection from \(C\). A class in its kernel can
be lifted to \(X(F)\) and its degree in \(I\) cancelled
by a local tuple. The remaining kernel consists exactly of
the degree-zero local tuples modulo the principal relations,
which proves exactness.

Finiteness of the middle group for this abelian field gives
finite index for the closed logarithm subgroup in \(Z^0\).
This free module has rank \(s-1\), since \(I\) has rank
one. The logarithms therefore span over \(\QQ_p\) and
contain the required independent list. The lattice spanned
by the selected columns has index \(p^{a_B}\), and its
quotient has exponent \(p^{e_B}\) by Smith reduction.
The same bounds hold after adjoining the remaining
logarithms. The order formula and multiplication of the
two annihilators in the exact sequence prove the two
inequalities.
\end{proof}

\subsection{An exact split contribution of order nine}\label{nav:I-80}
Let \(D=\QQ(\sqrt{5611})\). Since \(5611=31\cdot181\)
is squarefree and congruent to three modulo four, the
integer ring is \(\ZZ[\sqrt{5611}]\) and the fundamental
discriminant is \(22444\). The prime three splits into
two primes \(w,\bar w\). The following computation
measures the local logarithms at both primes.

\begin{theorem}[A split quadratic component on the prime-count fields]\leavevmode
\label{thm:splitnine319}
For every positive multiple \(r\) of \(18\),
\(\mathcal S_{3,r}(D)\simeq\ZZ/9\). For every positive
even \(r\) the group is cyclic and has order at least three.
For \(6\mid r\) it has order at least nine. For the
primitive quadratic character of conductor \(22444\),
the complementary groups over \(E_x\) fixed by
\(\Gal(E_x/D)\) therefore have total logarithmic order
on the twists divisible by \(18\) equal to
\[
       \left\lfloor{t(x)-1\over18}\right\rfloor\log 9
\]
on every common interval. Summing over all \(1\le r<t(x)\)
gives a lower bound
\(\bigl(\lfloor(t(x)-1)/2\rfloor+
         \lfloor(t(x)-1)/6\rfloor\bigr)\log3\).
\end{theorem}
\begin{proof}
First \(\Cl(\mathcal O_D[1/3])\{3\}=0\). Indeed the
Minkowski bound is \(\sqrt{5611}<75\). Among odd primes
less than \(75\), apart from \(3\) and the ramified
prime \(31\), the split primes are
\[
 5,\ 7,\ 11,\ 17,\ 19,\ 43,\ 47,\ 53,\ 61,\ 71 .
\]
This is checked by squaring residues modulo those primes
and testing the other primes \(13,23,29,37,41,59,67,73\).
The prime two is ramified. Inert primes are principal,
and ramified prime classes have order dividing two. For
each split prime its two classes are inverses. The rows
of the table are ordered so that its listed prime is the
only new split prime away from three in the norm.
\[
\begin{array}{c|r|l@{\qquad}c|r|l}
 q&a&|a^2-5611|&q&a&|a^2-5611|\\ \hline
 7&75&2\cdot7&5&74&3^3\cdot5\\
 61&72&7\cdot61&43&8&3\cdot43^2\\
 47&73&2\cdot3\cdot47&17&69&2\cdot5^2\cdot17\\
 19&62&3\cdot19\cdot31&11&56&3^2\cdot5^2\cdot11\\
 53&24&5\cdot19\cdot53&71&12&7\cdot11\cdot71
\end{array}
\]
For each odd split prime in a norm only one of the two
primes divides \(a+\sqrt{5611}\), since the two residue
roots are distinct. Each new exponent is a unit modulo
three. Reading the rows from left to right therefore
kills the class of each listed prime modulo three.
Minkowski's bound gives the assertion.

The fundamental unit is
\(\varepsilon=11954870+159597\sqrt{5611}\). For
verification the continued fraction of \(\sqrt{5611}\)
has initial entry \(74\) and period
\[
 1,9,1,2,2,2,1,1,1,2,2,2,1,9,1,148 .
\]
These entries follow from \(m'=da-m\),
\(d'=(5611-(m')^2)/d\), \(a'=\lfloor(74+m')/d'\rfloor\)
starting with \(m=0,d=1,a=74\). The convergent just
before the last period entry yields the displayed unit by
the continued-fraction theorem for Pell's equation.
Norm \(-1\) is impossible modulo four. The element
\(z=149+2\sqrt{5611}\) has norm \(-243=-3^5\). At one
prime over three it is a unit. The valuation vectors of
\(3,z\) at the two primes consequently generate a lattice
of index five in \(\ZZ^2\). Together with \(-1,\varepsilon\)
they generate a subgroup of the \(3\)-units of index
dividing five. Indeed the kernel of the valuation map
is the ordinary unit group.

Both completions are \(\QQ_3\). Identify local
cyclotomic degree with the logarithm of the unit part
divided by \(\log 4\), up to the common reciprocity sign.
The uniformizer \(3\) has degree zero. Both local images
equal \(\Gamma_D=1+3\ZZ_3\), written additively in this
coordinate. Thus \(Z^0\) in Theorem~\ref{thm:logunits319}
consists of pairs \((a,-a)\). At a unit \(u\) the
coordinate has valuation \(v_3(u^2-1)-1\), since
\(\log\) preserves valuation on \(1+3\ZZ_3\).
The two roots of \(5611\) modulo \(81\) are \(47,-47\).
At the first one \(\varepsilon\equiv53\bmod81\), and
at the second \(z\equiv55\bmod81\). Both
\(53^2-1\) and \(55^2-1\) have valuation three.
The two coordinates of each principal \(3\)-unit sum
to zero by reciprocity. Its logarithms are therefore
divisible by nine in these cases, with exact valuation
two. The generators above, including their index prime
to three, show that the closed logarithm lattice is
exactly \(9Z^0\). Theorem~\ref{thm:logunits319} now gives
\(\widetilde C(D)\simeq\ZZ/9\).

When \(2\cdot3^{h-1}\mid r\), the coefficient is trivial
modulo \(3^h\). Apply Theorem~\ref{thm:logstrict319} with
\(E=D,\ L=D(\mu_3)\), whose real subfield is \(D\).
The relative group acts trivially on the real class group
and the coefficient. The resulting \(3^h\)-torsion is
\(\operatorname{Hom}(\ZZ/9,\ZZ/3^h)\). For \(h=1\)
this shows cyclicity of the finite strict group and the
first lower bound. For \(h=2\) it gives the second lower
bound. For \(h=3\) the \(27\)-torsion has order nine,
so the full group has no element of order \(27\) and is
itself cyclic of order nine. Complex conjugation acts
trivially on cohomology over the real field \(D\), as it
is represented by an element of its absolute Galois
group. Finally \([E_x:D]\) is a power of two, so strict
descent identifies the invariant summand at three with
the one just computed. The complementary sign retains
the even twists. Counting multiples gives both sums.
\end{proof}

\subsection{Characteristic polynomials and capitulation in the real tower}\label{nav:classmatrix330}
Keep the interval and the prime \(p>2\) fixed. Use the ordinary
unramified module \(M_p\), the split quotient \(X_p\), and
\(D_p=\ker(M_p\to X_p)\) of Theorem~\ref{thm:decomploss330}.
The tame character and the field \(E=E_x\) vary along this interval.
Write \(H=\Gal(E_*(\mu_{p^\infty})/E_*(\mu_p))\), with generator
\(\gamma\) and cyclotomic value \(a=\kappa_p(\gamma)\). As before,
\(F_j\) denotes the layer of degree \(p^j\) above
\((E_*(\mu_p))^+\).

\begin{theorem}[The full order from the ordinary class module]\leavevmode
\label{thm:unramdet330}
The subgroup \(D_p\) is contained in the finite \(\ZZ_p\)-torsion
subgroup \(C_p\) of \(M_p\). Put \(W_p=M_p/C_p\) and
\(C'_p=C_p/D_p\). For every \(E_x\) let \(\Delta_x\) be the
prime-to-\(p\) torsion subgroup of \(G_x=\Gal(E_*(\mu_{p^\infty})/E_x)\).
For a character \(\tau\) of \(\Delta_x\) with values in \(\ZZ_p^\times\),
write \(e_\tau\) for its averaging idempotent and put
\[
 P_{x,\tau}(Z)=\det(ZI-\gamma\mid e_\tau W_p)\in\ZZ_p[Z].
\]
The determinant on a module of rank zero is one. For \(r\ge1\)
put \(\tau=\kappa_p^{-r}|_{\Delta_x}\) and \(b=a^{-r}\). Then
\(P_{x,\tau}(b)\ne0\) and the full complementary strict dual has order
\begin{equation}\label{eq:unramdet330}
 |A_{p,x,r}|=
   p^{v_p(P_{x,\tau}(b))}
   \bigl|\ker(\gamma-b\mid e_\tau C'_p)\bigr|.
\end{equation}
The same polynomial computes the free contribution to the ordinary
specialization \(U_{p,x,r}\). In its order formula the finite
group is \(e_\tau C_p\). Both free lattices and finite groups in
these formulas come from \(E_*\) for the whole interval.
\end{theorem}

\begin{proof}
Iwasawa's class module theorem identifies \(M_p\) with the inverse
limit under norm of the \(p\)-parts of the ordinary ideal class
groups of \(F_j\). It is a finitely generated torsion module over
\(\ZZ_p[[H]]\). The ground field \(F_0\) is abelian over \(\QQ\)
and \(F_\infty/F_0\) is its cyclotomic \(\ZZ_p\)-extension.
Ferrero and Washington's theorem therefore gives \(\mu=0\).
In the structure theorem over \(\ZZ_p[[H]]\), all elementary
factors then arise from powers of distinguished polynomials.
Each such quotient is finite free over \(\ZZ_p\) by polynomial
division. The kernel and cokernel in the structure theorem are
finite. It follows that \(M_p\) is finitely generated over
\(\ZZ_p\), \(C_p\) is finite, and \(W_p\) is finite free.
\cite{V-III-II-iwasawa15,V-III-II-fw13}

We show that \(\omega_j=\gamma^{p^j}-1\) is invertible on
\(W_p\otimes_{\ZZ_p}\QQ_p\) for every \(j\). Over \(F_j\), the
Galois group of the extension defining \(M_p\) is an extension
of \(H^{p^j}\simeq\ZZ_p\) by \(M_p\). A lift of a generator
splits this extension, since its closed procyclic subgroup maps
onto \(\ZZ_p\) with trivial kernel. The closed commutator group
is \(\omega_j M_p\). Consequently the kernel of the map from the
abelianization to \(H^{p^j}\) is \(M_p/\omega_j M_p\).

This abelianization is a quotient of the maximal abelian pro-\(p\)
Galois group of \(F_j\) unramified outside \(p\). That group
has \(\ZZ_p\)-rank one. To verify the rank, use the ray class
sequence with moduli supported at \(p\) and take its pro-\(p\)
limit. The product of the local principal unit groups has rank
\(\sum_{v\mid p}[F_{j,v}:\QQ_p]=[F_j:\QQ]\). Its image has finite
index, since the ideal class group is finite. Its kernel is
the closure of the diagonal global units that are locally
principal units. These form a subgroup of finite index in the
global units. Brumer's theorem on Leopoldt's conjecture for
abelian fields gives this closure rank \([F_j:\QQ]-1\), since
\(F_j\) is totally real and abelian over \(\QQ\).
\cite{V-III-brumer330} Thus \(M_p/\omega_j M_p\) has rank zero
after the rank-one quotient \(H^{p^j}\) is taken into account.
On the rational vector space, \(\omega_j\) is a surjective
endomorphism and is therefore invertible.

Some \(H^{p^j}\) fixes every prime above \(p\), since \(H\)
acts continuously on the finite set of these primes. The
corresponding Frobenius elements are then fixed, by the
description in Theorem~\ref{thm:decomploss330}. Their images
in \(W_p\) vanish because \(\omega_j\) has zero kernel there.
Thus \(D_p\subset C_p\). In particular the torsion subgroup of
\(X_p=M_p/D_p\) is exactly \(C'_p\), and its quotient by torsion
is \(W_p\).

The index of \(H\) in \(G_x\) is prime to \(p\). The factor
over \(E_*(\mu_p)/E_*\) divides \(p-1\), and \([E_*:E_x]\)
is a power of two. Since \(G_x\) is abelian it is the direct
product \(\Delta_x\times H\). Indeed, for \(m=[G_x:H]\) and
\(g\in G_x\), the element \(g^m\in H\) has a unique \(m\)-th
root in \(H\). Dividing \(g\) by this root gives a torsion lift
of its coset. The intersection with \(H\) is trivial.
The relations for \(\Delta_x\) select exactly the image of
\[
 e_\tau=|\Delta_x|^{-1}\sum_{\delta\in\Delta_x}\tau(\delta)^{-1}\delta .
\]
Indeed each relation kills this summand and \(z-e_\tau z\) is
a sum of relations. On the selected summand the relations for
\(H\) amount to \(\gamma-b\), since integral powers are dense
and the image of this endomorphism is compact.

Apply these facts to \(0\to C'_p\to X_p\to W_p\to0\).
The cokernel on \(e_\tau W_p\) is a quotient of the finite
group \(A_{p,x,r}=Q_{x,r}(X_p)\). Thus \(\gamma-b\) on this
free lattice has nonzero determinant, and has zero kernel.
If a torsion element lies in the image of \(\gamma-b\) on
\(e_\tau X_p\), its preimage maps to zero in the free lattice,
so it already belongs to \(e_\tau C'_p\). We obtain a short
exact sequence of the three cokernels. Smith normal form gives
the free cokernel order \(p^{v_p(P_{x,\tau}(b))}\). On the
finite group \(e_\tau C'_p\) an endomorphism has kernel and
cokernel of the same order. This proves \eqref{eq:unramdet330}.
Applying the same argument to \(0\to C_p\to M_p\to W_p\to0\)
proves the assertion for \(U_{p,x,r}\).
\end{proof}

\begin{proposition}[The finite term as a capitulation kernel]\leavevmode
\label{prop:capmodule330}
Let \(B_j=\widetilde C(F_j)\) be the degree-zero logarithmic
class group. For all sufficiently large \(j\), \(C'_p\)
identifies, equivariantly for the field action, with the subgroup
of \(B_j\) whose classes become zero under extension to
some \(F_k\), \(k\ge j\). For each such \(j\), one finite
\(k\) realizes this entire kernel for all \(E_x\) and all twists.
\end{proposition}

\begin{proof}
Use \(B_j=X_p/\omega_j X_p\), as in
Theorem~\ref{thm:strictdigits329}. For \(k>j\) write
\(\nu_{j,k}=\sum_{i=0}^{p^{k-j}-1}\gamma^{ip^j}\), so
\(\omega_k=\nu_{j,k}\omega_j\). Extension of a logarithmic
class from \(F_j\) to \(F_k\) takes \(z\bmod\omega_jX_p\) to
\(\nu_{j,k}z\bmod\omega_kX_p\). Indeed reciprocity sends
extension on ideles to group transfer. Representatives given
by the successive powers of a lift of \(\gamma^{p^j}\)
conjugate \(z\) by the powers in this sum. This computation
takes place in the group defining the split tower, hence passes
to its degree-zero abelianized kernel at each level.

Both \(\omega_k\) and \(\omega_j\) are invertible on
\(W_p\otimes\QQ_p\) by the preceding proof, so \(\nu_{j,k}\)
is injective on \(W_p\). If a class represented by \(z\) dies
at level \(k\), then \(\nu_{j,k}(z-\omega_j y)=0\) for some
\(y\in X_p\). Projecting to \(W_p\) shows that the class of
\(z\) has a representative in \(C'_p\).

Choose \(j_0\) so \(H^{p^{j_0}}\) acts trivially on the finite
group \(C'_p\), and choose \(l\ge0\) so \(p^l C'_p=0\).
If \(j\ge j_0\), then \(\omega_jX_p\cap C'_p=\omega_jC'_p=0\).
The first equality follows from injectivity of \(\omega_j\) on
\(W_p\). Thus \(C'_p\) injects into \(B_j\). For
\(k\ge j+\max(1,l)\) the operator \(\nu_{j,k}\) acts on \(C'_p\)
as multiplication by \(p^{k-j}\), which kills it. The two
inclusions give the asserted equality. All maps commute with
field automorphisms since the tower over \(\QQ\) is abelian.
\end{proof}

Formula \eqref{eq:unramdet330} separates two arithmetic quantities
in the remaining cohomology estimate. One needs a uniform bound
for the valuations of the ordinary characteristic polynomials at
the specified positive twists together with the finite kernels on
the capitulating classes. The proposition supplies a finite level
for each fixed prime and interval. It gives no uniform bound on
that level as the prime and interval vary.


\subsection{A determinant for full strict orders and the first class norm}\label{sec:normsample340}
An order bound from the first two ordinary class groups can be made for
all but an explicitly bounded number of twists. We use a determinant
which includes the finite subgroup of the cyclotomic module. All
valuations in this subsection give $p$ value one. For a finite
$\ZZ_p$-module $B$, write $\ell(B)=\log |B|/\log p$.

\begin{lemma}[Simultaneous specialization of a finite presentation]\leavevmode\label{lem:simuldet340}
Let $M$ be a finitely generated $\ZZ_p$-module, $G$ an endomorphism,
and $S$ a finite set of algebraic $p$-adic integers. Take a finite
extension $K/\QQ_p$ containing $S$, with ring of integers $O$ and
ramification index $e_K$. Suppose that for every $z\in S$ the quotient
$B_z=(M\otimes O)/(G-z)(M\otimes O)$ is finite. There is a monic
polynomial $P$ over the integers of a finite unramified extension of
$K$ such that, simultaneously,
\[
                 v_p(P(z))=\operatorname{length}_{O}(B_z)/e_K
                     \qquad(z\in S).
\]
One may take $\deg P=\dim_{\mathbf F_p}(M/pM)$.
\end{lemma}
\begin{proof}
Choose a minimal presentation of $M$ over $\ZZ_p$ in Smith form, with
$n$ generators and relation matrix $Q$. Its columns are the nonzero
powers of $p$ times the basis vectors belonging to the finite summands.
Lift $G$ to an $n$ by $n$ matrix $A$ on the generators. At $z$ the
presentation matrix is $[\,zI-A, Q\,]$. Since its cokernel is finite,
the minimum valuation of its maximal minors is the number on the
right, by Smith normal form over $O$. For $n=0$ take $P=1$.

Let $Z$ be a matrix of independent variables with as many rows as
$Q$ has columns and with $n$ columns. In $\det(zI-A+QZ)$, expand by
multilinearity in its columns. Each coefficient is zero or a signed
maximal minor of the presentation matrix. Conversely, choose any
$n$ columns of that matrix. Match the selected columns of $Q$ to the
omitted positions of $zI-A$. The monomial made from these entries of
$Z$ specifies precisely those replacements. Its coefficient is the
chosen minor up to sign. Thus the minimum valuation of the coefficients
is the required minimum. Divide the polynomial at each $z$ by a scalar
of that minimum valuation and reduce it in the residue field. The
product of the resulting polynomials is nonzero. Over the algebraic
closure of the residue field it has a point where it is nonzero, since
that field is infinite. The finitely many coordinates lie in a finite
extension of the residue field. Lift them to a matrix $Z_0$ in the
corresponding unramified extension of $O$. Taking
$P(X)=\det(XI-A+QZ_0)$ gives all the required equalities. Extension
of the residue field has preserved the normalized valuations.
\end{proof}

Fix an interval of Section~\ref{sec:missing}, its largest field $E_*$,
an odd prime $p$, and $a=\kappa_p(\gamma)$ as in
Theorem~\ref{thm:strictdigits329}. Set $d=v_p(a-1)$, $s_r=d+v_p(r)$.
Take an integer $t_*>1$ above all the weight cutoffs on this interval.
Keep the tame group $\Delta_0$, the idempotent $e$ removing the
invariants for $D'$, and the idempotent $e_\tau$ of
Theorem~\ref{thm:twolayers327}. Write $X=X_p$ for the split class
module of Theorem~\ref{thm:unramdet330}. For a character $\tau$ put
\[
 M=e_\tau eX,\quad c_\tau=\ell(C_{0,\tau}),\quad
 k_\tau=\ell\ker(C_{1,\tau}\longrightarrow C_{0,\tau}),\quad
 H=d+\lfloor\log(t_*)/\log p\rfloor .
\]
Here the map between the ordinary class groups is norm. Only weights
with $\kappa_p^{-r}|_{\Delta_0}=\tau$ are used in this component.

\begin{theorem}[Two independent bounds for exceptional twists]\leavevmode\label{thm:normsample340}
For each $\tau$ there are two sets $D^0_\tau,D^1_\tau$ of these weights
in $1\le r<t_*$, of cardinalities at most $\lfloor c_\tau/d\rfloor$
and $k_\tau$ respectively. The same sets work for every $E_x$ of the
interval. If $B_r=eA_{p,x,r}$ is the indicated summand of the full
complementary strict dual, then
\[
\begin{aligned}
 r\notin D^0_\tau &\quad\Longrightarrow\quad
        \ell(B_r)\le H c_\tau/d,\\
 r\notin D^1_\tau &\quad\Longrightarrow\quad
        \ell(p^{s_r}B_r)\le H k_\tau .
\end{aligned}
\]
In particular the second bound counts the full remaining length above
the base precision. It includes the finite classes that capitulate.
\end{theorem}
\begin{proof}
Put $b_r=a^{-r}$. The strict dual over $E_*$ in this component is
$M/(\gamma-b_r)M$, by Theorem~\ref{thm:unramdet330} and its proof.
Over $E_x$ the larger group of field automorphisms imposes further
relations, so its component is a quotient. All positive-twist
quotients over $E_*$ are finite. Moreover $M/(\gamma-1)M=C_{0,\tau}$.
Put $U=(\gamma-1)M$ and $N_p(Y)=1+Y+\cdots+Y^{p-1}$.
The kernel of the first class norm is $U/N_p(\gamma)U$, since
$\gamma^p-1=N_p(\gamma)(\gamma-1)$, by
Theorem~\ref{thm:twolayers327}. The image of $U$ in $B_r$ equals
$(b_r-1)B_r=p^{s_r}B_r$. The modules $M,U$ are finitely generated
over $\ZZ_p$ by Theorem~\ref{thm:unramdet330}. On their rational
spans $\gamma-b_r$ is injective, first on $M$ by finiteness and
then on its stable subspace spanned by $U$.

Apply Lemma~\ref{lem:simuldet340} to $M$ with the points $1,b_r$.
The roots of the monic polynomial are integral. Their distances to
$1$, measured by valuation and summed with multiplicity, total
$c_\tau$. A root $\alpha$ with $v_p(\alpha-1)<d$ contributes precisely
that same valuation at every $b_r$. There are at most
$\lfloor c_\tau/d\rfloor$ other roots. For each of these choose a
weight $r_0$ in the progression maximizing $v_p(b_{r_0}-\alpha)$,
provided the progression is nonempty. At every other weight the
ultrametric inequality and the binomial calculation of
Theorem~\ref{thm:strictdigits329} give
\[
 v_p(b_r-\alpha)\le v_p(b_r-b_{r_0})
            =d+v_p(r-r_0)\le H.
\]
Use the chosen weights as $D^0_\tau$. Since $H\ge d$, the sum over
all roots at a remaining weight is at most $Hc_\tau/d$.

For $U$ take as reference point a primitive $p$-th root $\zeta$ and
again include all the $b_r$. Its coefficient ring $O=\ZZ_p[\zeta]$
has ramification index $p-1$ and residue field $\mathbf F_p$.
Evaluating $Y$ at $\gamma$ identifies
\[
 (U\otimes O)/(\gamma-\zeta)(U\otimes O)
                   \simeq U/N_p(\gamma)U .
\]
Indeed $O=\ZZ_p[Y]/N_p(Y)$, and the further relation identifies
multiplication by $Y$ with $\gamma$ on generators of $U$. Thus the
reference valuation in the lemma is $k_\tau/(p-1)$. At $b_r\in\ZZ_p$
the normalized valuation equals the original $\ZZ_p$-length, by
flat scalar extension. Now $v_p(b_r-\zeta)=1/(p-1)$ for every
positive $r$. At most $k_\tau$ roots $\alpha$ satisfy
$v_p(\alpha-\zeta)\ge1/(p-1)$. For each such root select a closest $b_{r_0}$ as
above to form $D^1_\tau$. All other roots contribute their reference
valuation at every $b_r$. If that latter sum is $S$, the number
selected, with multiplicity, is at most $k_\tau-(p-1)S$. At an
unselected weight the valuation is at most
$S+H(k_\tau-(p-1)S)\le Hk_\tau$. Passing to the image of $U$
proves the second assertion.
\end{proof}

\begin{corollary}[A class norm bound on the common intervals]\leavevmode\label{cor:normbudget340}
Let $t=\lfloor x^{401}\rfloor$, $Y=x^{-3}\exp(x^\varepsilon/4)$,
$Z=x^{-3}\exp(x^\varepsilon/6)$, with $0<\varepsilon<1/52$.
For $p\le Z$ retain in $B_r=eA_{p,x,r}$ only terms with
$r\in D^0_\tau\cap D^1_\tau$. For $Z<p\le Y$ retain terms with
$r\in D^0_\tau$. Summed over these odd primes and $1\le r<t$, the
logarithmic orders of all the discarded $B_r$ are
$O_{\psi,\varepsilon}(x^{-2}\exp(x^\varepsilon/2))$.
At $p\le Z$ the exceptional count for each component is at most
$\min(\lfloor c_\tau/d\rfloor,k_\tau)$. Both ordinary class groups
and both selections use the same largest field throughout an interval.
\end{corollary}
\begin{proof}
For $r\notin D^0_\tau$ use the first bound of the theorem. For
$p\le Z$ and $r\notin D^1_\tau$, the base term $B_r/p^{s_r}B_r$
is a quotient of $C_{0,\tau}$ by Theorem~\ref{thm:twolayers327}.
Thus $c_\tau+Hk_\tau$ bounds the full length in this case. For
$n_*=[E_*:\QQ]$ there are at most $1+tn_* /(p-1)$ occurrences
of each $\tau$. Orthogonal tame idempotents give
$\sum_\tau c_\tau\log p\le\log|\mathrm{Cl}(F_0)|$ and
$\sum_\tau k_\tau\log p\le\log|\mathrm{Cl}(F_1)|$. The conductors
and the Minkowski estimate in Theorem~\ref{thm:alltrunc329} bound
these class number logarithms by $O(n_*p^{j+1}\log(2N_*n_*p^{j+1}))$
for $j=0,1$, where $N_*$ is the modulus for $E_*$. One may take
$t_*=2+\lceil V^{602}\rceil$, so $H\ll_\psi\log(2x)$ since
$p^{d-1}\le n_*$. For $p\le Y$ the conductor logarithm is
$O_{\psi,\varepsilon}(x^\varepsilon)$ and
$n_*\ll_\psi x^{3+12\varepsilon}$. Bounding the sums over primes
by sums over integers now gives a constant times
\[
 n_*x^\varepsilon\log(2x)
       \bigl(Y^2+tn_*Y+Z^3+tn_*Z^2\bigr).
\]
The $Y^2$ term is $O(x^{-3+13\varepsilon}\log(2x)\exp(x^\varepsilon/2))$.
The $Z^3$ term has $x^{-6+13\varepsilon}$ in its place. The remaining
powers of $x$ multiply at most $\exp(x^\varepsilon/3)$. These estimates
prove the stated bound. They concern the component cut out by $e$.
\end{proof}

\begin{proposition}[One exceptional order with fixed norm groups]\leavevmode\label{prop:normexample340}
In the category of finitely generated cyclotomic modules, the orders
of the norm groups at every layer, even with their transition maps on
underlying groups, give no uniform bound on a single exceptional
positive-twist order.
\end{proposition}
\begin{proof}
Fix $r_0\ge1$ and $a\in1+p\ZZ_p$ of infinite order. On $M=\ZZ_p$
let a topological generator $\gamma$ act by $a^{p^L-r_0}$, where
$p^L>r_0$ and $L>v_p(r_0)$. This defines a continuous action of
$\ZZ_p$. For every $j\ge0$, the ideal $(\gamma^{p^j}-1)M$ is
$p^{d+v_p(r_0)+j}M$, independent of $L$. The norm quotients are the
same cyclic groups with the same reduction maps. On the other hand
$\gamma-a^{-r_0}=a^{-r_0}(a^{p^L}-1)$ has valuation $d+L$. All
positive-twist quotients are finite, since $p^L-r_0>0$. Thus even
when every first difference of norm lengths equals one, the length
at $r_0$ can grow arbitrarily. This example concerns the module
argument. Deciding which such orders the arithmetic towers realize
requires information about their field actions.
\end{proof}


\subsection{A decomposition group that detects every strict class}\label{sec:localgeneration346}

Let $\ell$ be an odd prime and $b\ne\ell$ a prime congruent to one
modulo $\ell$. Let $\mathcal F$ be the maximal pro-$\ell$ extension
of $\QQ$ unramified away from $\ell,b$, and put
$\mathcal G=\Gal(\mathcal F/\QQ)$. The extensions considered here
may ramify at both primes. We use two elementary tests,
\begin{equation}\label{eq:localgeneration346}
 \ell^{(b-1)/\ell}\not\equiv1\pmod b,\qquad
 b^{\ell-1}\not\equiv1\pmod{\ell^2}.
\end{equation}
The second test is equivalent to $b\not\equiv1\bmod\ell^2$ under
our assumption on $b$. The two tests concern different places.

\begin{lemma}[The maximal elementary abelian quotient]\leavevmode\label{lem:localgeneration346}
The maximal elementary abelian quotient of $\mathcal G$ has
dimension two over $\mathbf F_\ell$. A decomposition group at
$\ell$ is all of $\mathcal G$ if and only if the first test in
\eqref{eq:localgeneration346} holds. A decomposition group at
$b$ is all of $\mathcal G$ if and only if the second test holds.
\end{lemma}
\begin{proof}
By Kronecker--Weber, every cyclic extension of $\QQ$ of degree
$\ell$ unramified away from these primes has character factoring
through the unit groups modulo a power of $\ell$ and a power of
$b$. At $\ell$ the principal units are topologically cyclic and
give one quotient of order $\ell$, already at modulus $\ell^2$.
At $b$ the principal units are pro-$b$, and
$\mathbf F_b^\times$ is cyclic of order divisible by $\ell$.
This gives one more quotient, already at modulus $b$. The two
degree-$\ell$ fields so obtained have relatively prime conductors,
and hence are linearly disjoint. For example, their intersection
would be an abelian extension of $\QQ$ unramified at every finite
prime and of odd degree, so has degree one. Their compositum
is the maximal elementary abelian subextension.
\cite[Chapters V and VI]{V-III-V-cft101}

In its Galois group inertia at $\ell$ spans the first factor.
The image of a decomposition group there in the second factor is
generated by the class of $\ell$ in
$\mathbf F_b^\times/(\mathbf F_b^\times)^\ell$. This class is
nontrivial precisely in the first test. At $b$ inertia spans
the second factor, and Frobenius on the first is the image of
$b$ in the degree-$\ell$ quotient of $(\ZZ/\ell^2\ZZ)^\times$.
It is nontrivial precisely in the second test. For $b=1+\ell u$,
the binomial expansion of $b^{\ell-1}$ modulo $\ell^2$ proves
the equivalence stated above.

A proper closed subgroup of a pro-$\ell$ group has proper image
in some finite quotient. Indeed a point outside the subgroup has
a coset neighborhood disjoint from it. A proper subgroup of the
resulting finite $\ell$-group lies in a maximal subgroup, which
is normal of index $\ell$. Pulling back that subgroup shows that
the original closed subgroup fails to surject onto some cyclic
quotient of order $\ell$. Decomposition groups are closed. Thus
surjectivity onto the two-dimensional elementary abelian quotient
is equivalent to equality with $\mathcal G$.
\end{proof}

\begin{theorem}[Detection over every finite layer]\leavevmode\label{thm:detectlayer346}
Suppose one of the tests in \eqref{eq:localgeneration346} holds,
and let $v$ be its corresponding place. Let $B$ be any finite
intermediate field of $\mathcal F/\QQ$, and let $w$ be its place
over $v$. It is unique. Let $O$ be the integers in a finite
extension $J/\QQ_\ell$ with uniformizer $\pi$. Suppose
$\vartheta:G_B\to O^\times$ is continuous, trivial modulo $\pi$,
and unramified away from the places over $\ell,b$.
On degree-one classes with coefficients $(J/O)(\vartheta)$
unramified away from these same places, localization at $w$
is injective.
\end{theorem}
\begin{proof}
By Lemma~\ref{lem:localgeneration346} the image of a local
decomposition group at $v$ in $\mathcal G$ is the whole group.
It acts transitively on the embeddings of $B$ into $\mathcal F$
and hence there is one place above $v$. Its subgroup fixing
the chosen embedding of $B$ surjects onto
$\Gal(\mathcal F/B)$. Equivalently the image of $G_{B_w}$
is that entire subgroup.

Let $z$ represent a class in the kernel of localization. For
some $y\in J/O$ one has
$z(h)=(\vartheta(h)-1)y$ on $G_{B_w}$. The image of $z$ is
finite by continuity, compactness of $G_B$ and discreteness
of $J/O$. Choose $m$ so that $y$ and every value of $z$ lie
in $A_m=(J/O)[\pi^m]$. The map
\[
      g\longmapsto (z(g),\,\vartheta(g)\bmod\pi^m)
\]
is a homomorphism to the semidirect product of $A_m$ with the
image of $\vartheta\bmod\pi^m$. This is a finite $\ell$-group,
since the character takes values in principal units modulo
$\pi^m$. Away from $\ell,b$ inertia acts trivially and a
coboundary on inertia is the zero cocycle. Thus the finite
Galois extension of $B$ cut out by this map is unramified there.

This finite extension is contained in $\mathcal F$ even if
$B/\QQ$ is not Galois. Indeed a finite Galois $\ell$-extension
inside $\mathcal F$ contains the normal closure of $B$. Base
change all conjugates of the finite extension just cut out by
the homomorphism to that Galois field. Each compositum with it
is Galois of $\ell$-power degree over it. Their joint compositum
is Galois of $\ell$-power degree over $\QQ$ and is unramified
at every prime outside $\ell,b$. It lies in $\mathcal F$.
Consequently every value of the displayed homomorphism occurs
already on $G_{B_w}$. The identity $z(g)=(\vartheta(g)-1)y$
depends only on that value, by the choice of $m$. It holds on
all of $G_B$, proving that the class is zero.
\end{proof}

\begin{corollary}[Full vanishing on the original period fields]\leavevmode\label{cor:periodblock346}
Write $\Gal(E_x/\QQ)=P\times H$ with $P$ its $\ell$-Sylow group,
$H$ of order prime to $\ell$, and let $B=E_x^H$. Suppose $B$
is unramified away from $\ell,b$ and at least one of the two
tests in \eqref{eq:localgeneration346} holds. For every character
$\beta$ of $H$, extended trivially on $P$, and every positive
$r$ with $\bar\beta^{-1}\bar\kappa_\ell^{-r}=1$ on $G_B$,
the entire $\beta$-component of $\mathcal S_{\ell,r}(E_x)$
over $\ZZ_\ell[\beta]$ vanishes. This holds at every $x$ of the
interval and at every exponent. If $B$ is unramified away from
$\ell$ alone, the same vanishing holds with no test at $b$.

For a scalar character $\alpha$ of $E_x$ with
$\bar\alpha^{-1}\bar\kappa_\ell^{-r}=1$ on $G_\QQ$ and conductor
supported on $\ell,b$, the stronger local group $S^{N_*}(\alpha,r)$
of Theorem~\ref{thm:levelstrict345} also vanishes under either
test. Under the first test already $S(\alpha,r)=0$. Under the
second test $S(\alpha,r)$ injects into the unramified local group
at $b$ of Lemma~\ref{lem:locallevel345}.
\end{corollary}
\begin{proof}
Restriction induces an isomorphism between the $\beta$-component
and the strict group over $B$ for
$\vartheta=\beta^{-1}\kappa_\ell^{-r}$ with divisible coefficients.
Indeed restriction to $E_x$ gives the invariants of the twisted
module by inflation--restriction, since $|H|$ is a unit. Untwisting
on $G_{E_x}$ selects the indicated character. At places above
$\ell$ and on inertia elsewhere the corresponding indices are
prime to $\ell$, so restriction also detects the local conditions.
This is the integral descent of Lemma~\ref{lem:des28}.
Since $G_\QQ/G_B$ has $\ell$-power order and the residual
character has order prime to $\ell$, triviality on $G_B$
implies triviality on $G_\QQ$. The prime-to-$\ell$ character
$\beta$ is therefore the Teichm\"uller lift of
$\bar\kappa_\ell^{-r}$, and is unramified away from $\ell$.

The field $B$ is an abelian $\ell$-extension in $\mathcal F$.
Under the first test apply Theorem~\ref{thm:detectlayer346} at
$\ell$, where a strict class is zero. Under the second apply it
at $b$. There the coefficient is unramified, and arithmetic
Frobenius acts by a root of unity times $(Nw)^{-r}$. This scalar
differs from one as $r>0$. Its difference from one acts
surjectively on the divisible coefficient. Thus the unramified
local cohomology is zero, also over a ramified field $B_w$.
In both cases injectivity proves the asserted full vanishing.
The relative degree of $E_x/K_\psi$ is a power of two, so $B$
is precisely the maximal $\ell$-subfield of the fixed $K_\psi$.
In particular the ramification condition on $B$ is independent
of $x$. If $B$ is unramified away from $\ell$ alone, use the
maximal pro-$\ell$ extension unramified away from $\ell$ alone.
Its elementary abelian quotient is the first cyclic factor in
Lemma~\ref{lem:localgeneration346}. Inertia at $\ell$ surjects
onto it. The same closed-subgroup and localization proofs apply.

For the scalar assertion use Theorem~\ref{thm:detectlayer346}
with base $\QQ$ and $\vartheta=\alpha^{-1}\kappa_\ell^{-r}$.
The first test gives injectivity at $\ell$. Under the second,
restriction at $b$ is injective and its image for $S(\alpha,r)$
is unramified. If $\alpha$ is ramified at $b$ then $b\mid N_*$
and $S^{N_*}$ imposes zero there. If $\alpha$ is unramified
there, the unramified local group already vanishes.
\end{proof}

\begin{proposition}[Split coefficient primes in the cubic example and beyond]\leavevmode\label{prop:splitexample346}
For each odd prime $\ell$, the primes $b\equiv1\bmod\ell$ for
which $\ell$ splits completely in the degree-$\ell$ subfield
of $\QQ(\zeta_b)$ and $b\not\equiv1\bmod\ell^2$ have natural
density $1/\ell^2$ among rational primes. In particular take
$\ell=3$, $b=61$, and a primitive cubic character $\psi$ of
conductor $61$. On every one of its period fields the full
$\Gal(E_x/K_\psi)$-invariant complementary group at three is
zero for every positive twist.
\end{proposition}
\begin{proof}
For $b\equiv1\bmod\ell$, complete splitting of $\ell$ in the
specified degree-$\ell$ field means that $\ell$ is an
$\ell$-th power modulo $b$. Away from $\ell$ this and the
congruence $b\equiv1\bmod\ell$ mean that $b$ splits completely
in $D=\QQ(\mu_\ell,\sqrt[\ell]{\ell})$, the splitting field of
$X^\ell-\ell$. Over $L=\QQ(\mu_\ell)$ the fields $D$ and
$C=\QQ(\mu_{\ell^2})$ are distinct cyclic extensions of
degree $\ell$. In the Kummer group $L^\times/L^{\times\ell}$
their defining classes are $\ell$ and $\zeta_\ell$. At the
unique prime above $\ell$, the first class has valuation
$\ell-1$ modulo $\ell$, and the second has valuation zero.
The second class is nonzero because $[C:L]=\ell$. They span
distinct lines, hence $D\cap C=L$. The group of $DC/\QQ$
has order $\ell^2(\ell-1)$. Exactly $\ell-1$ of its
elements restrict to the identity on $D$ and nontrivially to
$C$. This is a conjugacy-stable set and Chebotarev gives the
stated density. \cite[Chapter VI]{V-III-V-cft101}

For $b=61$ one has $3^5\equiv-1\bmod61$, whence
$3^{20}\equiv1\bmod61$, while $61\equiv7\bmod9$.
Thus $3$ splits in $K_\psi$ and the test at $61$ holds.
For even $r$ take $\beta=1$ in Corollary~\ref{cor:periodblock346}.
For odd $r$ the complementary sign is minus. It is zero on
the invariants from $K_\psi$, since this odd-degree Galois field
is totally real, so complex conjugation is the identity in its
Galois group over $\QQ$. The descent just proved identifies
these invariants with the asserted subgroup on $E_x$.
\end{proof}

The two local tests control the indicated components, including
twists for which the residue character is trivial and a uniform
bound just from the orders of successive class groups is
unavailable. The other character components and the comparison
of the prescribed integral theta periods remain to be bounded
in Section~\ref{sec:missing}.

\clearpage
\section{References}\label{nav:bibliography-V}
\begingroup\def\section*#1{}\begin{thebibliography}{99}\small
\bibitem{V-III-V-stick} L. C. Washington, \emph{Introduction to Cyclotomic Fields}, 2nd ed., Springer, 1997. \href{https://link.springer.com/book/10.1007/978-1-4612-1934-7}{Chapter 6, Stickelberger's theorem}.
\bibitem{V-III-V-kw24} A. V. Sutherland, \emph{The Kronecker--Weber theorem}, MIT 18.785, Lecture 20. \href{https://ocw.mit.edu/courses/18-785-number-theory-i-fall-2021/mit18_785f21_lec20.pdf}{Lecture notes}.
\bibitem{V-III-V-kummer30} R. Sharifi, \emph{Group and Galois Cohomology}, Section 2.5, Lemma 2.5.8. \href{https://www.math.ucla.edu/~sharifi/notes/groupcoh-ch02.html}{Kummer isomorphism}.
\bibitem{V-III-V-cft101} J. S. Milne, \emph{Class Field Theory}. \href{https://www.jmilne.org/math/CourseNotes/CFT.pdf}{Local and global Artin reciprocity and the correspondence with idele class quotients}.
\bibitem{V-III-V-ramare105} O. Ramar\'e, \emph{Approximate formulae for \(L(1,\chi)\)}, Acta Arith. \textbf{100} (2001), 245--266. \href{https://www.impan.pl/shop/en/publication/transaction/download/product/83550}{Corollary 1, p.~248}.
\bibitem{V-III-V-osburn115} R. Osburn, \emph{Vanishing of eigenspaces and cyclotomic fields}. \href{https://arxiv.org/abs/math/0412242}{Theorem 1.1, the rational cyclotomic case}.
\bibitem{V-III-V-mo138} T. Morisawa and R. Okazaki, \emph{Height and Weber's Class Number Problem}, J. Th\'eor. Nombres Bordeaux \textbf{28} (2016), 811--828. \href{https://www.numdam.org/item/10.5802/jtnb.965.pdf}{Corollary B}.
\bibitem{V-III-V-milne98} J. S. Milne, \emph{Arithmetic Duality Theorems}, 2nd ed., BookSurge, 2006. \href{https://www.jmilne.org/math/Books/ADTnot.pdf}{Part I, Sections 2 and 4}.
\bibitem{V-III-V-weibel98} C. A. Weibel, \emph{The \(K\)-book. An introduction to algebraic \(K\)-theory}, Graduate Studies in Mathematics 145, American Mathematical Society, 2013. \href{https://sites.math.rutgers.edu/~weibel/Kbook/Kbook.VI.pdf}{Chapter VI, Section 8}.
\bibitem{V-III-V-voight99} J. Voight, \emph{The Gauss higher relative class number problem}. \href{https://jvoight.github.io/articles/higherrel-062618.pdf}{Proposition 1.3 and its odd-prime form}.
\bibitem{V-III-V-fk191} T. Fukuda and K. Komatsu, \emph{On \(\ZZ_p\)-extensions of real quadratic fields}, J. Math. Soc. Japan \textbf{38} (1986), 95--102.
\href{https://www.jstage.jst.go.jp/article/jmath1948/38/1/38_1_95/_pdf}{Theorem 2, pp.~95--96, and the table on p.~99}.
\bibitem{V-III-weberFK524} T. Fukuda and K. Komatsu, \emph{Weber's class number problem in the cyclotomic $\mathbb Z_2$-extension of $\mathbb Q$, II}, J. Th\'eor. Nombres Bordeaux \textbf{22} (2010), 359--368, Corollary 1.1. \href{https://www.numdam.org/article/JTNB_2010__22_2_359_0.pdf}{Published paper}.
\bibitem{V-III-II-duality13} J. S. Milne, \emph{Arithmetic Duality Theorems}, 2nd ed., 2006, Chapters I and II. \href{https://www.jmilne.org/math/Books/ADTnot.pdf}{Book}.
\bibitem{V-III-II-gp13} C. Greither and C. D. Popescu, \emph{An equivariant main conjecture in Iwasawa theory and applications}, J. Algebraic Geom. \textbf{24} (2015), 629--692. \href{https://arxiv.org/abs/1103.3069}{Authors' paper}.
\bibitem{V-III-II-iwasawa15} K. Iwasawa, \emph{On \(\ZZ_\ell\)-extensions of algebraic number fields}, Ann. of Math. \textbf{98} (1973), 246--326.
\bibitem{V-III-II-fw13} B. Ferrero and L. C. Washington, \emph{The Iwasawa invariant \(\mu_p\) vanishes for abelian number fields}, Ann. of Math. \textbf{109} (1979), 377--395. \href{https://doi.org/10.2307/1971116}{Original paper}.
\bibitem{V-III-brumer330} A. Brumer, \emph{On the units of algebraic number fields}, Mathematika \textbf{14} (1967), 121--124. \href{https://doi.org/10.1112/S0025579300003703}{Original paper}.
\end{thebibliography}\endgroup
\end{document}
